\documentclass[10pt, letterpaper]{article}

% Inhaltsverzeichnis für Pakettypen (nur für Übersicht im Header, wird nicht im Dokument angezeigt)
% 1. Seitenlayout und Ränder
% 2. Sprache und Zeichensatz
% 3. Mathematik und Theorem-Umgebungen
% 4. Eigene Makros
% 5. Diagramme und Grafiken
% 6. Tabellen und Aufzählungen
% 7. Inhaltsverzeichnis
% 8. Abschnittsüberschriften
% 9. Abstrakt-Umgebung
% 10. Todos/Notizen
% 11. Rahmen/Box-Umgebungen
% 12. Python-Integration
% 13. Literaturverwaltung
% 14. Hyperlinks
% 15. Absatzeinstellungen
% 16. Umgebungen
% 17  Graphik
% 18  Extra
% 00. Titel und Autor

% --- 1. Seitenlayout und Ränder ---
\usepackage[margin=3cm]{geometry}

% --- 2. Sprache und Zeichensatz ---
\usepackage[english]{babel}
\usepackage[T1]{fontenc}
\usepackage[utf8]{inputenc}

% --- 3. Mathematik und Theorem-Umgebungen ---
\usepackage{amsmath, amssymb, amsthm}
\usepackage{mathrsfs}
\DeclareMathOperator{\WF}{WF}

% --- 4. Eigene Makros ---
\usepackage{xcolor}
\newcommand{\SKP}{\langle\cdot,\cdot\rangle}
\newcommand{\R}{\mathbb{R}}
\newcommand{\N}{\mathbb{N}}
\newcommand{\Q}{\mathbb{Q}}
\newcommand{\Z}{\mathbb{Z}}
\newcommand{\C}{\mathbb{C}}
\newcommand{\entwurf}[1]{\textcolor{red}{#1}}

% --- 5. Diagramme und Grafiken ---
\usepackage{graphicx}
\graphicspath{{../../Library/Images/}}
\usepackage[export]{adjustbox}
\usepackage{tikz}
\usetikzlibrary{decorations.pathreplacing, arrows.meta, positioning}
\usepackage{tikz-cd}

% --- 6. Tabellen und Aufzählungen ---
\usepackage{enumitem}
\setlist[itemize]{left=0.5cm}

\newenvironment{romanenum}[1][]
  {%
    \ifx&#1&
    \else
      \textbf{#1}\quad
    \fi
    \begin{enumerate}[label=\roman*)]
  }
  {%
    \end{enumerate}%
  }

% --- 7. Inhaltsverzeichnis ---
\usepackage{tocloft}
\renewcommand{\cftsecfont}{\footnotesize}
\renewcommand{\cftsubsecfont}{\footnotesize}
\renewcommand{\cftsubsubsecfont}{\footnotesize}
\renewcommand{\cftsecpagefont}{\footnotesize}
\renewcommand{\cftsubsecpagefont}{\footnotesize}
\renewcommand{\cftsubsubsecpagefont}{\footnotesize}
\usepackage{etoc}

% --- 8. Abschnittsüberschriften ---
\usepackage{titlesec}
\titleformat{\section}{\normalfont\large\bfseries}{\thesection}{1em}{}
\titleformat{\subsection}{\normalfont\normalsize\bfseries}{\thesubsection}{0.5em}{}
\titleformat{\subsubsection}{\normalfont\normalsize\bfseries}{\thesubsubsection}{0.5em}{}
\setcounter{secnumdepth}{4}

% --- 9. Abstrakt-Umgebung ---
\usepackage{changepage}
\renewenvironment{abstract}
  {
    \begin{adjustwidth}{1.5cm}{1.5cm}
    \small
    \textsc{Abstract. –}%
  }
  {
    \end{adjustwidth}
  }

% --- 10. Todos/Notizen ---
\usepackage{todonotes}

% --- 11. Rahmen/Box-Umgebungen ---
\usepackage{mdframed}
\usepackage{tcolorbox}
\colorlet{shadecolor}{gray!25}

\newenvironment{customTheorem}
  {\vspace{10pt}%
   \begin{mdframed}[
     backgroundcolor=gray!20,
     linewidth=0pt,
     innertopmargin=10pt,
     innerbottommargin=10pt,
     skipabove=\dimexpr\topsep+\ht\strutbox\relax,
     skipbelow=\topsep,
   ]}
  {\end{mdframed}
   \vspace{10pt}%
  }

% --- 12. Python-Integration ---
% (Deaktiviert in dieser Version, aktiviere bei Bedarf)
% \usepackage{pythontex}
% \usepackage[makestderr]{pythontex}

% --- 13. Literaturverwaltung ---
\usepackage{csquotes}
\usepackage[backend=biber, style=alphabetic, citestyle=alphabetic]{biblatex}
\addbibresource{bibliography.bib}

% --- 14. Hyperlinks ---
\usepackage{hyperref}
\hypersetup{
  colorlinks   = true,
  urlcolor     = blue,
  linkcolor    = blue,
  citecolor    = blue,
  frenchlinks  = true
}

% --- 15. Absatzeinstellungen ---
\usepackage[parfill]{parskip}
\sloppy

% --- 16. Umgebungen ---
\usepackage{thmtools}

\newcommand{\CustomHeading}[3]{%
  \par\medskip\noindent%
  \textbf{#1 #2} \textnormal{(#3)}.\enskip%
}

\newenvironment{DEF}[2]{\begin{unitbox}\CustomHeading{Definition}{#1}{#2}}{\end{unitbox}}
\newenvironment{PROP}[2]{\begin{unitbox}\CustomHeading{Proposition}{#1}{#2}}{\end{unitbox}}
\newenvironment{THEO}[2]{\begin{unitbox}\CustomHeading{Theorem}{#1}{#2}}{\end{unitbox}}
\newenvironment{LEM}[2]{\begin{unitbox}\CustomHeading{Lemma}{#1}{#2}}{\end{unitbox}}
\newenvironment{KORO}[2]{\begin{unitbox}\CustomHeading{Corollar}{#1}{#2}}{\end{unitbox}}
\newenvironment{REM}[2]{\begin{unitbox}\CustomHeading{Remark}{#1}{#2}}{\end{unitbox}}
\newenvironment{EXA}[2]{\begin{unitbox}\CustomHeading{Example}{#1}{#2}}{\end{unitbox}}
\newenvironment{STUD}[2]{\begin{unitbox}\CustomHeading{Study}{#1}{#2}}{\end{unitbox}}
\newenvironment{CONC}[2]{\begin{unitbox}\CustomHeading{Concept}{#1}{#2}}{\end{unitbox}}
\newenvironment{OTH}[2]{\begin{unitbox}\CustomHeading{Other}{#1}{#2}}{\end{unitbox}}
\newenvironment{EXE}[2]{\begin{unitbox}\CustomHeading{Exercise}{#1}{#2}}{\end{unitbox}}
\newenvironment{MOT}[2]{\begin{unitbox}\CustomHeading{Motivation}{#1}{#2}}{\end{unitbox}}
\newenvironment{PROOF}[2]{\begin{unitbox}\CustomHeading{Proof}{#1}{#2}}{\end{unitbox}}

% --- Unit Umgebung für Source-Inhalte ---
\usepackage{mdframed}
\newmdenv[
  linewidth=1pt,
  topline=false,
  bottomline=false,
  rightline=false,
  leftmargin=0cm,
  rightmargin=0cm,
  skipabove=10pt,
  skipbelow=10pt,
  innertopmargin=0.5\baselineskip,
  innerbottommargin=0.5\baselineskip,
  backgroundcolor=gray!10,
  linecolor=gray
]{unitbox}

\newenvironment{unit}[1]
  {\begin{unitbox}\textbf{Unit #1}\par\smallskip}
  {\end{unitbox}}

% --- 17. ... ---

% --- 18. Extras ---
\usepackage{stmaryrd}
\usepackage{bbold}  % falls du \mathbb{1} nutzen willst

% --- 00. Titel und Autor ---
\title{Mein Titel}
\author{Tim Jaschik}
\date{\today}


\def\AA{\mathring{\mathrm{A}}}

\begin{document}

\maketitle
\rule{\textwidth}{0.5pt}
\begin{abstract}
Kurze Beschreibung …
\end{abstract}
\rule{\textwidth}{0.5pt}
\vspace{0.5cm}

\tableofcontents

\pagebreak


\section*{Normtopologien}


\subsection*{Normierte Räume und Prähilberträume}



Bezeichnungen: Stets sei $\mathbb{K}=\mathbb{R}$ oder $\mathbb{K}=\mathbb{C}$ und $E$ ein $\mathbb{K}$-Vektorraum.


\subsection*{Definition}
i) Eine Abbildung $\|\cdot\|: E \rightarrow \mathbb{R}$ heißt Halbnorm auf einem $\mathbb{K}$-Vektorraum $E$, falls für alle $x, y \in E$ und $\lambda \in \mathbb{K}$ gilt:
$$
\begin{array}{ll}
(N 1) & \|\lambda x\|=|\lambda|\|x\| \quad \text { (Homogenität) } \\
(N 2) & \|x+y\| \leq\|x\|+\|y\| \quad \text { (Dreiecksungleichung) }
\end{array}
$$
Es folgt dann: $\left\|0_{E}\right\|=\left\|0_{\mathbb{K}} \cdot 0_{E}\right\| \stackrel{(N 1)}{=}\left|0_{\mathbb{K}}\right|\left\|0_{E}\right\|=0$ und damit $0=\|x+(-x)\| \stackrel{(N 2)}{\leq} \|x\|+\|-x\| \stackrel{(N 1)}{=} 2\|x\|$, also
$$
(N 3) \quad\|x\| \geq 0 \text { für alle } x \in E \text {. }
$$

ii) Eine Halbnorm $\|\cdot\|$ auf $E$ heißt Norm, falls zusätzlich für alle $x \in E$ gilt:
$$
\text { (N4) } \quad\|x\|=0 \Rightarrow x=0 \text { (Definitheit) }
$$
In diesem Fall heißt das Paar $(E,\|\cdot\|)$ normierter Raum.




\begin{DEF}{FA-A25-01-01}{Halbnorm, Norm und normierter Raum}
Sei $E$ ein $\mathbb{K}$-Vektorraum, wobei $\mathbb{K}\in\{\mathbb{R},\mathbb{C}\}$.

\begin{enumerate}
\item[(i)] Eine Abbildung 
\[
\|\cdot\|:E\longrightarrow\mathbb{R}
\]
heißt \textbf{Halbnorm} auf $E$, falls für alle $x,y\in E$ und $\lambda\in\mathbb{K}$ gilt:
\[
\begin{array}{ll}
\text{(N1)} & \|\lambda x\|=|\lambda|\,\|x\| \quad\text{(Homogenität)}\\[0.4em]
\text{(N2)} & \|x+y\|\le\|x\|+\|y\| \quad\text{(Dreiecksungleichung)}.
\end{array}
\]
Aus diesen Axiomen folgt unmittelbar:
\[
\|0_E\|=0
\quad\text{und damit}\quad
\|x\|\ge0\ \text{für alle }x\in E.
\tag{N3}
\]

\item[(ii)] Eine Halbnorm $\|\cdot\|$ auf $E$ heißt \textbf{Norm}, wenn zusätzlich gilt:
\[
\text{(N4)} \quad
\|x\|=0\ \Longrightarrow\ x=0.
\quad\text{(Definitheit)}
\]
\item[(iii)] Das Paar $(E,\|\cdot\|)$ heißt \textbf{normierter Raum}.
\end{enumerate}
\end{DEF}




\subsection*{Definition}


\begin{DEF}{FA-A25-01-02}{Äquivalenzrelation zwischen Halbnormen}
Zwei Halbnormen $\|\cdot\|_{1},\|\cdot\|_{2}$ auf $E$ heißen äquivalent, falls $a, b>0$ existieren mit
$$
a\|x\|_{1} \leq\|x\|_{2} \leq b\|x\|_{1} \text { für alle } x \in E \text {. }
$$
Wir schreiben dann: $\|\cdot\|_{1} \sim\|\cdot\|_{2}$.

Bemerkung: „\~{}" definiert Äquivalenzrelationen auf der Menge der Halbnormen und auf der Menge der Normen auf $E$.
\end{DEF}


\subsection*{Bemerkung und Beispiel}


i) $\operatorname{dim} E<\infty \Rightarrow$ Alle Normen auf $E$ sind äquivalent (bekannt aus dem Grundstudium).\\
ii) $E=\mathbb{K}^{N}$. Für $x \in \mathbb{K}^{N}$ und $p \in[1, \infty]$ sei

$$
|x|_{p}:= \begin{cases}\left(\left|x_{1}\right|^{p}+\ldots+\left|x_{N}\right|^{p}\right)^{\frac{1}{p}}, & p<\infty \\ \max _{1 \leq i \leq N}\left|x_{i}\right| & p=\infty\end{cases}
$$

$\underline{\text { Bekannt aus der Analysis II: }}|\cdot|_{p}: E \rightarrow \mathbb{R}$ ist eine Norm auf $E$. Dabei gilt die Höldersche Ungleichung:\\
$\sum_{i=1}^{N}\left|x_{i} y_{i}\right| \leq|x|_{p}|y|_{p^{\prime}}$ für alle $x, y \in \mathbb{K}^{N}$ mit $\quad p^{\prime}:=\left\{\begin{array}{lll}\frac{p}{p-1} & \text { falls } & 1<p<\infty \\ \infty & \text { falls } & p=1 \\ 1 & \text { falls } & p=\infty\end{array}\right.$\\
$p^{\prime}$ heißt konjugierter Exponent zu $p$. Beachte: $p, q \in(1, \infty)$ konjugiert $\Longleftrightarrow \frac{1}{p}+\frac{1}{q}=1$\\
iii) Sei $M$ eine beliebige Menge, und sei $\ell^{\infty}(M, \mathbb{K})$ der $\mathbb{K}$-Vektorraum der beschränkten Funktionen $u: M \rightarrow \mathbb{K}$ mit punktweisen linearen Operationen.\\
Dann ist $\ell^{\infty}(M, \mathbb{K})$ ein normierter $\mathbb{K}$-Vektorraum mit der Supremumsnorm $\|\cdot\|_{\infty}$ definiert durch $\|u\|_{\infty}:=\sup _{z \in M}|u(z)|$.\\
Der Beweis der Normeigenschaften ist sehr einfach.\\
Ist $M=\mathbb{N}$, so notiert man die Funktionswerte von $u \in \ell^{\infty}(\mathbb{N}, \mathbb{K})$ als Folge, d. h. man schreibt $u=\left(u_{k}\right)_{k}$ mit $u_{k}=u(k)$. Dann ist $\|u\|_{\infty}=\sup _{k \in \mathbb{N}}\left|u_{k}\right|$.

Kurzschreibweisen: $\ell^{\infty}(M)$ statt $\ell^{\infty}(M, \mathbb{K})$, $\ell^{\infty}$ statt $\ell^{\infty}(\mathbb{N})=\ell^{\infty}(\mathbb{N}, \mathbb{K})$.







\begin{EXA}{FA-A25-01-03}{Äquivalenz von Normen in endlichdimensionalen Räumen}
Ist $E$ ein endlichdimensionaler $\mathbb{K}$-Vektorraum, d.\,h.\ $\dim E<\infty$, so sind alle Normen auf $E$ \emph{äquivalent}.  
Das heißt: Für zwei beliebige Normen $\|\cdot\|_1$, $\|\cdot\|_2$ auf $E$ existieren Konstanten $c,C>0$ mit
\[
c\,\|x\|_1 \le \|x\|_2 \le C\,\|x\|_1
\quad\text{für alle }x\in E.
\]
Insbesondere induzieren alle Normen auf $E$ dieselbe Topologie.
\end{EXA}




\begin{EXA}{FA-A25-01-04}{$p$-Normen auf $\mathbb{K}^N$}
Für $E=\mathbb{K}^N$ und $p\in[1,\infty]$ sei für $x=(x_1,\dots,x_N)\in\mathbb{K}^N$ definiert:
\[
|x|_p :=
\begin{cases}
\displaystyle \Bigl(\sum_{i=1}^N |x_i|^p\Bigr)^{1/p}, & p<\infty,\\[1em]
\displaystyle \max_{1\le i\le N} |x_i|, & p=\infty.
\end{cases}
\]
Dann ist $|\cdot|_p$ eine Norm auf $\mathbb{K}^N$.

\smallskip
Dabei gilt die \textbf{Hölder’sche Ungleichung:}
\[
\sum_{i=1}^N |x_i y_i|
\;\le\;
|x|_p\,|y|_{p'},
\quad
x,y\in\mathbb{K}^N,
\]
wobei $p'$ der zu $p$ \emph{konjugierte Exponent} ist, definiert durch
\[
p' :=
\begin{cases}
\frac{p}{p-1}, & 1<p<\infty,\\
\infty, & p=1,\\
1, & p=\infty,
\end{cases}
\qquad\text{so daß}\qquad
\frac{1}{p}+\frac{1}{p'}=1.
\]
\end{EXA}




\begin{EXA}{FA-A25-01-05}{$\ell^\infty(M,\mathbb{K})$ mit Supremumsnorm}
Sei $M$ eine beliebige Menge.  
Dann ist
\[
\ell^{\infty}(M,\mathbb{K})
:=\bigl\{u:M\to\mathbb{K}\mid u\text{ ist beschränkt}\bigr\}
\]
ein $\mathbb{K}$-Vektorraum (mit punktweiser Addition und Skalarmultiplikation).

Auf $\ell^{\infty}(M,\mathbb{K})$ definieren wir die \textbf{Supremumsnorm}
\[
\|u\|_{\infty} := \sup_{z\in M}|u(z)|.
\]
Dann ist $(\ell^{\infty}(M,\mathbb{K}),\|\cdot\|_\infty)$ ein normierter Raum.

\smallskip
Für $M=\mathbb{N}$ schreiben wir $u=(u_k)_{k\in\mathbb{N}}$ mit $u_k=u(k)$ und
\[
\|u\|_{\infty} = \sup_{k\in\mathbb{N}}|u_k|.
\]
Kurzschreibweise:
\[
\ell^{\infty}(M):=\ell^{\infty}(M,\mathbb{K}),
\qquad
\ell^{\infty}:=\ell^{\infty}(\mathbb{N}).
\]
\end{EXA}






Wir wollen im Folgenden weitere unendlich-dimensionale normierte Räume durch ein allgemeines Prinzip konstruieren. Dieses Prinzip liefert auch einen (alternativen) Beweis der Halbnormeigenschaften von $|\cdot|_{p}$ aus $1.3($ ii $)$.

\subsection*{Definition}
(a) Eine Teilmenge $A \subset E$ heißt\\
(i) konvex, wenn für alle $x, y \in A, \lambda \in[0,1]$ gilt: $(1-\lambda) x+\lambda y \in A$;\\
(ii) symmetrisch, wenn gilt: $x \in A \Longrightarrow \theta x \in A$ für alle $\theta \in \mathbb{K}$ mit $|\theta|=1$;\\
(iii) absolut konvex, wenn $A$ konvex und symmetrisch ist.\\
(iv) absorbierend, wenn für jedes $x \in E$ ein $s>0$ existiert mit $s x \in A$.\\
(b) Eine auf einer konvexen Teilmenge $A \subset E$ definierte Funktion

$$
f: A \rightarrow \mathbb{R} \cup\{\infty\}
$$

heißt konvex, wenn $f((1-\lambda) x+\lambda y) \leq(1-\lambda) f(x)+\lambda f(y)$ für alle $x, y \in A$, $\lambda \in[0,1]$ gilt.\\
Ist $A$ absolut konvex, so heißt $f$ absolut konvex, wenn $f$ konvex und gerade ist. Dabei heißt $f$ gerade, wenn $f(\theta x)=f(x)$ für alle $x \in A$ und $\theta \in \mathbb{K}$ mit $|\theta|=1$ gilt.



\begin{DEF}{FA-A25-01-06}{Konvexität, Symmetrie und absolute Konvexität}
Sei $E$ ein $\mathbb{K}$-Vektorraum, wobei $\mathbb{K}\in\{\mathbb{R},\mathbb{C}\}$.

\begin{enumerate}
\item[(a)] Eine Teilmenge $A\subset E$ heißt
\begin{enumerate}
\item[(i)] \textbf{konvex}, falls für alle $x,y\in A$ und alle $\lambda\in[0,1]$ gilt:
\[
(1-\lambda)x+\lambda y \in A.
\]

\item[(ii)] \textbf{symmetrisch}, falls für alle $\theta\in\mathbb{K}$ mit $|\theta|=1$ gilt:
\[
x\in A \ \Longrightarrow\ \theta x\in A.
\]

\item[(iii)] \textbf{absolut konvex}, falls $A$ sowohl konvex als auch symmetrisch ist.

\item[(iv)] \textbf{absorbierend}, falls für jedes $x\in E$ ein $s>0$ existiert mit
\[
s x \in A.
\]
\end{enumerate}

\item[(b)] Sei $A\subset E$ konvex und
\[
f:A\longrightarrow\mathbb{R}\cup\{\infty\}.
\]
Dann heißt $f$ \textbf{konvex}, wenn für alle $x,y\in A$ und $\lambda\in[0,1]$ gilt:
\[
f\bigl((1-\lambda)x+\lambda y\bigr)\le (1-\lambda)f(x)+\lambda f(y).
\]

Ist $A$ \emph{absolut konvex}, so heißt $f$ \textbf{absolut konvex}, wenn $f$ konvex und zusätzlich \textbf{gerade} ist, d.\,h.
\[
f(\theta x)=f(x)\quad\text{für alle }x\in A\text{ und }\theta\in\mathbb{K}\text{ mit }|\theta|=1.
\]
\end{enumerate}
\end{DEF}




\subsection*{Bemerkung und Beispiel}
i). Ist $\|\cdot\|$ eine Halbnorm auf $E$ und $r>0$, so sind die Mengen $U_{r}(0):=\{x \in E$ : $\|x\|<r\}$ und $B_{r}(0):=\{x \in E:\|x\| \leq r\}$ absolut konvex und absorbierend.\\
ii). $A \subset E$ ist absolut konvex genau dann, wenn $\lambda x+\mu y \in A$ für alle $x, y \in A$, $\lambda, \mu \in \mathbb{K}$ mit $|\lambda|+|\mu| \leq 1$ gilt.\\
iii). Ist $A \subset E$ konvex, so gilt $\sum_{i=1}^{n} \lambda_{i} x_{i} \in A$ für alle $n \in \mathbb{N}, x_{1}, \ldots, x_{n} \in A$ und $\lambda_{1}, \ldots, \lambda_{n} \in[0, \infty)$ mit $\sum_{i=1}^{n} \lambda_{i}=1$.\\
iv). Ist $A \subset E$ absolut konvex, so gilt $\sum_{i=1}^{n} \lambda_{i} x_{i} \in A$ für alle $n \in \mathbb{N}, x_{1}, \ldots, x_{n} \in A$ und $\lambda_{1}, \ldots, \lambda_{n} \in \mathbb{K}$ mit $\sum_{i=1}^{n}\left|\lambda_{i}\right| \leq 1$.\\
v). Sei $I \subset \mathbb{R}$ ein Intervall. Dann ist $f \in C^{1}(I)$ genau dann konvex, wenn $f^{\prime}$ monoton wachsend ist.

Die Beweise von i)-v) sind allesamt recht einfach.



Grundfakten zu (absolut) konvexen Mengen und konvexen Funktionen


\begin{PROP}{FA-A25-01-07}{Elementare Eigenschaften zu (absolut) konvexen Mengen und konvexen Funktionen}
Sei $E$ ein Vektorraum über $\mathbb{K}\in\{\mathbb{R},\mathbb{C}\}$.
\begin{enumerate}
\item Ist $\|\cdot\|$ eine Halbnorm auf $E$ und $r>0$, so sind
\[
U_{r}(0):=\{x\in E:\|x\|<r\}
\quad\text{und}\quad
B_{r}(0):=\{x\in E:\|x\|\le r\}
\]
absolut konvex und absorbierend.
\item Eine Menge $A\subset E$ ist genau dann absolut konvex, wenn
\[
\lambda x+\mu y\in A \quad\text{für alle }x,y\in A,\ \lambda,\mu\in\mathbb{K}\text{ mit }|\lambda|+|\mu|\le 1.
\]
\item Ist $A\subset E$ konvex, so gilt
\[
\sum_{i=1}^{n}\lambda_i x_i\in A
\quad\text{für alle }n\in\mathbb{N},\ x_i\in A,\ \lambda_i\ge 0,\ \sum_{i=1}^{n}\lambda_i=1.
\]
\item Ist $A\subset E$ absolut konvex, so gilt
\[
\sum_{i=1}^{n}\lambda_i x_i\in A
\quad\text{für alle }n,\ x_i\in A,\ \lambda_i\in\mathbb{K},\ \sum_{i=1}^{n}|\lambda_i|\le 1.
\]
\item Sei $I\subset\mathbb{R}$ ein Intervall und $f\in C^{1}(I)$. 
Dann ist $f$ konvex genau dann, wenn $f'$ monoton wachsend ist.
\end{enumerate}
\end{PROP}





\begin{PROOF}{FA-A25-01-08}{P: Elementare Eigenschaften zu (absolut) konvexen Mengen und konvexen Funktionen}
\textbf{(i)} (Absolut konvex) Sei $x,y$ in $U_r(0)$ (bzw.\ $B_r(0)$) und $|\lambda|+|\mu|\le 1$.
Dann
\[
\|\lambda x+\mu y\|\le |\lambda|\,\|x\|+|\mu|\,\|y\| < (|\lambda|+|\mu|)\,r \le r
\]
(bzw.\ $\le r$), also $\lambda x+\mu y\in U_r(0)$ (bzw.\ $B_r(0)$). Damit ist die Menge absolut konvex.
(Absorbierend) Für beliebiges $x\in E$ und $r>0$ wähle $t:=\max\{1,\|x\|/r\}>0$. Dann
$\|x\| \le rt$ und somit $x\in t B_r(0)$; ebenso mit strikt $<$ für $U_r(0)$.
\medskip

\textbf{(ii)} „$\Rightarrow$“: Ist $A$ absolut konvex (konvex und balanciert), so ist für $|\lambda|,|\mu|\le 1$
die Kombination $\lambda x+\mu y\in A$ eine Folge von Balanciertheit (Skalierung) und Konvexität:
$\lambda x,\,\mu y\in A$ und für $t:=\frac{|\lambda|}{|\lambda|+|\mu|}$ gilt
$\lambda x+\mu y=(|\lambda|+|\mu|)\big(t\,\operatorname{sgn}(\lambda)x+(1-t)\,\operatorname{sgn}(\mu)y\big)\in A$.
„$\Leftarrow$“: Setze in der Bedingung $(\lambda,\mu)=(t,1-t)$ mit $t\in[0,1]$; dann ist $A$ konvex.
Setze ferner $y=0$ und $\mu=0$, $|\lambda|\le 1$; dann folgt balanciert: $\lambda x\in A$ für $x\in A$. 

\medskip
\textbf{(iii)} Induktion nach $n$. Für $n=2$ ist es die Definition. 
Sei die Aussage für $n$ richtig. Für $n+1$ setze
$t:=\sum_{i=1}^{n}\lambda_i\in[0,1]$ und $y:=\frac{1}{t}\sum_{i=1}^{n}\lambda_i x_i\in A$ (Induktionsannahme).
Dann ist $\sum_{i=1}^{n+1}\lambda_i x_i=t\,y+(1-t)\,x_{n+1}\in A$ (Konvexität). 

\medskip
\textbf{(iv)} Schreibe
\[
\sum_{i=1}^{n}\lambda_i x_i=\Big(\sum_{i=1}^{n}|\lambda_i|\Big)\,
\sum_{i=1}^{n}\frac{|\lambda_i|}{\sum_{j}|\lambda_j|}\,\operatorname{sgn}(\lambda_i)\,x_i.
\]
Da $A$ balanciert ist, liegen $\operatorname{sgn}(\lambda_i)\,x_i\in A$. 
Da $A$ konvex ist, liegt die konvexe Kombination in $A$. Schließlich skaliert man 
mit dem Faktor $\sum_i|\lambda_i|\le 1$ erneut per Balanciertheit. 

\medskip
\textbf{(v)} „$\Rightarrow$“: Für $x<y<z$ gilt bei Konvexität die Sekantensteigungs-Monotonie
\[
\frac{f(y)-f(x)}{y-x}\le \frac{f(z)-f(y)}{z-y}.
\]
Mit dem Mittelwertsatz existieren $\xi\in(x,y)$ und $\eta\in(y,z)$ mit
$f'(\xi)=\frac{f(y)-f(x)}{y-x}$ und $f'(\eta)=\frac{f(z)-f(y)}{z-y}$; also $f'(\xi)\le f'(\eta)$.
Daraus folgt, dass $f'$ monoton wächst.

„$\Leftarrow$“: Sei $f'$ monoton wachsend. Dann sind die Sekantensteigungen 
$x<y<z\Rightarrow \frac{f(y)-f(x)}{y-x}\le \frac{f(z)-f(y)}{z-y}$.
Dies ist äquivalent zur Konvexität (z.\,B.\ per Integralformel
$f(y)=f(x)+\int_{x}^{y}f'(t)\,dt$ und Monotonie von $f'$).
\end{PROOF}







\subsection*{Satz}



\begin{EXA}{FA-A25-01-09}{Induzierte Halbnorm auf Vektorraum durch absorbierende und absolut konvexe Menge}
Sei $A \subset E$ absolut konvex und absorbierend. Dann wird durch
$$
\|x\|_{A}:=\inf \left\{t>0: \frac{x}{t} \in A\right\}=\inf \{t>0: x \in t A\} \quad \text { für } x \in E
$$
eine Halbnorm auf $E$ definiert.


Beweis. Da $A$ absorbierend ist, erhält man $\|x\|_{A}<\infty$ für alle $x \in E$.\\
Zur Homogenität: Für $x \in E$ und $\lambda \in \mathbb{K}$ gilt
$$
\begin{aligned}
\|\lambda x\|_{A} & =\inf \{t>0: \lambda x \in t A\}=\inf \{t>0:|\lambda| x \in t A\} \\
& =\inf \{|\lambda| s: s>0, x \in s A\}=|\lambda|\|x\|_{A} .
\end{aligned}
$$
Zur Dreiecksungleichung: Seien $x, y \in E$, und seien $s, t>0$ mit $\frac{x}{s}, \frac{y}{t} \in A$. Mit $\lambda=\frac{t}{t+s}$ (also $1-\lambda=\frac{s}{t+s}$ ) ist dann
$$
\frac{x+y}{s+t}=(1-\lambda) \frac{x}{s}+\lambda \frac{y}{t} \in A
$$
aufgrund der Konvexität von $A$, und somit folgt $\|x+y\|_{A} \leq s+t$. Durch Infimumbildung in $s$ und $t$ folgt $\|x+y\|_{A} \leq\|x\|_{A}+\|y\|_{A}$. \(\square\)
\end{EXA}




\subsection*{Definition und Satz}



\begin{EXA}{FA-A25-01-10}{$\ell^{p}(\mathbb{N}, \mathbb{K}$ als normierter Vektorraum bzgl $p$-Norm}
Für $p \in[1, \infty)$ sei $\ell^{p}(\mathbb{N}, \mathbb{K})$ die Menge aller Folgen $x=\left(x_{k}\right)_{k}$ in $\mathbb{K}$ mit $\sum_{k \in \mathbb{N}}\left|x_{k}\right|^{p}<\infty$. $\ell^{p}(\mathbb{N}, \mathbb{K}$ ) ist ein normierter $\mathbb{K}$-Vektorraum (bzgl. komponentenweiser Addition und Skalarmultiplikation) mit der Norm $\|\cdot\|_{p}$ definiert durch
$$
\|x\|_{p}=\left(\sum_{k \in \mathbb{N}}\left|x_{k}\right|^{p}\right)^{\frac{1}{p}} \quad \text { für } x=\left(x_{k}\right)_{k} \in \ell^{p}(\mathbb{N}, \mathbb{K}) .
$$


Beweis. Offensichtlich ist die Nullfolge $0 \in \ell^{p}(\mathbb{N}, \mathbb{K})$. Seien $x=\left(x_{k}\right)_{k}, y=\left(y_{k}\right)_{k} \in \ell^{p}(\mathbb{N}, \mathbb{K})$ und $\lambda \in \mathbb{K}$. Dann ist
$$
\sum_{k \in \mathbb{N}}\left|\lambda x_{k}\right|^{p}=|\lambda|^{p} \sum_{k}\left|x_{k}\right|^{p}<\infty,
$$
und somit $\lambda x=\left(\lambda x_{k}\right)_{k} \in \ell^{p}(\mathbb{N}, \mathbb{K})$. Da die Funktion $(0, \infty) \rightarrow \mathbb{R}, t \mapsto t^{p}$ gemäß 1.5 v ) konvex ist, gilt ferner
$$
\sum_{k \in \mathbb{N}}\left|x_{k}+y_{k}\right|^{p} \leq \sum_{k \in \mathbb{N}} 2^{p}\left(\frac{\left|x_{k}\right|}{2}+\frac{\left|y_{k}\right|}{2}\right)^{p} \leq 2^{p-1} \sum_{k \in \mathbb{N}}\left(\left|x_{k}\right|^{p}+\left|y_{k}\right|^{p}\right)<\infty
$$
und somit $x+y=\left(x_{k}+y_{k}\right)_{k} \in \ell^{p}(\mathbb{N}, \mathbb{K})$. Also ist $\ell^{p}(\mathbb{N}, \mathbb{K})$ ein $\mathbb{K}$-Vektorraum. Zudem ist
$$
A:=\left\{x \in \ell^{p}(\mathbb{N}, \mathbb{K}): \sum_{k \in \mathbb{N}}\left|x_{k}\right|^{p} \leq 1\right\}
$$
offensichtlich eine absorbierende Teilmenge von $\ell^{p}(\mathbb{N}, \mathbb{K})$, und aufgrund der Konvexität der Funktion $t \mapsto t^{p}$ ist $A$ auch absolut konvex. Somit definiert
$$
\|x\|_{A}=\inf \left\{t>0: \sum_{k \in \mathbb{N}}\left|\frac{x_{k}}{t}\right|^{p} \leq 1\right\}=\inf \left\{t>0: \sum_{k \in \mathbb{N}}\left|x_{k}\right|^{p} \leq t^{p}\right\}=\left(\sum_{k \in \mathbb{N}}\left|x_{k}\right|^{p}\right)^{\frac{1}{p}}=\|x\|_{p} .
$$
eine Halbnorm auf $\ell^{p}(\mathbb{N}, \mathbb{K})$ gemäß 1.6.

Zum Beweis der Definitheit (N4) beachte man:
$$
\|x\|_{p}=0 \Longrightarrow \sum_{k \in \mathbb{N}}\left|x_{k}\right|^{p}=0 \Longrightarrow x_{k}=0 \text { für alle } k \in \mathbb{N} \Longrightarrow x=0 \text {. }
$$ \(\square\)
\end{EXA}




2.8 Bemerkung. (a) Analog definiert man $\ell^{p}(\mathcal{I}, \mathbb{K})$ für eine beliebige höchstens abzählbare Indexmenge $\mathcal{I}$ anstelle von $\mathbb{N}$, z.B. $\mathcal{I}=\mathbb{Z}$. Kurzschreibweisen: $\ell^{p}(\mathcal{I})$ statt $\ell^{p}(\mathcal{I}, \mathbb{K})$ und $\ell^{p}$ statt $\ell^{p}(\mathbb{N})=\ell^{p}(\mathbb{N}, \mathbb{K})$. Man beachte: Der Fall $\mathcal{I}=\{1, \ldots, N\}$ führt wieder zurück auf $\mathbb{K}^{N}$.




\begin{DEF}{FA-A25-01-11}{Verallgemeinerte $\ell^{p}$-Räume über beliebigen Indexmengen}
Sei $I$ eine \textbf{höchstens abzählbare Indexmenge} und $\mathbb{K}\in\{\mathbb{R},\mathbb{C}\}$.

\begin{enumerate}
\item Für $1\le p<\infty$ definieren wir
\[
\ell^{p}(I,\mathbb{K})
:=\Bigl\{\,x=(x_{i})_{i\in I}\subset\mathbb{K}
\;\Big|\;
\sum_{i\in I}|x_{i}|^{p}<\infty
\Bigr\},
\quad
\|x\|_{p}:=\Bigl(\sum_{i\in I}|x_{i}|^{p}\Bigr)^{1/p}.
\]

\item Für $p=\infty$ setzen wir
\[
\ell^{\infty}(I,\mathbb{K})
:=\Bigl\{\,x=(x_{i})_{i\in I}\subset\mathbb{K}
\;\Big|\;
\sup_{i\in I}|x_{i}|<\infty
\Bigr\},
\quad
\|x\|_{\infty}:=\sup_{i\in I}|x_{i}|.
\]
\end{enumerate}

\textbf{Kurzschreibweisen:}
\[
\ell^{p}(I):=\ell^{p}(I,\mathbb{K}), 
\quad
\ell^{p}:=\ell^{p}(\mathbb{N},\mathbb{K}).
\]

\textbf{Bemerkung.}
Im Spezialfall $I=\{1,\ldots,N\}$ gilt
\[
\ell^{p}(I,\mathbb{K})=\mathbb{K}^{N}
\quad\text{mit}\quad
\|x\|_{p}=\Bigl(\sum_{i=1}^{N}|x_{i}|^{p}\Bigr)^{1/p},
\]
sodass die endliche Dimension als Spezialfall der $\ell^{p}$-Räume erscheint.
\end{DEF}




(b) Sind $p, q \in[1, \infty]$ konjugierte Exponenten, so gilt die Höldersche Ungleichung

$$
\sum_{k \in \mathcal{I}}\left|x_{k} y_{k}\right| \leq\|x\|_{p}\|y\|_{q} \quad \text { für } x=\left(x_{k}\right)_{k} \in \ell^{p}(\mathcal{I}, \mathbb{K}), y=\left(y_{k}\right)_{k} \in \ell^{q}(\mathcal{I}, \mathbb{K})
$$

Im Fall $p=\infty, q=1$ sieht man dies direkt:

$$
\sum_{k \in \mathcal{I}}\left|x_{k} y_{k}\right| \leq \sum_{k \in \mathcal{I}}\|x\|_{\infty}\left|y_{k}\right|=\|x\|_{\infty}\|y\|_{1}
$$

Im Fall $p, q \in(1, \infty)$ verwenden wir die Youngsche Ungleichung

$$
a b \leq \frac{a^{p}}{p}+\frac{b^{q}}{q} \quad \text { für } a, b \geq 0 .
$$

Diese Ungleichung beweist man leicht durch Minimierung der Funktion $x \mapsto \frac{x^{p}}{p}-x b$ in $[0, \infty)$ für festes $b \geq 0$. Unter Verwendung von (2.1) folgt nun mit $s=\|x\|_{p}$, $t=\|y\|_{q}:$

$$
\begin{aligned}
\sum_{k \in \mathcal{I}}\left|x_{k} y_{k}\right|=s t \sum_{k \in \mathcal{I}}\left|\frac{x_{k}}{s} \cdot \frac{y_{k}}{t}\right| & \leq s t \sum_{k \in \mathcal{I}}\left(\frac{1}{p}\left|\frac{x_{k}}{s}\right|^{p}+\frac{1}{q}\left|\frac{y_{k}}{t}\right|^{q}\right) \\
& =s t\left(\left.\frac{1}{p}| | \frac{x}{s}\right|_{p} ^{p}+\left.\frac{1}{q}| | \frac{y}{t}\right|_{q} ^{q}\right)=s t\left(\frac{1}{p}+\frac{1}{q}\right)=s t
\end{aligned}
$$




\begin{THEO}{FA-A25-01-12}{Höldersche Ungleichung in allg. $\ell^p$-Räumen}
Seien $p,q\in[1,\infty]$ \textbf{konjugierte Exponenten}, das heißt $\tfrac{1}{p}+\tfrac{1}{q}=1$, und sei $\mathcal{I}$ eine höchstens abzählbare Indexmenge.  
Dann gilt für alle Folgen 
\[
x=(x_k)_{k\in\mathcal{I}}\in\ell^{p}(\mathcal{I},\mathbb{K}), 
\quad
y=(y_k)_{k\in\mathcal{I}}\in\ell^{q}(\mathcal{I},\mathbb{K})
\]
die \textbf{Höldersche Ungleichung}:
\[
\boxed{
\sum_{k\in\mathcal{I}}|x_k y_k|
\;\le\;
\|x\|_{p}\,\|y\|_{q}.
}
\]
\end{THEO}



\begin{PROOF}{FA-A25-01-13}{P: Höldersche Ungleichung in allg. $\ell^p$-Räumen}
\textbf{Fall 1: $p=\infty$, $q=1$.}  
Dann gilt unmittelbar
\[
\sum_{k\in\mathcal{I}}|x_k y_k|
\le
\sum_{k\in\mathcal{I}}\|x\|_{\infty}\,|y_k|
=\|x\|_{\infty}\sum_{k\in\mathcal{I}}|y_k|
=\|x\|_{\infty}\|y\|_{1}.
\]

\medskip
\textbf{Fall 2: $p,q\in(1,\infty)$.}  
Wir verwenden die \textbf{Youngsche Ungleichung}
\[
ab\le \frac{a^{p}}{p}+\frac{b^{q}}{q},\qquad a,b\ge 0.
\]
Diese folgt durch Minimierung der Funktion 
$f(x)=\tfrac{x^{p}}{p}-xb$ auf $[0,\infty)$ bei festem $b\ge0$.

Setze $s=\|x\|_{p}$ und $t=\|y\|_{q}$.  
Dann gilt für $s,t>0$:
\[
\sum_{k\in\mathcal{I}}|x_k y_k|
=s\,t\sum_{k\in\mathcal{I}}\Big|\frac{x_k}{s}\cdot\frac{y_k}{t}\Big|
\le s\,t\sum_{k\in\mathcal{I}}
\Big(
\frac{1}{p}\Big|\frac{x_k}{s}\Big|^{p}
+\frac{1}{q}\Big|\frac{y_k}{t}\Big|^{q}
\Big).
\]
Da $\|\frac{x}{s}\|_{p}=1$ und $\|\frac{y}{t}\|_{q}=1$, folgt
\[
\sum_{k\in\mathcal{I}}|x_k y_k|
\le s\,t\Big(\frac{1}{p}+\frac{1}{q}\Big)
=s\,t=\|x\|_{p}\|y\|_{q}.
\]

\medskip
\textbf{Fall 3: $p=1$, $q=\infty$.}  
Analog zum ersten Fall ergibt sich
\[
\sum_{k\in\mathcal{I}}|x_k y_k|
\le
\|x\|_{1}\|y\|_{\infty}.
\]

Damit sind alle Fälle abgedeckt und die Aussage bewiesen.
\end{PROOF}


2.9 Satz. Sei $p \in[1, \infty)$.\\
(a) Für $q \in(p, \infty]$ gilt $\ell^{p} \subseteq \ell^{q}, \ell^{p} \neq \ell^{q}$, und $\|x\|_{q} \leq\|x\|_{p}$ für alle $x \in \ell^{p}$, aber $\|\cdot\|_{q}$ und $\|\cdot\|_{p}$ sind nicht äquivalent auf $\ell^{p}$.\\
(b) $\underset{\substack{q \rightarrow \infty \\ q \geq p}}{\lim }\|x\|_{q}=\|x\|_{\infty}$ für alle $x \in \ell^{p}$.

Beweis. Übung.\\




\begin{PROP}{FA-A25-01-14}{Einbettung und Grenzfall für $\ell^p$-Normen}
Sei $p\in[1,\infty)$.

\textbf{(a)} Für jedes $q\in(p,\infty]$ gilt
\[
\ell^{p}\subseteq \ell^{q},\qquad \|x\|_{q}\le \|x\|_{p}\quad\text{für alle }x\in \ell^{p},
\]
und die Inklusion ist echt, d.\,h.\ $\ell^{p}\ne \ell^{q}$. Zudem sind die Normen $\|\cdot\|_{p}$ und
$\|\cdot\|_{q}$ auf $\ell^{p}$ \emph{nicht} äquivalent.

\textbf{(b)} Für jedes $x\in \ell^{p}$ gilt
\[
\lim_{\substack{q\to\infty\\ q\ge p}}\|x\|_{q} = \|x\|_{\infty}.
\]
\end{PROP}


\begin{PROOF}{FA-A25-01-15}{P: Einbettung und Grenzfall für $\ell^p$-Normen}
(a) \emph{Inklusion und Normabschätzung.} Sei $x=(x_k)\in \ell^{p}$ und $q>p$.
Zerlege die Indexmenge in $A:=\{k:\,|x_k|\ge 1\}$ und $A^{c}$.
Da $\sum_k |x_k|^{p}<\infty$, ist $A$ endlich (sonst divergiert $\sum_{k\in A}|x_k|^{p}\ge \sum_{k\in A}1=\infty$).
Für $k\in A^{c}$ gilt $|x_k|^{q}\le |x_k|^{p}$ (weil $|x_k|<1$ und $q>p$), für $k\in A$ gilt
$|x_k|^{q}\le |x_k|^{p}$ (da $A$ endlich und $|x_k|<\infty$).
Somit
\[
\sum_k |x_k|^{q}\le \sum_k |x_k|^{p}<\infty,\quad\text{also }x\in \ell^{q}\ \text{und}\ \|x\|_{q}\le \|x\|_{p}.
\]

\emph{Echte Inklusion.} Wähle eine Folge $x_k:=k^{-\alpha}$ mit
\[
\frac{1}{q}<\alpha<\frac{1}{p}.
\]
Dann ist $\sum_k k^{-\alpha q}$ konvergent ($\alpha q>1$), aber
$\sum_k k^{-\alpha p}$ divergent ($\alpha p\le 1$). Also $x\in \ell^{q}\setminus \ell^{p}$.

\emph{Nichtäquivalenz der Normen auf $\ell^{p}$.} Setze für $n\in\mathbb{N}$
\[
x^{(n)}:=\big(\underbrace{n^{-1/p},\dots,n^{-1/p}}_{n\text{ Einträge}},0,0,\dots\big).
\]
Dann $\|x^{(n)}\|_{p}=(n\cdot n^{-1})^{1/p}=1$, aber
\[
\|x^{(n)}\|_{q}=\big(n\cdot n^{-q/p}\big)^{1/q}=n^{\,1/q-1/p}\xrightarrow[n\to\infty]{}0
\quad(\text{da }1/q-1/p<0).
\]
Damit kann keine Konstante $c>0$ existieren mit $c\|x\|_{p}\le \|x\|_{q}$ für alle $x\in \ell^{p}$;
also sind $\|\cdot\|_{p}$ und $\|\cdot\|_{q}$ nicht äquivalent auf $\ell^{p}$.

\smallskip
(b) Setze $M:=\|x\|_{\infty}$. Für alle $q\ge p$ gilt $M\le \|x\|_{q}\le \|x\|_{p}$, also ist
$\big(\|x\|_{q}\big)_{q\ge p}$ monoton fallend und nach unten durch $M$ beschränkt.
Sei $\varepsilon>0$. Wähle eine \emph{endliche} Menge $F\subset \mathbb{N}$ so, dass
\[
\sum_{k\notin F}|x_k|^{p}<\varepsilon^{p}.
\]
Dann für alle $q\ge p$:
\[
\|x\|_{q}
=\Big(\sum_{k\in F}|x_k|^{q}+\sum_{k\notin F}|x_k|^{q}\Big)^{\!1/q}
\le \Big(\sum_{k\in F}|x_k|^{q}\Big)^{\!1/q}+\Big(\sum_{k\notin F}|x_k|^{p}\Big)^{\!1/q}
\le \Big(\sum_{k\in F}|x_k|^{q}\Big)^{\!1/q}+\varepsilon.
\]
Da $F$ endlich ist, gilt
\[
\lim_{q\to\infty}\Big(\sum_{k\in F}|x_k|^{q}\Big)^{1/q}=\max_{k\in F}|x_k|\le M,
\]
also $\limsup_{q\to\infty}\|x\|_{q}\le M+\varepsilon$. Da $\varepsilon>0$ beliebig war,
$\limsup_{q\to\infty}\|x\|_{q}\le M$.
Andererseits ist stets $\|x\|_{q}\ge M$ (denn in der Summe steht mindestens ein Term
$\ge M^{q}$ bis auf Approximation der Supremumsstelle), also $\liminf_{q\to\infty}\|x\|_{q}\ge M$.
Zusammen folgt $\lim_{q\to\infty,\,q\ge p}\|x\|_{q}=M=\|x\|_{\infty}$.
\end{PROOF}




2.10 Definition. (a) Eine Abbildung $h: E \times E \rightarrow \mathbb{K}$ heißt hermitesche Form auf $E$, falls gilt:\\
(H1) $h(\cdot, y): E \rightarrow \mathbb{K}$ ist $\mathbb{K}$-linear für alle $y \in E$;\\
(H2) $h(x, y)=\overline{h(y, x)}$ für alle $x, y \in E$. Insbesondere folgt: $h(x, x) \in \mathbb{R}$ für alle $x \in E$.\\
(b) Eine hermitesche Form $h$ heißt positiv, falls $h(x, x) \geq 0$ für alle $x \in E$ gilt.\\
(c) Eine positive hermitesche Form heißt Skalarprodukt, wenn für alle $x \in E$ die Implikation $h(x, x)=0 \Longrightarrow x=0$ gilt. In diesem Fall schreiben wir oft $\langle x, y\rangle$ statt $h(x, y)$. Das Paar $(E,\langle\cdot, \cdot\rangle)$ heißt dann Prähilbertraum.



\begin{DEF}{FA-A25-01-16}{Hermitesche Form und Prähilbertraum}
Sei $E$ ein $\mathbb{K}$-Vektorraum mit $\mathbb{K}\in\{\mathbb{R},\mathbb{C}\}$.

\begin{enumerate}[label=(\alph*)]

\item Eine Abbildung
\[
h: E\times E\to\mathbb{K}
\]
heißt \textbf{hermitesche Form} auf $E$, falls gilt:
\begin{description}
\item[(H1)] Für jedes feste $y\in E$ ist $h(\cdot,y):E\to\mathbb{K}$ \textbf{$\mathbb{K}$-linear}.
\item[(H2)] Es gilt die \textbf{Hermitizität}:
\[
h(x,y)=\overline{h(y,x)}\qquad\forall\,x,y\in E.
\]
\end{description}
Insbesondere folgt aus (H2): $h(x,x)\in\mathbb{R}$ für alle $x\in E$.

\item Eine hermitesche Form $h$ heißt \textbf{positiv (semi-definit)}, falls
\[
h(x,x)\ge 0\qquad\forall\,x\in E.
\]

\item Eine positive hermitesche Form $h$ heißt \textbf{Skalarprodukt}, wenn zusätzlich gilt:
\[
h(x,x)=0\ \Longrightarrow\ x=0.
\]
In diesem Fall schreiben wir oft
\[
\langle x,y\rangle := h(x,y)
\]
und nennen das Paar $(E,\langle\cdot,\cdot\rangle)$ einen \textbf{Prähilbertraum}.
\end{enumerate}
\end{DEF}


2.11 Bemerkung. 

Ist $h$ eine hermitesche Form auf $E$, so gilt
$$
h\left(x, \alpha y_{1}+\beta y_{2}\right)=\bar{\alpha} h\left(x, y_{1}\right)+\bar{\beta} h\left(x, y_{2}\right) \quad \text { für alle } x, y_{1}, y_{2} \in E, \alpha, \beta \in \mathbb{K} \text {. }
$$
Ist $\mathbb{K}=\mathbb{R}$, so ist $h$ genau dann eine hermitesche Form, wenn es eine symmetrische Bilinearform ist.\\



\begin{PROP}{FA-A25-01-17}{Linearitäts- und Symmetrieeigenschaften hermitescher Formen}
Sei $h:E\times E\to\mathbb{K}$ eine hermitesche Form auf einem $\mathbb{K}$-Vektorraum $E$.

\begin{enumerate}[label=(\alph*)]
\item Dann gilt für alle $x,y_{1},y_{2}\in E$ und $\alpha,\beta\in\mathbb{K}$:
\[
h\big(x,\alpha y_{1}+\beta y_{2}\big)
=\overline{\alpha}\,h(x,y_{1})+\overline{\beta}\,h(x,y_{2}).
\]
Das heißt: eine hermitesche Form ist \textbf{linear im ersten Argument} und \textbf{konjugiert-linear im zweiten Argument}.

\item Im reellen Fall $\mathbb{K}=\mathbb{R}$ gilt:
\[
h(x,y)=h(y,x)\quad\forall\,x,y\in E,
\]
also ist $h$ genau dann eine hermitesche Form, wenn $h$ eine \textbf{symmetrische Bilinearform} ist.
\end{enumerate}
\end{PROP}


\begin{PROOF}{FA-A25-01-18}{P: Linearitäts- und Symmetrieeigenschaften hermitescher Formen}
(a) Für feste $x\in E$ gilt nach Definition (H1): $h(x,\cdot)$ ist \emph{linear}.  
Nach (H2) folgt für jedes $y\in E$:
\[
h(x,y)=\overline{h(y,x)}.
\]
Setzt man hier $y=\alpha y_{1}+\beta y_{2}$ ein und verwendet die Linearität von $h(\cdot,y)$ im ersten Argument, erhält man
\[
\begin{aligned}
h(x,\alpha y_{1}+\beta y_{2})
&=\overline{h(\alpha y_{1}+\beta y_{2},x)}
=\overline{\alpha h(y_{1},x)+\beta h(y_{2},x)}\\
&=\overline{\alpha}\,\overline{h(y_{1},x)}+\overline{\beta}\,\overline{h(y_{2},x)}
=\overline{\alpha}\,h(x,y_{1})+\overline{\beta}\,h(x,y_{2}).
\end{aligned}
\]

(b) Im Fall $\mathbb{K}=\mathbb{R}$ entfällt die Konjugation, und (H2) reduziert sich zu $h(x,y)=h(y,x)$, also Symmetrie. Damit ist jede hermitesche Form eine symmetrische Bilinearform und umgekehrt.
\end{PROOF}




2.12 Satz. Sei h eine positive hermitesche Form auf E. Dann gilt:\\
(a) $|h(x, y)|^{2} \leq h(x, x) h(y, y)$ für alle $x, y \in E$ (Cauchy-Schwarz-Ungleichung)\\
(b) Durch $\|x\|:=\sqrt{h(x, x)}$ ist eine Halbnorm auf $E$ definiert. Diese ist genau dann eine Norm, wenn h ein Skalarprodukt ist.

Beweis. Übung.



\begin{THEO}{FA-A25-01-19}{Cauchy--Schwarz Ungleichung (Positive Hermitesche Form)}
Sei $E$ ein $\mathbb{K}$-Vektorraum ($\mathbb{K}=\mathbb{R}$ oder $\mathbb{C}$) und $h:E\times E\to\mathbb{K}$ eine \emph{positive hermitesche Form}, d.\,h.\ hermitesch und $h(x,x)\ge 0$ für alle $x\in E$. Dann gilt für alle $x,y\in E$:
\[
|h(x,y)|^{2}\ \le\ h(x,x)\,h(y,y).
\]
\end{THEO}

\begin{PROOF}{FA-A25-01-20}{P: Cauchy--Schwarz Ungleichung (Positive Hermitesche Form)}
Ist $y=0$, so ist die Behauptung trivial. Sei also $y\neq 0$. Für $\lambda\in\mathbb{K}$ ist wegen der Positivität
\[
0\ \le\ h(x-\lambda y,\,x-\lambda y).
\]
Benutze Linearität in der ersten und konjugiert-lineare Abhängigkeit in der zweiten Variablen:
\[
\begin{aligned}
h(x-\lambda y,x-\lambda y)
&= h(x,x) - \overline{\lambda}\,h(x,y) - \lambda\,h(y,x) + |\lambda|^{2} h(y,y)\\
&= h(x,x) - 2\Re\big(\overline{\lambda}\,h(x,y)\big) + |\lambda|^{2}h(y,y).
\end{aligned}
\]
Wähle
\[
\lambda\ :=\ \frac{h(x,y)}{h(y,y)}\qquad(\text{wohldefiniert, da }h(y,y)>0 \text{ für } y\neq 0).
\]
Dann folgt
\[
0\ \le\ h(x,x)\ -\ \frac{|h(x,y)|^{2}}{h(y,y)}.
\]
Nach Umstellen ergibt sich $|h(x,y)|^{2}\le h(x,x)\,h(y,y)$.
\end{PROOF}


Halbnorm-Eigenschaft von positiven hermiteschen Formen

\begin{PROP}{FA-A25-01-21}{Halbnorm-Eigenschaft von positiven hermiteschen Formen}
Sei $E$ ein $\mathbb{K}$-Vektorraum ($\mathbb{K}=\mathbb{R}$ oder $\mathbb{C}$) und $h:E\times E\to\mathbb{K}$ eine \emph{positive hermitesche Form}, d.\,h.\ hermitesch und $h(x,x)\ge 0$ für alle $x\in E$. Dann definiert die Vorschrift
\[
\|x\|\ :=\ \sqrt{h(x,x)}\qquad(x\in E)
\]
eine \emph{Halbnorm} auf $E$. Sie ist genau dann eine \emph{Norm}, wenn $h$ ein Skalarprodukt ist (d.\,h.\ $h(x,x)=0\Rightarrow x=0$).
\end{PROP}

\begin{PROOF}{FA-A25-01-22}{P: Halbnorm-Eigenschaft von positiven hermiteschen Formen}
Für $\alpha\in\mathbb{K}$ gilt
\[
\|\alpha x\|^{2}=h(\alpha x,\alpha x)=\alpha\,\overline{\alpha}\,h(x,x)=|\alpha|^{2}\|x\|^{2}
\quad\Rightarrow\quad \|\alpha x\|=|\alpha|\|x\|.
\]
Weiter
\[
\|x+y\|^{2}=h(x+y,x+y)=\|x\|^{2}+\|y\|^{2}+2\Re\,h(x,y)
\le \|x\|^{2}+\|y\|^{2}+2|h(x,y)|
\]
und mit (a)
\[
\|x+y\|^{2}\ \le\ \|x\|^{2}+\|y\|^{2}+2\|x\|\|y\|=(\|x\|+\|y\|)^{2},
\]
also $\|x+y\|\le \|x\|+\|y\|$. Positivität $\|x\|\ge 0$ ist klar. Damit ist $\|\cdot\|$ eine Halbnorm.

Ist $h$ \emph{definit} (Skalarprodukt), so folgt aus $\|x\|=0$ die Implikation $x=0$; die Halbnorm ist dann eine Norm.
Umgekehrt: Ist $\|\cdot\|$ eine Norm, so impliziert $\|x\|=0$ stets $x=0$, also ist $h$ definit.
\end{PROOF}





2.13 Beispiel. Sei $\mathcal{I}$ eine höchstens abzählbare Indexmenge. Dann ist $\ell^{2}(\mathcal{I})$ ein Prähilbertraum mit Skalarprodukt $\langle\cdot, \cdot\rangle$ definiert durch

$$
\langle x, y\rangle:=\sum_{k \in \mathcal{I}} x_{k} \bar{y}_{k} \text { für } x=\left(x_{k}\right)_{k \in \mathcal{I}}, y=\left(y_{k}\right)_{k \in \mathcal{I}} \in \ell^{2}(\mathcal{I}) \text {. }
$$

Die Konvergenz der obigen Reihe folgt aus dem Majorantenkriterium, denn es gilt

$$
\left|x_{k} \bar{y}_{k}\right|=\left|x_{k}\right|\left|y_{k}\right| \leq \frac{1}{2}\left(\left|x_{k}\right|^{2}+\left|y_{k}\right|^{2}\right) \quad \text { für alle } k \in \mathbb{N},
$$

und die Reihen $\sum_{k}\left|x_{k}\right|^{2}$ und $\sum_{k}\left|y_{k}\right|^{2}$ konvergieren.


\begin{EXA}{FA-A25-01-23}{Standardskalarprodukt auf $\ell^{2}(\mathcal{I})$ (Prähilbertraum)}
Sei $\mathcal{I}$ eine höchstens abzählbare Indexmenge und $\mathbb{K}\in\{\mathbb{R},\mathbb{C}\}$.
Dann ist
\[
\ell^{2}(\mathcal{I})
=\Bigl\{x=(x_{k})_{k\in\mathcal{I}}:\ \sum_{k\in\mathcal{I}}|x_{k}|^{2}<\infty\Bigr\}
\]
ein Prähilbertraum mit Skalarprodukt
\[
\langle x,y\rangle
:=\sum_{k\in\mathcal{I}} x_{k}\,\overline{y_{k}},
\qquad x=(x_{k}),\ y=(y_{k})\in \ell^{2}(\mathcal{I}).
\]

\textbf{Konvergenz der Reihe.}
Für alle $k$ gilt die Abschätzung
\[
|x_{k}\,\overline{y_{k}}|=|x_{k}|\,|y_{k}|
\le \tfrac{1}{2}\big(|x_{k}|^{2}+|y_{k}|^{2}\big),
\]
so dass die Reihe $\sum_{k}|x_{k}\overline{y_{k}}|$ vom konvergenten Majoranten
$\tfrac{1}{2}\sum_{k}\big(|x_{k}|^{2}+|y_{k}|^{2}\big)$ beherrscht wird; damit konvergiert
$\sum_{k}x_{k}\overline{y_{k}}$ absolut.

\textbf{Eigenschaften.}
Für $x,y,z\in \ell^{2}(\mathcal{I})$ und $\alpha\in\mathbb{K}$:
\[
\langle \alpha x+y, z\rangle=\alpha \langle x,z\rangle+\langle y,z\rangle,\qquad
\langle x,y\rangle=\overline{\langle y,x\rangle},\qquad
\langle x,x\rangle=\sum_{k}|x_{k}|^{2}\ge 0.
\]
Insbesondere ist $h(x,y):=\langle x,y\rangle$ eine positive hermitesche Form, und
\[
\|x\|:=\sqrt{\langle x,x\rangle}
=\Big(\sum_{k\in\mathcal{I}}|x_{k}|^{2}\Big)^{1/2}
\]
ist die zugehörige Norm (Cauchy--Schwarz $\Rightarrow$ Dreiecksungleichung),
womit $(\ell^{2}(\mathcal{I}),\langle\cdot,\cdot\rangle)$ ein Prähilbertraum ist.
\end{EXA}


\pagebreak




\subsection*{Topologische Eigenschaften normierter Räume}


Wir wiederholen zunächst die grundlegenden topologischen Begriffe in metrischen Räumen.


2.14 Definition. 


Sei $X$ eine Menge. Eine Metrik auf $X$ ist eine Abbildung $d: X \times X \rightarrow$ $[0, \infty)$, so dass für $x, y, z \in X$ gilt:\\
(M1) $d(x, y)=0 \Longleftrightarrow x=y$,\\
$(\mathrm{M} 2) d(x, y)=d(y, x)$,\\
(M3) $d(x, z) \leq d(x, y)+d(y, z)$.\\
Das Paar ( $X, d$ ) heißt metrischer Raum.\\


\begin{DEF}{FA-A25-01-24}{Metrischer Raum}
Sei $X$ eine Menge.  
Eine Abbildung
\[
d: X \times X \longrightarrow [0,\infty)
\]
heißt \textbf{Metrik} auf $X$, falls für alle $x,y,z\in X$ gilt:

\begin{enumerate}[label=(M\arabic*)]
\item \textbf{(Trennbarkeit)}  
      $d(x,y)=0 \ \Longleftrightarrow\ x=y$;
\item \textbf{(Symmetrie)}  
      $d(x,y)=d(y,x)$;
\item \textbf{(Dreiecksungleichung)}  
      $d(x,z)\le d(x,y)+d(y,z)$.
\end{enumerate}

Das Paar $(X,d)$ heißt dann ein \textbf{metrischer Raum}.  
Für $x\in X$ und $r>0$ bezeichnet man die Menge
\[
B_{r}(x):=\{\,y\in X:\ d(x,y)<r\,\}
\]
als die \textbf{offene Kugel} mit Mittelpunkt $x$ und Radius $r$.
\end{DEF}


2.15 Definition und Bemerkung. (a) Für $A \subseteq X, x^{*} \in X$ und $\varepsilon>0$ setzen wir:

$$
\begin{aligned}
\operatorname{diam} A & :=\sup \{d(x, y) \mid x, y \in A\} & & (\text { Durchmesser von A), } \\
\operatorname{dist}\left(x^{*}, A\right) & :=\inf \left\{d\left(x^{*}, y\right) \mid y \in A\right\} & & \left(\text { Abstand von } x^{*} z u A\right), \\
U_{\varepsilon}\left(x^{*}\right) & :=\left\{y \in X \mid d\left(y, x^{*}\right)<\varepsilon\right\} & & \left(\text { offene } \varepsilon \text {-Kugel um } x^{*}\right), \\
B_{\varepsilon}\left(x^{*}\right) & :=\left\{y \in X \mid d\left(y, x^{*}\right) \leq \varepsilon\right\} & & \left(\text { abgeschlossene } \varepsilon \text {-Kugel um } x^{*}\right), \\
U_{\varepsilon}(A) & :=\{y \in X \mid \operatorname{dist}(y, A)<\varepsilon\} & & (\text { offene } \varepsilon \text {-Umgebung von A) }, \\
B_{\varepsilon}(A) & :=\{y \in X \mid \operatorname{dist}(y, A) \leq \varepsilon\} & & (\text { abgeschl. } \varepsilon \text {-Umgebung von A), } \\
\bar{A} & :=\{x \in X \mid \operatorname{dist}(x, A)=0\} & & (\text { Abschluss von A), } \\
\operatorname{int}(A):=\AA & :=\left\{x \in A \mid \exists \varepsilon>0: U_{\varepsilon}(x) \subseteq A\right\} & & (\text { offener Kern von A), } \\
\partial A & :=\bar{A} \backslash \AA & & (\text { Rand von A). }
\end{aligned}
$$

(b) $A$ heißt

\begin{itemize}
  \item abgeschlossen, falls $A=\bar{A}$;
  \item offen, falls $A=\AA$;
  \item Umgebung von $x$, falls $x \in \AA$;
  \item dicht in $X$, falls $\bar{A}=X$;
  \item beschränkt, falls $\operatorname{diam} A<\infty$.
\end{itemize}



\begin{DEF}{FA-A25-01-25}{Grundbegriffe der metrischen Topologie}
Sei $(X,d)$ ein metrischer Raum, $A\subseteq X$, $x^{*}\in X$ und $\varepsilon>0$.  
Dann definieren wir:

\begin{align*}
\operatorname{diam}A
&:=\sup\{\,d(x,y)\mid x,y\in A\,\}
&&\text{(Durchmesser von $A$)},\\[0.5em]
\operatorname{dist}(x^{*},A)
&:=\inf\{\,d(x^{*},y)\mid y\in A\,\}
&&\text{(Abstand des Punktes $x^{*}$ von $A$)},\\[0.5em]
U_{\varepsilon}(x^{*})
&:=\{\,y\in X\mid d(y,x^{*})<\varepsilon\,\}
&&\text{(offene $\varepsilon$-Kugel um $x^{*}$)},\\[0.5em]
B_{\varepsilon}(x^{*})
&:=\{\,y\in X\mid d(y,x^{*})\le\varepsilon\,\}
&&\text{(abgeschlossene $\varepsilon$-Kugel um $x^{*}$)},\\[0.5em]
U_{\varepsilon}(A)
&:=\{\,y\in X\mid \operatorname{dist}(y,A)<\varepsilon\,\}
&&\text{(offene $\varepsilon$-Umgebung von $A$)},\\[0.5em]
B_{\varepsilon}(A)
&:=\{\,y\in X\mid \operatorname{dist}(y,A)\le\varepsilon\,\}
&&\text{(abgeschlossene $\varepsilon$-Umgebung von $A$)},\\[0.5em]
\overline{A}
&:=\{\,x\in X\mid \operatorname{dist}(x,A)=0\,\}
&&\text{(Abschluss von $A$)},\\[0.5em]
\operatorname{int}(A)=\mathring{A}
&:=\{\,x\in A\mid \exists\,\varepsilon>0:\ U_{\varepsilon}(x)\subseteq A\,\}
&&\text{(offener Kern von $A$)},\\[0.5em]
\partial A
&:=\overline{A}\setminus \mathring{A}
&&\text{(Rand von $A$)}.
\end{align*}

\textbf{(b) Terminologie.}  
Eine Teilmenge $A\subseteq X$ heißt
\begin{itemize}
  \item \textbf{abgeschlossen}, falls $A=\overline{A}$;
  \item \textbf{offen}, falls $A=\mathring{A}$;
  \item \textbf{Umgebung} eines Punktes $x\in X$, falls $x\in\mathring{A}$;
  \item \textbf{dicht in $X$}, falls $\overline{A}=X$;
  \item \textbf{beschränkt}, falls $\operatorname{diam}A<\infty$.
\end{itemize}
\end{DEF}




Es gilt:


$U_{\varepsilon}\left(x^{*}\right), U_{\varepsilon}(A)$ sind offen, $B_{\varepsilon}\left(x^{*}\right), B_{\varepsilon}(A)$ sind abgeschlossen;

$A$ offen $\Longleftrightarrow X \backslash A$ abgeschlossen;

$X, \varnothing$ sind offen und abgeschlossen;

$\AA$ ist offen, $\bar{A}$ ist abgeschlossen;

$\partial A$ ist abgeschlossen.\\

(c) Seien $A_{i} \subseteq X, i \in \mathcal{I}$ ( $\mathcal{I}$ beliebige Indexmenge). Dann gilt:
  
  Sind alle $A_{i}, i \in \mathcal{I}$, offen, so auch $\bigcup_{i \in \mathcal{I}} A_{i}$ und, falls $\mathcal{I}$ endlich ist, auch $\bigcap_{i \in \mathcal{I}} A_{i}$.
  
  Sind alle $A_{i}, i \in \mathcal{I}$, abgeschlossen, so auch $\bigcap_{i \in \mathcal{I}} A_{i}$ und, falls $\mathcal{I}$ endlich ist, auch $\bigcup_{i \in \mathcal{I}} A_{i}$.\\
(d) Sei $a \in X$ und $\left(x_{k}\right)_{k} \subseteq X$ eine Folge. Falls $d\left(x_{k}, a\right) \rightarrow 0$ für $k \rightarrow \infty$, so nennen wir $a$ Grenzwert der Folge $\left(x_{k}\right)_{k}$ und schreiben $a=\lim _{k \rightarrow \infty} x_{k}$ bzw. $x_{k} \rightarrow a$ für $k \rightarrow \infty$. Leicht zu sehen: Durch diese Eigenschaft ist $a$ eindeutig bestimmt.\\
(e) Ist $E$ ein $\mathbb{K}$-Vektorraum, so induziert jede Norm $\|\cdot\|$ auf $E$ eine Metrik $d$, gegeben durch $d(x, y):=\|x-y\|$ für alle $x, y \in E$. Ist ( $X, d$ ) ein metrischer Raum und $A \subseteq X$, so ist auch ( $A, d$ ) ein metrischer Raum. Insbesondere ist jede Teilmenge eines normierten Raumes ein metrischer Raum mit der durch die Norm induzierten Metrik.\\

(f) Die Eigenschaften „offen" und „abgeschlossen" hängen vom umgebenden metrischen Raum ab. Ist ( $X, d$ ) ein metrischer Raum und $M \subseteq X$, so gilt für $A \subseteq M$ :
  
   $A$ offen in $(M, d) \Longleftrightarrow$ es existiert $A^{\prime} \subseteq X, A^{\prime}$ offen in $X$, mit $A=A^{\prime} \cap M$;
  
   $A$ abgeschlossen in $(M, d) \Longleftrightarrow$ es existiert $A^{\prime} \subseteq X, A^{\prime}$ abgeschlossen in $X$, mit $A=A^{\prime} \cap M$.\\

(g) Sei $E$ ein $\mathbb{K}$-Vektorraum. Zwei äquivalente Normen $\|\cdot\|_{1}$ und $\|\cdot\|_{2}$ auf $E$ erzeugen die gleiche Topologie, d.h. das gleiche System offener Teilmengen. Für $A \subseteq E$ gilt also, dass $A$ genau dann offen bzgl. $\|\cdot\|_{1}$ ist, wenn $A$ offen bzgl. $\|\cdot\|_{2}$ ist. Gleichermaßen invariant bleiben die Eigenschaften „abgeschlossen", „beschränkt" und „dicht", die Definitionen $\bar{A}$ und $\AA$ und die Konvergenz von Folgen. Insbesondere sind diese Begriffe im Vektorraum $\mathbb{K}^{N}$, unabhängig von der Wahl einer Norm, wohldefiniert.\\



\begin{PROP}{FA-A25-01-26}{Grundlegende topologische Eigenschaften metrischer Räume}
Sei $(X,d)$ ein metrischer Raum.

\textbf{(a) Eigenschaften von Kugeln, Offenheit und Abgeschlossenheit.}
\begin{itemize}
  \item Die Mengen $U_{\varepsilon}(x^{*})$ und $U_{\varepsilon}(A)$ sind offen, während $B_{\varepsilon}(x^{*})$ und $B_{\varepsilon}(A)$ abgeschlossen sind.
  \item $A$ ist offen $\Longleftrightarrow$ $X\setminus A$ ist abgeschlossen.
  \item $X$ und $\varnothing$ sind sowohl offen als auch abgeschlossen.
  \item Der offene Kern $\mathring{A}$ ist offen, der Abschluss $\overline{A}$ ist abgeschlossen.
  \item Der Rand $\partial A=\overline{A}\setminus \mathring{A}$ ist abgeschlossen.
\end{itemize}

\textbf{(b) Stabilität von Offenheit und Abgeschlossenheit unter Mengenoperationen.}\\
Seien $(A_i)_{i\in\mathcal{I}}$ Teilmengen von $X$ für eine beliebige Indexmenge $\mathcal{I}$.
\begin{itemize}
  \item Sind alle $A_i$ offen, so ist auch $\bigcup_{i\in\mathcal{I}}A_i$ offen, und falls $\mathcal{I}$ endlich ist, ist auch $\bigcap_{i\in\mathcal{I}}A_i$ offen.
  \item Sind alle $A_i$ abgeschlossen, so ist auch $\bigcap_{i\in\mathcal{I}}A_i$ abgeschlossen, und falls $\mathcal{I}$ endlich ist, ist auch $\bigcup_{i\in\mathcal{I}}A_i$ abgeschlossen.
\end{itemize}

\textbf{(c) Grenzwerte von Folgen.}\\
Für $a\in X$ und eine Folge $(x_k)_{k\in\mathbb{N}}\subseteq X$ heißt
\[
d(x_k,a)\to 0\quad\text{für }k\to\infty
\]
ein \textbf{Grenzwertverhalten}, und man schreibt $x_k\to a$ bzw. $a=\lim_{k\to\infty}x_k$.
Der Grenzwert $a$ ist eindeutig bestimmt.

\textbf{(d) Von Normen induzierte Metriken.}
\begin{itemize}
  \item Ist $E$ ein $\mathbb{K}$-Vektorraum und $\|\cdot\|$ eine Norm auf $E$, so definiert
        \[
        d(x,y):=\|x-y\|
        \]
        eine Metrik auf $E$.
  \item Ist $(X,d)$ ein metrischer Raum und $A\subseteq X$, so ist auch $(A,d)$ ein metrischer Raum (mit der Einschränkung der Metrik).
  \item Insbesondere ist jede Teilmenge eines normierten Raumes mit der induzierten Metrik ein metrischer Raum.
\end{itemize}

\textbf{(e) Relative Offenheit und Abgeschlossenheit.}\\
Sei $(X,d)$ ein metrischer Raum und $M\subseteq X$. Dann gilt für $A\subseteq M$:
\begin{itemize}
  \item $A$ ist offen in $(M,d)$ $\Longleftrightarrow$ es existiert eine offene Menge $A'\subseteq X$ mit $A=A'\cap M$;
  \item $A$ ist abgeschlossen in $(M,d)$ $\Longleftrightarrow$ es existiert eine abgeschlossene Menge $A'\subseteq X$ mit $A=A'\cap M$.
\end{itemize}

\textbf{(f) Äquivalente Normen.}\\
Sei $E$ ein $\mathbb{K}$-Vektorraum.  
Sind $\|\cdot\|_1$ und $\|\cdot\|_2$ äquivalente Normen auf $E$, so induzieren sie dieselbe Topologie:
\begin{itemize}
  \item $A\subseteq E$ ist offen bzgl.\ $\|\cdot\|_1$ $\Longleftrightarrow$ $A$ ist offen bzgl.\ $\|\cdot\|_2$;
  \item Ebenso bleiben die Eigenschaften „abgeschlossen“, „beschränkt“ und „dicht“ invariant, ebenso die Definitionen von $\overline{A}$, $\mathring{A}$ und die Konvergenz von Folgen.
\end{itemize}
Insbesondere sind diese Begriffe in $\mathbb{K}^N$ unabhängig von der Wahl der Norm wohldefiniert.
\end{PROP}



2.16 Definition. 

Der metrische Raum ($X,d$) heißt separabel, falls $X$ eine abzählbare dichte Teilmenge besitzt. Hier und im Folgenden steht „abzählbar" für „endlich oder abzählbar unendlich".


\begin{DEF}{FA-A25-01-27}{Separable metrische Räume}
Ein metrischer Raum $(X,d)$ heißt \textbf{separabel},  
wenn es eine \textbf{abzählbare dichte Teilmenge} $D\subseteq X$ gibt, d.\,h.
\[
\overline{D}=X.
\]
Hierbei bedeutet „abzählbar“ stets „endlich oder abzählbar unendlich“.  

Äquivalent gilt: $(X,d)$ ist separabel $\Longleftrightarrow$  
es existiert eine Folge $(x_n)_{n\in\mathbb{N}}\subseteq X$, deren Abschluss ganz $X$ ist, also  
\[
\forall\,x\in X,\ \forall\,\varepsilon>0:\ \exists\,n\in\mathbb{N}\text{ mit }d(x,x_n)<\varepsilon.
\]
\end{DEF}



2.17 Beispiel. (a) $\mathbb{K}^{N}$ ist separabel, denn
  
  $\mathbb{Q}^{N}$ ist abzählbar und dicht in $\mathbb{R}^{N}$
  
  $(\mathbb{Q}+\mathrm{i} \mathbb{Q})^{N}$ ist abzählbar und dicht in $\mathbb{C}^{N}$.\\

(b) $\mathbb{R} \backslash \mathbb{Q}$ ist separabel, denn $\sqrt{2}+\mathbb{Q}$ ist eine abzählbare und dichte Teilmenge von $\mathbb{R} \backslash \mathbb{Q}$.




\begin{EXA}{FA-A25-01-28}{Beispiele separabler metrischer Räume}
\begin{enumerate}
\item \textbf{Der Raum $\mathbb{K}^{N}$ ist separabel.}\\[0.3em]
Für $\mathbb{K}=\mathbb{R}$ gilt:  
$\mathbb{Q}^{N}$ ist abzählbar und dicht in $\mathbb{R}^{N}$, da für jedes $x\in\mathbb{R}^{N}$ und jedes $\varepsilon>0$ ein $q\in\mathbb{Q}^{N}$ mit $\|x-q\|<\varepsilon$ existiert.

Für $\mathbb{K}=\mathbb{C}$ gilt analog:  
$(\mathbb{Q}+\mathrm{i}\mathbb{Q})^{N}$ ist abzählbar und dicht in $\mathbb{C}^{N}$, da jede komplexe Zahl $z=a+\mathrm{i}b$ durch rationale Approximationen $a_n,b_n\in\mathbb{Q}$ beliebig gut angenähert werden kann.

\item \textbf{Der irrationale Zahlenraum $\mathbb{R}\setminus\mathbb{Q}$ ist separabel.}\\[0.3em]
Die Menge
\[
\sqrt{2}+\mathbb{Q}=\{\sqrt{2}+q\mid q\in\mathbb{Q}\}
\]
ist abzählbar, da sie das Bild einer abzählbaren Menge unter einer bijektiven Abbildung ist,  
und sie ist dicht in $\mathbb{R}\setminus\mathbb{Q}$, da für jedes irrationale $x$ und jedes $\varepsilon>0$ ein $q\in\mathbb{Q}$ existiert mit $|x-(\sqrt{2}+q)|<\varepsilon$.
\end{enumerate}
\end{EXA}



2.18 Satz. 

($X,d$) separabel, $A \subseteq X \Longrightarrow(A, d)$ separabel.



\begin{PROP}{FA-A25-01-29}{Vererbung der Separabilität auf Teilräume}
Sei $(X,d)$ ein separabler metrischer Raum und $A\subseteq X$ eine Teilmenge.  
Dann ist auch der Teilraum $(A,d)$ separabel.
\end{PROP}

Beweis. 

\begin{PROOF}{FA-A25-01-30}{P: Vererbung der Separabilität auf Teilräume}
Nach Voraussetzung existiert eine abzählbare Menge $Y=\left\{y_{n} \mid n \in \mathbb{N}\right\} \subseteq X$, welche in $X$ dicht liegt. Wir setzen
$$
H:=\left\{(m, n) \in \mathbb{N} \times \mathbb{N} \left\lvert\, B_{\frac{1}{m}}\left(y_{n}\right) \cap A \neq \varnothing\right.\right\}
$$
und wählen $b_{m, n} \in B_{\frac{1}{m}}\left(y_{n}\right) \cap A$ für jedes $(m, n) \in H$. Da $H$ abzählbar ist, ist auch die Menge
$$
B:=\left\{b_{m, n} \mid(m, n) \in H\right\} \subseteq A
$$
abzählbar. Ferner ist $B$ dicht in $A$ (Übung!). Es folgt, dass $A$ separabel ist.
\end{PROOF}


2.19 Definition. 

Sei ($E,\|\cdot\|$) ein normierter Raum. Eine Teilmenge $M \subseteq E$ heißt total, falls span $M=E$ gilt. Hier und im Folgenden bezeichne "span $M^{\prime \prime}$ das Vektorraumerzeugnis von $M$ in $E$, welches als der Unterraum aller Linearkombinationen von Elementen aus $M$ bzw. äquivalent als der Schnitt aller Unterräume von $E$, welche $M$ enthalten, gegeben ist.



\begin{DEF}{FA-A25-01-31}{Totale Menge und lineare Hülle}
Sei $(E,\|\cdot\|)$ ein normierter (bzw.\ linearer) Raum und $M\subseteq E$.

\begin{enumerate}[label=(\alph*)]
\item Die \textbf{lineare Hülle} (oder \textbf{span}) von $M$ ist
\[
\operatorname{span} M
:= \bigl\{\textstyle\sum_{j=1}^{n}\alpha_j x_j \;\big|\; n\in\mathbb{N},\ \alpha_j\in\mathbb{K},\ x_j\in M\bigr\}.
\]
Äquivalent ist
\[
\operatorname{span} M=\bigcap\{\,U\subseteq E \mid U\text{ Unterraum von }E,\ M\subseteq U\,\}.
\]

\item Eine Teilmenge $M$ heißt \textbf{total} in $E$, falls
\[
\operatorname{span} M = E.
\]
\end{enumerate}

\textbf{Bemerkung.}
Man unterscheidet gelegentlich von der algebraischen Totalität auch die \emph{topologische} Variante:
$M$ heißt \emph{fundamental} (oder \emph{topologisch total}), wenn $\overline{\operatorname{span} M}=E$.
Jede Hamel-Basis ist total; eine Orthonormalbasis in einem Hilbertraum ist fundamental (i.\,A.\ nicht total im algebraischen Sinne).
\end{DEF}


2.20 Satz. 

\begin{PROP}{FA-A25-01-32}{Normierter Raum ist separabel $\Leftrightarrow$ $\exists$ Abzählbare totale Teilmenge}
Sei $(E,\|\cdot\|)$ ein normierter $\mathbb{K}$-Vektorraum. Dann ist $E$ genau dann separabel, wenn $E$ eine abzählbare totale Teilmenge enthält.
\end{PROP}

Beweis. 

\begin{PROOF}{FA-A25-01-33}{P: Normierter Raum ist separabel $\Leftrightarrow$ $\exists$ Abzählbare totale Teilmenge}
„$\Longrightarrow$ ": Sei $E$ separabel. Dann existiert eine abzählbare Teilmenge $M \subseteq E$ mit $\bar{M}=E$. Es folgt $E=\bar{M} \subseteq \overline{\operatorname{span} M} \subseteq E$, also $\overline{\operatorname{span} M}=E$. Demnach ist $M$ total.

„$\Longleftarrow$": Sei $M \subseteq E$ abzählbar und total. Wir setzen
$$
D:=\left\{\sum_{k=1}^{n} a_{k} v_{k} \mid n \in \mathbb{N}, a_{1}, \ldots, a_{n} \in R, v_{1}, \ldots, v_{n} \in M\right\},
$$
wobei
$$
R:=\left\{\begin{array}{lll}
\mathbb{Q}, & \text { falls } & \mathbb{K}=\mathbb{R} \\
\mathbb{Q}+i \mathbb{Q}, & \text { falls } & \mathbb{K}=\mathbb{C}
\end{array}\right.
$$
sei. $D$ ist abzählbar, da $M$ und $R$ abzählbar sind.

Seien nun $x \in E$ und $\varepsilon>0$ gegeben. Da $M$ total ist, existieren $v_{1}, \ldots, v_{n} \in M$ und $\lambda_{1}, \ldots, \lambda_{n} \in \mathbb{K}$ mit
$$
\left\|x-\sum_{k=1}^{n} \lambda_{k} v_{k}\right\|<\frac{\varepsilon}{2}
$$
Sei $C:=\sum_{k=1}^{n}\left\|v_{k}\right\|$. Da $R$ dicht in $\mathbb{K}$ ist, existieren $a_{1}, \ldots, a_{n} \in R$ mit
$$
\left|\lambda_{k}-a_{k}\right|<\frac{\varepsilon}{2 C} \quad \text { für } k=1, \ldots, n \text {. }
$$
Es folgt
$$
\left\|x-\sum_{k=1}^{n} a_{k} v_{k}\right\| \stackrel{(2.2)}{\leq} \frac{\varepsilon}{2}+\left\|\sum_{k=1}^{n}\left(\lambda_{k}-a_{k}\right) v_{k}\right\| \leq \frac{\varepsilon}{2}+\sum_{k=1}^{n}\left|\lambda_{k}-a_{k}\right|\left\|v_{k}\right\|<\frac{\varepsilon}{2}+\frac{\varepsilon}{2}=\varepsilon
$$
Da $\varepsilon>0$ beliebig gewählt war, folgt $x \in \bar{D}$. Somit ist $\bar{D}=E$, und $E$ ist separabel.
\end{PROOF}


\begin{DEF}{FA-A25-01-34}{Finite Folgen in Körpern}
Eine Folge $\left(x_{k}\right)_{k} \subseteq \mathbb{K}$ heißt finit, falls $x_{k} \neq 0$ für höchstens endlich viele $k \in \mathbb{N}$ gilt.
\end{DEF}


Sei $\mathcal{F}$ der $\mathbb{K}$-Vektorraum der finiten Folgen. Es gilt dann $\mathcal{F} \subseteq \ell^{p}$ für alle $p \in[1, \infty]$. Eine Basis $B$ von $\mathcal{F}$ ist gegeben durch die Einheitsvektoren $e_{k}:=\left(\delta_{k j}\right)_{j}=(0, \ldots, 0,1,0, \ldots), k \in \mathbb{N}$. Für $p \in[1, \infty)$ ist $\mathcal{F}$ dicht in $\ell^{p}$ (Übung), also $B$ total in $\ell^{p}$. Es folgt mit Satz 2.20: $\ell^{p}$ ist für $p \in[1, \infty)$ separabel.\\
(b) $\ell^{\infty}$ ist nicht separabel (Übung).





\begin{EXA}{FA-A25-01-35}{Separabilität der $\ell^{p}$-Räume und Nicht-Separabilität von $\ell^{\infty}$}
Sei $\mathcal{F}\subset \mathbb{K}^{\mathbb{N}}$ der $\mathbb{K}$-Vektorraum der \emph{finiten Folgen}, d.\,h.
\[
\mathcal{F}:=\{x=(x_k)_{k\in\mathbb{N}}:\ \exists N\ \forall k>N:\ x_k=0\}.
\]
Für $k\in\mathbb{N}$ bezeichne $e_k:=(\delta_{kj})_{j\in\mathbb{N}}$ den $k$-ten Einheitsvektor.
Dann gilt:

\begin{enumerate}[label=(\alph*)]
\item Für jedes $p\in[1,\infty]$ ist $\mathcal{F}\subset \ell^{p}$ und $\{e_k:k\in\mathbb{N}\}$ ist eine (algebraische) Basis von $\mathcal{F}$.

\item Für $p\in[1,\infty)$ ist $\mathcal{F}$ dicht in $\ell^{p}$. Insbesondere ist $\ell^{p}$ separabel.

\item Der Raum $\ell^{\infty}$ ist \emph{nicht} separabel.
\end{enumerate}

(a) Offensichtlich ist jede finite Folge in $\ell^{p}$ (die Norm wird durch eine endliche Summe/ ein endliches Supremum bestimmt).  
Die Einheitsvektoren erzeugen jede finite Folge eindeutig, also bilden $(e_k)$ eine Hamel-Basis von $\mathcal{F}$.

\smallskip
(b) \emph{Dichtheit von $\mathcal{F}$ in $\ell^{p}$ für $1\le p<\infty$.}
Sei $x=(x_k)\in \ell^{p}$ und definiere die \emph{Trunkate}
\[
T_N x := \sum_{k=1}^{N} x_k e_k\in \mathcal{F}.
\]
Dann
\[
\|x-T_Nx\|_{p}^{p}=\sum_{k>N}|x_k|^{p}\xrightarrow[N\to\infty]{}0,
\]
weil $\sum_{k=1}^{\infty}|x_k|^{p}<\infty$. Also $T_Nx\to x$ in $\ell^{p}$, d.\,h.\ $\overline{\mathcal{F}}=\ell^{p}$.  
Da $\mathcal{F}$ einen abzählbaren totalen Erzeuger besitzt (z.\,B.\ die rationalen Linearkombinationen der $e_k$), ist $\ell^{p}$ separabel.

\smallskip
(c) \emph{$\ell^{\infty}$ ist nicht separabel.}
Betrachte die (überabzählbare) Menge
\[
\mathcal{A}:=\{\,\chi_{S}=(\mathbf{1}_{\{k\in S\}})_{k\in\mathbb{N}} \in \{0,1\}^{\mathbb{N}} \subset \ell^{\infty}\mid S\subseteq \mathbb{N}\,\}.
\]
Für $S\neq T$ gilt $\|\chi_{S}-\chi_{T}\|_{\infty}=1$. Somit sind die offenen Kugeln
\[
B_{1/2}(\chi_{S})=\{y\in\ell^{\infty}:\ \|y-\chi_{S}\|_{\infty}<1/2\},\qquad S\subseteq\mathbb{N},
\]
paarweise \emph{disjunkt}. Wäre $D\subset\ell^{\infty}$ abzählbar und dicht, so müsste jede dieser Kugeln ein Element aus $D$ enthalten; damit wäre $D$ unzählbar—ein Widerspruch. Folglich besitzt $\ell^{\infty}$ keine abzählbare dichte Teilmenge und ist nicht separabel.
\end{EXA}





\pagebreak


\subsection*{Stetige Abbildungen}


Seien im Folgenden ( $X, d_{X}$ ) und ( $Y, d_{Y}$ ) metrische Räume.\\
2.22 Definition und Bemerkung. Seien $f: X \rightarrow Y$ eine Abbildung und $a \in X$.\\
(a) $f$ heißt stetig in $a$, wenn $\lim _{n \rightarrow \infty} f\left(x_{n}\right)=f(a)$ für jede Folge $\left(x_{n}\right)_{n}$ in $X$ mit $\lim _{n \rightarrow \infty} x_{n}=a$ gilt. Dies ist bekanntermaßen äquivalent zu jeder der folgenden Eigenschaften:

\begin{itemize}
  \item Für alle $\varepsilon>0$ existiert $\delta>0$ mit $f\left(U_{\delta}(a)\right) \subseteq U_{\varepsilon}(f(a))$.
  \item Für jede Umgebung $V$ von $f(a)$ in $Y$ ist $f^{-1}(V)$ eine Umgebung von $a$ in $X$.\\
(b) $f$ heißt stetig, wenn $f$ in jedem Punkt aus $X$ stetig ist. Dies ist äquivalent zu jeder der folgenden Eigenschaften:
  \item Für jede in $Y$ offene Teilmenge $V$ von $Y$ ist $f^{-1}(V)$ offen in $X$.
  \item Für jede in $Y$ abgeschlossene Teilmenge $V$ von $Y$ ist $f^{-1}(V)$ abgeschlossen in $X$.
  \item Für alle $A \subseteq X$ gilt $f(\bar{A}) \subseteq \overline{f(A)}$.\\
(c) Wir setzen
\end{itemize}

$$
C(X, Y):=\{f: X \rightarrow Y \mid f \text { ist stetig }\}
$$

und

$$
C_{\mathrm{b}}(X, Y):=\{f \in C(X, Y) \mid f(X) \text { ist beschränkt }\} .
$$

Ist speziell $Y=\mathbb{K}$ (mit $\mathbb{K}=\mathbb{R}$ oder $\mathbb{K}=\mathbb{C}$ ), so schreiben wir kurz $C(X)$ anstelle von $C(X, \mathbb{K})$ und $C_{\mathrm{b}}(X)$ anstelle von $C_{\mathrm{b}}(X, \mathbb{K})$.



\begin{DEF}{FA-A25-01-36}{Stetigkeit in metrischen Räumen; Funktionsräume}
Seien $(X,d_X)$ und $(Y,d_Y)$ metrische Räume, $f:X\to Y$ und $a\in X$.

\textbf{(a) Punktweise Stetigkeit.}
Die Abbildung $f$ heißt \textbf{stetig in $a$}, wenn eine (äquivalente) der folgenden Bedingungen erfüllt ist:
\begin{enumerate}
\item \textbf{Folgenkriterium:} Für jede Folge $(x_n)_n\subset X$ mit $x_n\to a$ gilt $f(x_n)\to f(a)$.
\item \textbf{$\varepsilon$-$\delta$-Kriterium:} Für alle $\varepsilon>0$ existiert $\delta>0$ mit
      \[
      f\big(U_\delta(a)\big)\subset U_\varepsilon\big(f(a)\big).
      \]
\item \textbf{Umgebungskriterium:} Für jede Umgebung $V$ von $f(a)$ in $Y$ ist $f^{-1}(V)$ eine Umgebung von $a$ in $X$.
\end{enumerate}

\textbf{(b) Globale Stetigkeit.}
$f$ heißt \textbf{stetig}, wenn $f$ in jedem Punkt $a\in X$ stetig ist. Dies ist äquivalent zu jeder der folgenden Aussagen:
\begin{enumerate}
\item Für jede offene Menge $V\subset Y$ ist $f^{-1}(V)$ offen in $X$.
\item Für jede abgeschlossene Menge $V\subset Y$ ist $f^{-1}(V)$ abgeschlossen in $X$.
\item Für alle $A\subset X$ gilt die Abschlussabbildungseigenschaft
      \[
      f(\overline{A})\subset \overline{f(A)}.
      \]
\end{enumerate}

\textbf{(c) Funktionsräume.}
\[
C(X,Y):=\{\,f:X\to Y\mid f\ \text{stetig}\,\},\qquad
C_b(X,Y):=\{\,f\in C(X,Y)\mid f(X)\ \text{beschränkt}\,\}.
\]
Für $Y=\mathbb{K}$ (mit $\mathbb{K}\in\{\mathbb{R},\mathbb{C}\}$) schreiben wir kurz
\[
C(X):=C(X,\mathbb{K}),\qquad C_b(X):=C_b(X,\mathbb{K}).
\]
\end{DEF}



2.23 Bemerkung. (a) Die Komposition stetiger Abbildungen ist stetig\\
(b) Sind $f_{n}: X \rightarrow Y$ stetig für $n \in \mathbb{N}$ und konvergiert die Funktionenfolge $\left(f_{n}\right)_{n}$ gleichmäßig gegen $f: X \rightarrow Y$ (d.h. $\sup _{x \in X} d_{Y}\left(f_{n}(x), f(x)\right) \rightarrow 0$ für $n \rightarrow \infty$ ), so ist auch $f$ stetig.\\
(c) Sind $f, g: X \rightarrow Y$ stetig und ist $\{x \in X \mid f(x)=g(x)\}$ dicht in $X$, so ist $f \equiv g$.\\
(d) Die Mengen $C(X)$ und $C_{\mathrm{b}}(X)$ sind $\mathbb{K}$-Vektorräume. Genauer ist $C_{\mathrm{b}}(X)$ ein abgeschlossener Unterraum von $\ell^{\infty}(X)$ bzgl. der Norm $\|\cdot\|_{\infty}$. Dies folgt aus (b), da die Konvergenz bzgl. $\|\cdot\|_{\infty}$ genau der gleichmäßigen Konvergenz entspricht. Falls nicht explizit anders bemerkt, betrachten wir $C_{\mathrm{b}}(X)$ im Folgenden stets mit der Norm $\|\cdot\|_{\infty}$\\



\begin{CONC}{FA-A25-01-37}{Eigenschaften stetiger Abbildungen und der Räume $C(X)$, $C_b(X)$}
\textbf{(a) Kompositionsstetigkeit.}  
Sind $f:X\to Y$ und $g:Y\to Z$ stetig, so ist auch $g\circ f:X\to Z$ stetig.

\smallskip
\textbf{(b) Grenzstetigkeit bei gleichmäßiger Konvergenz.}  
Sei $(f_n)_{n\in\mathbb{N}}\subset C(X,Y)$ eine Folge stetiger Abbildungen, die \emph{gleichmäßig} gegen $f:X\to Y$ konvergiert, d.\,h.
\[
\sup_{x\in X} d_Y\big(f_n(x),f(x)\big)\xrightarrow[n\to\infty]{}0.
\]
Dann ist $f$ stetig.

\smallskip
\textbf{(c) Eindeutigkeit durch dichte Übereinstimmung.}  
Sind $f,g\in C(X,Y)$ und stimmt $f$ mit $g$ auf einer dichten Teilmenge $D\subset X$ überein,
\[
f(x)=g(x)\quad\text{für alle }x\in D,
\]
so gilt bereits $f\equiv g$ auf ganz $X$.

\smallskip
\textbf{(d) Algebraische und topologische Struktur.}  
\begin{itemize}
\item Sowohl $C(X)$ als auch $C_b(X)$ sind $\mathbb{K}$-Vektorräume.
\item $C_b(X)$ ist ein \textbf{abgeschlossener Unterraum} von $\ell^{\infty}(X)$ bezüglich der Supremumsnorm
      \[
      \|f\|_{\infty}:=\sup_{x\in X}|f(x)|.
      \]
      Die Abgeschlossenheit folgt aus (b), da $\|\cdot\|_{\infty}$-Konvergenz genau gleichmäßiger Konvergenz entspricht.
\item Wenn nicht anders angegeben, wird $C_b(X)$ stets mit der Norm $\|\cdot\|_{\infty}$ betrachtet.
\end{itemize}
\end{CONC}


2.24 Definition. Sei $f: X \rightarrow Y$ eine Abbildung\\
(a) $f$ heißt gleichmäßig stetig, falls für alle $\varepsilon>0$ ein $\delta>0$ existiert, so dass für alle $x_{1}, x_{2} \in X$ die Implikation gilt: $d_{X}\left(x_{1}, x_{2}\right)<\delta \Longrightarrow d_{Y}\left(f\left(x_{1}\right), f\left(x_{2}\right)\right)<\varepsilon$.\\
(b) $f$ heißt Lipschitz-stetig, falls $L>0$ existiert mit

$$
d_{Y}\left(f\left(x_{1}\right), f\left(x_{2}\right)\right) \leq L d_{X}\left(x_{1}, x_{2}\right) \quad \text { für alle } x_{1}, x_{2} \in X .
$$

(c) $f$ heißt Isometrie, falls $d_{Y}\left(f\left(x_{1}\right), f\left(x_{2}\right)\right)=d_{X}\left(x_{1}, x_{2}\right)$ für alle $x_{1}, x_{2} \in X$.\\
(d) $f$ heißt Homöomorphismus, falls $f$ bijektiv ist und $f$ sowie $f^{-1}$ stetig sind. Existiert solch ein $f$, so heißen $X$ und $Y$ homöomorph. Wir schreiben in diesem Fall $X \approx Y$.\\
2.25 Bemerkung. (a) Für eine Abbildung $f: X \rightarrow Y$ gelten folgende Implikationen: $f$ Isometrie $\Longrightarrow f$ Lipschitz-stetig $\Longrightarrow f$ gleichmäßig stetig $\Longrightarrow f$ stetig.\\
(b) Es gelte $\varnothing \neq A \subseteq X$. Dann ist die Abbildung $X \rightarrow \mathbb{R}, x \mapsto \operatorname{dist}(x, A)$ Lipschitzstetig (mit Lipschitz-Konstante 1).



\begin{DEF}{FA-A25-01-38}{Gleichmäßige, Lipschitz-Stetigkeit, Isometrien und Homöomorphismen}
Seien $(X,d_X)$ und $(Y,d_Y)$ metrische Räume und $f:X\to Y$ eine Abbildung.

\textbf{(a) Gleichmäßige Stetigkeit.}  
$f$ heißt \textbf{gleichmäßig stetig}, wenn gilt:
\[
\forall\,\varepsilon>0\ \exists\,\delta>0:\quad 
d_X(x_1,x_2)<\delta \ \Rightarrow\ d_Y\big(f(x_1),f(x_2)\big)<\varepsilon
\quad\forall\,x_1,x_2\in X.
\]
Der Wert $\delta$ hängt dabei nur von $\varepsilon$, nicht aber von $x_1,x_2$ ab.

\textbf{(b) Lipschitz-Stetigkeit.}  
$f$ heißt \textbf{Lipschitz-stetig}, wenn ein $L>0$ existiert mit
\[
d_Y\big(f(x_1),f(x_2)\big)\le L\, d_X(x_1,x_2)
\quad\forall\,x_1,x_2\in X.
\]
Das kleinste mögliche $L$ wird als \emph{Lipschitz-Konstante} von $f$ bezeichnet.

\textbf{(c) Isometrien.}  
$f$ heißt \textbf{Isometrie}, wenn
\[
d_Y\big(f(x_1),f(x_2)\big)=d_X(x_1,x_2)
\quad\forall\,x_1,x_2\in X.
\]
Eine Isometrie erhält also sämtliche Abstände exakt.

\textbf{(d) Homöomorphismen.}  
$f$ heißt \textbf{Homöomorphismus}, wenn $f$ bijektiv ist und sowohl $f$ als auch $f^{-1}$ stetig sind.  
In diesem Fall heißen $(X,d_X)$ und $(Y,d_Y)$ \textbf{homöomorph}, und wir schreiben
\[
X \approx Y.
\]
\end{DEF}




\begin{CONC}{FA-A25-01-39}{Zusammenhang der Stetigkeitsbegriffe}
Für eine Abbildung $f:X\to Y$ gilt die Implikationskette
\[
f\text{ Isometrie}\ \Longrightarrow\ f\text{ Lipschitz-stetig}\ \Longrightarrow\ f\text{ gleichmäßig stetig}\ \Longrightarrow\ f\text{ stetig.}
\]
\end{CONC}


\begin{EXA}{FA-A25-01-40}{Distanzfunktion ist Lipschitz-Stetig}
Für eine nichtleere Teilmenge $A\subset X$ ist die Abbildung
\[
f_A:X\to\mathbb{R},\quad f_A(x):=\operatorname{dist}(x,A):=\inf_{a\in A} d_X(x,a)
\]
\textbf{Lipschitz-stetig} mit Lipschitz-Konstante $L=1$, d.\,h.
\[
|\operatorname{dist}(x_1,A)-\operatorname{dist}(x_2,A)|\le d_X(x_1,x_2)
\quad\forall\,x_1,x_2\in X.
\]
\end{EXA}



Seien im Folgenden $\left(E,\|\cdot\|_{E}\right),\left(F,\|\cdot\|_{F}\right)$ normierte $\mathbb{K}$-Vektorräume.\\
2.26 Satz. Sei $T: E \rightarrow F \mathbb{K}$-linear.\\
(a) Äquivalent sind:

\begin{itemize}
  \item $T$ ist stetig;
  \item $T$ ist stetig in 0 ;
  \item $T\left(B_{1}(0)\right)$ ist in $F$ beschränkt (Erinnerung: $B_{1}(0):=\left\{x \in E \mid\|x\|_{E} \leq 1\right\}$ );
  \item $T$ ist Lipschitz-stetig.\\
(b) Ist $\operatorname{dim} E<\infty$, so ist $T$ stetig.
\end{itemize}



\begin{THEO}{FA-A25-01-41}{Äquivalente Stetigkeitskriterien für lineare Abbildungen}
Seien $(E,\|\cdot\|_{E})$ und $(F,\|\cdot\|_{F})$ normierte $\mathbb{K}$-Vektorräume und $T:E\to F$ eine $\mathbb{K}$-lineare Abbildung.

\textbf{(a) Äquivalenzen.}
Folgende Aussagen sind äquivalent:
\begin{enumerate}
\item[(i)] $T$ ist stetig (in jedem Punkt).
\item[(ii)] $T$ ist stetig in $0$.
\item[(iii)] $T(B_1(0))$ ist in $F$ beschränkt, wobei $B_1(0):=\{x\in E:\|x\|_{E}\le 1\}$.
\item[(iv)] $T$ ist Lipschitz-stetig.
\end{enumerate}

\textbf{(b) Endlicher Rang des Urbildraums.}
Ist $\dim E<\infty$, so ist jede lineare Abbildung $T:E\to F$ stetig.
\end{THEO}







Beweis. 

\begin{PROOF}{FA-A25-01-42}{P: Äquivalente Stetigkeitskriterien für lineare Abbildungen}
(a) bekannt und einfach.

(b) Man sieht direkt, dass eine Norm $\|\cdot\|_{T}$ auf $E$ definiert ist durch $\|x\|_{T}:=$ $\|x\|_{E}+\|T x\|_{F}$. Nach Bemerkung und Beispiel 2.3(a) ist $\|\cdot\|_{T}$ äquivalent zu $\|\cdot\|_{E}$, da $E$ endlichdimensional ist. Somit existiert $C>0$ mit
$$
\|x\|_{T} \leq C\|x\|_{E} \quad \text { für alle } x \in E,
$$
also insbesondere
$$
\|T x\|_{F} \leq\|x\|_{T} \leq C\|x\|_{E} \leq C \quad \text { für alle } x \in B_{1}(0) .
$$
\end{PROOF}


2.27 Definition. 

(a) Wir setzen
$$
\mathcal{L}(E, F):=\{T: E \rightarrow F \mid \mathrm{T} \text { ist linear und stetig }\} .
$$
Man rechnet leicht nach, dass $\mathcal{L}(E, F)$ mit der Definition
$$
\|T\|:=\|T\|_{\mathcal{L}(E, F)}:=\sup _{x \in B_{1}(0)}\|T x\|_{F}=\sup _{x \in E \backslash\{0\}} \frac{\|T x\|_{F}}{\|x\|_{E}}
$$
ein normierter Raum ist.

(b) Der Vektorraum $E^{\prime}:=\mathcal{L}(E, \mathbb{K})$ heißt topologischer Dualraum von $E$.\\

(c) Wir setzen $\mathcal{L}(E):=\mathcal{L}(E, E)$ (Raum der stetigen Endomorphismen von $E$ ).\\

(d) Eine bijektive $\mathbb{K}$-lineare Abbildung $T: E \rightarrow F$ heißt topologischer Isomorphismus, falls $T$ und $T^{-1}$ stetig sind. Existiert eine solche Abbildung $T$, so heißen $\left(E,\|\cdot\|_{E}\right)$ und $\left(F,\|\cdot\|_{F}\right)$ topologisch isomorph. Wir setzen
$$
\operatorname{Iso}(E, F):=\{T: E \rightarrow F \mid T \text { ist ein topologischer Isomorphismus }\}
$$
und schreiben kurz $\operatorname{Iso}(E)$ anstelle von $\operatorname{Iso}(E, E)$.\\

(e) Die Identität id: $E \rightarrow E$ schreiben wir meist als $I:=\mathrm{id}$, so wie in der Operatorentheorie üblich.



\begin{DEF}{FA-A25-01-43}{Stetige lineare Abbildungen, Dualraum und topologische Isomorphismen}
Seien $(E,\|\cdot\|_E)$ und $(F,\|\cdot\|_F)$ normierte $\mathbb{K}$-Vektorräume.

\textbf{(a) Raum stetiger linearer Abbildungen.}  
Wir definieren
\[
\mathcal{L}(E,F):=\{T:E\to F\mid T\text{ ist linear und stetig}\}.
\]
Für $T\in\mathcal{L}(E,F)$ wird die \textbf{Operatornorm} definiert durch
\[
\|T\|:=\|T\|_{\mathcal{L}(E,F)}:=\sup_{x\in B_1(0)}\|Tx\|_F
=\sup_{x\in E\setminus\{0\}}\frac{\|Tx\|_F}{\|x\|_E},
\]
wobei $B_1(0):=\{x\in E:\|x\|_E\le1\}$ die Einheitskugel in $E$ bezeichnet.  
Mit dieser Norm ist $\mathcal{L}(E,F)$ ein normierter Raum.

\textbf{(b) Topologischer Dualraum.}  
Der Raum
\[
E':=\mathcal{L}(E,\mathbb{K})
\]
heißt \textbf{topologischer Dualraum} von $E$.  
Seine Elemente heißen \emph{stetige Linearformen} oder \emph{(kontinuierliche) Funktionale} auf $E$.

\textbf{(c) Stetige Endomorphismen.}  
Wir setzen
\[
\mathcal{L}(E):=\mathcal{L}(E,E),
\]
also den Raum der stetigen Endomorphismen von $E$.

\textbf{(d) Topologische Isomorphismen.}  
Eine bijektive $\mathbb{K}$-lineare Abbildung $T:E\to F$ heißt \textbf{topologischer Isomorphismus},  
falls sowohl $T$ als auch $T^{-1}$ stetig sind.  
Existiert eine solche Abbildung, so heißen $(E,\|\cdot\|_E)$ und $(F,\|\cdot\|_F)$ \textbf{topologisch isomorph}.  
Man definiert
\[
\operatorname{Iso}(E,F):=\{T:E\to F\mid T \text{ ist topologischer Isomorphismus}\},
\qquad
\operatorname{Iso}(E):=\operatorname{Iso}(E,E).
\]

\textbf{(e) Identität.}  
Die Identitätsabbildung $\mathrm{id}_E:E\to E$ wird im Folgenden kurz mit
\[
I:=\mathrm{id}_E
\]
bezeichnet, entsprechend der üblichen Notation in der Operatorentheorie.
\end{DEF}




2.28 Bemerkung. (a) Die Elemente von $\mathcal{L}(E, F)$ nennt man auch beschränkte lineare Operatoren, und $\|\cdot\|=\|\cdot\|_{\mathcal{L}(E, F)}$ nennt man die Operatornorm.\\
(b) Ist $T \in \mathcal{L}(E, F)$, so ist $c=\|T\|$ die kleinste nichtnegative Zahl mit der Eigenschaft

$$
\|T x\|_{F} \leq c\|x\|_{E} \quad \text { für alle } x \in E .
$$

(c) Ist $\left(G,\|\cdot\|_{G}\right)$ ein weiterer normierter Raum und sind $T \in \mathcal{L}(E, F)$ und $S \in \mathcal{L}(F, G)$ gegeben, so gilt

$$
\|S T x\|_{G} \leq\|S\|\|T x\|_{F} \leq\|S\|\|T\|\|x\|_{E} \quad \text { für alle } x \in E,
$$

also

$$
S T \in \mathcal{L}(E, G) \quad \text { und } \quad\|S T\| \leq\|S\|\|T\| \quad \text { (Submultiplikativität der Norm). }
$$

Insbesondere gilt $\left\|T^{n}\right\| \leq\|T\|^{n}$ für $T \in \mathcal{L}(E)$ und $n \in \mathbb{N}$.



\begin{CONC}{FA-A25-01-44}{Operatornorm und Submultiplikativität}
\textbf{(a) Beschränkte lineare Operatoren.}  
Die Elemente von $\mathcal{L}(E,F)$ heißen \textbf{beschränkte lineare Operatoren},  
und die Norm $\|\cdot\|=\|\cdot\|_{\mathcal{L}(E,F)}$ wird \textbf{Operatornorm} genannt.

\textbf{(b) Charakterisierung der Operatornorm.}  
Für $T\in\mathcal{L}(E,F)$ ist
\[
c=\|T\|=\sup_{x\neq 0}\frac{\|Tx\|_F}{\|x\|_E}
\]
die kleinste nichtnegative Zahl mit der Eigenschaft
\[
\|Tx\|_F\le c\,\|x\|_E\quad\forall\,x\in E.
\]
Das heißt, $\|T\|$ ist die optimale Lipschitz-Konstante von $T$.

\textbf{(c) Komposition und Submultiplikativität.}  
Sei $(G,\|\cdot\|_G)$ ein weiterer normierter Raum und seien $T\in\mathcal{L}(E,F)$ sowie $S\in\mathcal{L}(F,G)$.  
Dann gilt für alle $x\in E$:
\[
\|STx\|_G \le \|S\|\|Tx\|_F \le \|S\|\|T\|\|x\|_E.
\]
Daraus folgt:
\[
ST \in \mathcal{L}(E,G)
\quad\text{und}\quad
\|ST\| \le \|S\|\|T\|.
\]
Diese Eigenschaft heißt \textbf{Submultiplikativität der Operatornorm}.

Insbesondere gilt für jedes $T\in\mathcal{L}(E)$ und jedes $n\in\mathbb{N}$:
\[
\|T^n\| \le \|T\|^n.
\]
\end{CONC}



2.29 Beispiel. 

(a) Sei der $\mathbb{K}$-Vektorraum $\mathcal{F}$ der finiten Folgen (siehe Definition und Bemerkung 2.21) versehen mit $\|\cdot\|_{\infty}$, d.h. $\|x\|_{\infty}=\sup _{k \in \mathbb{N}}\left|x_{k}\right|=\max _{k \in \mathbb{N}}\left|x_{k}\right|$ für $x \in \mathcal{F}$. Sei ferner $T: \mathcal{F} \rightarrow \mathcal{F}$ definiert durch $T\left(x_{k}\right)_{k}=\left(\frac{1}{k} x_{k}\right)_{k}$. Man sieht leicht, dass $T$ bijektiv ist. Da ferner $\|T x\|_{\infty} \leq\|x\|_{\infty}$ für alle $x \in \mathcal{F}$ gilt, ist $T$ auch stetig. Allerdings ist $T^{-1}$ nicht stetig, da

$$
\frac{1}{k} e_{k} \rightarrow 0 \text { in }\left(\mathcal{F},\|\cdot\|_{\infty}\right) \text { und }\left\|T^{-1} \frac{1}{k} e_{k}\right\|_{\infty}=\left\|e_{k}\right\|_{\infty}=1
$$

für alle $k$ gilt. Hier sei $e_{k}:=\left(\delta_{n k}\right)_{n} \in \mathcal{F}$ wie in Definition und Bemerkung 2.21 der $k$-te Einheitsvektor in $\mathcal{F}$.



\begin{EXA}{FA-A25-01-45}{Nicht stetigen inverse Operator trotz Bijektivität (Raum finiter Folgen)}
\textbf{(a) Konstruktion.}  
Sei $\mathcal{F}$ der $\mathbb{K}$-Vektorraum der \emph{finiten Folgen}, d.\,h.
\[
\mathcal{F} := \{x=(x_k)_{k\in\mathbb{N}} \mid x_k=0 \text{ für alle bis auf endlich viele } k\}.
\]
Wir versehen $\mathcal{F}$ mit der Norm
\[
\|x\|_\infty := \sup_{k\in\mathbb{N}} |x_k| = \max_{k\in\mathbb{N}} |x_k|,
\]
und betrachten die lineare Abbildung
\[
T:\mathcal{F}\to\mathcal{F}, \quad T(x_k)_k := \left(\frac{1}{k}x_k\right)_k.
\]

\textbf{(b) Eigenschaften.}  
Offensichtlich ist $T$ bijektiv, da für jedes $y=(y_k)\in\mathcal{F}$ gilt:
\[
x := (k y_k)_k \in \mathcal{F}, \quad T(x)=y.
\]
Ferner gilt für alle $x\in\mathcal{F}$:
\[
\|Tx\|_\infty = \sup_k \left|\frac{1}{k}x_k\right| \le \sup_k |x_k| = \|x\|_\infty.
\]
Also ist $T$ linear und stetig mit $\|T\| \le 1$.

\textbf{(c) Nichtstetigkeit des Inversen.}  
Für das Inverse gilt $T^{-1}(x_k)_k = (k x_k)_k$.  
Betrachten wir die Folge $\frac{1}{k} e_k$, wobei $e_k := (\delta_{nk})_n$ der $k$-te Einheitsvektor in $\mathcal{F}$ ist.
Dann gilt
\[
\frac{1}{k} e_k \to 0 \quad \text{in } (\mathcal{F},\|\cdot\|_\infty),
\]
denn $\left\|\frac{1}{k} e_k\right\|_\infty = \frac{1}{k} \to 0$.

Andererseits:
\[
T^{-1}\left(\frac{1}{k} e_k\right) = e_k,
\quad\text{und somit}\quad
\left\|T^{-1}\frac{1}{k} e_k\right\|_\infty = \|e_k\|_\infty = 1 \quad \forall k.
\]
Also konvergiert $\frac{1}{k} e_k \to 0$, während $T^{-1}\left(\frac{1}{k} e_k\right)\not\to 0$.  
Damit ist $T^{-1}$ \emph{nicht stetig}.

\textbf{(d) Fazit.}  
$T$ ist also ein Beispiel für einen \emph{bijektiven, stetigen linearen Operator}, dessen Inverses \emph{nicht stetig} ist.  
Folglich ist $T$ \emph{kein topologischer Isomorphismus}.
\end{EXA}


(b) Sei $1 \leq p<q \leq \infty$. Dann ist die identische Abbildung $T:\left(\ell^{p},\|\cdot\|_{p}\right) \rightarrow\left(\ell^{p},\|\cdot\|_{q}\right)$ stetig (da $\|x\|_{q} \leq\|x\|_{p}$ für alle $x \in \ell^{p}$ gilt), aber $T^{-1}$ ist nicht stetig, da $\|\cdot\|_{q}$ und $\|\cdot\|_{p}$ auf $\ell^{p}$ nicht äquivalent sind.\\


\begin{EXA}{FA-A25-01-46}{Identität zwischen $\ell^p$ mit unterschiedlichen Normen ist nicht topologisch isomorph}
Sei $1\le p<q\le\infty$ und
\[
T:\big(\ell^{p},\|\cdot\|_{p}\big)\longrightarrow \big(\ell^{p},\|\cdot\|_{q}\big),\qquad T(x)=x
\]
die identische Abbildung. Dann gilt:
\begin{enumerate}
\item $T$ ist stetig (sogar Lipschitz mit Konstante $1$).
\item $T^{-1}:\big(\ell^{p},\|\cdot\|_{q}\big)\to\big(\ell^{p},\|\cdot\|_{p}\big)$ ist \emph{nicht} stetig. Insbesondere sind die Normen $\|\cdot\|_{p}$ und $\|\cdot\|_{q}$ auf $\ell^{p}$ nicht äquivalent; folglich ist $T$ kein topologischer Isomorphismus.
\end{enumerate}

\emph{Beweis.}
(1) Für alle $x\in\ell^p$ gilt $\|x\|_{q}\le \|x\|_{p}$ (Inklusion $\ell^p\subset \ell^q$ für $p<q$), also
\[
\|Tx\|_{q}=\|x\|_{q}\le \|x\|_{p},
\]
d.\,h. $T$ ist Lipschitz mit Konstante $1$ und damit stetig.

\smallskip
(2) Wir zeigen, dass $T^{-1}$ nicht stetig ist, indem wir eine Folge $(x^{(N)})_{N}$ konstruieren, die in $\|\cdot\|_{q}$ gegen $0$ konvergiert, deren $\|\cdot\|_{p}$-Norm jedoch konstant $1$ bleibt.

\textbf{Fall $q<\infty$:} Definiere $x^{(N)}\in\ell^p$ durch
\[
x^{(N)}_k:=\begin{cases}
N^{-1/p}, & 1\le k\le N,\\
0,& k>N.
\end{cases}
\]
Dann ist
\[
\|x^{(N)}\|_{p}=\Big(\sum_{k=1}^{N} |N^{-1/p}|^{p}\Big)^{1/p}
=(N\cdot N^{-1})^{1/p}=1,
\]
während
\[
\|x^{(N)}\|_{q}=\Big(\sum_{k=1}^{N} |N^{-1/p}|^{q}\Big)^{1/q}
= \big(N\cdot N^{-q/p}\big)^{1/q} = N^{\frac{1}{q}-\frac{1}{p}}\xrightarrow[N\to\infty]{}0,
\]
da $\frac{1}{q}-\frac{1}{p}<0$. Also $x^{(N)}\to 0$ in $\|\cdot\|_{q}$, aber $\|x^{(N)}\|_{p}=1$ für alle $N$. Somit ist $T^{-1}$ nicht stetig.

\textbf{Fall $q=\infty$:} Mit derselben Folge gilt
\[
\|x^{(N)}\|_{p}=1,\qquad \|x^{(N)}\|_{\infty}=\sup_{1\le k\le N} N^{-1/p}=N^{-1/p}\xrightarrow[N\to\infty]{}0.
\]
Also wieder $x^{(N)}\to 0$ in $\|\cdot\|_{\infty}$, aber nicht in $\|\cdot\|_{p}$.

In beiden Fällen ist $T^{-1}$ nicht stetig.
\end{EXA}



(c) (Integraloperatoren) Seien $[a, b],[c, d] \subseteq \mathbb{R}$ kompakte Intervalle, und sei $k:[c, d] \times$ $[a, b] \rightarrow \mathbb{C}$ stetig. Wir schreiben $C[a, b]:=C([a, b]), C[c, d]:=C([c, d])$ und definieren die lineare Abbildung $K: C[a, b] \rightarrow C[c, d]$ durch
$$
(K u)(t):=\int_{a}^{b} k(t, s) u(s) \mathrm{d} s \quad \text { für } t \in[c, d]
$$
Offensichtlich ist $K$ wohldefiniert und linear. $K$ ist stetig, denn für alle $u \in C[a, b]$\\
und $t \in[c, d]$ gilt
$$
|(K u)(t)| \leq \int_{a}^{b}|k(t, s)||u(s)| \mathrm{d} s \leq M\|u\|_{\infty}
$$
mit
$$
M:=\max _{c \leq t \leq d} \int_{a}^{b}|k(t, s)| \mathrm{d} s<\infty
$$
also $\|K u\|_{\infty} \leq M\|u\|_{\infty}$ für $u \in C[a, b]$. Es folgt somit $K \in \mathcal{L}(C[a, b], C[c, d])$ mit $\|K\| \leq M$. Wir werden nun sehen, dass sogar $\|K\|=M$ gilt. Sei dazu $t_{0} \in[c, d]$ ein Punkt, wo das Maximum in (2.3) angenommen wird. Für $\varepsilon>0$ sei ferner
$$
u_{\varepsilon} \in C[a, b] \quad \text { definiert durch } \quad u_{\varepsilon}(s)=\frac{\overline{k\left(t_{0}, s\right)}}{\left|k\left(t_{0}, s\right)\right|+\varepsilon}
$$
$\mathrm{Da}\left\|u_{\varepsilon}\right\|_{\infty} \leq 1$ ist, unabhängig von $\varepsilon>0$, gilt
$$
\|K\| \geq\left\|K u_{\varepsilon}\right\|_{\infty} \geq\left|\left(K u_{\varepsilon}\right)\left(t_{0}\right)\right|=\int_{a}^{b} \frac{\left|k\left(t_{0}, s\right)\right|^{2}}{\left|k\left(t_{0}, s\right)\right|+\varepsilon} \mathrm{d} s
$$
wobei das letzte Integral für $\varepsilon \rightarrow 0$ nach dem Satz von der majorisierten Konvergenz gegen $\int_{a}^{b}\left|k\left(t_{0}, s\right)\right| d s=M$ konvergiert. Also ist auch $\|K\| \geq M$, und insgesamt folgt Gleichheit. 

Spezielles Beispiel: Sei $K \in \mathcal{L}(C[0,1], C[0,1])$ der Lösungsoperator zum Problem (1.2) aus Kapitel 1, d.h. es gelte
$$
K f(t)=\int_{0}^{1} G(t, s) f(s) \mathrm{d} s \quad \text { für } f \in C[0,1], t \in[0,1]
$$
mit der Greenfunktion
$$
G(t, s)= \begin{cases}(1-t) s, & s \leq t \\ (1-s) t, & s \geq t\end{cases}
$$
Dann ist
$$
\|K\|=\max _{t \in[0,1]} h(t)
$$
mit
$$
\begin{aligned}
h(t) & =\int_{0}^{1}|G(t, s)| \mathrm{d} s=(1-t) \int_{0}^{t} s \mathrm{~d} s+t \int_{t}^{1}(1-s) \mathrm{d} s \\
& =(1-t) \int_{0}^{t} s \mathrm{~d} s+t \int_{0}^{1-t} s \mathrm{~d} s=\frac{1}{2}\left((1-t) t^{2}+t(1-t)^{2}\right)=\frac{t(1-t)}{2}
\end{aligned}
$$
Es folgt $\|K\|=\frac{1}{8}$.



\begin{EXA}{FA-A25-01-47}{Operatornorm eines Integraloperators auf $C$-Räumen}
Seien $[a,b],[c,d]\subset\mathbb{R}$ kompakte Intervalle und $k:[c,d]\times[a,b]\to\mathbb{C}$ stetig. Definiere
\[
K:C[a,b]\longrightarrow C[c,d],\qquad (Ku)(t):=\int_a^b k(t,s)u(s)\,ds.
\]
Dann gilt:
\begin{enumerate}
\item[(i)] $K$ ist wohldefiniert, linear und stetig mit
\[
\|Ku\|_\infty \le M\,\|u\|_\infty\quad\text{für alle }u\in C[a,b],
\quad\text{wobei}\quad
M:=\max_{t\in[c,d]}\int_a^b |k(t,s)|\,ds<\infty.
\]
Insbesondere $K\in\mathcal{L}\!\left(C[a,b],C[c,d]\right)$ und $\|K\|\le M$.
\item[(ii)] Es gilt sogar Gleichheit $\|K\|=M$.
\end{enumerate}

\textbf{Beweis.}
(i) Stetigkeit von $k$ auf dem Kompaktum $[c,d]\times[a,b]$ impliziert Beschränktheit. Für $u\in C[a,b]$ und $t\in[c,d]$:
\[
|(Ku)(t)|=\Big|\int_a^b k(t,s)u(s)\,ds\Big|
\le \int_a^b |k(t,s)|\,|u(s)|\,ds
\le \|u\|_\infty \int_a^b |k(t,s)|\,ds.
\]
Maximieren nach $t$ liefert $\|Ku\|_\infty\le M\|u\|_\infty$ und damit die Stetigkeit sowie $\|K\|\le M$.

(ii) Wähle $t_0\in[c,d]$ mit $\int_a^b |k(t_0,s)|\,ds=M$. Für $\varepsilon>0$ setze
\[
u_\varepsilon(s):=\frac{\overline{k(t_0,s)}}{|k(t_0,s)|+\varepsilon}\in C[a,b],
\qquad \|u_\varepsilon\|_\infty\le 1.
\]
Dann
\[
\|K\|\ \ge\ \|Ku_\varepsilon\|_\infty\ \ge\ |(Ku_\varepsilon)(t_0)|
=\int_a^b \frac{|k(t_0,s)|^2}{|k(t_0,s)|+\varepsilon}\,ds
\ \xrightarrow[\varepsilon\to 0]{}\ \int_a^b |k(t_0,s)|\,ds=M,
\]
wobei der letzte Schritt mit dem Satz von der majorisierten Konvergenz (Majorante $|k(t_0,s)|$) folgt. Also $\|K\|\ge M$. Zusammen mit (i) folgt $\|K\|=M$. \qed
\end{EXA}



\begin{EXA}{FA-A25-01-48}{Norm des Green-Integraloperators auf $C[0,1]$}
Betrachte $K\in\mathcal{L}(C[0,1],C[0,1])$ gegeben durch
\[
(Kf)(t)=\int_0^1 G(t,s)f(s)\,ds,\qquad
G(t,s)=\begin{cases}(1-t)s,& s\le t,\\ (1-s)t,& s\ge t.\end{cases}
\]
Dann ist
\[
\|K\|=\max_{t\in[0,1]} h(t),\qquad
h(t):=\int_0^1 |G(t,s)|\,ds.
\]
Da $G\ge 0$, ist $|G|=G$ und
\[
h(t)=(1-t)\int_0^t s\,ds\ +\ t\int_t^1(1-s)\,ds
=(1-t)\frac{t^2}{2}+t\frac{(1-t)^2}{2}
=\frac{t(1-t)}{2}.
\]
Das Maximum von $t(1-t)$ auf $[0,1]$ liegt bei $t=\tfrac12$ mit Wert $\tfrac14$. Also
\[
\|K\|=\max_{t\in[0,1]}\frac{t(1-t)}{2}=\frac{1}{8}.
\]
\end{EXA}



\pagebreak

\subsection*{Vollständigkeit und Reihen}


2.30 Definition. (a) Eine Folge $\left(x_{n}\right)_{n}$ in $X$ heißt Cauchyfolge, wenn $\lim _{n \rightarrow \infty} \sup _{m \geq n} d\left(x_{m}, x_{n}\right)=0$ gilt.\\
(b) ( $X, d$ ) heißt vollständig, wenn jede Cauchyfolge in $X$ konvergiert.\\
(c) Ein vollständiger normierter Raum ( $E,\|\cdot\|$ ) heißt Banachraum.\\
(d) Ein vollständiger Prähilbertraum ( $E,\langle\cdot, \cdot\rangle$ ) heißt Hilbertraum.

In (c) und (d) bezieht sich die Vollständigkeit dabei auf die induzierte Metrik.\\



\begin{DEF}{FA-A25-01-49}{Cauchyfolgen, Vollständigkeit und Banach-/Hilberträume}
\begin{enumerate}
\item[(a)] Eine Folge $\left(x_{n}\right)_{n}$ in einem metrischen Raum $(X,d)$ heißt \textbf{Cauchyfolge}, falls gilt
\[
\lim_{n\to\infty}\ \sup_{m\ge n} d(x_m,x_n)=0,
\]
d.\,h. für jedes $\varepsilon>0$ existiert $N\in\mathbb{N}$ mit
\[
d(x_m,x_n)<\varepsilon \quad\text{für alle } m,n\ge N.
\]

\item[(b)] Ein metrischer Raum $(X,d)$ heißt \textbf{vollständig}, wenn jede Cauchyfolge in $X$ konvergiert, d.\,h. es existiert für jede Cauchyfolge $(x_n)$ ein $x\in X$ mit $x_n\to x$.

\item[(c)] Ein \textbf{Banachraum} ist ein vollständiger normierter Raum $(E,\|\cdot\|)$, d.\,h. jeder Cauchyfolge $(x_n)$ in $E$ existiert ein Grenzwert $x\in E$ mit $\|x_n - x\| \to 0$.

\item[(d)] Ein \textbf{Hilbertraum} ist ein vollständiger Prähilbertraum $(E,\langle\cdot,\cdot\rangle)$, d.\,h. ein innerer Produktraum, dessen induzierte Norm
\[
\|x\|:=\sqrt{\langle x,x\rangle}
\]
den Raum vollständig macht.

\smallskip
\textit{Anmerkung:} In (c) und (d) bezieht sich die Vollständigkeit jeweils auf die von der Norm induzierte Metrik $d(x,y):=\|x-y\|$.
\end{enumerate}
\end{DEF}


2.31 Bemerkung. Sei $\left(x_{n}\right)_{n} \subseteq X$ eine Folge.\\
(a) Ist $\left(x_{n}\right)_{n}$ konvergent, so ist $\left(x_{n}\right)_{n}$ eine Cauchyfolge.\\
(b) Ist $\left(x_{n}\right)_{n}$ eine Cauchyfolge, so ist $\left(x_{n}\right)_{n}$ in $X$ beschränkt, d.h. die Menge $\left\{x_{n} \mid n \in \mathbb{N}\right\} \subseteq X$ ist beschränkt.\\
(c) Ist $\left(x_{n}\right)_{n}$ eine Cauchyfolge und besitzt $\left(x_{n}\right)_{n}$ eine gegen $a \in X$ konvergente Teilfolge, so konvergiert bereits die Folge $\left(x_{n}\right)_{n}$ selbst gegen $a$.\\
(d) Ist $\left(x_{n}\right)_{n}$ eine Cauchyfolge, $(Y, \tilde{d})$ ein weiterer metrischer Raum und $f: X \rightarrow Y$ gleichmäßig stetig, so ist auch $\left(f\left(x_{n}\right)\right)_{n}$ eine Cauchyfolge (in $Y$ ).\\



\begin{CONC}{FA-A25-01-50}{Eigenschaften von Cauchyfolgen}
Sei $(X,d)$ ein metrischer Raum und $\left(x_{n}\right)_{n}\subseteq X$ eine Folge.
\begin{enumerate}
\item[(a)] \textbf{Konvergenz impliziert Cauchy-Eigenschaft.}  
Ist $\left(x_{n}\right)_{n}$ konvergent, d.\,h. $x_{n}\to a$ für ein $a\in X$, so ist $\left(x_{n}\right)_{n}$ eine Cauchyfolge.  
\emph{Begründung:} Für jedes $\varepsilon>0$ existiert $N$ mit $d(x_n,a)<\varepsilon/2$ für alle $n\ge N$. Dann gilt für $m,n\ge N$:
\[
d(x_m,x_n)\le d(x_m,a)+d(a,x_n)<\varepsilon.
\]

\item[(b)] \textbf{Cauchyfolgen sind beschränkt.}  
Ist $\left(x_{n}\right)_{n}$ eine Cauchyfolge, so existiert $R>0$ mit 
\[
d(x_n,x_m)\le R \quad\text{für alle }m,n\in\mathbb{N}.
\]
Damit ist die Menge $\{x_n : n\in\mathbb{N}\}$ beschränkt.  
\emph{Begründung:} Wähle $N$ mit $d(x_n,x_m)<1$ für $n,m\ge N$. Dann gilt für beliebiges $n$:  
\[
d(x_n,x_N)\le 1+\max_{1\le k\le N}d(x_k,x_N),
\]
also beschränkt.

\item[(c)] \textbf{Cauchyfolge mit konvergenter Teilfolge ist konvergent.}  
Besitzt eine Cauchyfolge $(x_n)$ eine Teilfolge $(x_{n_k})$, die gegen $a\in X$ konvergiert, so gilt bereits $x_n\to a$.  
\emph{Begründung:} Für jedes $\varepsilon>0$ existieren $N_1,N_2$ mit  
$d(x_n,x_m)<\varepsilon/2$ für $n,m\ge N_1$ und $d(x_{n_k},a)<\varepsilon/2$ für $k\ge N_2$.  
Dann gilt für $n\ge N:=\max(N_1,n_{N_2})$:  
\[
d(x_n,a)\le d(x_n,x_{n_k})+d(x_{n_k},a)<\varepsilon.
\]

\item[(d)] \textbf{Gleichmäßig stetige Abbildungen erhalten Cauchyfolgen.}  
Sei $(Y,\tilde d)$ ein weiterer metrischer Raum und $f:(X,d)\to(Y,\tilde d)$ gleichmäßig stetig.  
Dann ist $(f(x_n))_n$ eine Cauchyfolge in $Y$.  
\emph{Begründung:} Für jedes $\varepsilon>0$ existiert $\delta>0$ mit  
$d(x_m,x_n)<\delta\Rightarrow \tilde d(f(x_m),f(x_n))<\varepsilon$.  
Da $(x_n)$ Cauchy ist, existiert $N$ mit $d(x_m,x_n)<\delta$ für $m,n\ge N$, womit die Behauptung folgt. \qed
\end{enumerate}
\end{CONC}


2.32 Satz. Sei $A \subseteq X$.\\
(a) Ist ( $X, d$ ) vollständig und $A \subseteq X$ abgeschlossen, so ist auch ( $A, d$ ) vollständig.\\
(b) Ist ( $A, d$ ) vollständig, so ist $A$ abgeschlossen in $X$.



\begin{PROP}{FA-A25-01-51}{Vollständigkeit und Abgeschlossenheit}
Sei $(X,d)$ ein metrischer Raum und $A\subseteq X$.
\begin{enumerate}
\item[(a)] Ist $(X,d)$ vollständig und $A$ abgeschlossen, so ist auch der Teilraum $(A,d)$ vollständig.

\item[(b)] Ist $(A,d)$ vollständig, so ist $A$ abgeschlossen in $X$.
\end{enumerate}
\end{PROP}



Beweis. 

\begin{PROOF}{FA-A25-01-52}{P: Vollständigkeit und Abgeschlossenheit}
(a) Sei $\left(x_{n}\right)_{n}$ eine Cauchyfolge in $A$. Da $X$ vollständig ist, existiert $x=$ $\lim _{n \rightarrow \infty} x_{n} \in X$, wobei $x \in \bar{A}=A$ nach Voraussetzung gilt. Also konvergiert $\left(x_{n}\right)_{n}$ in $A$.

(b) Sei $x \in \bar{A}$. Dann ist $x=\lim _{n \rightarrow \infty} x_{n}$ für eine Folge $\left(x_{n}\right)_{n}$ in $A$. Nach Bemerkung 2.31(a) ist $\left(x_{n}\right)_{n}$ eine Cauchyfolge in $A$, und damit konvergiert $\left(x_{n}\right)_{n}$ nach Voraussetzung in $A$. Es folgt $x \in A$. Insgesamt folgt $A=\bar{A}$.
\end{PROOF}


2.33 Bemerkung. Sei ( $Y, \tilde{d}$ ) ein weiterer metrischer Raum und $f: X \rightarrow Y$ eine Abbildung.\\
(a) Ist $f$ eine Isometrie, so ist neben $f$ auch $f^{-1}: f(X) \rightarrow X$ gleichmäßig stetig, und aus Bemerkung 2.31(d) folgt:
$$
(X, d) \text { vollständig } \Longleftrightarrow(f(X), \tilde{d}) \text { vollständig }
$$
Gilt dies, so ist $f(X)$ abgeschlossen in $Y$ nach Satz 2.32.

(b) Ist $f$ ein Homöomorphismus, so erhält $f$ nicht notwendigerweise die Vollständigkeit: Seien z.B. $X=\mathbb{R}$ und $Y=\left(-\frac{\pi}{2}, \frac{\pi}{2}\right)$, jeweils versehen mit der Betragsmetrik $d(x, y)=|x-y|$. Sei ferner der Homöomorphismus $f: X \rightarrow Y$ gegeben durch $f(x)=\arctan x$. Dann ist $X$ vollständig und $Y$ nicht, da $Y \subseteq X$ nicht abgeschlossen ist.


\begin{CONC}{FA-A25-01-53}{Isometrien, Homöomorphismen und Vollständigkeit}
Seien $(X,d)$ und $(Y,\tilde d)$ metrische Räume und $f:X\to Y$.

\begin{enumerate}
\item[(a)] \textbf{Isometrien erhalten Vollständigkeit (auf dem Bild).}
Ist $f$ eine Isometrie, d.\,h. $\tilde d(f(x_1),f(x_2))=d(x_1,x_2)$ für alle $x_1,x_2\in X$, so ist $f$ insbesondere bijektiv auf $f(X)$ und sowohl $f$ als auch $f^{-1}:f(X)\to X$ sind gleichmäßig stetig. Aus Bemerkung~\ref{ANA-L02-31}~(d) folgt die Äquivalenz
\[
(X,d)\ \text{vollständig}\ \Longleftrightarrow\ (f(X),\tilde d)\ \text{vollständig}.
\]
Ist $(X,d)$ vollständig, so ist auch $(f(X),\tilde d)$ vollständig und damit $f(X)$ nach Satz~\ref{ANA-L02-32} in $Y$ abgeschlossen.

\item[(b)] \textbf{Homöomorphismen erhalten Vollständigkeit im Allgemeinen nicht.}
Sei $X=\mathbb{R}$ mit $d(x,y)=|x-y|$ und $Y=\bigl(-\frac{\pi}{2},\frac{\pi}{2}\bigr)$ mit derselben Betragsmetrik. Der Homöomorphismus
\[
f:X\to Y,\qquad f(x)=\arctan x
\]
ist stetig bijektiv mit stetiger Umkehrabbildung $\tan:Y\to X$. Dennoch ist $X$ vollständig, wohingegen $Y$ nicht vollständig ist, da $Y$ in $X$ nicht abgeschlossen ist (z.\,B.\ ist $(y_n)$ mit $y_n=\frac{\pi}{2}-\frac1n$ eine Cauchyfolge in $Y$, die in $\mathbb{R}$ gegen $\frac{\pi}{2}\notin Y$ konvergiert).
\end{enumerate}
\end{CONC}


2.34 Bemerkungen und Beispiele. (a) Sind $\|\cdot\|_{1},\|\cdot\|_{2}$ äquivalente Normen auf einem $\mathbb{K}$-Vektorraum $E$, so gilt:

$$
\left(E,\|\cdot\|_{1}\right) \text { Banachraum } \quad \Longleftrightarrow \quad\left(E,\|\cdot\|_{2}\right) \text { Banachraum }
$$

Allgemeiner gilt für zwei topologisch isomorphe $\mathbb{K}$-Vektorräume ( $E,\|\cdot\|_{E}$ ) und $\left(F,\|\cdot\|_{F}\right)$ :

$$
E \text { Banachraum } \quad \Longleftrightarrow \quad F \text { Banachraum }
$$



\begin{PROP}{FA-A25-01-54}{Erhaltung von Vollständigkeit unter Norm-Äquivalenz und Isomorphismus}
Seien $\|\cdot\|_{1}$ und $\|\cdot\|_{2}$ zwei äquivalente Normen auf einem $\mathbb{K}$-Vektorraum $E$.

\begin{enumerate}
\item (a) Dann gilt
\[
\left(E,\|\cdot\|_{1}\right) \text{ ist vollständig (Banachraum)}
\quad \Longleftrightarrow \quad
\left(E,\|\cdot\|_{2}\right) \text{ ist vollständig (Banachraum)}.
\]

\item (b) Allgemeiner gilt: Sind $(E,\|\cdot\|_{E})$ und $(F,\|\cdot\|_{F})$ topologisch isomorph, d.\,h.\ es existiert ein topologischer Isomorphismus $T\in\operatorname{Iso}(E,F)$, dann gilt
\[
E \text{ ist Banachraum } \quad \Longleftrightarrow \quad F \text{ ist Banachraum.}
\]
\end{enumerate}
\end{PROP}


\begin{PROOF}{FA-A25-01-55}{P: Erhaltung von Vollständigkeit unter Norm-Äquivalenz und Isomorphismus}
Aus der Äquivalenz der Normen existieren Konstanten $c,C>0$ mit
\[
c\|x\|_{1}\le \|x\|_{2}\le C\|x\|_{1}\quad \text{für alle } x\in E.
\]
Damit gilt: Eine Folge $(x_n)$ ist genau dann Cauchy bzgl.\ $\|\cdot\|_{1}$, wenn sie Cauchy bzgl.\ $\|\cdot\|_{2}$ ist.  
Somit konvergiert $(x_n)$ in einer Norm genau dann, wenn sie in der anderen konvergiert.  
Daher ist $(E,\|\cdot\|_{1})$ vollständig $\iff (E,\|\cdot\|_{2})$ vollständig. \qed

Ist $E$ vollständig und $T:E\to F$ ein stetiger, bijektiver Operator mit stetiger Umkehrabbildung, so ist das Bild einer Cauchyfolge in $E$ wiederum eine Cauchyfolge in $F$, und jede konvergente Folge in $E$ wird auf eine konvergente Folge in $F$ abgebildet.  
Die Umkehrung folgt analog unter Anwendung von $T^{-1}$.
\end{PROOF}





(b) $\mathbb{K}^{n}$ ist ein Banachraum (bzgl. jeder Norm). Ist ferner $E$ irgendein normierter Raum mit $n:=\operatorname{dim} E<\infty$, so ist $E$ topologisch isomorph zu $\mathbb{K}^{n}$ und damit ein Banachraum nach (a).\\


\begin{EXA}{FA-A25-01-56}{Endlichdimensionale Banachräume}
\begin{enumerate}
\item[(a)] Der Raum $\mathbb{K}^{n}$ ist bezüglich jeder Norm vollständig, also ein \textbf{Banachraum}.  
Denn alle Normen auf einem endlichdimensionalen Raum sind äquivalent, und die Vollständigkeit gilt insbesondere für die euklidische Norm.

\item[(b)] Ist $E$ ein normierter Raum mit $\dim E = n < \infty$, so existiert ein topologischer Isomorphismus
\[
T : E \longrightarrow \mathbb{K}^{n}.
\]
Nach Satz~\ref{ANA-L02-34} folgt daraus:
\[
E \text{ ist Banachraum } \quad \Longleftrightarrow \quad \mathbb{K}^{n} \text{ ist Banachraum}.
\]
Da $\mathbb{K}^{n}$ vollständig ist, ist somit auch jeder endlichdimensionale normierte Raum $E$ ein Banachraum.
\end{enumerate}
\end{EXA}



(c) Jeder endlich-dimensionale Unterraum eines normierten Raums ist abgeschlossen nach (b) und Satz 2.32.\\


\begin{LEM}{FA-A25-01-57}{Endlichdimensionale Unterräume sind Abgeschlossen}
Sei $(E,\|\cdot\|)$ ein normierter Raum. Dann ist jeder endlichdimensionale Unterraum $U \subseteq E$ abgeschlossen.
\end{LEM}


\begin{PROOF}{FA-A25-01-58}{P: Endlichdimensionale Unterräume sind Abgeschlossen}
Nach Beispiel ist $U$ endlichdimensional und damit ein Banachraum bezüglich der von $E$ induzierten Norm.  
Da $E$ selbst ein metrischer Raum ist und $(U,\|\cdot\|)$ als Teilraum von $E$ vollständig ist, folgt aus Satz~\ref{ANA-L02-32}, dass $U$ in $E$ abgeschlossen ist.
\end{PROOF}



(d) 

$\ell^{p}$ ist ein Banachraum für $p \in[1, \infty]$. Wir beweisen dies zunächst im Fall $p<\infty$ : Sei $\left(x_{n}\right)_{n}$ eine Cauchyfolge in $\ell^{p}$, wobei wir $x_{n}=\left(\xi_{n k}\right)_{k}$ notieren. Sei $k \in \mathbb{N}$ fest; für $m, n \in \mathbb{N}$ gilt dann
$$
\left|\xi_{m k}-\xi_{n k}\right|^{p} \leq \sum_{j \in \mathbb{N}}\left|\xi_{m j}-\xi_{n j}\right|^{p}=\left\|x_{m}-x_{n}\right\|_{p}^{p}
$$
und somit ist die Folge $\left(\xi_{n k}\right)_{n} \subseteq \mathbb{K}$ eine Cauchyfolge. Da $\mathbb{K}$ vollständig ist, existiert $\xi_{k}:=\lim _{n \rightarrow \infty} \xi_{n k} \in \mathbb{K}$. Setze $x:=\left(\xi_{k}\right)_{k}$. Für alle $\ell \in \mathbb{N}$ gilt dann
$$
\sum_{k=1}^{\ell}\left|\xi_{k}\right|^{p}=\lim _{n \rightarrow \infty} \sum_{k=1}^{\ell}\left|\xi_{n k}\right|^{p} \leq \limsup _{n \rightarrow \infty}\left\|x_{n}\right\|_{p}^{p} \stackrel{\text { Bemerkung 2.31(b) }}{<} \infty
$$
Also ist $\sum_{k \in \mathbb{N}}\left|\xi_{k}\right|^{p}<\infty$ und somit $x \in \ell^{p}$. Für alle $n, \ell \in \mathbb{N}$ gilt zudem
$$
\sum_{k=1}^{\ell}\left|\xi_{k}-\xi_{n k}\right|^{p}=\lim _{m \rightarrow \infty} \sum_{k=1}^{\ell}\left|\xi_{m k}-\xi_{n k}\right|^{p} \leq \sup _{m \geq n}\left\|x_{m}-x_{n}\right\|_{p}^{p}=: c_{n}
$$
und somit $\left\|x-x_{n}\right\|_{p}^{p} \leq c_{n}$. Nach Voraussetzung gilt zudem $\lim _{n \rightarrow \infty} c_{n}=0$ und somit $\lim _{n \rightarrow \infty} x_{n}=x$ in $\ell^{p}$; dies war zu zeigen. 

Der Beweis für $p=\infty$ ist ähnlich, nur einfacher. 





\begin{EXA}{FA-A25-01-59}{$\ell^{p}$ ist vollständig}
Für jedes $p\in[1,\infty]$ ist $(\ell^{p},\|\cdot\|_{p})$ ein Banachraum.

\begin{proof}
\textbf{Fall $1\le p<\infty$.}
Sei $(x_n)_{n}$ eine Cauchyfolge in $\ell^{p}$, mit $x_n=(\xi_{nk})_{k\in\mathbb{N}}$.
Dann gilt für jedes feste $k\in\mathbb{N}$ und alle $m,n$:
\[
|\xi_{mk}-\xi_{nk}|^{p}\le \sum_{j=1}^{\infty}|\xi_{mj}-\xi_{nj}|^{p}=\|x_m-x_n\|_{p}^{\,p}.
\]
Also ist $(\xi_{nk})_n$ eine Cauchyfolge in $\mathbb{K}$; da $\mathbb{K}$ vollständig ist, existiert
\[
\xi_k:=\lim_{n\to\infty}\xi_{nk}\in\mathbb{K}.
\]
Definiere $x=(\xi_k)_{k\in\mathbb{N}}$. Für jedes $\ell\in\mathbb{N}$ folgt aus Grenzübergang und Monotonie
\[
\sum_{k=1}^{\ell}|\xi_k|^{p}
=\lim_{n\to\infty}\sum_{k=1}^{\ell}|\xi_{nk}|^{p}
\le \limsup_{n\to\infty}\|x_n\|_{p}^{\,p}<\infty,
\]
wobei die Beschränktheit aus der Cauchy-Eigenschaft von $(x_n)$ folgt (vgl.\ Bem.~2.31(b)). 
Damit ist $\sum_{k=1}^{\infty}|\xi_k|^{p}<\infty$, also $x\in\ell^{p}$.

Weiter sei
\[
c_n:=\sup_{m\ge n}\|x_m-x_n\|_{p}^{\,p}\xrightarrow[n\to\infty]{}0.
\]
Für beliebige $n,\ell\in\mathbb{N}$ gilt, unter Benutzung von $\xi_k=\lim_{m\to\infty}\xi_{mk}$ und Fatou/Monotonie,
\[
\sum_{k=1}^{\ell}|\xi_k-\xi_{nk}|^{p}
=\lim_{m\to\infty}\sum_{k=1}^{\ell}|\xi_{mk}-\xi_{nk}|^{p}
\le \sup_{m\ge n}\sum_{k=1}^{\infty}|\xi_{mk}-\xi_{nk}|^{p}
= c_n.
\]
Lässt man $\ell\to\infty$ folgen (Monotone Konvergenz), erhält man $\|x-x_n\|_{p}^{\,p}\le c_n$ und daher
$\|x-x_n\|_{p}\to0$. Also konvergiert $(x_n)$ in $\ell^{p}$ gegen $x$.

\medskip
\textbf{Fall $p=\infty$.}
Sei $(x_n)_n$ Cauchy in $\ell^{\infty}$, $x_n=(\xi_{nk})_{k}$. Dann gibt es $M$ mit
$\|x_n\|_{\infty}\le M$ für alle $n$, also $|\xi_{nk}|\le M$ für alle $n,k$.
Ferner: Für jedes feste $k$ ist $(\xi_{nk})_n$ Cauchy in $\mathbb{K}$ (da
$|\xi_{mk}-\xi_{nk}|\le \|x_m-x_n\|_{\infty}$) und konvergiert gegen ein $\xi_k\in\mathbb{K}$.
Setze $x=(\xi_k)_k$. Dann gilt
\[
\|x\|_{\infty}=\sup_{k}|\xi_k|\le \sup_{k}\liminf_{n}|\xi_{nk}|
\le \liminf_{n}\sup_{k}|\xi_{nk}|=\liminf_{n}\|x_n\|_{\infty}\le M,
\]
also $x\in\ell^{\infty}$. Schließlich
\[
\|x-x_n\|_{\infty}=\sup_{k}|\xi_k-\xi_{nk}|
=\sup_{k}\lim_{m\to\infty}|\xi_{mk}-\xi_{nk}|
\le \limsup_{m\to\infty}\|x_m-x_n\|_{\infty}\xrightarrow[n\to\infty]{}0,
\]
da $(x_n)$ Cauchy ist. Somit konvergiert $(x_n)$ gegen $x$ in $\ell^{\infty}$.

In beiden Fällen ist $\ell^{p}$ vollständig.
\end{proof}
\end{EXA}





(e) 


Ist $M$ eine beliebige Menge, so ist $\ell^{\infty}(M)$ ein Banachraum (Übung).



\begin{EXA}{FA-A25-01-60}{Vollständigkeit von $\ell^{\infty}(M)$}
Sei $M$ eine beliebige Menge und
\[
\ell^{\infty}(M):=\{f:M\to\mathbb{K}\mid \|f\|_{\infty}:=\sup_{x\in M}|f(x)|<\infty\}.
\]
Dann ist $(\ell^{\infty}(M),\|\cdot\|_{\infty})$ ein Banachraum.

\begin{proof}
Sei $(f_n)_{n}$ eine Cauchyfolge in $(\ell^{\infty}(M),\|\cdot\|_{\infty})$. Dann gilt:
\[
\forall \varepsilon>0\ \exists N\ \forall m,n\ge N:\ \sup_{x\in M}|f_m(x)-f_n(x)|<\varepsilon.
\]
Insbesondere ist für jedes feste $x\in M$ die skalare Folge $\bigl(f_n(x)\bigr)_n$ Cauchy in $\mathbb{K}$ und besitzt einen Grenzwert; setze
\[
f(x):=\lim_{n\to\infty} f_n(x)\quad (x\in M).
\]
Wir zeigen $f\in\ell^{\infty}(M)$ und $f_n\to f$ in $\|\cdot\|_{\infty}$.  
Aus der Cauchy-Eigenschaft folgt die Gleichmäßigkeitsbeschränktheit: Es gibt $C>0$ mit $\|f_n\|_{\infty}\le C$ für alle $n$. Für jedes $x$ erhält man dann $|f(x)|=\lim_{n}|f_n(x)|\le C$, also $\|f\|_{\infty}\le C<\infty$.

Für die Konvergenz: Sei $\varepsilon>0$. Wähle $N$ wie oben. Für $n\ge N$ und jedes $x\in M$ gilt
\[
|f(x)-f_n(x)|=\lim_{m\to\infty}|f_m(x)-f_n(x)|\le \sup_{m\ge N}\sup_{y\in M}|f_m(y)-f_n(y)|<\varepsilon,
\]
folglich $\|f-f_n\|_{\infty}<\varepsilon$. Also $f_n\to f$ in $\ell^{\infty}(M)$, und der Raum ist vollständig.
\end{proof}
\end{EXA}


\begin{EXA}{FA-A25-01-61}{$C_{\mathrm{b}}(X)$ ist Banachraum}
Sei $(X,d)$ ein metrischer Raum und $C_{\mathrm{b}}(X):=\{f\in C(X)\mid \|f\|_\infty<\infty\}$ mit der Supremumsnorm $\|f\|_\infty:=\sup_{x\in X}|f(x)|$. Dan ist $C_{\mathrm{b}}(X)$ ein abgeschlossener Unterraum von $\ell^\infty(X)$ und damit ein Banachraum.

\emph{Beweis.}  
Nach Bemerkung 2.23(d) ist $C_{\mathrm{b}}(X)$ ein linearer Unterraum von $\ell^\infty(X)$ und unter gleichmäßiger Konvergenz abgeschlossen: Konvergiert $(f_n)\subset C_{\mathrm{b}}(X)$ gleichmäßig gegen $f$, so ist $f$ beschränkt und stetig, also $f\in C_{\mathrm{b}}(X)$. 

Mit Theorem~\ref{ANA-L02-38} ist $\ell^\infty(X)$ ein Banachraum. Nach Satz~2.32(a) ist jeder abgeschlossene Unterraum eines vollständigen metrischen Raums vollständig. Folglich ist $C_{\mathrm{b}}(X)$ vollständig und damit (mit $\|\cdot\|_\infty$) ein Banachraum. \qed
\end{EXA}



(f) Nach Bemerkung 2.23(d) ist der Unterraum $C_{\mathrm{b}}(X)$ abgeschlossen in $\ell^{\infty}(X)$ und damit ein Banachraum nach (e) und Satz 2.32(a).




(g) $\ell^{2}$ ist ein Hilbertraum mit dem Skalarprodukt $\langle\cdot, \cdot\rangle$ aus Beispiel 2.13.\\



\begin{EXA}{FA-A25-01-62}{$\ell^{2}$ ist ein Hilbertraum}
Definiere für $x=(x_k)_{k\in\mathbb{N}},\,y=(y_k)_{k\in\mathbb{N}}\in\ell^{2}$ das Skalarprodukt
\[
\langle x,y\rangle := \sum_{k=1}^{\infty} x_k\,\overline{y_k}.
\]
Dann gilt:
\begin{enumerate}
\item[(i)] Die Reihe ist absolut konvergent (Cauchy–Schwarz) und $\langle\cdot,\cdot\rangle$ ist eine hermitesche Form im Sinn von Def.~2.10: linear im ersten Argument, $\langle x,y\rangle=\overline{\langle y,x\rangle}$, sowie positiv definit, d.\,h. $\langle x,x\rangle>0$ für $x\neq 0$.
\item[(ii)] Die durch $\|x\|:=\sqrt{\langle x,x\rangle}$ induzierte Norm stimmt mit der $\ell^{2}$-Norm überein:
\[
\|x\|=\left(\sum_{k=1}^{\infty}|x_k|^{2}\right)^{1/2}=\|x\|_{2}.
\]
\item[(iii)] $(\ell^{2},\|\cdot\|_{2})$ ist vollständig (vgl.\ Theorem~\ref{ANA-L02-37} für $p=2$).
\end{enumerate}
Aus (i)–(iii) folgt, dass $(\ell^{2},\langle\cdot,\cdot\rangle)$ ein Hilbertraum ist. 
Insbesondere bildet die Standardbasis $(e_k)_{k\in\mathbb{N}}$ mit $e_k=(\delta_{jk})_{j}$ ein Orthonormalsystem.
\end{EXA}




2.35 Satz. 

\begin{EXA}{FA-A25-01-63}{Raum der beschränkten linearen Operatoren zwischen Banachräumen ist vollständig}
Seien $E, F$ normierte Räume. Ist $F$ ein Banachraum, so ist auch $\mathcal{L}(E, F)$ ein Banachraum. Insbesondere ist der Dualraum $E^{\prime}$ eines normierten Raumes $E$ stets ein Banachraum.

\begin{proof} Sei $\left(T_{n}\right)_{n}$ eine Cauchyfolge in $\mathcal{L}(E, F)$ und $c_{n}:=\sup _{m \geq n}\left\|T_{m}-T_{n}\right\|$ für $n \in \mathbb{N}$. Dann gilt $\lim _{n \rightarrow \infty} c_{n}=0$. Für alle $x \in E$ folgt
$$
\sup _{m \geq n}\left\|T_{m} x-T_{n} x\right\| \leq \sup _{m \geq n}\left\|T_{m}-T_{n}\right\|\|x\|=c_{n}\|x\| \rightarrow 0 \quad \text { für } n \rightarrow \infty
$$
d.h. $\left(T_{n} x\right)_{n}$ ist eine Cauchyfolge in $F$. Aus der Voraussetzung folgt somit die Existenz von
$$
T x:=\lim _{n \rightarrow \infty} T_{n} x \in F \quad \text { für alle } x \in E .
$$
Aus der Linearität der Abbildungen $T_{n}$ und der Stetigkeit der linearen Operationen in $F$ folgt direkt, dass auch die Zuordnung $x \mapsto T x$ linear ist. Ferner ist $\|\cdot\|$ stetig und daher
$$
\|T x\|=\lim _{n \rightarrow \infty}\left\|T_{n} x\right\| \leq\left(\sup _{n \in \mathbb{N}}\left\|T_{n}\right\|\right)\|x\| \quad \text { für alle } x \in E,
$$
d.h. $T$ ist stetig. Man beachte dabei, dass $\left(\left\|T_{n}\right\|\right)_{n}$ wegen Bemerkung 2.31(b) beschränkt ist. Schließlich ist
$$
\begin{aligned}
\left\|\left(T-T_{n}\right) x\right\|=\lim _{m \rightarrow \infty}\left\|\left(T_{m}-T_{n}\right) x\right\| & \\
& \leq\left(\sup _{m \geq n}\left\|T_{m}-T_{n}\right\|\right)\|x\|=c_{n}\|x\| \quad \text { für alle } x \in E
\end{aligned}
$$
d.h. $\left\|T-T_{n}\right\| \leq c_{n} \rightarrow 0$ für $n \rightarrow \infty$. Dies zeigt die Vollständigkeit von $\mathcal{L}(E,F)$. \end{proof}
\end{EXA}



2.36 Definition. 

Seien $(E,\|\cdot\|)$ ein normierter Raum und $\left(x_{n}\right)_{n} \subseteq E$ eine Folge.\\
(a) Wenn $S:=\lim _{n \rightarrow \infty} \sum_{k=1}^{n} x_{k} \in E$ existiert, so nennen wir die Reihe $\sum_{k} x_{k}$ konvergent und setzen $\sum_{k=1}^{\infty} x_{k}:=S$.\\
(b) Die Reihe $\sum_{k} x_{k}$ heißt absolut konvergent, wenn $\sum_{k=1}^{\infty}\left\|x_{k}\right\|<\infty$ ist.\\


\begin{DEF}{FA-A25-01-64}{Konvergenz und absolute Konvergenz von Reihen in normierten Räumen}
Sei $(E,\|\cdot\|)$ ein normierter Raum und $(x_n)_{n\in\mathbb{N}}\subseteq E$ eine Folge.

\begin{enumerate}
\item[(a)] \textbf{(Konvergenz einer Reihe)}  
Die Reihe $\sum_{k=1}^{\infty} x_k$ heißt \emph{konvergent}, falls die Folge der Partialsummen
\[
S_n := \sum_{k=1}^{n} x_k
\]
in $E$ konvergiert, d.\,h.
\[
S := \lim_{n\to\infty} S_n \in E \quad \text{existiert}.
\]
In diesem Fall setzen wir
\[
\sum_{k=1}^{\infty} x_k := S.
\]

\item[(b)] \textbf{(Absolute Konvergenz)}  
Die Reihe $\sum_{k=1}^{\infty} x_k$ heißt \emph{absolut konvergent}, wenn
\[
\sum_{k=1}^{\infty} \|x_k\| < \infty
\]
gilt.
\end{enumerate}
\end{DEF}



2.37 Satz. 

\begin{PROP}{FA-A25-01-65}{Charakterisierung von Banachräumen durch absolute Konvergenz von Reihen}
Sei $(E,\|\cdot\|)$ ein normierter Raum. Dann gilt: $E$ ist genau dann ein Banachraum, wenn jede absolut konvergente Reihe in $E$ auch konvergent ist.
\end{PROP}


Beweis. 

\begin{PROOF}{FA-A25-01-66}{P: Charakterisierung von Banachräumen durch absolute Konvergenz von Reihen}
„ $\Longrightarrow$ ": Genau wie für $E=\mathbb{R}$ bzw. $\mathbb{C}$ („Cauchykriterium"), man muss nur alle Beträge durch Normen ersetzen.

„$\Longleftarrow$": Sei $\left(x_{n}\right)_{n} \subseteq E$ eine Cauchyfolge. Dann existiert für alle $k \in \mathbb{N}$ ein $\tau_{k} \in \mathbb{N}$ mit
$$
\sup _{m \geq n}\left\|x_{m}-x_{n}\right\|<2^{-k} \quad \text { für } n \geq \tau_{k} \text {. }
$$
Wähle nun induktiv $n_{1} \geq \tau_{1}$ und $n_{k+1} \geq \max \left\{\tau_{k+1}, n_{k}+1\right\}$ für $k \in \mathbb{N}$. Dann gilt
$$
\left\|x_{n_{k+1}}-x_{n_{k}}\right\| \leq 2^{-k} \quad \text { für alle } k,
$$
also insbesondere
$$
\sum_{k=1}^{\infty}\left\|x_{n_{k+1}}-x_{n_{k}}\right\| \leq \sum_{k=1}^{\infty} 2^{-k}=1
$$
dh. die Reihe $\sum_{k}\left(x_{n_{k+1}}-x_{n_{k}}\right)$ konvergiert absolut. Nach Voraussetzung konvergiert sie damit auch, d.h. es existiert
$$
x:=\lim _{k \rightarrow \infty} x_{n_{k}}=x_{n_{1}}+\lim _{k \rightarrow \infty} \sum_{j=1}^{k}\left(x_{n_{j+1}}-x_{n_{j}}\right) \quad \in E
$$
Nach Bemerkung 2.31(d) konvergiert die Folge damit insgesamt gegen $x$, was zu zeigen war.
\end{PROOF}




\pagebreak


\subsection*{Kompaktheit}
Im Folgenden sei ( $X, d$ ) ein metrischer Raum.

2.38 Definition und Bemerkung. (a) ( $X, d$ ) heißt kompakt, falls gilt: Ist $\mathcal{I}$ eine beliebige Indexmenge und sind $U_{i} \subseteq X, i \in \mathcal{I}$, offene Mengen mit $X=\bigcup_{i \in \mathcal{I}} U_{i}$, so existiert eine endliche Teilmenge $\mathcal{J} \subseteq \mathcal{I}$ mit $X=\bigcup_{i \in \mathcal{J}} U_{i}$. Mit anderen Worten: Jede offene Überdeckung von $X$ lässt sich auf eine endliche Teilüberdeckung reduzieren.\\
(b) ( $X, d$ ) heißt präkompakt, wenn für alle $\varepsilon>0$ eine endliche Teilmenge $M \subseteq X$ existiert mit $X=B_{\varepsilon}(M)$ (hier könnte man äquivalent auch $U_{\varepsilon}(M)$ nehmen; überlegen!).\\
(c) $A \subseteq X$ heißt kompakt (bzw. präkompakt), wenn der metrische Teilraum $(A, d)$ kompakt (bzw. präkompakt) ist. Somit gilt:

 $A$ ist kompakt genau dann, wenn für jedes Familie $U_{i} \subseteq X, i \in \mathcal{I}$ offener Mengen mit $A \subseteq \bigcup_{i \in \mathcal{I}} U_{i}$ eine endliche Teilmenge $\mathcal{J} \subseteq \mathcal{I}$ existiert mit $A \subseteq$ $\bigcup_{i \in \mathcal{J}} U_{i}$.
 
 $A$ ist präkompakt genau dann für jedes $\varepsilon>0$ eine endliche Teilmenge $M \subseteq A$ existiert mit $A \subseteq B_{\varepsilon}(M)$ (wobei $B_{\varepsilon}(M)$ hier die Umgebung von $M$ in $X$ bezeichne).\\

(d) $A \subseteq X$ heißt relativ kompakt, wenn $\bar{A}$ kompakt ist.\\



\begin{DEF}{FA-A25-01-67}{Kompaktheit, Präkompaktheit und relative Kompaktheit in metrischen Räumen}
Sei $(X,d)$ ein metrischer Raum.

\begin{enumerate}
\item[(a)] \textbf{(Kompaktheit)}  
$(X,d)$ heißt \emph{kompakt}, wenn gilt:  
Ist $(U_i)_{i\in\mathcal{I}}$ eine Familie offener Mengen $U_i\subseteq X$ mit
\[
X = \bigcup_{i\in\mathcal{I}} U_i,
\]
so existiert eine endliche Teilmenge $\mathcal{J}\subseteq\mathcal{I}$, sodass bereits
\[
X = \bigcup_{i\in\mathcal{J}} U_i.
\]
Das heißt: jede offene Überdeckung von $X$ besitzt eine endliche Teilüberdeckung.

\item[(b)] \textbf{(Präkompaktheit)}  
$(X,d)$ heißt \emph{präkompakt}, wenn für jedes $\varepsilon>0$ eine endliche Teilmenge $M\subseteq X$ existiert mit
\[
X = B_{\varepsilon}(M) := \bigcup_{x\in M} B_{\varepsilon}(x).
\]
(Äquivalent: $X=U_{\varepsilon}(M)$ mit offenen $\varepsilon$-Kugeln.)

\item[(c)] \textbf{(Kompakte bzw. präkompakte Teilmengen)}  
Eine Teilmenge $A\subseteq X$ heißt \emph{kompakt} (bzw. \emph{präkompakt}), wenn der metrische Teilraum $(A,d)$ kompakt (bzw. präkompakt) ist.  
Äquivalent gilt:
\begin{itemize}
  \item $A$ ist kompakt $\Longleftrightarrow$  
  jede Familie $(U_i)_{i\in\mathcal{I}}$ offener Mengen in $X$ mit $A\subseteq\bigcup_{i\in\mathcal{I}}U_i$ besitzt eine endliche Teilfamilie $\mathcal{J}\subseteq\mathcal{I}$ mit $A\subseteq\bigcup_{i\in\mathcal{J}}U_i$;
  \item $A$ ist präkompakt $\Longleftrightarrow$  
  für jedes $\varepsilon>0$ existiert eine endliche Teilmenge $M\subseteq A$ mit $A\subseteq B_{\varepsilon}(M)$.
\end{itemize}

\item[(d)] \textbf{(Relative Kompaktheit)}  
Eine Menge $A\subseteq X$ heißt \emph{relativ kompakt}, wenn ihr Abschluss $\bar{A}$ kompakt ist.
\end{enumerate}
\end{DEF}



2.39 Satz (Cantor). 

\begin{THEO}{FA-A25-01-68}{Cantor}
Sei  ($X,d$) vollständig, und seien $A_{n} \subseteq X, n \in \mathbb{N}$, abgeschlossene, nichtleere Mengen mit $A_{n+1} \subseteq A_{n}$ für alle $n$ und $\lim _{n \rightarrow \infty} \operatorname{diam} A_{n}=0$. Dann ist $A:=$ $\bigcap_{n \in \mathbb{N}} A_{n}$ nichtleer.
\end{THEO}


Beweis. 

\begin{PROOF}{FA-A25-01-69}{P: Cantor}
Wähle $x_{n} \in A_{n}$ für $n \in \mathbb{N}$. Nach Voraussetzung ist dann
$$
x_{m} \in A_{n} \quad \text { für } m \geq n .
$$
Also folgt
$$
\sup _{m \geq n} d\left(x_{m}, x_{n}\right) \leq \operatorname{diam} A_{n} \rightarrow 0 \quad \text { für } n \rightarrow \infty,
$$
d.h. $\left(x_{n}\right)_{n}$ ist eine Cauchyfolge in $X$. Da ( $X, d$ ) vollständig ist, existiert $x:=\lim _{n \rightarrow \infty} x_{n}$. Wegen (2.4) ist $x \in \bar{A}_{n}=A_{n}$ für alle $n \in \mathbb{N}$; also $x \in A$.

Nachbemerkung: Es gilt sogar $A=\{x\}$, da $\operatorname{diam} A \leq \lim _{n \rightarrow \infty} \operatorname{diam} A_{n}=0$ ist.
\end{PROOF}


2.40 Satz. 


Für einen metrischen Raum ( $X, d$ ) sind folgende Eigenschaften äquivalent.
(i) $X$ ist kompakt.\\
(ii) Jede Folge $\left(x_{k}\right)_{k}$ in $X$ besitzt eine in $X$ konvergente Teilfolge.\\
(iii) $X$ ist präkompakt und vollständig.



\begin{THEO}{FA-A25-01-70}{Äquivalenz von Kompaktheit, Folgenkompaktheit und Präkompaktheit+Vollständigkeit}
Für einen metrischen Raum $(X,d)$ sind äquivalent:
\begin{enumerate}
\item[(i)] $X$ ist kompakt.
\item[(ii)] Jede Folge $(x_k)_{k\in\mathbb{N}}$ in $X$ besitzt eine in $X$ konvergente Teilfolge (d.\,h.\ $X$ ist folgenkompakt/sequentiell kompakt).
\item[(iii)] $X$ ist präkompakt (totale Beschränktheit) und vollständig.
\end{enumerate}
\end{THEO}


Beweis. 

\begin{PROOF}{FA-A25-01-71}{P: Äquivalenz von Kompaktheit, Folgenkompaktheit und Präkompaktheit+Vollständigkeit}
„(i) $\Longrightarrow$ (ii)": Sei $X$ kompakt, und sei $\left(x_{k}\right)_{k}$ eine Folge in $X$. Ohne Einschränkung sei $M:=\left\{x_{k} \mid k \in \mathbb{N}\right\}$ eine unendliche Menge. Angenommen, kein $a \in X$ ist Grenzwert einer Teilfolge von $\left(x_{k}\right)_{k}$. Dann existiert für jedes $a \in X$ ein $\varepsilon_{a}>0$ so, dass $U_{\varepsilon_{a}}(a) \cap M$ endlich ist. Da $X=\bigcup_{a \in X} U_{\varepsilon_{a}}(a)$ kompakt ist, existieren $a_{1}, \ldots, a_{n} \in X$ mit $X=$ $\bigcup_{i=1}^{n} U_{\varepsilon_{a_{i}}}\left(a_{i}\right)$. Es folgt: 
$$M=M \cap X=\bigcup_{i=1}^{n} M \cap U_{\varepsilon_{a_{i}}}\left(a_{i}\right)$$ 
ist endlich. Widerspruch.

„(ii) $\Longrightarrow$ (iii)": Sei $\left(x_{k}\right)_{k}$ eine Cauchyfolge in $X$. Diese besitzt nach Voraussetzung eine in $X$ konvergente Teilfolge. Nach Bemerkung 2.31(c) ist dann aber schon $\left(x_{k}\right)_{k}$ selbst konvergent. Demnach ist $X$ vollständig.

Wir zeigen nun die Präkompaktheit von $X$. Sei dazu $\varepsilon>0$. Für einen Widerspruchsbeweis nehmen wir an, es gäbe keine endliche Menge $M \subseteq X$ mit $X=B_{\varepsilon}(M)$. Wir wählen dann $a_{1} \in X$ beliebig.
$$
\begin{aligned}
& \text { Da } X \nsubseteq U_{\varepsilon}\left(a_{1}\right) \text {, existiert } a_{2} \in X \backslash U_{\varepsilon}\left(a_{1}\right) ; \\
& \text { da } X \nsubseteq U_{\varepsilon}\left(a_{1}\right) \cup U_{\varepsilon}\left(a_{2}\right) \text {, existiert } a_{3} \in X \backslash\left(U_{\varepsilon}\left(a_{1}\right) \cap U_{\varepsilon}\left(a_{2}\right)\right) ; \\
& \quad \ldots \quad \quad \ldots
\end{aligned}
$$
$$
\text { da } X \nsubseteq \bigcup_{j=1}^{n} U_{\varepsilon}\left(a_{j}\right), \text { existiert } a_{n+1} \in X \backslash \bigcup_{j=1}^{n} U_{\varepsilon}\left(a_{j}\right) .
$$
Diese Konstruktion definiert induktiv eine Folge $\left(a_{n}\right)_{n}$ in $X$ derart, dass
$$
d\left(a_{n}, a_{m}\right) \geq \varepsilon \quad \text { für alle } m, n \in \mathbb{N} \text { mit } m \neq n
$$
gilt. Nach Voraussetzung muss $\left(a_{n}\right)_{n}$ aber eine konvergente Teilfolge $\left(a_{n_{j}}\right)_{j}$ besitzen. Diese ist dann eine Cauchyfolge, im Widerspruch zu (2.5). Es folgt die Präkompaktheit von $X$. 

„(iii) $\Longrightarrow$ (i)": Seien $U_{i} \subseteq X$ für $i \in \mathcal{I}$ offene Teilmengen mit
$$
X=\bigcup_{i \in \mathcal{I}} U_{i}
$$
Da $X$ präkompakt ist, existiert für jedes $n \in \mathbb{N}$ eine endliche Menge $M_{n} \subseteq X$ mit
$$
X=B_{\frac{1}{n}}\left(M_{n}\right)=\bigcup_{x \in M_{n}} B_{\frac{1}{n}}(x)
$$
Sei nun angenommen, dass $X$ nicht von endlich vielen der Mengen $U_{i}$ überdeckt wird. Da $M_{1}$ endlich ist, existiert $x_{1} \in M_{1}$ so, dass $A_{1}:=B_{1}\left(x_{1}\right)$ nicht von endlich vielen $U_{i}$ überdeckt wird. Da ferner $M_{2}$ endlich ist und $A_{1}=\bigcup_{x \in M_{2}} A_{1} \cap B_{\frac{1}{2}}(x)$ gilt, existiert $x_{2} \in M_{2}$ so, dass $A_{2}:=A_{1} \cap B_{\frac{1}{2}}\left(x_{2}\right)$ nicht von endlich vielen $U_{i}$ überdeckt wird. Induktiv findet man: Es gibt $x_{n} \in M_{n}$, so dass $A_{n}:=A_{n-1} \cap B_{\frac{1}{n}}\left(x_{n}\right)$ nicht von endlich vielen $U_{i}$ überdeckt wird.

Dies liefert eine Folge abgeschlossener Mengen mit $A_{n+1} \subseteq A_{n}$ und diam $A_{n} \leq \frac{2}{n}$ für alle $n \in \mathbb{N}$. Da $X$ vollständig ist, folgt mit Satz 2.39, dass ein $a \in \bigcap_{n \in \mathbb{N}} A_{n}$ existiert. Ferner existiert $i \in \mathcal{I}$ mit $a \in U_{i}$. Da $U_{i}$ offen ist, existiert auch ein $\delta>0$ mit $B_{\delta}(a) \subseteq U_{i}$. Für $n>\frac{2}{\delta}$ ist $\operatorname{diam} A_{n}<\delta$ und $a \in A_{n}$, also $A_{n} \subseteq B_{\delta}(a) \subseteq U_{i}$ im Widerspruch zu (2.7). Der Widerspruch zeigt, dass eine endliche Teilmenge $J \subseteq \mathcal{I}$ existiert mit $X=\bigcup_{i \in J} U_{i}$. Es folgt die Kompaktheit von $X$.
\end{PROOF}


2.41 Bemerkung. 


(a) Ist $X$ kompakt und $A \subseteq X$ abgeschlossen, so ist $A$ kompakt.\\


\begin{LEM}{FA-A25-01-72}{Abgeschlossene Teilmengen kompakter metrischer Räume sind kompakt}
Sei $(X,d)$ ein kompakter metrischer Raum und $A \subseteq X$ eine abgeschlossene Teilmenge.  
Dann ist auch $A$ kompakt.
\end{LEM}


\begin{PROOF}{FA-A25-01-73}{P: Abgeschlossene Teilmengen kompakter metrischer Räume sind kompakt}
Sei $(U_i)_{i \in \mathcal{I}}$ eine offene Überdeckung von $A$, d.\,h.
\[
A \subseteq \bigcup_{i \in \mathcal{I}} U_i.
\]
Da $A$ abgeschlossen ist, ist $X \setminus A$ offen in $X$. Somit ist
\[
(X \setminus A) \cup \bigcup_{i \in \mathcal{I}} U_i
\]
eine offene Überdeckung von $X$. Wegen der Kompaktheit von $X$ existiert eine endliche Teilmenge $\mathcal{J} \subseteq \mathcal{I}$ mit
\[
X = (X \setminus A) \cup \bigcup_{i \in \mathcal{J}} U_i.
\]
Durch Schneiden mit $A$ folgt
\[
A = \bigcup_{i \in \mathcal{J}} (A \cap U_i) = \bigcup_{i \in \mathcal{J}} U_i,
\]
da $A \cap (X \setminus A) = \varnothing$. Somit besitzt jede offene Überdeckung von $A$ eine endliche Teilüberdeckung, also ist $A$ kompakt.
\end{PROOF}




(b) Ist $X$ kompakt, $(Y, \tilde{d})$ ein weiterer metrischer Raum und $f: X \rightarrow Y$ stetig, so ist auch $f(X) \subseteq Y$ kompakt, und $f$ ist gleichmäßig stetig.



\begin{LEM}{FA-A25-01-74}{Stetiges Bild eines kompakten Raums ist kompakt}
Sei $(X,d)$ ein kompakter metrischer Raum, $(Y,\tilde d)$ ein metrischer Raum und
$f:X\to Y$ stetig. Dann ist $f(X)$ kompakt (als Teilraum von $Y$).
\end{LEM}




\begin{PROOF}{FA-A25-01-75}{P: Stetiges Bild eines kompakten Raums ist kompakt}
Sei $(V_i)_{i\in\mathcal I}$ eine offene Überdeckung von $f(X)$ in $Y$. Dann ist
$(f^{-1}(V_i))_{i\in\mathcal I}$ eine offene Überdeckung von $X$, denn $f$ ist stetig.
Wegen der Kompaktheit von $X$ existiert eine endliche Teilfamilie $\mathcal J\subset\mathcal I$
mit $X=\bigcup_{i\in\mathcal J} f^{-1}(V_i)$. Anwenden von $f$ liefert
\[
f(X)=\bigcup_{i\in\mathcal J} V_i,
\]
also besitzt jede offene Überdeckung von $f(X)$ eine endliche Teilüberdeckung. Damit ist
$f(X)$ kompakt.
\end{PROOF}





\begin{PROP}{FA-A25-01-76}{Heine--Cantor}
Sei $(X,d)$ ein kompakter metrischer Raum, $(Y,\tilde d)$ ein metrischer Raum und
$f:X\to Y$ stetig. Dann ist $f$ gleichmäßig stetig.
\end{PROP}


\begin{PROOF}{FA-A25-01-77}{P: Heine--Cantor}
Angenommen, $f$ sei \emph{nicht} gleichmäßig stetig. Dann existiert $\varepsilon_0>0$
und für jedes $n\in\mathbb N$ Punkte $x_n,y_n\in X$ mit
$d(x_n,y_n)<\tfrac1n$ und $\tilde d\bigl(f(x_n),f(y_n)\bigr)\ge\varepsilon_0$.
Da $X$ kompakt ist, besitzt $(x_n)$ eine konvergente Teilfolge $(x_{n_k})$
mit Grenzwert $x\in X$. Wegen $d(x_{n_k},y_{n_k})\to0$ folgt auch $y_{n_k}\to x$.
Aus der Stetigkeit von $f$ folgt $f(x_{n_k})\to f(x)$ und $f(y_{n_k})\to f(x)$,
also $\tilde d\bigl(f(x_{n_k}),f(y_{n_k})\bigr)\to0$ im Widerspruch zu
$\tilde d\bigl(f(x_{n_k}),f(y_{n_k})\bigr)\ge \varepsilon_0$. Also ist $f$
gleichmäßig stetig.
\end{PROOF}


(c) (Satz von Heine-Borel) 


\begin{PROP}{FA-A25-01-78}{Heine-Borel}
Sei $E$ ein endlichdimensionaler normierter Raum und $A \subseteq E$. Dann ist $A$ genau dann kompakt, wenn $A$ beschränkt und abgeschlossen ist.
\end{PROP}

Beweis. 


\begin{PROOF}{FA-A25-01-79}{P: Heine-Borel}
Bekannt aus der Analysis II.
\end{PROOF}




2.42 Bemerkung.

(a) Ist $X$ präkompakt, so ist $X$ offensichtlich beschränkt. Ferner ist $X$ dann auch separabel, denn für jedes $n \in \mathbb{N}$ existiert eine endliche Teilmenge $M_{n} \subseteq X$ mit $X=B_{\frac{1}{n}}\left(M_{n}\right)$. Somit ist die Menge $M:=\bigcup_{n \in \mathbb{N}} M_{n}$ abzählbar und dicht in $X$.\\


\begin{PROP}{FA-A25-01-80}{Präkompakte metrische Räume sind beschränkt und separabel}
Sei $(X,d)$ ein metrischer Raum. Ist $X$ präkompakt (totale Beschränktheit), so gilt:
\begin{enumerate}
\item[(i)] $X$ ist beschränkt.
\item[(ii)] $X$ ist separabel.
\end{enumerate}
\end{PROP}


\begin{PROOF}{FA-A25-01-81}{P: Präkompakte metrische Räume sind beschränkt und separabel}
(i) Beschränktheit. Präkompaktheit bedeutet: Für jedes $\varepsilon>0$ gibt es eine \emph{endliche} Menge
$M_\varepsilon\subset X$ mit $X\subset \bigcup_{a\in M_\varepsilon} B_\varepsilon(a)$.
Wähle $\varepsilon=1$ und schreibe $M_1=\{a_1,\dots,a_N\}$.
Für beliebige $x,y\in X$ existieren daher Indizes $i,j$ mit $d(x,a_i)\le 1$ und $d(y,a_j)\le 1$. Dann
\[
d(x,y)\le d(x,a_i)+d(a_i,a_j)+d(a_j,y)\le 1+\max_{1\le i,j\le N} d(a_i,a_j)+1.
\]
Damit ist $\operatorname{diam}X\le 2+\max_{i,j} d(a_i,a_j)<\infty$, also $X$ beschränkt.

(ii) Separabilität. Für jedes $n\in\mathbb N$ existiert eine endliche Menge $M_n\subset X$
mit $X\subset \bigcup_{a\in M_n} B_{1/n}(a)$. Setze
\[
M:=\bigcup_{n=1}^\infty M_n.
\]
Dann ist $M$ abzählbar (als abzählbare Vereinigung endlicher Mengen). Wir zeigen, dass $M$ dicht in $X$ ist.
Sei $x\in X$ und $\varepsilon>0$. Wähle $n\in\mathbb N$ mit $1/n<\varepsilon$.
Dann gibt es $a\in M_n\subset M$ mit $d(x,a)<1/n<\varepsilon$.
Also liegt jede $\varepsilon$-Kugel um $x$ im Abschluss von $M$, d.\,h.\ $\overline{M}=X$.
Damit ist $X$ separabel.
\end{PROOF}



(b) Für $A \subseteq X$ sind äquivalent:

\begin{itemize}
  \item $A$ ist präkompakt
  \item Für jedes $\varepsilon>0$ existiert eine präkompakte Menge $C \subseteq X$ mit $A \subseteq B_{\varepsilon}(C)$.
  \item Für jedes $\varepsilon>0$ existieren endlich viele Teilmengen $A_{1}, \ldots, A_{n}$ von $A$ mit $A \subseteq \bigcup_{i=1}^{n} A_{i}$ und diam $A_{i} \leq \varepsilon$ für $i=1, \ldots, n$.
  \item $\bar{A}$ ist präkompakt.
\end{itemize}

Beweis als Übung.\\



\begin{PROP}{FA-A25-01-82}{Äquivalenzen zur Präkompaktheit von metrischen Teilräumen}
Sei $(X,d)$ ein metrischer Raum und $A\subseteq X$. Die folgenden Aussagen sind äquivalent:
\begin{enumerate}
\item[(i)] $A$ ist präkompakt (totale Beschränktheit), d.\,h. für jedes $\varepsilon>0$ gibt es eine \emph{endliche}
Menge $M\subset X$ mit $A\subset \bigcup_{x\in M} B_\varepsilon(x)$.
\item[(ii)] Für jedes $\varepsilon>0$ existiert eine präkompakte Menge $C\subseteq X$ mit
$A\subseteq B_\varepsilon(C):=\{x\in X:\operatorname{dist}(x,C)<\varepsilon\}$.
\item[(iii)] Für jedes $\varepsilon>0$ gibt es endlich viele Teilmengen $A_1,\dots,A_n\subseteq A$ mit
\[
A\subseteq \bigcup_{i=1}^n A_i
\qquad\text{und}\qquad
\operatorname{diam}(A_i)\le \varepsilon\ \ (i=1,\dots,n).
\]
\item[(iv)] Der Abschluss $\overline{A}$ ist präkompakt.
\end{enumerate}
\end{PROP}


\begin{PROOF}{FA-A25-01-83}{P: Äquivalenzen zur Präkompaktheit von metrischen Teilräumen}
\noindent (i)$\Rightarrow$(ii): Sei $\varepsilon>0$. Wähle ein endliches $\varepsilon$-Netz $M\subset X$ für $A$, also
$A\subset \bigcup_{x\in M} B_\varepsilon(x)$. Setze $C:=M$. Endliche Mengen sind präkompakt, und
$A\subseteq B_\varepsilon(C)$.

\smallskip
\noindent (ii)$\Rightarrow$(iii): Sei $\varepsilon>0$. Wähle nach (ii) eine präkompakte Menge $C\subset X$ mit
$A\subset B_{\varepsilon/2}(C)$. Da $C$ präkompakt ist, existiert ein endliches $(\varepsilon/2)$-Netz
$\{c_1,\dots,c_n\}\subset X$ für $C$. Setze
\[
A_i:=A\cap B_{\varepsilon/2}(c_i)\quad (i=1,\dots,n).
\]
Dann überdecken die $A_i$ die Menge $A$, und für $x,y\in A_i$ gilt
$d(x,y)\le d(x,c_i)+d(c_i,y)\le \varepsilon$, also $\operatorname{diam}(A_i)\le \varepsilon$.

\smallskip
\noindent (iii)$\Rightarrow$(i): Sei $\varepsilon>0$. Wähle nach (iii) eine endliche Überdeckung $A\subset\bigcup_{i=1}^n A_i$
mit $\operatorname{diam}(A_i)\le \varepsilon$. Wähle in jedem nichtleeren $A_i$ einen Punkt $x_i\in A_i$.
Dann gilt für $x\in A_i$: $d(x,x_i)\le \operatorname{diam}(A_i)\le \varepsilon$, also
$A_i\subset B_\varepsilon(x_i)$ und damit $A\subset \bigcup_{i=1}^n B_\varepsilon(x_i)$. Also ist $A$ präkompakt.

\smallskip
\noindent (i)$\Rightarrow$(iv): Ist $A$ präkompakt, so ist $A$ total beschränkt. Der Abschluss einer total
beschränkten Menge ist wieder total beschränkt (denn jedes $\varepsilon$-Netz für $A$ ist auch
ein $\varepsilon$-Netz für $\overline{A}$). Also ist $\overline{A}$ präkompakt.

\smallskip
\noindent (iv)$\Rightarrow$(i): Ist $\overline{A}$ präkompakt, so besitzt $\overline{A}$ für jedes $\varepsilon>0$ ein endliches
$\varepsilon$-Netz; dieses überdeckt dann insbesondere $A\subset \overline{A}$. Also ist $A$ präkompakt.

\smallskip
Damit sind (i)–(iv) äquivalent.
\end{PROOF}


(c) Für $A \subseteq X$ gilt wegen Satz 2.32(a), Satz 2.40 und (b):

\begin{itemize}
  \item $A$ relativ kompakt $\Longrightarrow \bar{A}$ kompakt $\Longrightarrow \bar{A}$ präkompakt $\Longrightarrow A$ präkompakt;
  \item $(X, d)$ vollständig und $A$ präkompakt $\Longrightarrow \bar{A}$ vollständig und präkompakt $\Longrightarrow$ $\bar{A}$ kompakt $\Longrightarrow A$ relativ kompakt,\\
d.h. in einem vollständigen metrischen Raum sind Teilmengen genau dann präkompakt, wenn sie relativ kompakt sind.\\
\end{itemize}


\begin{PROP}{FA-A25-01-84}{Präkompakt vs. relativ kompakt; Äquivalenz in vollständigen Räumen}
Sei $(X,d)$ ein metrischer Raum und $A\subseteq X$.
\begin{enumerate}
\item[(i)] Es gilt stets die Implikationskette
\[
A\ \text{relativ kompakt}
\ \Longrightarrow\ 
\overline{A}\ \text{kompakt}
\ \Longrightarrow\
\overline{A}\ \text{präkompakt}
\ \Longrightarrow\
A\ \text{präkompakt}.
\]
\item[(ii)] Ist $(X,d)$ vollständig, so sind folgende Aussagen äquivalent:
\[
A\ \text{präkompakt}
\quad\Longleftrightarrow\quad
A\ \text{relativ kompakt}.
\]
\end{enumerate}
\end{PROP}


\begin{PROOF}{FA-A25-01-85}{P: Präkompakt vs. relativ kompakt; Äquivalenz in vollständigen Räumen}
\begin{description}
\item[(i)] Per Definition bedeutet „relativ kompakt“, dass $\overline{A}$ kompakt ist. Kompakte metrische Räume sind präkompakt, also ist $\overline{A}$ präkompakt. Aus $\overline{A}$ präkompakt folgt wegen $\overline{A}\supset A$ unmittelbar $A$ präkompakt (jedes $\varepsilon$-Netz für $\overline{A}$ überdeckt auch $A$).

\item[(ii) „$\Rightarrow$“] Sei $A$ präkompakt. Dann ist $\overline{A}$ ebenfalls präkompakt (Abschluss einer total beschränkten Menge ist total beschränkt) und zudem nach Satz~2.32(a) vollständig, da $\overline{A}$ in dem vollständigen Raum $X$ abgeschlossen ist. Mit Satz~2.40 (Charakterisierung der Kompaktheit als „präkompakt + vollständig“) folgt, dass $\overline{A}$ kompakt ist, also $A$ relativ kompakt.

\smallskip
\noindent „$\Leftarrow$“ Ist $A$ relativ kompakt, so ist $\overline{A}$ kompakt und damit präkompakt; daraus folgt wieder $A$ präkompakt. 
\end{description}
\end{PROOF}


2.43 Satz. 


\begin{PROP}{FA-A25-01-86}{Charakterisierung von Präkompaktheit durch durch Beschränktheit und $\epsilon$-Umgebgungsbedingung}
Sei $(E,\|\cdot\|)$ ein normierter Raum und $A \subseteq E$. Dann gilt: $A$ ist präkompakt genau dann, wenn $A$ beschränkt ist und für jedes $\varepsilon>0$ ein endlichdimensionaler Unterraum $F \subseteq E$ existiert mit $A \subseteq B_{\varepsilon}(F)$.
\end{PROP}


Beweis. 

\begin{PROOF}{FA-A25-01-87}{P: Charakterisierung von Präkompaktheit durch durch Beschränktheit und $\epsilon$-Umgebgungsbedingung}
„$\Longrightarrow$ ": Sei $A$ präkompakt. Wie bereits bemerkt, ist $A$ dann auch beschränkt. Sei ferner $\varepsilon>0$, und sei $M \subseteq A$ endlich mit $A \subseteq B_{\varepsilon}(M)$. Setze $F:=\operatorname{span} M$. Dann ist $\operatorname{dim} F<\infty$ und $A \subseteq B_{\varepsilon}(F)$.

„$\Longleftarrow$": Nach Voraussetzung ist $r:=\sup _{x \in A}\|x\|<\infty$. Sei $\varepsilon>0$, und sei $F \subseteq E$ ein endlichdimensionaler Unterraum mit $A \subseteq B_{\varepsilon}(F)$. Setze $C:=\{x \in F \mid\|x\| \leq r+\varepsilon\}$. Nach Bemerkung 2.41(c) ist $C$ dann kompakt, also auch präkompakt. Ferner existiert für jedes $x \in A$ ein $y \in F$ mit $\|x-y\|<\varepsilon$, also insbesondere $\|y\| \leq\|x\|+\varepsilon \leq r+\varepsilon$ und somit $y \in C$. Also ist $A \subseteq B_{\varepsilon}(C)$, und mit Bemerkung 2.42(b) folgt die Behauptung.
\end{PROOF}



2.44 Lemma (Rieszsches Lemma). 


\begin{LEM}{FA-A25-01-88}{Rieszsches Lemma}
Sei ($E,\|\cdot\|$) ein normierter Raum und $F\subseteq E$ ein echter abgeschlossener Unterraum. Seien ferner $r>0$ und $\varepsilon \in(0,1)$ gegeben. Dann existiert ein $x \in E$ mit $\|x\|=r$ und $\operatorname{dist}(x, F) \geq(1-\varepsilon) r$.
\end{LEM}

Beweis. 

\begin{PROOF}{FA-A25-01-89}{P: Rieszsches Lemma}
Da $F=\bar{F} \neq E$ gilt, existiert $z \in E \backslash F$. Es folgt $\alpha:=\operatorname{dist}(z, F)>0$. Ferner existiert $y \in F$ mit $\|z-y\|<\frac{\alpha}{1-\varepsilon}$. Seien nun $\lambda:=\frac{r}{\|z-y\|}$ und $x:=\lambda(z-y)$. Dann ist $\|x\|=r$, und für alle $w \in F$ gilt:
$$
\|x-w\|=\lambda\|z-\underbrace{\left(y+\lambda^{-1} w\right)}_{\in F}\| \geq \lambda \alpha=\frac{r \alpha}{\|z-y\|}>(1-\varepsilon) r .
$$
Es folgt $\operatorname{dist}(x, F) \geq(1-\varepsilon) r$.
\end{PROOF}



2.45 Satz. 



\begin{THEO}{FA-A25-01-90}{Heine--Borel-Charakterisierung von endlich dim. normierten Räumen}
Für einen normierten Raum $(E,\|\cdot\|)$ sind äquivalent:
\begin{enumerate}
\item[(i)] $\dim E<\infty$.
\item[(ii)] Jede abgeschlossene und beschränkte Teilmenge von $E$ ist kompakt.
\item[(iii)] Die Einheitskugel $B_1(0)=\{x\in E:\|x\|\le 1\}$ ist präkompakt (totale Beschränktheit).
\end{enumerate}
\end{THEO}


Für einen normierten Raum ($E,\|\cdot\|$) sind äquivalent:

(i) $\operatorname{dim} E<\infty$\\
(ii) Jede abgeschlossene und beschränkte Teilmenge von E ist kompakt\\
(iii) $B_{1}(0) \subseteq E$ ist präkompakt.

Beweis. 

\begin{PROOF}{FA-A25-01-91}{P: Heine--Borel-Charakterisierung von endlich dim. normierten Räumen}
„(i) $\Longrightarrow$ (ii)": siehe Bemerkung 2.41(c).

„(ii) $\Longrightarrow$ (iii)": $B_{1}(0)$ ist abgeschlossen und beschränkt, also kompakt nach Voraussetzung und damit auch präkompakt.

„(iii) $\Longrightarrow$ (i)": Nach Satz 2.43 existiert ein endlichdimensionaler (und damit nach Bemerkungen und Beispiele 2.34(c) abgeschlossener) Unterraum $F \subseteq E$ mit
$$
B_{1}(0) \subseteq B_{\frac{1}{3}}(F)
$$
Ist $F \neq E$, so existiert nach Lemma 2.44 zu $\varepsilon=\frac{1}{3}$ und $r=1$ ein $x \in E \backslash F$ mit $\|x\|=1$ und dist $(x, F) \geq\left(1-\frac{1}{3}\right)=\frac{2}{3}>\frac{1}{3}$, also $x \notin B_{\frac{1}{3}}(F)$ im Widerspruch zu (2.8). Es folgt also $E=F$; somit ist $E$ endlichdimensional.
\end{PROOF}




\pagebreak

\section{Vollständigkeit}

\subsection{Der Satz von Baire}

Sei stets ($X,d$) ein metrischer Raum.


3.1 Satz (von Baire). 

\begin{PROP}{FA-A25-02-01}{Baire}
Sei ($X,d$) ein vollständiger metrischer Raum, und seien $D_{n} \subseteq X$ offene und dichte Teilmengen für $n \in \mathbb{N}$. Dann ist auch $D:=\bigcap_{n \in \mathbb{N}} D_{n}$ dicht in $X$.
\end{PROP}

Beweis. 

\begin{PROOF}{FA-A25-02-02}{P: Baire}
Es reicht, zu zeigen, dass $D \cap U \neq \varnothing$ für jede nichtleere offene Teilmenge $U$ von $X$ gilt. Sei also $U \subseteq X$ offen und nichtleer. Da $D_{1} \subseteq X$ offen und dicht ist, ist $D_{1} \cap U$ offen und nichtleer. Wähle $x_{1} \in D_{1} \cap U$ und $r_{1} \in\left(0, \frac{1}{2}\right)$ mit $A_{1}:=B_{r_{1}}\left(x_{1}\right) \subseteq D_{1} \cap U$. Dann gilt:
$$
\varnothing \neq \AA_{1} \subseteq A_{1}=\bar{A}_{1} \subseteq D_{1} \cap U \quad \text { und } \quad \operatorname{diam} A_{1} \leq 1
$$
Konstruiere nun sukzessive Mengen $A_{n} \subseteq X, n \geq 2$ mit
$$
\varnothing \neq \AA_{n} \subseteq A_{n}=\bar{A}_{n} \subseteq D_{n} \cap A_{n-1} \quad \text { und } \quad \operatorname{diam} A_{n} \leq \frac{1}{n}
$$
Sei dazu $n \geq 2$, und seien $A_{1}, \ldots, A_{n-1}$ bereits konstruiert. Da $\varnothing \neq \AA_{n-1}$ offen in $X$ und $D_{n}$ offen und dicht ist, muss $D_{n} \cap \AA_{n-1}$ offen und nichtleer sein. Wähle $x_{n} \in D_{n} \cap \AA_{n-1}$ und $r_{n} \in\left(0, \frac{1}{2 n}\right)$ mit $A_{n}:=B_{r_{n}}\left(x_{n}\right) \subseteq D_{n} \cap \AA_{n-1}$. Mit dieser Wahl von $A_{n}$ gilt (3.2). Anwendung von Satz 2.39 auf die Mengen $A_{n}, n \in \mathbb{N}$ liefert schließlich $A:=\bigcap_{n \in \mathbb{N}} A_{n} \neq \varnothing$. Aus (3.1) und (3.2) folgt aber
$$
A_{n} \subseteq\left(D_{1} \cap \cdots \cap D_{n}\right) \cap U \quad \text { für alle } n \in \mathbb{N}, \quad \text { also } A \subseteq D \cap U .
$$
Es folgt $D \cap U \neq \varnothing$.
\end{PROOF}


3.2 Korollar. 

Sei ($X,d$) vollständig, und sei $A_{n} \subseteq X$ abgeschlossen für $n \in \mathbb{N}$. Ferner sei $A:=\bigcup_{n \in \mathbb{N}} A_{n}$. Dann gilt:

(a) Ist $\AA_{n}=\varnothing$ für alle $n \in \mathbb{N}$, so ist auch $\AA=\varnothing$.\\
(b) Enthält $A$ eine offene Kugel $U_{r}(a)$ mit $a \in X, r>0$, dann enthält mindestens eine der Mengen $A_{n}$ eine abgeschlossene Kugel der Form $B_{s}(b)$ mit $b \in X, s>0$.




\begin{KORO}{FA-A25-02-03}{Vereinigung abgeschlossener Mengen in vollständigen Räumen}
Sei $(X,d)$ vollständig und seien $A_n\subseteq X$ abgeschlossen für $n\in\mathbb{N}$. Setze $A:=\bigcup_{n\in\mathbb{N}}A_n$. Dann gilt:
\begin{enumerate}\itemsep0.2em
\item[(a)] Gilt $\operatorname{int}(A_n)=\varnothing$ für alle $n\in\mathbb{N}$, so ist auch $\operatorname{int}(A)=\varnothing$.
\item[(b)] Enthält $A$ eine offene Kugel $U_r(a)$ mit $a\in X$, $r>0$, so existieren $n\in\mathbb{N}$ sowie $b\in X$, $s>0$ mit $B_s(b)\subseteq A_n$.
\end{enumerate}
\end{KORO}


Beweis. 

\begin{PROOF}{FA-A25-02-04}{P: Vereinigung abgeschlossener Mengen in vollständigen Räumen}
(a) Nach Voraussetzung ist $D_{n}:=X \backslash A_{n}$ offen und dicht für alle $n \in \mathbb{N}$. Nach Satz 3.1 ist somit $X \backslash A=\bigcap_{n \in \mathbb{N}} D_{n}$ dicht in $X$. Dies liefert $\AA=X \backslash \overline{X \backslash A}=\varnothing$.

(b) folgt direkt aus (a).
\end{PROOF}


3.3 Definition. 

Sei $A \subseteq X$.

(a) $A$ heißt nirgends dicht (in $X$ ), wenn $\operatorname{int}(\bar{A})=\varnothing$ gilt.\\
(b) $A$ heißt mager (oder von erster Kategorie), falls sich $A$ als abzählbare Vereinigung nirgends dichter Mengen schreiben lässt.



\begin{DEF}{FA-A25-02-05}{Nirgends dichte und magere Mengen}
Sei $(X,d)$ ein metrischer Raum und $A\subseteq X$.
\begin{enumerate}\itemsep0.2em
\item[(a)] $A$ heißt \emph{nirgends dicht} (in $X$), wenn das Innere des Abschlusses leer ist, d.\,h.
\[
\operatorname{int}(\overline{A})=\varnothing.
\]
\item[(b)] $A$ heißt \emph{mager} (oder \emph{von erster Kategorie}), wenn $A$ als abzählbare Vereinigung nirgends dichter Teilmengen dargestellt werden kann, d.\,h. es existieren Mengen $A_k\subseteq X$ mit
\[
A=\bigcup_{k=1}^{\infty} A_k
\quad\text{und}\quad
\operatorname{int}(\overline{A_k})=\varnothing\ \text{für alle }k\in\mathbb{N}.
\]
\end{enumerate}
\end{DEF}



3.4 Bemerkungen und Beispiele. 

(a) 

Die Vereinigung abzählbar vieler magerer Mengen ist mager.




\begin{EXA}{FA-A25-02-06}{Abzählbare Vereinigung magerer Mengen ist mager}
Sei $(X,d)$ ein metrischer Raum. Ist $A_m\subseteq X$ für jedes $m\in\mathbb{N}$ mager, so ist auch
\[
A:=\bigcup_{m=1}^\infty A_m
\]
mager.

\begin{proof}
Für jedes $m$ existiert eine Folge $(N_{m,k})_{k\in\mathbb{N}}$ nirgends dichter Mengen mit
$A_m=\bigcup_{k=1}^\infty N_{m,k}$. Dann gilt
\[
A=\bigcup_{m=1}^\infty A_m=\bigcup_{m=1}^\infty\bigcup_{k=1}^\infty N_{m,k}.
\]
Da $\mathbb{N}\times\mathbb{N}$ abzählbar ist, existiert eine Bijektion
$\varphi:\mathbb{N}\to\mathbb{N}\times\mathbb{N}$, so dass
\[
A=\bigcup_{n=1}^\infty N_{\varphi(n)}
\]
eine abzählbare Vereinigung nirgends dichter Mengen ist. Also ist $A$ mager.
\end{proof}
\end{EXA}




(b) 


Für $z \in X$ ist äquivalent:
$$
\{z\} \text { nirgends dicht in } X \Longleftrightarrow \operatorname{dist}(z, X \backslash\{z\})=0
$$
Gilt dies für alle $z \in X$, so ist jede abzählbare Teilmenge von $X$ mager. Dies ist insbesondere dann der Fall, wenn $\{0\} \neq X$ ein normierter Raum (mit induzierter Metrik) ist.



\begin{EXA}{FA-A25-02-07}{Punkte als nirgends dichte Mengen und Magerkeit abzählbarer Mengen}
Sei $(X,d)$ ein metrischer Raum und $z\in X$. Dann sind die folgenden Aussagen äquivalent:
\[
\{z\}\ \text{ist nirgends dicht in }X
\quad\Longleftrightarrow\quad
\operatorname{dist}\bigl(z,\,X\setminus\{z\}\bigr)=0.
\]

\begin{proof}
\emph{„$\Rightarrow$“:} Ist $\{z\}$ nirgends dicht, so hat $\overline{\{z\}}=\{z\}$ kein Inneres, d.\,h.
es gibt kein $r>0$ mit $U_r(z)\subseteq\{z\}$. Damit gilt: Zu jedem $r>0$ existiert ein $y\in X\setminus\{z\}$
mit $d(y,z)<r$, also $\operatorname{dist}(z,X\setminus\{z\})=0$.

\emph{„$\Leftarrow$“:} Gilt $\operatorname{dist}(z,X\setminus\{z\})=0$, so enthält jede offene Kugel um $z$
auch Punkte aus $X\setminus\{z\}$. Somit besitzt $\{z\}$ kein Inneres, d.\,h. $\operatorname{int}(\overline{\{z\}})=\varnothing$,
also ist $\{z\}$ nirgends dicht.
\end{proof}

\textbf{Folgerung.}  
Gilt diese Äquivalenz für alle $z\in X$, so ist jede abzählbare Teilmenge $A=\{z_1,z_2,\ldots\}$ von $X$
eine abzählbare Vereinigung nirgends dichter Mengen und daher mager:
\[
A=\bigcup_{n=1}^{\infty}\{z_n\}.
\]

\textbf{Spezialfall.}  
Ist $X$ ein normierter Raum mit $\{0\}\neq X$, so gilt für jedes $z\in X$,
\[
\operatorname{dist}(z, X\setminus\{z\})=0,
\]
da es stets Folgen $(x_n)\subset X\setminus\{z\}$ mit $\|x_n-z\|\to 0$ gibt.
Also ist jede abzählbare Teilmenge eines normierten Raumes mager.
\end{EXA}



(c) 

Seien $(E,\|\cdot\|)$ ein normierter Raum und $F \subseteq E$ ein Unterraum, welcher nicht dicht ist. Dann ist $F$ auch nirgends dicht, denn: $F$ nicht dicht bedeutet, dass $v \in E \backslash \bar{F}$ existiert. Da $\bar{F}$ ein Unterraum von $E$ ist (Übung), können wir dabei $0 \neq\|v\|<1$ annehmen. Für alle $x \in \bar{F}$ und $\delta>0$ ist dann $x+\delta v \in U_{\delta}(x) \backslash \bar{F}$. Es folgt $\operatorname{int}(\bar{F})=\varnothing$.




\begin{EXA}{FA-A25-02-08}{Nicht dichter Unterraum ist nirgends dicht}
Sei $(E,\|\cdot\|)$ ein normierter Raum und $F\subseteq E$ ein (linearer) Unterraum. Ist $F$ nicht dicht in $E$, so ist $F$ nirgends dicht in $E$, d.\,h.
\[
\operatorname{int}\bigl(\overline{F}\bigr)=\varnothing.
\]

\begin{proof}
Da $F$ nicht dicht ist, existiert $v\in E\setminus \overline{F}$. Die Menge $\overline{F}$ ist ein abgeschlossener Unterraum von $E$ (denn $F$ ist linear und die Abschlussbildung erhält lineare Strukturen). Insbesondere ist $\overline{F}$ translations- und skalierungsinvariant innerhalb von $F$, aber $v\notin\overline{F}$.

Wähle (durch Normierung) $v$ so, dass $0<\|v\|<1$. Zeige nun: Für jedes $x\in\overline{F}$ und jedes $\delta>0$ gilt
\[
x+\delta v\ \in\ U_\delta(x)\setminus \overline{F}.
\]
Tat­sächlich ist $d(x,x+\delta v)=\|\delta v\|=\delta\|v\|<\delta$, also $x+\delta v\in U_\delta(x)$. Wäre zudem $x+\delta v\in\overline{F}$, so folgte wegen Linearität und Abgeschlossenheit von $\overline{F}$, dass
\[
(x+\delta v)-x=\delta v\in \overline{F}
\quad\Longrightarrow\quad
v=\tfrac{1}{\delta}(\delta v)\in \overline{F},
\]
im Widerspruch zu $v\notin \overline{F}$. Also liegt \emph{jede} offene Kugel $U_\delta(x)$ in $E$ nicht vollständig in $\overline{F}$.

Damit besitzt $\overline{F}$ keinen inneren Punkt, also $\operatorname{int}(\overline{F})=\varnothing$. Folglich ist $F$ nirgends dicht.
\end{proof}
\end{EXA}


\begin{EXA}{FA-A25-02-09}{Echte abgeschlossene Unterräume sind nirgends dicht}
Ist $F$ ein \emph{echter} abgeschlossener Unterraum von $(E,\|\cdot\|)$, so ist $F$ nirgends dicht. (Denn ein abgeschlossener, nicht dichter Unterraum erfüllt die obige Voraussetzung.)
\end{EXA}



In normierten Räumen gilt sogar: Hat ein Unterraum $F$ nichtleeres Inneres, so ist $F=E$. (Denn enthält $F$ eine offene Kugel $U_r(x_0)$, so folgt wegen Linearität $U_r(0)\subset F$ und damit $F=E$.)


3.5 Satz. 

\begin{PROP}{FA-A25-02-10}{$X\backslash A$ ist dicht in vollständigen Raum für $A$ mager}
Seien ($X,d$) vollständig und $A \subseteq X$ mager. Dann ist $X \backslash A$ dicht in $X$. Insbesondere ist $X$ nicht mager in sich, falls $X \neq \varnothing$ ist.
\end{PROP}

Beweis. 

\begin{PROOF}{FA-A25-02-11}{P: $X\backslash A$ ist dicht in vollständigen Raum für $A$ mager}
Sei $A=\bigcup_{n \in \mathbb{N}} A_{n}$ mit $\operatorname{int}\left(\bar{A}_{n}\right)=\varnothing$ für alle $n$. Nach Korollar 3.2 gilt dann $\AA \subseteq \dot{B}=\varnothing$ für $B:=\bigcup_{n \in \mathbb{N}} \overline{A_{n}}$. Es folgt $\overline{X \backslash A}=X \backslash \AA=X$.
\end{PROOF}


3.6 Korollar. 

\begin{KORO}{FA-A25-02-12}{Basis von $\infty$-dim Banachräumen sind überabzählbar}
Sei $E$ ein unendlichdimensionaler Banachraum. Dann ist jede Basis von $E$ (im Sinne der linearen Algebra) überabzählbar.
\end{KORO}

Beweis. 

\begin{PROOF}{FA-A25-02-13}{P: Basis von $\infty$-dim Banachräumen sind überabzählbar}
Angenommen, $E$ besäße eine abzählbare Basis $\left\{v_{1}, v_{2}, \ldots,\right\}$. Dann ist $E=$ $\bigcup_{n \in \mathbb{N}} V_{n}$ mit $V_{n}:=\operatorname{span}\left\{v_{1}, \ldots, v_{n}\right\}$. Nach Voraussetzung und Bemerkungen und Beispiele [Endlichdimensionale Unterräume sind Abgeschlossen] ist $V_{n}=\bar{V}_{n} \neq E$ für alle $n$, und somit ist $V_{n}$ nirgends dicht in $E$ nach Bemerkungen und Beispiele [Nicht dichter Unterraum ist nirgends dicht]. Dann ist $E$ mager in sich, im Widerspruch zu [$X\backslash A$ ist dicht in vollständigen Raum für $A$ mager].
\end{PROOF}


3.7 Bemerkung. Wir werden bald sehen, dass der Satz von Baire wichtige Konsequenzen für die Struktur der Menge stetiger linearer Abbildungen hat (Satz von der gleichmäßigen Beschränktheit, Satz von der offenen Abbildung).


\pagebreak

\subsection{Gleichmäßige Beschränktheit}

Stets seien $E, F$ normierte $\mathbb{K}$-Vektorräume


3.8 Hauptsatz (von der gleichmäßigen Beschränktheit). 

Sei $E$ ein Banachraum, und sei $\mathcal{H} \subseteq \mathcal{L}(E, F)$ punktweise beschränkt, d.h.,
$$
\text { für jedes } x \in E \text { existiere } C_{x}>0 \text { mit }\|T x\|_{F} \leq C_{x} \text { für alle } T \in \mathcal{H} \text {. }
$$
Dann ist $\mathcal{H}$ beschränkt in $\mathcal{L}(E, F)$, d.h. es existiert $C>0$ mit $\|T\| \leq C$ für alle $T \in \mathcal{H}$.



\begin{THEO}{FA-A25-02-14}{Prinzip der gleichmäßigen Beschränktheit (Banach--Steinhaus)}
Sei $E$ ein Banachraum, $F$ ein normierter Raum und $\mathcal{H}\subseteq \mathcal{L}(E,F)$ eine Familie stetiger linearer Operatoren, die \emph{punktweise beschränkt} ist, d.\,h.
\[
\forall x\in E\ \exists C_x>0\ \text{mit}\ \sup_{T\in\mathcal{H}}\|Tx\|_F\le C_x.
\]
Dann ist $\mathcal{H}$ in $\mathcal{L}(E,F)$ \emph{gleichmäßig beschränkt}, d.\,h. es existiert $C>0$ mit
\[
\|T\|\le C\qquad\text{für alle }T\in\mathcal{H}.
\]
\end{THEO}


Beweis. 

\begin{PROOF}{FA-A25-02-15}{P: Prinzip der gleichmäßigen Beschränktheit (Banach--Steinhaus)}
Für $n \in \mathbb{N}$ sei
$$
A_{n}:=\left\{x \in E \mid\|T x\|_{F} \leq n \text { für alle } T \in \mathcal{H}\right\}=\bigcap_{T \in \mathcal{H}} T^{-1}(\underbrace{B_{n}(0)}_{\subseteq F}) .
$$
Dann ist $A_{n}$ abgeschlossen in $E$ für alle $n \in \mathbb{N}$, und nach Voraussetzung gilt $E=\bigcup_{n \in \mathbb{N}} A_{n}$. Da $E$ vollständig ist, existieren nach Korollar 3.2(b) ein $n \in \mathbb{N}, z \in E$ und $s>0$ mit $B_{s}(z) \subseteq A_{n}$, d.h.
$$
\|T x\|_{F} \leq n \quad \text { für alle } x \in B_{s}(z), T \in \mathcal{H} .
$$
Sei nun $v \in E \backslash\{0\}$ und $T \in \mathcal{H}$ beliebig. Dann ist $x_{v}:=z+s \frac{v}{\|v\|_{E}} \in B_{s}(z)$ und $v=\frac{\|v\|_{E}}{s}\left(x_{v}-z\right)$. Es folgt
$$
\|T v\|_{F}=\frac{\|v\|_{E}}{s}\left\|T x_{v}-T z\right\|_{F} \leq \frac{\|v\|_{E}}{s}\left(\left\|T x_{v}\right\|_{F}+\|T z\|_{F}\right) \stackrel{(3.3)}{\leq} \frac{2 n}{s}\|v\|_{E}
$$
Insgesamt folgt $\sup _{T \in \mathcal{H}}\|T\| \leq \frac{2 n}{s}<\infty$.
\end{PROOF}


3.9 Korollar. 


Sei $E$ ein Banachraum, und sei $\left(T_{n}\right)_{n}$ eine Folge in $\mathcal{L}(E, F)$.\\
(a) Gilt $\lim \sup _{n \rightarrow \infty}\left\|T_{n}\right\|=\infty$, so existiert $x \in E$ derart, dass die Folge $\left(T_{n} x\right)_{n}$ in $F$ nicht konvergiert.\\
(b) Existiert $T x:=\lim _{n \rightarrow \infty} T_{n} x \in F$ für alle $x \in E$, so ist ( $\left\|T_{n}\right\|$ ) beschränkt und es gilt $T \in \mathcal{L}(E, F)$ mit

$$
\|T\| \leq \liminf _{n \rightarrow \infty}\left\|T_{n}\right\|
$$

(c) Falls $F$ auch ein Banachraum ist, und falls gilt, dass $\left(\left\|T_{n}\right\|\right)$ beschränkt ist und $\left(T_{n} x\right)$ lediglich für alle $x$ in einer in $E$ dichten Teilmenge konvergiert, dann konvergiert ( $T_{n} x$ ) für alle $x \in E$ und (b) ist anwendbar.



\begin{KORO}{FA-A25-02-16}{Konvergenz linearer Operatorfolgen in Banachräumen}
Sei $E$ ein Banachraum, $F$ ein normierter Raum und $\bigl(T_n\bigr)_{n\in\mathbb{N}}\subset\mathcal{L}(E,F)$ eine Folge stetiger linearer Operatoren. Dann gelten folgende Aussagen:

\begin{enumerate}[label=(\alph*)]
\item
Gilt
\[
\limsup_{n\to\infty}\|T_n\|=\infty,
\]
so existiert ein $x\in E$, für das die Folge $(T_nx)_{n}$ in $F$ \emph{nicht konvergiert}.

\item
Existiert für jedes $x\in E$ der Grenzwert
\[
Tx:=\lim_{n\to\infty}T_nx\in F,
\]
so ist $(\|T_n\|)$ beschränkt, und der Grenzoperator $T$ liegt in $\mathcal{L}(E,F)$ mit
\[
\|T\|\le \liminf_{n\to\infty}\|T_n\|.
\]

\item
Ist $F$ zusätzlich vollständig (d.\,h. ein Banachraum), $(\|T_n\|)$ beschränkt und konvergiert $(T_nx)$ nur für alle $x$ aus einer \emph{dichten Teilmenge} $D\subset E$, so konvergiert $(T_nx)$ für alle $x\in E$, und (b) ist anwendbar.
\end{enumerate}
\end{KORO}




Beweis. 

\begin{PROOF}{FA-A25-02-17}{P: Konvergenz linearer Operatorfolgen in Banachräumen}
(a) Wenn $\lim _{n \rightarrow \infty} T_{n} x \in F$ für alle $x \in E$ existiert, so ist die Menge $\left\{T_{n} \mid n \in \mathbb{N}\right\} \subseteq \mathcal{L}(E, F)$ punktweise beschränkt im Sinne von [Prinzip der gleichmäßigen Beschränktheit (Banach--Steinhaus)] Dieser liefert dann ein $c>0$ mit $\left\|T_{n}\right\| \leq c$ für alle $n \in \mathbb{N}$. Damit folgt (a).

(b) Für $T$ wie in (b) gilt
$$
\|T x\|=\lim _{n \rightarrow \infty}\left\|T_{n} x\right\|=\liminf _{n \rightarrow \infty}\left\|T_{n} x\right\| \leq \liminf _{n \rightarrow \infty}\left\|T_{n}\right\|\|x\| \quad \text { für alle } x \in E
$$
und (a) liefert $\liminf _{n \rightarrow \infty}\left\|T_{n}\right\|<\infty$. Somit ist $T$ (als offensichtlich lineare Abbildung) stetig und es folgt ().

(c) Sei $D$ ein dichte Teilmenge von $E$, so dass ( $T_{n} x$ ) für alle $x \in D$ konvergiert. Wir setzen $C:=\sup \left\|T_{n}\right\|$ und zeigen, dass ( $T_{n} x$ ) für alle $x \in E$ konvergiert. Dafür halten wir $x \in E$ fest und betrachten eine Folge $\left(x_{k}\right) \subseteq D$ mit $x_{k} \rightarrow x$. Sei $\varepsilon>0$. Es existiert $k$ mit $\left\|x-x_{k}\right\| \leq \varepsilon /(3 C)$, und es existiert $n_{0}$ mit $\left\|T_{m} x_{k}-T_{n} x_{k}\right\| \leq \varepsilon / 3$ für alle $m, n \geq n_{0}$. Dann folgt
$$
\begin{aligned}
\left\|T_{m} x-T_{n} x\right\| \leq\left\|T_{m} x-T_{m} x_{k}\right\|+\left\|T_{m} x_{k}-T_{n} x_{k}\right\|+\left\|T_{n} x_{k}-T_{n} x\right\| & \\
& \leq C\left\|x-x_{k}\right\|+\varepsilon / 3+C\left\|x-x_{k}\right\| \leq \varepsilon
\end{aligned}
$$
für alle $m, n \geq n_{0}$, das heißt, ( $T_{n} x$ ) ist eine Cauchyfolge in $F$ und konvergiert demnach.
\end{PROOF}




3.10 Bemerkung (Satz von Banach-Steinhaus). Aus Korollar 3.9 folgt für Banachräume $E$ und $F$ und eine Folge $\left(T_{n}\right) \subseteq \mathcal{L}(E, F)$ die Äquivalenz der Aussagen\\
(i) $\left(T_{n} x\right)$ konvergiert in $F$ für alle $x \in E$;\\
(ii) $\left\|T_{n}\right\|_{\mathcal{L}(E, F)}$ ist beschränkt und ( $T_{n} x$ ) konvergiert für alle $x$ in einer dichten Teilmenge von $E$.

Ist eine dieser Aussagen erfüllt, dann definiert $T x:=\lim _{n \rightarrow \infty} T_{n} x$ ein $T \in \mathcal{L}(E, F)$.



\begin{PROP}{FA-A25-02-18}{Satz von Banach--Steinhaus (Folgenversion)}
Seien $E$ und $F$ Banachräume und $\left(T_n\right)_{n\in\mathbb{N}}\subseteq\mathcal{L}(E,F)$ eine Folge stetiger linearer Operatoren. Dann sind die folgenden Aussagen äquivalent:

\begin{enumerate}[label=(\roman*)]
\item
Für alle $x\in E$ konvergiert die Folge $(T_nx)$ in $F$.
\item
Die Folge $(\|T_n\|_{\mathcal{L}(E,F)})$ ist beschränkt, und $(T_nx)$ konvergiert für alle $x$ in einer \emph{dichten Teilmenge} von $E$.
\end{enumerate}

Ist eine (und damit beide) dieser Aussagen erfüllt, so definiert
\[
Tx:=\lim_{n\to\infty}T_nx
\]
einen wohldefinierten Operator $T\in\mathcal{L}(E,F)$, d.\,h.
\[
T_nx\longrightarrow Tx\quad\text{für alle }x\in E,
\qquad\text{und}\qquad
\|T\|\le\liminf_{n\to\infty}\|T_n\|.
\]
\end{PROP}



\begin{PROOF}{FA-A25-02-19}{P: Satz von Banach--Steinhaus (Folgenversion)}
Die Äquivalenz folgt unmittelbar aus [Konvergenz linearer Operatorfolgen in Banachräumen]:

\begin{itemize}
\item[(i)$\Rightarrow$(ii)] 
Wenn $(T_nx)$ für alle $x\in E$ konvergiert, dann ist $\{T_n\}$ punktweise beschränkt. Nach dem Prinzip der gleichmäßigen Beschränktheit ist $(\|T_n\|)$ beschränkt, und die Konvergenz gilt trivially auf jeder Teilmenge, also insbesondere auf einer dichten.

\item[(ii)$\Rightarrow$(i)]
Ist $(\|T_n\|)$ beschränkt und konvergiert $(T_nx)$ auf einer dichten Teilmenge $D\subset E$, so folgt aus der Vollständigkeit von $F$ und der linearen Stetigkeit der $T_n$, dass $(T_nx)$ für alle $x\in E$ konvergiert (Korollar~3.9(c)).

\end{itemize}
Der Grenzoperator $T$ ist dann stetig, linear und erfüllt die obige Normabschätzung.
\end{PROOF}


\pagebreak


\subsection{Der Satz von der offenen Abbildung}
Stets seien $E, F$ normierte $\mathbb{K}$-Vektorräume\\


3.11 Definition und Bemerkung. Seien $X, Y$ metrische Räume. Eine Abbildung $f: X \rightarrow Y$ heißt offen, falls für $A \subseteq X$ gilt:

$$
A \text { ist offen in } X \quad \Longrightarrow \quad f(A) \text { ist offen in } Y \text {. }
$$

Man beachte: Ist $f: X \rightarrow Y$ offen und bijektiv, so ist $f^{-1}$ stetig (vgl. Definition und Bemerkung 2.22(b))\\



\begin{DEF}{FA-A25-02-20}{Offene Abbildungen zwischen metrischen Räumen}
Seien $(X,d_X)$ und $(Y,d_Y)$ metrische Räume.  
Eine Abbildung $f:X\to Y$ heißt \textbf{offen}, falls gilt:
\[
A\subseteq X \text{ offen in } X \ \Longrightarrow\ f(A)\text{ ist offen in } Y.
\]

\textbf{Bemerkung.}  
Ist $f:X\to Y$ eine \emph{offene} und \emph{bijektive} Abbildung, so ist die Umkehrabbildung
\[
f^{-1}:Y\to X
\]
\emph{stetig}.  
Damit ist eine offene Bijektion automatisch ein \textbf{Homöomorphismus}.
(Vgl. Definition und Bemerkung~2.22(b).)
\end{DEF}






3.12 Hauptsatz (von der offenen Abbildung). 

\begin{THEO}{FA-A25-02-21}{Hauptsatz von der offenen Abbildung}
Sei $E$ ein Banachraum, und sei $T \in$ $\mathcal{L}(E, F)$ derart, dass $\operatorname{Bild}(T)$ nicht mager in $F$ ist. Dann ist $T$ surjektiv und offen.
\end{THEO}


Beweis von Hauptsatz 3.12. 

\begin{PROOF}{FA-A25-02-22}{P: Hauptsatz von der offenen Abbildung}
Zur Abkürzung verwenden wir hier für $r>0$ und $x \in E$ die Notation $B_{r}(x ; E)$ für die abgeschlossene Kugel in $E, B_{r} E:=B_{r}(0 ; E)$, und analoge Schreibweisen in $F$.

Zuerst zeigen wir, dass $r>0$ existiert mit
$$
B_{r} F \subseteq \overline{T\left(B_{1} E\right)}
$$
Für $n \in \mathbb{N}$ setzen wir $A_{n}:=\overline{T\left(B_{n} E\right)}$. Weil
$$
\operatorname{Bild}(T)=\bigcup_{n=1}^{\infty} T\left(B_{n} E\right)
$$
nicht mager ist, existiert $n_{0}$ mit $\operatorname{int}\left(A_{n_{0}}\right) \neq \varnothing$. [Abschluss, Inneres, Bilder und absolute Konvexität] liefert jetzt, dass auch
$$
\operatorname{int}\left(A_{1}\right)=\operatorname{int}\left(\frac{1}{n_{0}} A_{n_{0}}\right)=\frac{1}{n_{0}} \operatorname{int}\left(A_{n_{0}}\right) \neq \varnothing
$$
gilt und dass $y_{0} \in F$ und $r>0$ existieren, so dass $B_{r}\left(y_{0} ; F\right) \subseteq A_{1}$ gilt. Wegen [Abschluss, Inneres, Bilder und absolute Konvexität] ist $A_{1}$ absolut konvex. Daraus folgt $B_{r}\left(-y_{0} ; F\right) \subseteq A_{1}$. Für $x \in B_{r} F$ gilt also $-y_{0}+x, y_{0}+x \in A_{1}$ und somit wegen der Konvexität von $A_{1}$ :
$$
x=\frac{1}{2}\left(-y_{0}+x+y_{0}+x\right) \in A_{1} .
$$
Damit haben wir (3.5) gezeigt.

Im zweiten Schritt zeigen wir für $r$ aus (3.5), dass
$$
B_{r / 2} F \subseteq T\left(B_{1} E\right)
$$
gilt. Sei $y \in B_{r / 2} F$ gewählt. Dann zeigt (3.5), dass $z_{1} \in B_{1 / 2} E$ existiert, so dass $\| y-$ $T z_{1} \| \leq r / 4$ gilt. Mit (3.5) erhalten wir wieder $z_{2} \in B_{1 / 4} E$, so dass $\left\|y-T z_{1}-T z_{2}\right\| \leq r / 8$ ist. Per Induktion konstruieren wir eine Folge $\left(z_{n}\right) \subseteq E$, so dass gilt:
$$
\left\|z_{n}\right\| \leq\left(\frac{1}{2}\right)^{n}, \quad\left\|y-\sum_{k=1}^{n} T z_{k}\right\| \leq r\left(\frac{1}{2}\right)^{n+1}
$$
Wir setzen nun $x_{n}:=\sum_{k=1}^{n} z_{k}$. Dann folgt
$$
\left\|x_{n}\right\| \leq \sum_{k=1}^{n}\left\|z_{k}\right\| \leq \sum_{k=1}^{n}\left(\frac{1}{2}\right)^{k} \leq 1
$$
Wegen (3.7) konvergiert $\sum z_{k}$ absolut. Da $E$ vollständig ist, liefert [Charakterisierung von Banachräumen durch absolute Konvergenz von Reihen], dass $\left(x_{n}\right)$ gegen ein $x \in B_{1} E$ konvergiert. Unter Benutzung von (3.7) prüft man leicht, dass $y=T x \in T\left(B_{1} E\right)$ gilt. Damit ist (3.6) gezeigt. Es folgt sofort
$$
F=\bigcup_{n=1}^{\infty} n B_{r / 2} F \subseteq \bigcup_{n=1}^{\infty} n T\left(B_{1} E\right)=\bigcup_{n=1}^{\infty} T\left(n B_{1} E\right)=\operatorname{Bild}(T)
$$
d.h., $T$ ist surjektiv.

Im letzten Schritt nehmen wir an, dass $U \subseteq E$ offen ist. Für beliebiges $y \in T(U)$ existiert $x \in U$ mit $y=T x$. Weil $U$ offen ist, existiert $t>0$ mit $x+B_{t} E \subseteq U$. Aus (3.6) folgt $B_{t r / 2} F \subseteq T\left(B_{t} E\right)$. Dann gilt $y+B_{t r / 2} F \subseteq y+T\left(B_{t} E\right)=T\left(x+B_{t} E\right) \subseteq T(U)$, und $y+B_{t r / 2}$ ist daher eine Umgebung von $y$, die in $T(U)$ enthalten ist. Da $y \in T(U)$ beliebig war, ist $T(U)$ offen.
\end{PROOF}




Bevor wir den Beweis bringen, betrachten wir erst einmal einige Folgerungen aus diesem zentralen Resultat.\\


3.13 Beispiel. Sei $1 \leq p<q \leq \infty$. Nach Satz 2.9 ist die Inklusion $i:\left(\ell^{p},\|\cdot\|_{p}\right) \hookrightarrow$ $\left(\ell^{q},\|\cdot\|_{q}\right)$ stetig. Da $\ell^{p}$ ein Banachraum und $i$ nicht surjektiv ist, folgt mit Hauptsatz 3.12, dass $\ell^{p}$ in $\left(\ell^{q}\|\cdot\|_{q}\right)$ mager ist.



\begin{EXA}{FA-A25-02-26}{Magerkeit der Inklusion $\ell^p \hookrightarrow \ell^q$}
Sei $1 \leq p < q \leq \infty$.  
Nach [Einbettung und Grenzfall für $\ell^p$-Normen] ist die kanonische Inklusion
\[
i:\bigl(\ell^{p},\|\cdot\|_{p}\bigr)\hookrightarrow \bigl(\ell^{q},\|\cdot\|_{q}\bigr), 
\qquad i(x)=x,
\]
\textbf{stetig}, da für alle $x\in\ell^{p}$ gilt
\[
\|x\|_{q} \le \|x\|_{p}.
\]
Da $\ell^{p}$ vollständig (also ein \textbf{Banachraum}) ist und $i$ \emph{nicht surjektiv} ist
($\ell^{p}\subsetneq \ell^{q}$ für $p<q$), können wir den \textbf{Hauptsatz von Baire} 
[Hauptsatz von der offenen Abbildung] anwenden.

\medskip
\textbf{Folgerung:}  
Das Bild $i(\ell^{p})=\ell^{p}$ ist in $\bigl(\ell^{q},\|\cdot\|_{q}\bigr)$ eine 
\emph{magere Teilmenge}.  

\[
\boxed{\ell^{p} \text{ ist in } \bigl(\ell^{q},\|\cdot\|_{q}\bigr) \text{ mager.}}
\]
\end{EXA}



3.14 Korollar. Seien $E, F$ Banachräume und $T \in \mathcal{L}(E, F)$. Dann gilt:\\
(a) Ist $T$ surjektiv, so ist $T$ offen.\\
(b) Ist $T$ bijektiv, so ist $T$ ein topologischer Isomorphismus.


\begin{KORO}{FA-A25-02-27}{Offenheit surjektiver Operatoren und topologischer Isomorphismus}
Seien $E,F$ Banachräume und $T\in\mathcal{L}(E,F)$. Dann gilt:
\begin{enumerate}[label=(\alph*)]
\item Ist $T$ surjektiv, so ist $T$ eine \emph{offene Abbildung}.
\item Ist $T$ bijektiv, so ist $T$ ein \emph{topologischer Isomorphismus}, d.\,h. $T^{-1}\in\mathcal{L}(F,E)$.
\end{enumerate}
\end{KORO}



Beweis. 

\begin{PROOF}{FA-A25-02-28}{P: Offenheit surjektiver Operatoren und topologischer Isomorphismus}
(a) Nach Voraussetzung und [$X\backslash A$ ist dicht in vollständigen Raum für $A$ mager] ist Bild $T=F$ nicht mager in $F$. Also ist $T$ offen nach [Hauptsatz von der offenen Abbildung]

(b) Aus (a) und Definition von offener Abbildung folgt die Stetigkeit von $T^{-1}$.
\end{PROOF}


3.15 Korollar. Seien $\|\cdot\|_{1},\|\cdot\|_{2}$ Normen auf einem $\mathbb{K}$-Vektorraum $V$ derart, dass $\left(V,\|\cdot\|_{1}\right)$, ( $V,\|\cdot\|_{2}$ ) Banachräume sind. Ferner möge $C>0$ existieren mit $\|x\|_{1} \leq C\|x\|_{2}$ für alle $x \in V$. Dann sind $\|\cdot\|_{1}$ und $\|\cdot\|_{2}$ äquivalent.

Beweis. Nach Voraussetzung ist $I:\left(V,\|\cdot\|_{2}\right) \rightarrow\left(V,\|\cdot\|_{1}\right)$ stetig, also ein topologischer Isomorphismus nach Korollar 3.14(b). Insbesondere ist $I:\left(V,\|\cdot\|_{1}\right) \rightarrow\left(V,\|\cdot\|_{2}\right)$ stetig, d.h. es existiert $D>0$ mit $\|x\|_{2} \leq D\|x\|_{1}$ für alle $x \in V$. Die Behauptung folgt.


\begin{KORO}{FA-A25-02-29}{Äquivalenz von Normen auf vollständigen Räumen}
Seien $\|\cdot\|_{1},\|\cdot\|_{2}$ zwei Normen auf einem $\mathbb{K}$-Vektorraum $V$ derart, 
dass $(V,\|\cdot\|_{1})$ und $(V,\|\cdot\|_{2})$ Banachräume sind. 
Angenommen, es existiert $C>0$ mit
\[
\|x\|_{1}\le C\,\|x\|_{2}\quad\text{für alle }x\in V.
\]
Dann sind $\|\cdot\|_{1}$ und $\|\cdot\|_{2}$ äquivalent; d.\,h. es existiert $D>0$ mit
\[
\|x\|_{2}\le D\,\|x\|_{1}\quad\text{für alle }x\in V.
\]
\end{KORO}



\begin{PROOF}{FA-A25-02-30}{P: Äquivalenz von Normen auf vollständigen Räumen}
Betrachte die Identität $I:(V,\|\cdot\|_{2})\to (V,\|\cdot\|_{1})$, $I(x)=x$. 
Aus $\|x\|_{1}\le C\|x\|_{2}$ folgt die Stetigkeit von $I$. 
Da sowohl Definitions- als auch Zielraum Banachräume sind und $I$ bijektiv ist, 
liefert [Offenheit surjektiver Operatoren und topologischer Isomorphismus] die Stetigkeit der Inversen 
$I^{-1}:(V,\|\cdot\|_{1})\to (V,\|\cdot\|_{2})$. Also existiert $D>0$ mit $\|x\|_{2}\le D\|x\|_{1}$ für alle $x\in V$. 
Damit sind die Normen äquivalent.
\end{PROOF}



Für den Beweis von Hauptsatz 3.12 vereinbaren wir zunächst ein paar praktische Bezeichnungen.


3.16 Notation. Seien $x \in E, \alpha \in \mathbb{K} \backslash\{0\}$ und $A, B \subseteq E$. Wir bezeichnen

$$
\begin{aligned}
x+A & :=\{x+y \mid y \in A\} \\
\alpha A & :=\{\alpha y \mid y \in A\} \\
A+B & :=\{y+z \mid y \in A, z \in B\}
\end{aligned}
$$


Operationen auf Mengen in Vektorräumen

\begin{DEF}{FA-A25-02-23}{Operationen auf Mengen in Vektorräumen}
Sei $(E,\|\cdot\|)$ ein Vektorraum über dem Körper $\mathbb{K}$, 
und seien $x\in E$, $\alpha\in\mathbb{K}\setminus\{0\}$ sowie $A,B\subseteq E$.  
Dann definieren wir folgende Mengenoperationen:

\[
\begin{aligned}
x + A &:= \{\,x + y \mid y \in A\,\}
&&\text{(Translation von $A$ um den Vektor $x$),}\\[0.5em]
\alpha A &:= \{\,\alpha y \mid y \in A\,\}
&&\text{(Skalierung von $A$ mit dem Faktor $\alpha$),}\\[0.5em]
A + B &:= \{\,y + z \mid y \in A,\ z \in B\,\}
&&\text{(Minkowski-Summe der Mengen $A$ und $B$).}
\end{aligned}
\]
\end{DEF}



3.17 Lemma. Seien $x \in E, \alpha \in \mathbb{K} \backslash\{0\}$ und $A \subseteq E$. Dann gelten:\\
(a) $\overline{\alpha A}=\alpha \bar{A}$ und $\operatorname{int}(\alpha A)=\alpha \operatorname{int}(A)$.\\
(b) $\overline{x+A}=x+\bar{A}$ und $\operatorname{int}(x+A)=x+\operatorname{int}(A)$.\\
(c) Für einen Vektorraum $F$ und eine lineare Abbildung $T: E \rightarrow F$ ist $T(A)$ absolut konvex, falls $A$ absolut konvex ist.\\
(d) Wenn $A$ absolut konvex ist, dann ist auch $\bar{A}$ absolut konvex.

Beweis. Übung.\\



\begin{LEM}{FA-A25-02-24}{Abschluss, Inneres, Bilder und absolute Konvexität}
Seien $E$ ein normierter (oder allgemeiner: topologischer) Vektorraum über $\mathbb{K}\in\{\mathbb{R},\mathbb{C}\}$, $x\in E$, $\alpha\in\mathbb{K}\setminus\{0\}$ und $A\subseteq E$. Dann gilt:
\begin{enumerate}\itemsep0.25em
\item[(a)] $\overline{\alpha A}=\alpha\,\overline{A}$ und $\operatorname{int}(\alpha A)=\alpha\,\operatorname{int}(A)$.
\item[(b)] $\overline{x+A}=x+\overline{A}$ und $\operatorname{int}(x+A)=x+\operatorname{int}(A)$.
\item[(c)] Ist $F$ ein Vektorraum und $T:E\to F$ linear, so ist $T(A)$ absolut konvex, falls $A$ absolut konvex ist.
\item[(d)] Ist $A$ absolut konvex, so ist auch $\overline{A}$ absolut konvex.
\end{enumerate}
\end{LEM}



\begin{PROOF}{FA-A25-02-25}{P: Abschluss, Inneres, Bilder und absolute Konvexität}
\textbf{(a)} Die Abbildung $m_\alpha:E\to E$, $m_\alpha(y)=\alpha y$, ist ein Homöomorphismus mit inversem $m_{\alpha^{-1}}$.
Daher gilt für jede Menge $M\subset E$:
\[
\overline{m_\alpha(M)}=m_\alpha(\overline{M})
\qquad\text{und}\qquad
\operatorname{int}(m_\alpha(M))=m_\alpha(\operatorname{int}(M)).
\]
Mit $M=A$ folgt die Behauptung.

\medskip
\textbf{(b)} Die Translation $\tau_x:E\to E$, $\tau_x(y)=x+y$, ist ein Homöomorphismus mit Inverser $\tau_{-x}$. Wie in (a) folgt
\[
\overline{\tau_x(A)}=\tau_x(\overline{A})
\qquad\text{und}\qquad
\operatorname{int}(\tau_x(A))=\tau_x(\operatorname{int}(A)),
\]
also die Behauptung.

\medskip
\textbf{(c)} Erinnern wir: „absolut konvex“ bedeutet: $A$ ist konvex und balanciert (kreisförmig), d.\,h.
$\lambda A\subset A$ für alle $\lambda\in\mathbb{K}$ mit $|\lambda|\le 1$.
Ist $A$ konvex, so ist $T(A)$ konvex, denn für $u=T(a_1), v=T(a_2)$ und $t\in[0,1]$ gilt
\[
t u+(1-t)v=T(ta_1+(1-t)a_2)\in T(A).
\]
Ist $A$ balanciert und $|\lambda|\le 1$, so folgt
\[
\lambda\,T(A)=T(\lambda A)\subset T(A).
\]
Damit ist $T(A)$ balanciert und somit absolut konvex.

\medskip
\textbf{(d)} Abschluss erhält sowohl Konvexität als auch Balanciertheit:
\begin{itemize}\itemsep0.15em
\item Konvexität: Sind $y_1,y_2\in\overline{A}$, existieren Folgen $(a_n),(b_n)\subset A$ mit $a_n\to y_1$, $b_n\to y_2$. Für $t\in[0,1]$ ist $t a_n+(1-t)b_n\in A$ (Konvexität von $A$) und $t a_n+(1-t)b_n\to t y_1+(1-t)y_2$. Also $t y_1+(1-t)y_2\in\overline{A}$.
\item Balanciertheit: Sei $|\lambda|\le 1$ und $y\in\overline{A}$. Dann gibt es $a_n\in A$ mit $a_n\to y$. Da $A$ balanciert ist, gilt $\lambda a_n\in A$, und $\lambda a_n\to \lambda y$. Also $\lambda y\in\overline{A}$.
\end{itemize}
Damit ist $\overline{A}$ absolut konvex.
\end{PROOF}



\pagebreak


\subsection{Der Satz vom abgeschlossenen Graphen}


Seien stets $E$ und $F$ normierte Räume über $\mathbb{K}$.


3.18 Definition und Bemerkung. Auf dem $\mathbb{K}$-Vektorraum $E \times F$ sei die Norm $\|\cdot\|$ definiert durch $\|(x, y)\|:=\|x\|_{E}+\|y\|_{F}$. Dann gilt:

$$
(E \times F,\|\cdot\|) \text { ist Banachraum } \quad \Longleftrightarrow \quad E \text { und } F \text { sind Banachräume }
$$


\begin{EXA}{FA-A25-02-31}{Produkt von Banachräumen}
Seien $E,F$ normierte Räume und definiere auf $E\times F$ die Norm
\[
\|(x,y)\|\;:=\;\|x\|_E+\|y\|_F.
\]
Dann gilt
\[
(E\times F,\|\cdot\|)\text{ ist Banach} \quad\Longleftrightarrow\quad
E\text{ und }F\text{ sind Banachräume.}
\]


\begin{proof}
\emph{``$\Rightarrow$'':} Angenommen, $(E\times F,\|\cdot\|)$ ist vollständig.
Sei $(x_n)$ eine Cauchy-Folge in $E$. Dann ist auch $\bigl((x_n,0)\bigr)$ eine
Cauchy-Folge in $E\times F$, denn
\[
\|(x_n,0)-(x_m,0)\|=\|x_n-x_m\|_E\to 0.
\]
Nach Vollständigkeit existiert $(x, y)\in E\times F$ mit $(x_n,0)\to (x,y)$.
Aus $\|(x_n,0)-(x,y)\|=\|x_n-x\|_E+\|y\|_F\to 0$ folgt $y=0$ und $x_n\to x$ in $E$.
Also ist $E$ vollständig. Das gleiche Argument mit $(0,y_n)$ zeigt, dass auch $F$
vollständig ist.

\medskip
\emph{``$\Leftarrow$'':} Angenommen, $E$ und $F$ sind vollständig.
Sei $\bigl((x_n,y_n)\bigr)$ eine Cauchy-Folge in $E\times F$. Dann gilt
\[
\|x_n-x_m\|_E \le \|(x_n,y_n)-(x_m,y_m)\|\to 0,\qquad
\|y_n-y_m\|_F \le \|(x_n,y_n)-(x_m,y_m)\|\to 0,
\]
also sind $(x_n)$ in $E$ und $(y_n)$ in $F$ Cauchy und konvergieren wegen der
Vollständigkeit gegen $x\in E$ bzw.\ $y\in F$. Dann
\[
\|(x_n,y_n)-(x,y)\|=\|x_n-x\|_E+\|y_n-y\|_F\to 0,
\]
also konvergiert $\bigl((x_n,y_n)\bigr)$ in $E\times F$. Damit ist $E\times F$
vollständig.
\end{proof}
\end{EXA}



3.19 Satz (vom abgeschlossenen Graphen). 


\begin{THEO}{FA-A25-02-32}{Satz vom abgeschlossenen Graphen}
Seien $E, F$ Banachräume, und sei $T: E \rightarrow F$ eine lineare Abbildung, deren Graph $\Gamma:=\{(x, y) \in E \times F \mid y=T x\}$ in $E \times F$ abgeschlossen ist. Dann ist $T \in \mathcal{L}(E, F)$, also stetig.
\end{THEO}

Beweis. 

\begin{PROOF}{FA-A25-02-33}{P: Satz vom abgeschlossenen Graphen}
Nach Voraussetzung und [{FA-A25-02-31}{Produkt von Banachräumen}] ist ($E \times F,\|\cdot\|$) ein Banachraum und somit auch $(\Gamma,\|\cdot\|)$, da $\Gamma \subseteq E \times F$ abgeschlossen ist. Sei $G: \Gamma \rightarrow E$ definiert durch $G(x, y)=x$. Da 
$$\|G(x, y)\|_{E}=\|x\|_{E} \leq\|(x, y)\|$$ 
für $(x, y) \in \Gamma$ ist $G \in$ $\mathcal{L}(\Gamma, E)$. Ferner ist $G$ bijektiv mit $G^{-1}(x)=(x, T x)$ für $x \in E$. Gemäß [{FA-A25-02-27}{Offenheit surjektiver Operatoren und topologischer Isomorphismus}] ist nun $G^{-1} \in \mathcal{L}(E, \Gamma)$; also existiert $C>0$ mit
$$
\|x\|_{E}+\|T x\|_{F}=\left\|G^{-1}(x)\right\| \leq C\|x\|_{E} \quad \text { für alle } x \in E,
$$
und somit $\|T x\|_{F} \leq C\|x\|_{E}$ für alle $x \in E$. Es folgt die Stetigkeit von $T$.
\end{PROOF}


3.20 Bemerkung. Sei $T: E \rightarrow F$  $\mathbb{K}$-linear. Offensichtlich sind dann folgende Eigenschaften äquivalent:

\begin{itemize}
  \item Der Graph von $T$ ist abgeschlossen.
  \item Für jede Folge $\left(x_{n}\right)_{n}$ in $E$, für die Grenzwerte
\end{itemize}

$$
x:=\lim _{n \rightarrow \infty} x_{n} \in E \quad \text { und } \quad y:=\lim _{n \rightarrow \infty} T x_{n} \in F
$$

existieren, gilt $y=T x$.

Für jede Nullfolge $\left(x_{n}\right)_{n}$ in $E$, für die der Grenzwert $y:=\lim _{n \rightarrow \infty} T x_{n} \in F$ existiert, gilt $y=0$.
  
  


\begin{CONC}{FA-A25-02-34}{Charakterisierung abgeschlossener Graphen linearer Abbildungen}
Sei $T:E\to F$ eine $\mathbb{K}$-lineare Abbildung zwischen normierten Räumen. Dann sind die folgenden Aussagen äquivalent:
\begin{enumerate}[label=(\alph*)]
\item Der \textbf{Graph} von $T$, 
\[
\Gamma(T):=\{(x,Tx)\in E\times F\},
\]
ist abgeschlossen in $E\times F$.
\item Für jede Folge $(x_n)\subseteq E$, für die 
\[
x_n\to x\in E \quad\text{und}\quad Tx_n\to y\in F
\]
gilt, folgt $y=Tx$.
\item Für jede Nullfolge $(x_n)\subseteq E$, für die $Tx_n\to y\in F$ existiert, gilt $y=0$.
\end{enumerate}


\begin{proof}
(a) $\Rightarrow$ (b): Ist der Graph abgeschlossen und gilt $x_n\to x$, $Tx_n\to y$, so folgt $(x_n,Tx_n)\to (x,y)\in E\times F$. 
Da der Grenzwert eines konvergenten Graphenpunktes wieder im Graphen liegt, gilt $(x,y)\in\Gamma(T)$, also $y=Tx$.

\smallskip
(b) $\Rightarrow$ (c): Setze $x=0$. Dann folgt aus $x_n\to 0$ und $Tx_n\to y$, dass $y=T0=0$.

\smallskip
(c) $\Rightarrow$ (a): Sei $(x_n,Tx_n)\to(x,y)$ in $E\times F$. 
Dann gilt $x_n-x\to 0$ und $T(x_n-x)=Tx_n-Tx\to y-Tx$. 
Nach (c) folgt $y-Tx=0$, also $(x,y)\in\Gamma(T)$ und der Graph ist abgeschlossen.
\end{proof}
\end{CONC}
  

3.21 Bemerkung. Seien $V, W$ Teilräume von $E$ mit $E=V \oplus W$. Aus der linearen Algebra ist bekannt, dass genau eine Projektion $P: E \rightarrow E$ (d.h. $P$ linear und $P^{2}=P$ ) mit $V=\operatorname{Kern} P$ und $W=\operatorname{Bild} P$ existiert. Wir nennen $P$ die Projektion auf $W$ längs $V$. Dann ist $I-P: E \rightarrow E$ die Projektion auf $V$ längs $W$.


\begin{DEF}{FA-A25-02-35}{Projektionen bei direkter Summe von Teilräumen}
Seien $V,W$ Unterräume eines Vektorraums $E$ mit direkter Summe 
\[
E = V \oplus W.
\]
Dann existiert genau eine lineare Abbildung $P:E\to E$ mit den Eigenschaften
\[
P^2 = P, \qquad \ker P = V, \qquad \operatorname{im} P = W.
\]
Diese Abbildung $P$ heißt \textbf{Projektion auf $W$ längs $V$}.
Ebenso ist $I - P$ die \textbf{Projektion auf $V$ längs $W$}, also
\[
(I-P)^2 = I-P, \qquad \ker (I-P) = W, \qquad \operatorname{im}(I-P) = V.
\]

\begin{proof}[Begründung]
Da $E=V\oplus W$ gilt, besitzt jedes $x\in E$ eine eindeutige Zerlegung 
$x = v + w$ mit $v\in V$, $w\in W$. 
Wir definieren 
\[
P(x) := w.
\]
Dann ist $P$ offensichtlich linear, $P^2 = P$, und 
$\ker P = V$, $\operatorname{im} P = W$. 
Die Eindeutigkeit folgt daraus, dass die Darstellung $x=v+w$ eindeutig ist.
\end{proof}
\end{DEF}



3.22 Definition. Seien $V, W$ Teilräume von $E$ mit $E=V \oplus W$, und sei $P$ die Projektion auf $W$ längs $V$.\\

(a) Offensichtlich ist $P$ genau dann stetig, wenn $I-P$ stetig ist. Gilt dies, so schreiben wir $E=V \oplus_{\text {top }} W$ (topologische direkte Summe) und nennen $W$ ein topologisches Komplement von $V$ (bzw. umgekehrt). Insbesondere sind in diesem Fall $V$ und $W$ als Kerne stetiger linearer Operatoren abgeschlossen in $E$.\\

(b) Ein Unterraum $V \subseteq E$ heißt stetig projiziert, falls es eine Projektion $P \in \mathcal{L}(E)$ mit $V=$ Bild $P$ gibt, d.h. falls $V$ ein topologisches Komplement in $E$ besitzt.\\



\begin{DEF}{FA-A25-02-36}{Topologische direkte Summe und stetige Projektionen}
Seien $V, W$ Teilräume eines normierten Raumes $E$ mit direkter Summe 
\[
E = V \oplus W,
\]
und sei $P:E\to E$ die Projektion auf $W$ längs $V$, d.h.\ $P^2=P$, $\ker P=V$, $\operatorname{im}P=W$.

\begin{enumerate}[label=(\alph*)]
\item
Die Projektion $P$ ist genau dann stetig, wenn auch $I-P$ stetig ist. 
Ist dies der Fall, so schreibt man
\[
E = V \oplus_{\mathrm{top}} W
\]
und nennt diese Zerlegung eine \textbf{topologische direkte Summe}.
In diesem Fall heißen $V$ und $W$ \textbf{topologische Komplemente} voneinander.
Insbesondere sind $V$ und $W$ als Kerne stetiger linearer Operatoren abgeschlossen in $E$.

\item
Ein Unterraum $V \subseteq E$ heißt \textbf{stetig projiziert}, 
falls eine stetige Projektion $P \in \mathcal{L}(E)$ existiert mit
\[
V = \operatorname{im} P.
\]
Äquivalent dazu besitzt $V$ ein topologisches Komplement $W$ in $E$, 
d.h.\ $E = V \oplus_{\mathrm{top}} W$.
\end{enumerate}
\end{DEF}


3.23 Beispiel. 

(a) Sei $E=\left(\mathbb{R}^{2},|\cdot|_{\infty}\right), V=\operatorname{span}\{(1,0)\}$ und $W_{n}=\operatorname{span}\left\{\left(1, \frac{1}{n}\right)\right\}$ für $n \in \mathbb{N}$. Dann ist $E=V \oplus W_{n}$. Sei $P_{n}$ die Projektion mit Bild $P=W_{n}$ und Kern $P=V$. Da ( 0,1 ) = - ( $n, 0)+(n, 1) \in V \oplus W_{n}$ gilt, ist $P_{n}(0,1)=(n, 1)$, also $\left|P_{n}(0,1)\right|_{\infty}=n=n|(0,1)|_{\infty}$. Es folgt $\left\|P_{n}\right\| \geq n$.\\


\begin{EXA}{FA-A25-02-37}{Nichtstetige Projektionen bei variierenden Komplementen}
Sei 
\[
E = (\mathbb{R}^2, \|\cdot\|_\infty), 
\quad V = \operatorname{span}\{(1,0)\},
\quad W_n = \operatorname{span}\!\left\{\left(1, \tfrac{1}{n}\right)\right\}, \quad n\in\mathbb{N}.
\]
Dann gilt für jedes $n$: 
\[
E = V \oplus W_n.
\]
Sei $P_n:E\to E$ die lineare Projektion auf $W_n$ längs $V$, d.h.\ $\operatorname{im}P_n=W_n$ und $\ker P_n=V$.

Da 
\[
(0,1) = -(n,0) + (n,1) \in V \oplus W_n,
\]
folgt 
\[
P_n(0,1) = (n,1).
\]
Damit ergibt sich
\[
\|P_n(0,1)\|_\infty = n = n\,\|(0,1)\|_\infty,
\]
also 
\[
\|P_n\| \ge n.
\]
\textbf{Folgerung:} Die Normen der Projektionen $P_n$ wachsen unbeschränkt, d.h.\ die Folge $(P_n)$ ist nicht gleichmäßig beschränkt. Insbesondere sind die Projektionen $P_n$ für große $n$ nicht stetig in gleichmäßigem Sinne. Dies zeigt, dass die Existenz einer direkten Summe $E = V \oplus W_n$ allein keine topologische direkte Summe garantiert.
\end{EXA}



(b) Sei nun $E=\mathcal{F}$ der Raum der finiten Folgen mit Norm $\|\cdot\|_{\infty}$. Dann ist $E=V \oplus W$ mit
$$
V=\left\{x=\left(x_{j}\right)_{j} \in \mathcal{F} \mid x_{j}=0 \text { für } j \text { gerade }\right\} \quad \text { und } \quad W=\operatorname{span}\left\{d_{n} \mid n \in \mathbb{N}\right\},
$$
wobei $d_{n}=\left(0, \ldots, 0,1, \frac{1}{n}, 0, \ldots\right)=e_{2 n-1}+\frac{1}{n} e_{2 n}$ für $n \in \mathbb{N}$ sei. Sei $P$ die Projektion auf $W$ längs $V$. Wie in (a) sieht man
$$
e_{2 n}=-n e_{2 n-1}+n d_{n} \in V \oplus W \quad \text { für alle } n \in \mathbb{N},
$$
also $\left\|P e_{2 n}\right\|_{\infty}=\left\|n d_{n}\right\|_{\infty}=n=n\left\|e_{2 n}\right\|_{\infty}$. Es folgt, dass $P$ nicht stetig ist.



\begin{EXA}{FA-A25-02-38}{Nichtstetige Projektion im Raum der finiten Folgen}
Sei $E=\mathcal{F}$ der Raum der finiten Folgen über $\mathbb{K}$ mit der Supremumsnorm 
$\|x\|_\infty=\sup_j |x_j|$. Definiere
\[
V:=\{x=(x_j)_j\in\mathcal{F}\mid x_j=0 \text{ für alle geraden } j\}, 
\qquad 
d_n:=e_{2n-1}+\tfrac{1}{n}e_{2n}\ (n\in\mathbb{N}),
\]
und
\[
W:=\operatorname{span}\{d_n\mid n\in\mathbb{N}\}.
\]
Dann gilt $E=V\oplus W$. Sei $P:E\to E$ die Projektion auf $W$ längs $V$, d.\,h.\ $\operatorname{im}P=W$ und $\ker P=V$.

Für jedes $n\in\mathbb{N}$ gilt die Zerlegung
\[
e_{2n}=-n\,e_{2n-1} + n\,d_n \in V\oplus W,
\]
also
\[
P(e_{2n}) = n\,d_n \quad\Longrightarrow\quad 
\|P(e_{2n})\|_\infty=\|n\,d_n\|_\infty = n = n\,\|e_{2n}\|_\infty.
\]
Daraus folgt
\[
\|P\|\;\ge\; \frac{\|P(e_{2n})\|_\infty}{\|e_{2n}\|_\infty} \;=\; n
\quad\text{für alle } n\in\mathbb{N},
\]
also ist $\|P\|=\infty$ und $P$ \emph{nicht stetig}.

\textbf{Fazit:} Obwohl $E=V\oplus W$ algebraisch gilt, ist dies \emph{keine} topologische direkte Summe; das Bild-Komplement $W$ ist hier kein topologisches Komplement von $V$.
\end{EXA}



3.24 Satz. Seien $V, W$ Teilräume von $E$ mit $E=V \oplus W$. Dann gilt:\\

(a) Ist $V$ endlichdimensional und $W$ in $E$ abgeschlossen, so ist $E=V \oplus_{\text {top }} W$.\\

(b) Ist $E$ ein Banachraum und sind $V$ und $W$ abgeschlossen in $E$, so ist $E=V \oplus_{\text {top }} W$.

Beweis. Sei $P: E \rightarrow E$ die Projektion mit Bild $P=V$ und Kern $P=W$.

(a) Angenommen, $P$ ist nicht stetig. Dann existiert eine Folge $\left(x_{n}\right)_{n} \subseteq E \backslash\{0\}$ mit $\left\|P x_{n}\right\| \geq n\left\|x_{n}\right\|$ für alle $n \in \mathbb{N}$. Setze $y_{n}:=\frac{x_{n}}{\left\|P x_{n}\right\|}$. Dann gilt $y_{n} \rightarrow 0$ für $n \rightarrow \infty$, aber $P y_{n} \in S_{V}:=\{z \in V \mid\|z\|=1\}$ für alle $n$. Da $\operatorname{dim} V<\infty$, ist $S_{V}$ kompakt. Also existiert $z \in S_{V}$ und eine Teilfolge $\left(y_{n_{\ell}}\right)_{\ell}$ mit $\lim _{\ell \rightarrow \infty} P y_{n_{\ell}}=z$. Somit ist
$$
-z=\lim _{\ell \rightarrow \infty} \underbrace{(I-P) y_{n_{\ell}}}_{\in W} \in W,
$$
da $W$ abgeschlossen ist. Es folgt $z \in S_{V} \cap W$, im Widerspruch zu $V \cap W=\{0\}$. Der Widerspruch zeigt die Stetigkeit von $P$.

(b) Nach Satz 3.19 reicht es, zu zeigen, dass $P$ einen abgeschlossenen Graphen besitzt. Zum Beweis verwenden wir Bemerkung 3.20. Sei also $\left(x_{n}\right)_{n}$ eine Nullfolge in $E$ derart, dass $y:=\lim _{n \rightarrow \infty} P x_{n} \in E$ existiert. Da $V=\operatorname{Bild} P$ in $E$ abgeschlossen ist, folgt $y \in V$. Da ferner $W=\operatorname{Bild}(I-P)$ in $E$ abgeschlossen ist, folgt $y=-\lim _{n \rightarrow \infty}(I-P) x_{n} \in W$. Somit ist $y \in V \cap W=\{0\}$, d.h. $y=0$. Gemäß Bemerkung 3.20 folgt die Behauptung.



\begin{PROP}{FA-A25-02-39}{Topologische direkte Summe unter Endlichdimensionalität bzw. Abgeschlossenheit}
Seien $V,W$ Teilräume eines normierten Raums $E$ mit $E=V\oplus W$ (algebraische direkte Summe), und sei $P:E\to E$ die Projektion mit $\operatorname{im}P=V$ und $\ker P=W$. Dann gilt:
\begin{enumerate}[label=(\alph*)]
\item Ist $V$ endlichdimensional und $W$ in $E$ abgeschlossen, so ist $P$ stetig; insbesondere $E=V\oplus_{\mathrm{top}} W$.
\item Ist $E$ ein Banachraum und $V,W$ in $E$ abgeschlossen, so ist $P$ stetig; insbesondere $E=V\oplus_{\mathrm{top}} W$.
\end{enumerate}
\end{PROP}



\begin{PROOF}{FA-A25-02-40}{P: Topologische direkte Summe unter Endlichdimensionalität bzw. Abgeschlossenheit}
\textbf{(a)} Angenommen, $P$ sei \emph{nicht} stetig. Dann gibt es eine Folge $(x_n)\subset E\setminus\{0\}$ mit
\[
\|Px_n\|\;\ge n\,\|x_n\|\qquad(n\in\mathbb{N}).
\]
Setze $y_n:=x_n/\|Px_n\|$. Dann $\|y_n\|\le 1/n\to 0$, aber $Py_n\in S_V:=\{v\in V:\|v\|=1\}$ für alle $n$.
Da $V$ endlichdimensional ist, ist $S_V$ kompakt; es existieren also $z\in S_V$ und eine Teilfolge $(n_\ell)$ mit $Py_{n_\ell}\to z$.
Aus $y_{n_\ell}=Py_{n_\ell}+(I-P)y_{n_\ell}$ und $y_{n_\ell}\to 0$ folgt $(I-P)y_{n_\ell}\to -z$.
Da $W=\ker P$ abgeschlossen ist und $(I-P)y_{n_\ell}\in W$, folgt $-z\in W$, also $z\in V\cap W=\{0\}$—im Widerspruch zu $z\in S_V$.
Also ist $P$ stetig, d.\,h. $E=V\oplus_{\mathrm{top}} W$.

\medskip
\textbf{(b)} Sei nun $E$ ein Banachraum und $V,W$ abgeschlossen. Nach dem Satz vom abgeschlossenen Graphen genügt es zu zeigen, dass der Graph von $P$ abgeschlossen ist.
Nach Bemerkung über abgeschlossene Graphen reicht es zu verifizieren:

\emph{Ist $(x_n)$ eine Nullfolge in $E$ und existiert $y=\lim_{n\to\infty}Px_n$ in $E$, so folgt $y=0$.}

Tatsächlich: Aus $Px_n\to y$ und der Abgeschlossenheit von $V=\operatorname{im}P$ folgt $y\in V$. Weiter gilt
\[
(I-P)x_n\;\longrightarrow\; -y,
\]
da $x_n\to 0$. Wegen der Abgeschlossenheit von $W=\ker P=\operatorname{im}(I-P)$ liegt $-y\in W$, also $y\in V\cap W=\{0\}$.
Damit ist der Graph von $P$ abgeschlossen, also $P$ stetig, und somit $E=V\oplus_{\mathrm{top}} W$.
\end{PROOF}


\pagebreak

\section{Räume messbarer Funktionen}


\subsection{Präkompaktheit für Mengen stetiger Abbildungen}

Seien stets $X, Y$ metrische Räume und $E$ ein normierter Raum.\\
4.1 Bemerkung. Ist $X$ kompakt, so ist $C(X, Y)=C_{\mathrm{b}}(X, Y)$ nach Bemerkung 2.41(b) und Bemerkung 2.42(a). Insbesondere ist dann also $C(X)=C_{\mathrm{b}}(X)$ ein Banachraum mit Norm $\|\cdot\|_{\infty}$ nach Bemerkungen und Beispiele 2.34(f).\\
4.2 Definition. Eine Teilmenge $H \subseteq C(X, Y)$ heißt\\
(a) gleichgradig stetig in $a \in X$, wenn für jedes $\varepsilon>0$ ein $\delta>0$ existiert mit $f\left(U_{\delta}(a)\right) \subseteq$ $U_{\varepsilon}(f(a))$ für alle $f \in H$.\\
(b) gleichgradig stetig, wenn (a) für alle $a \in X$ gilt.

Eine Folge $\left(f_{n}\right)_{n}$ in $C(X, Y)$ heißt gleichgradig stetig (in $a \in X$ ), wenn dies für die Menge $\left\{f_{n} \mid n \in \mathbb{N}\right\}$ gilt.\\
4.3 Bemerkung und Beispiel. (a) Gleichgradige Stetigkeit vererbt sich auf Teilmengen von gleichgradig stetigen Mengen von Funktionen.\\
(b) Sei $L>0$ fest gewählt und

$$
H:=\{f \in C(X, Y) \mid d(f(x), f(y)) \leq L \mathrm{~d}(x, y) \quad \text { für alle } x, y \in X\} .
$$

Dann ist $H$ gleichgradig stetig.\\
(c) Sei $X=[0,1], Y=\mathbb{R}$ und sei die Funktionenfolge $\left(f_{n}\right)_{n}$ in $C(X, Y)$ definiert durch $f_{n}(x)=x^{n}$ für $n \in \mathbb{N}$. Wie man leicht sieht, ist diese Folge gleichgradig stetig in jedem Punkt $x \in[0,1)$. Sie ist allerdings nicht gleichgradig stetig in $x=1$ : Für festes $\varepsilon \in\left(0, \frac{1}{2}\right)$ gilt nämlich

$$
\varepsilon^{\frac{1}{n}} \rightarrow 1 \quad \text { für } n \rightarrow \infty
$$

aber

$$
\left|f_{n}\left(\varepsilon^{\frac{1}{n}}\right)-f_{n}(1)\right|=1-\varepsilon>\varepsilon \quad \text { für alle } n \in \mathbb{N} \text {. }
$$

Also existiert kein $\delta>0$ derart, dass für alle $n \in \mathbb{N}$ die Implikation

$$
|x-1|<\delta \quad \Longrightarrow \quad\left|f_{n}(x)-f_{n}(1)\right|<\varepsilon
$$

gilt.\\
4.4 Definition und Satz. Sei $\ell^{\infty}(X, E):=\{f: X \rightarrow E \mid f(X) \subseteq E$ beschränkt $\}$. Dann gilt:\\
(a) $\ell^{\infty}(X, E)$ ist ein normierter Raum mit Norm $\|\cdot\|_{\infty}$ definiert durch $\|f\|_{\infty}:=$ $\sup _{x \in X}\|f(x)\|$ für $f \in \ell^{\infty}(X, E)$.\\
(b) $C_{\mathrm{b}}(X, E)$ ist ein abgeschlossener Teilraum von $\ell^{\infty}(X, E)$.\\
(c) Ist $X$ kompakt, so ist $C(X, E)=C_{\mathrm{b}}(X, E)$.\\
(d) Es gelten folgende Äquivalenzen:

\begin{verbatim}
$E$ Banachraum $\Longleftrightarrow \ell^{\infty}(X, E)$ Banachraum $\Longleftrightarrow C_{\mathrm{b}}(X, E)$ Banachraum
\end{verbatim}

Beweis. (a)-(c) wie für $E=\mathbb{K}$. (d): leichte Übung.\\
$\qquad$ Ende des Inhalts für 2024-05-07\\
4.5 Satz (von Arzelà und Ascoli). Ist $X$ kompakt, so ist eine Teilmenge $H \subseteq C(X, E)$ genau dann präkompakt, wenn $H$ gleichgradig stetig ist und wenn für alle $a \in X$ die Menge $H(a):=\{f(a) \mid f \in H\}$ in $E$ präkompakt ist.

Beweis. „": Sei $\varepsilon>0$. Da $H$ gleichgradig stetig ist, existiert zu jedem $a \in X$ ein $\delta_{a}>0$ mit $f\left(U_{\delta_{a}}(a)\right) \subseteq U_{\frac{\varepsilon}{3}}(f(a))$ für alle $f \in H$. Da $X=\bigcup_{a \in X} U_{\delta_{a}}(a)$ kompakt ist, existieren $a_{1}, \ldots, a_{n}$ derart, dass

$$
X=\bigcup_{i=1}^{n} U_{i} \quad \text { mit } U_{i}:=U_{\delta_{a_{i}}}\left(a_{i}\right) \text { gilt. }
$$

Da die Mengen $H\left(a_{i}\right) \subseteq E, i=1, \ldots, n$ präkompakt sind, existiert eine endliche Menge $M \subseteq E$ mit $\bigcup_{i=1}^{n} H\left(a_{i}\right) \subseteq B_{\frac{\varepsilon}{6}}(M)$. Definiert man für $v=\left(v_{1}, \ldots, v_{n}\right) \in M^{n}$ also

$$
L_{v}:=\left\{f \in H \left\lvert\,\left\|f\left(a_{i}\right)-v_{i}\right\| \leq \frac{\varepsilon}{6}\right. \text { für } i=1, \ldots, n\right\},
$$

so folgt nach Konstruktion $H \subseteq \bigcup_{v \in M^{n}} L_{v}$. Da es sich hier um eine endliche Vereinigung handelt, reicht es gemäß Bemerkung 2.42(b) nun, zu zeigen:

$$
\operatorname{diam} L_{v} \leq \varepsilon \text { in } C(X, E) \text { bzgl. }\|\cdot\|_{\infty} \text { für alle } v \in M^{n} .
$$

Seien dazu $f, g \in L_{v}$, und sei $x \in X$ beliebig. Dann ist $x \in U_{i}$ für ein $i \in\{1, \ldots, n\}$, und somit

$$
\|f(x)-g(x)\| \leq \underbrace{\left\|f(x)-f\left(a_{i}\right)\right\|}_{\leq \frac{\varepsilon}{3}}+\underbrace{\left\|f\left(a_{i}\right)-v_{i}\right\|}_{\leq \frac{\varepsilon}{6}}+\underbrace{\left\|v_{i}-g\left(a_{i}\right)\right\|}_{\leq \frac{\varepsilon}{6}}+\underbrace{\left\|g\left(a_{i}\right)-g(x)\right\|}_{\leq \frac{\varepsilon}{3}} \leq \varepsilon .
$$

Es folgt $\|f-g\|_{\infty} \leq \varepsilon$, und damit ist $\operatorname{diam} L_{v} \leq \varepsilon$.\\
„ $\Longrightarrow$ ": Seien $a \in X$ und $\varepsilon>0$. Nach Voraussetzung existiert eine endliche Teilmenge $M \subseteq H$ mit

$$
H \subseteq \bigcup_{f \in M} B_{\frac{\varepsilon}{4}}(f) \quad \text { in } C(X, E) \text { bzgl. }\|\cdot\|_{\infty}
$$

Da $M$ endlich ist, existiert zudem $\delta>0$ so, dass $f\left(U_{\delta}(a)\right) \subseteq U_{\frac{\varepsilon}{3}}(f(a))$ für alle $f \in M$ gilt. Sei nun $g \in H$ beliebig. Dann existiert $f \in M$ mit $\|g-f\|_{\infty} \leq \frac{\varepsilon}{4}$, und somit gilt für $y \in U_{\delta}(a)$ :

$$
\begin{aligned}
\|g(y)-g(a)\| & \leq\|g(y)-f(y)\|+\|f(y)-f(a)\|+\|f(a)-g(a)\| \\
& \leq 2\|f-g\|_{\infty}+\|f(y)-f(a)\| \leq 2 \frac{\varepsilon}{4}+\frac{\varepsilon}{3}<\varepsilon .
\end{aligned}
$$

Es folgt $g\left(U_{\delta}(a)\right) \subseteq U_{\varepsilon}(g(a))$, und dies liefert die gleichgradige Stetigkeit von $H$ in $a$. Nach Definition von $\|\cdot\|_{\infty}$ liefert (4.2) aber auch direkt die Inklusion

$$
H(a) \subseteq \bigcup_{f \in M} B_{\frac{\varepsilon}{4}}(f(a)) \quad \text { in } E \text { bzgl. }\|\cdot\|
$$

also ist $H(a)$ in $E$ präkompakt.\\
4.6 Korollar. Seien $X$ kompakt und $\left(f_{n}\right)_{n}$ eine beschränkte und gleichgradig stetige Folge in $C(X)$. Dann hat $\left(f_{n}\right)_{n}$ eine gleichmäßig konvergente Teilfolge.

Beweis. Sei $H:=\left\{f_{n} \mid n \in \mathbb{N}\right\} \subseteq C(X)$. Dann ist $H$ gleichgradig stetig und $H(a)=$ $\{f(a) \mid f \in H\} \subseteq \mathbb{K}$ beschränkt und damit auch präkompakt für alle $a \in X$. Mit Satz 4.5 folgt, dass $H$ in $C(X)$ präkompakt und somit auch relativ kompakt ist (nach Bemerkung 2.42(c), da $C(X)$ vollständig ist). Also hat $\left(f_{n}\right)_{n}$ eine in $\bar{H}$ konvergente Teilfolge.

\subsection*{Der Satz von Stone-Weierstraß}
Im Folgenden sei ( $X, d$ ) stets ein kompakter metrischer Raum.\\
4.7 Bemerkung. Im Folgenden wollen wir ein hinreichendes Kriterium dafür herleiten, dass eine Teilmenge im normierten Raum ( $C(X),\|\cdot\|_{\infty}$ ) dicht liegt. Dabei nutzen wir die Tatsache, dass dieser Raum auch eine kommutative $\mathbb{K}$-Algebra mit Eins ist, d.h.:

\begin{itemize}
  \item $C(X)$ ist abgeschlossen unter der (punktweisen) Multiplikation • von Funktionen, und diese ist assoziativ, distributiv (bzgl. der Addition in $C(X)$ ) und kommutativ.
  \item Die konstante Einsfunktion $1=1_{X}$ ist ein Einselement für die Multiplikation.
  \item Die Skalarmultiplikation verträgt sich mit der Multiplikation von Funktionen, d.h. es gilt $\lambda(f \cdot g)=(\lambda f) \cdot g=f \cdot(\lambda g)$ für $f, g \in C(X)$ und $\lambda \in \mathbb{K}$.
\end{itemize}

Ein Unterraum $A \subseteq C(X)$ heißt Teilalgebra von $C(X)$ (mit 1), wenn gilt:

\begin{itemize}
  \item $1 \in A$
  \item Sind $f, g \in A$, so ist auch $f \cdot g \in A$.\\
4.8 Satz (von Stone-Weierstraß). Sei A eine Teilalgebra von $C(X)=C(X, \mathbb{K})$ mit folgenden Eigenschaften.\\
(i) A trennt die Punkte von $X$, d.h. für je zwei Punkte $x, y \in X$ mit $x \neq y$ existiert eine Funktion $f \in A$ mit $f(x) \neq f(y)$.\\
(ii) Im Fall $\mathbb{K}=\mathbb{C}$ ist $A$ abgeschlossen unter komplexer Konjugation, d.h. für $f \in A$ ist auch $\bar{f} \in A$.\\
Dann liegt $A$ dicht in $C(X)$.\\
Um dies zu beweisen, benötigen wir zunächst zwei Hilfssätze:\\
4.9 Hilfssatz (Spezialfall der binomischen Reihe). Für $s \in[-1,1]$ gilt
\end{itemize}

$$
\sqrt{1-s}=\sum_{n=0}^{\infty} a_{n} s^{n}
$$

mit

$$
a_{n}=(-1)^{n}\binom{1 / 2}{n}=(-1)^{n} \prod_{j=1}^{n} \frac{\frac{3}{2}-j}{j} \quad \text { für } n \in \mathbb{N}_{0}
$$

und die Reihe konvergiert absolut.\\
Beweis. Die aus der Analysis bekannte binomische Reihenentwicklung zur Potenz $\frac{1}{2}$ liefert (4.3) zunächst für $s \in(-1,1)$. Wir zeigen nun die Abschätzung

$$
\left|a_{n}\right| \leq n^{-\frac{3}{2}} \quad \text { für } n \geq 1 \text {. }
$$

Dies ist klar für $a_{1}=-\frac{1}{2}$. Für $n \geq 2$ gilt zudem

$$
\frac{\left|a_{n}\right|}{\left|a_{n-1}\right|}=\left|\frac{\frac{3}{2}-n}{n}\right|=1-\frac{3}{2 n} \leq\left(1-\frac{1}{n}\right)^{\frac{3}{2}}=\left(\frac{n}{n-1}\right)^{-\frac{3}{2}} .
$$

Hier haben wir die Konvexität der Funktion $(-1, \infty) \rightarrow \mathbb{R}, t \mapsto(1+t)^{\frac{3}{2}}$ verwendet (diese liefert die Ungleichung $(1+t)^{\frac{3}{2}} \geq 1+\frac{3}{2} t$ für $t \in(-1, \infty)$ ). Per Induktion folgt also für $n \geq 2$ :

$$
\left|a_{n}\right| \leq\left(\frac{n}{n-1}\right)^{-\frac{3}{2}}\left|a_{n-1}\right| \leq\left(\frac{n}{n-1}\right)^{-\frac{3}{2}}(n-1)^{-\frac{3}{2}}=n^{-\frac{3}{2}}
$$

Insbesondere konvergiert die Reihe in (4.3) für alle $s \in[-1,1]$ absolut, und mit dem Abelschen Grenzwertsatz folgt, dass Hilfssatz 4.9 auch für $s= \pm 1$ gilt. (Randbemerkung: Es gilt $a_{n}<0$ für $n \geq 1$.)\\
4.10 Hilfssatz. Sei A eine Teilalgebra von C ( $X, \mathbb{R}$ ). Dann gilt:\\
(a) Der Abschluss $\bar{A}$ von $A$ ist auch eine Teilalgebra von $C(X, \mathbb{R})$.\\
(b) Ist $A$ abgeschlossen in $C(X, \mathbb{R})$, so gilt:\\
(i) Ist $f \in A$ mit $f \geq 0$, so ist auch $\sqrt{f} \in A$.\\
(ii) Sind $f, g \in A$, so sind auch die Funktionen $\min \{f, g\}, \max \{f, g\} \in A$.

Beweis. (a) Wegen $\|f g\|_{\infty} \leq\|f\|_{\infty}\|g\|_{\infty}$ für $f, g \in C(X)$ ist die Multiplikation in $C(X)$ stetig (überlegen!). Das liefert $f g \in \bar{A}$ für $f, g \in \bar{A}$ (man approximiere $f$ und $g$ aus $A$ heraus!). Da ferner $1 \in A \subseteq \bar{A}$ gilt, folgt die Behauptung.\\
(b)(i) Ohne Einschränkung sei $0 \leq f \leq 1$ (sonst normalisiere man $f$ ). Dann gilt $g=1-f \in A$ und $0 \leq g \leq 1$. Gemäß Hilfssatz 4.9 gilt

$$
\sqrt{f(x)}=\sqrt{1-g(x)}=\sum_{n=0}^{\infty} a_{n} g^{n}(x) \quad \text { für } x \in X
$$

mit den Binomialkoeffizienten aus (4.4). Da $A$ abgeschlossen ist und weil $\sum_{n=0}^{N} a_{n} g^{n} \in A$ gilt für $N \in \mathbb{N}$, reicht es, zu zeigen, dass die Reihe in (4.6) bzgl. der Norm $\|\cdot\|_{\infty}$ von $C(X, \mathbb{R})$ konvergiert; dann ist $\sqrt{f} \in A$. Wegen $\left\|g^{n}\right\|_{\infty} \leq 1$ für alle $n \in \mathbb{N}_{0}$ ist aber

$$
\sum_{n=0}^{\infty}\left\|a_{n} g^{n}\right\|_{\infty} \leq \sum_{n=0}^{\infty}\left|a_{n}\right|\|g\|_{\infty}^{n} \leq \sum_{n=0}^{\infty}\left|a_{n}\right|<\infty
$$

gemäß Hilfssatz 4.9, d.h. die Reihe $\sum_{n=0}^{\infty} a_{n} g^{n}$ konvergiert absolut im Banachraum $C(X, \mathbb{R})$. Nach Satz 2.37 konvergiert sie dann auch.\\
(b)(ii) Wegen (b)(i) liegt mit $f$ auch $f^{2}$ und somit $|f|=\sqrt{f^{2}}$ in $A$. Mit $f, g \in A$ sind also auch

$$
\min \{f, g\}=\frac{1}{2}(f+g-|f-g|) \quad \text { und } \quad \max \{f, g\}=-\min \{-f,-g\}
$$

Funktionen in $A$.

Beweis von Satz 4.8. Wir betrachten zunächst den Fall $\mathbb{K}=\mathbb{R}$. Sei $f \in C(X, \mathbb{R})$ und $\varepsilon>0$ gegeben. Es reicht, zu zeigen, dass $g \in \bar{A}$ existiert mit $\|f-g\|_{\infty} \leq \varepsilon$. Zur Konstruktion der Funktion $g$ betrachten wir zunächst beliebige $x, y \in X$ mit $x \neq y$. Nach Voraussetzung existiert eine Funktion $h=h_{x, y} \in A$ mit $h(x) \neq h(y)$. Mit dieser Wahl definieren wir nun $g_{x, y} \in A$ durch

$$
g_{x, y}(z):=f(y)+(f(x)-f(y)) \frac{h(z)-h(y)}{h(x)-h(y)} \quad \text { für } z \in X
$$

dann gilt $g_{x, y}(x)=f(x)$ und $g_{x, y}(y)=f(y)$. Ferner definieren wir die konstante Funktion $g_{x, x} \in A, g_{x, x} \equiv f(x)$ für $x \in X$. Sei nun $x \in X$ fest gewählt. Für $y \in X$ betrachten wir\\
dann die offene Umgebung

$$
U_{y}:=\left\{z \in X \mid g_{x, y}(z)<f(z)+\varepsilon\right\}
$$

von $y$. Aufgrund der Kompaktheit von $X$ existieren $y_{1}, \ldots, y_{n} \in X$ mit $X=\bigcup_{j=1}^{n} U_{\underline{y_{j}}}$. Ferner liegt wegen Hilfssatz 4.10 die Funktion $g_{x}:=\min \left\{g_{x, y_{i}} \mid i=1, \ldots, n\right\}$ in $\bar{A}$. Dabei gilt $g_{x}(x)=f(x)$ und $g_{x}(z)<f(z)+\varepsilon$ für alle $z \in X$. Daher ist $V_{x}:=$ $\left\{z \in X \mid g_{x}(z)>f(z)-\varepsilon\right\}$ eine offene Umgebung von $x$. Wiederum aufgrund der Kompaktheit von $X$ existieren $x_{1}, \ldots, x_{m} \in X$ mit $X=\bigcup_{j=1}^{m} V_{x_{j}}$. Ferner ist auch $g:=\max \left\{g_{x_{j}} \mid j=1, \ldots, m\right\} \in \bar{A}$, und es gilt

$$
f(z)-\varepsilon<g<f(z)+\varepsilon \quad \text { für alle } z \in X .
$$

Es folgt $\|f-g\|_{\infty} \leq \varepsilon$, wie benötigt.\\
Wir betrachten schließlich den Fall $\mathbb{K}=\mathbb{C}$. Für $A_{0}:=\{\operatorname{Re} f \mid f \in A\}$ gilt dann: 1 . Für $f \in A$ ist auch $\operatorname{Re} f=\frac{1}{2}(f+\bar{f}) \in A$, und somit folgt $A_{0} \subseteq A$. 2. Für $f \in A$ ist $\operatorname{Im} f=\operatorname{Re}(-\mathrm{i} f) \in A_{0}$. Aus 1. und 2. folgt die Mengengleichheit

$$
A=A_{0}+\mathrm{i} A_{0}
$$

und diese impliziert insbesondere, dass $A_{0}$ die Punkte von $X$ trennt. Ferner ist $A_{0}$ eine Teilalgebra von $C(X, \mathbb{R})$, denn: Für $\operatorname{Re} f, \operatorname{Re} g \in A_{0}$ mit $f, g \in A$ ist

$$
\operatorname{Re} f \cdot \operatorname{Re} g=\operatorname{Re}(f \operatorname{Re} g) \in A_{0}
$$

da nach 1. auch Re $g$ und somit $f \operatorname{Re} g$ in $A$ liegt. Ferner ist offensichtlich $1 \in A_{0}$. Wie bereits gezeigt, liegt $A_{0}$ also dicht in $C(X, \mathbb{R})$, und somit liegt $A$ wegen (4.7) auch dicht in $C(X, \mathbb{C})$.\\
4.11 Korollar. (a) Sei $X \subseteq \mathbb{R}^{N}$ kompakt. Dann liegen die Polynomfunktionen in $x_{1}, \ldots, x_{N}$ (mit Koeffizienten in $\mathbb{K}$ ) dicht in $C(X, \mathbb{K})$.\\
(b) Sei $X \subseteq \mathbb{C}^{N}$ kompakt. Dann liegen die Polynomfunktionen in $z_{1}, \ldots, z_{N}, \bar{z}_{1}, \ldots, \bar{z}_{N}$ mit komplexen Koeffizienten dicht in $C(X, \mathbb{C})$.

Beweis. (a) Offensichtlich bilden die genannten Funktionen eine Teilalgebra von $C(X, \mathbb{K})$, welche die Punkte von $X$ trennt, da zu je zwei verschiedenen Punkten $x, y \in \mathbb{R}^{N}$ eine affin lineare Funktion $g: \mathbb{R}^{N} \rightarrow \mathbb{R}$ existiert mit $g(x) \neq g(y)$ (z.B. die Funktion $\left.z \mapsto g(z)=(z-y) \cdot(x-y), z \in \mathbb{R}^{N}\right)$. Also folgt die Behauptung aus den Satz von StoneWeierstraß.\\
(b) Die Behauptung folgt wie (a) aus Satz 4.8, wenn man zusätzlich beachtet, dass die hier betrachteten Funktionen eine Teilalgebra bilden, welche auch unter der komplexen Konjugation abgeschlossen ist.\\
4.12 Definition und Korollar. Sei $C_{2 \pi}(\mathbb{R}, \mathbb{C})$ der Vektorraum der stetigen $2 \pi$ periodischen Funktionen $\mathbb{R} \rightarrow \mathbb{C}$. Dann liegt der Teilraum $T$ der trigonometrischen Polynome dicht in $C_{2 \pi}(\mathbb{R}, \mathbb{C})$ bzgl. $\|\cdot\|_{\infty}$. Dabei heißt eine Funktion $f \in C_{2 \pi}(\mathbb{R}, \mathbb{C})$ trigonometrisches Polynom, wenn es $n \in \mathbb{N}$ und $a_{k} \in \mathbb{C}, k=-n, \ldots, n$, gibt mit $f(t)=\sum_{k=-n}^{n} a_{k} \mathrm{e}^{\mathrm{i} k t}$ für $t \in \mathbb{R}$.

Beweis. Sei $S^{1}:=\{z \in \mathbb{C}| | z \mid=1\} \subseteq \mathbb{C}$ der Einheitskreis. Dann ist die Abbildung

$$
T: C\left(S^{1}, \mathbb{C}\right) \rightarrow C_{2 \pi}(\mathbb{R}, \mathbb{C}), \quad(T f)(t):=f\left(\mathrm{e}^{\mathrm{i} t}\right)
$$

ist ein isometrischer Isomorphismus (bzgl. $\|\cdot\|_{\infty}$ ), welcher Polynome in $z$ und $\bar{z}$ genau auf trigonometrische Polynome abbildet. Die Behauptung folgt also durch Anwendung von Korollar 4.11(b) auf $X:=S^{1} \subseteq \mathbb{C}$.\\
4.13 Korollar. Ist ( $X, d$ ) ein kompakter metrischer Raum, so ist $C(X)$ separabel.

Beweis. Da $X$ separabel ist (siehe Bemerkung 2.42), existieren $x_{n} \in X, n \in \mathbb{N}$ derart, dass $\left\{x_{n} \mid n \in \mathbb{N}\right\}$ in $X$ dicht liegt. Für $n \in \mathbb{N}$ sei nun $\varphi_{n} \in C(X)$ definiert durch $\varphi_{n}(x):=d\left(x, x_{n}\right)$. Sei ferner $A \subseteq C(X)$ die kleinste Teilalgebra, welche die Funktionen $\varphi_{n}, n \in \mathbb{N}$ enthält. Als Vektorraum ist $A=\operatorname{span} M$, wobei $M$ die Menge der endlichen Produkte $\varphi_{n_{1}} \cdots \varphi_{n_{m}}$ mit $m \in \mathbb{N}_{0}$ und $n_{l} \in \mathbb{N}, l=1, \ldots, m$ ist (leeres Produkt ist die Einsfunktion). Offensichtlich ist $M$ abzählbar. Ferner trennt die Algebra $A$ die Punkte von $X$ : Sind nämlich $x, y \in X$ mit $x \neq y$, so existiert $n \in \mathbb{N}$ mit $\varphi_{n}(x)=d\left(x, x_{n}\right)<\frac{d(x, y)}{2}$. Es folgt dann

$$
\varphi_{n}(y)=d\left(y, x_{n}\right)>d(y, x)-d\left(x, x_{n}\right)>\frac{d(x, y)}{2}>\varphi_{n}(x)
$$

Im Fall $\mathbb{K}=\mathbb{C}$ ist $A$ zudem abgeschlossen unter komplexer Konjugation, da die Funktionen $\varphi_{n}, n \in \mathbb{N}$ reellwertig sind. Also folgt mit dem Satz von Stone-Weierstraß, dass $A$ in $C(X)$ dicht liegt. Somit ist $M$ eine abzählbare totale Teilmenge von $C(X)$, und damit ist $C(X)$ separabel nach Satz 2.20.

\section*{3. $L^{p}$-Räume}
Stets sei im Folgenden ( $\Omega, \mathcal{M}, \mu$ ) ein (reeller) Maßraum. Ferner verwenden wir die Vereinbarung $\frac{r}{\infty}=0$ für $r \in \mathbb{R}$ und setzen $\overline{\mathbb{R}}:=\mathbb{R} \cup\{+\infty,-\infty\}$. Wir wiederholen zunächst zentrale Begriffe der Maß- und Integrationstheorie und Konvergenzsätze, welche größtenteils aus der Integrationstheorie bekannt sind.\\
4.14 Definition. Seien ( $Y, d$ ) ein metrischer Raum und $f: \Omega \rightarrow Y$ eine Funktion.

\begin{itemize}
  \item $f$ heißt $\mu$-messbar (oder einfach messbar), wenn für jede offene Teilmenge $U \subseteq Y$ die Menge $f^{-1}(U) \subseteq \Omega \mu$-messbar ist.
  \item Ist $f \mu$-messbar und nimmt $f$ nur endlich viele Werte an, so nennt man $f$ eine Stufenfunktion (auch Treppenfunktion oder Elementarfunktion).
\end{itemize}

Wir schreiben:

$$
\begin{aligned}
\mathcal{M}(\Omega, Y) & :=\{f: \Omega \rightarrow Y \mid f \mu \text {-messbar }\} \\
\mathcal{E}(\Omega, Y) & :=\{f: \Omega \rightarrow Y \mid f \text { Stufenfunktion }\}
\end{aligned}
$$

Im Fall $Y=\mathbb{K}(\mathbb{K}=\mathbb{R}$ oder $\mathbb{K}=\mathbb{C})$ schreiben wir auch $\mathcal{M}(\Omega)$ bzw. $\mathcal{E}(\Omega)$. Im Fall $Y=\mathbb{R}$ gilt insbesondere

$$
f \in \mathcal{M}(\Omega, \mathbb{R}) \Longleftrightarrow\{x \in \Omega \mid f(x)>c\} \text { ist messbar für alle } c \in \mathbb{R} \text {. }
$$

Wir werden im Folgenden (nicht-endliche) Funktionen $f: \Omega \rightarrow \overline{\mathbb{R}} \mu$-messbar nennen, wenn sie die Bedingung auf der rechten Seite von (4.8) erfüllen. Auch für solche Funktionen schreiben wir dann $f \in \mathcal{M}(\Omega, \mathbb{R})$ bzw. $f \in \mathcal{M}(\Omega)$.\\
4.15 Bemerkung. Für $f \in \mathcal{E}(\Omega)$ gilt $f=\sum_{y \in \mathbb{K}} y 1_{\{f=y\}}$, wobei wir hier und im Folgenden die Kurzschreibweise $\{f=y\}:=f^{-1}(y)=\{x \in \Omega \mid f(x)=y\} \subseteq \Omega$ verwenden.\\
4.16 Definition und Bemerkung. (a) Für eine $\mu$-messbare Funktion $f: \Omega \rightarrow \overline{\mathbb{R}}$ gilt offenbar:

$$
\int_{\Omega} f^{ \pm} \mathrm{d} \mu<\infty \quad \Longleftrightarrow \quad \int_{\Omega}|f| \mathrm{d} \mu<\infty
$$

Hier seien $f^{+}:=\max \{f, 0\}$ bzw. $f^{-}:=-\min \{f, 0\}$ der Positiv- bzw. Negativteil von $f$. Es gilt also u.a. $f^{ \pm} \geq 0$ und $f=f^{+}-f^{-}$. Die Funktion $f$ heißt $\mu$-integrierbar, falls die Bedingungen in (4.9) gelten. In diesem Fall setzt man

$$
\int_{\Omega} f \mathrm{~d} \mu:=\int_{\Omega} f^{+} \mathrm{d} \mu-\int_{\Omega} f^{-} \mathrm{d} \mu \in \mathbb{R} .
$$

(b) Eine $\mu$-messbare Funktion $f: \Omega \rightarrow \mathbb{C}$ heißt $\mu$-integrierbar, wenn $\operatorname{Re} f$ und $\operatorname{Im} f$ $\mu$-integrierbar sind. In diesem Fall setzt man

$$
\int_{\Omega} f \mathrm{~d} \mu:=\int_{\Omega} \operatorname{Re} f \mathrm{~d} \mu+\mathrm{i} \int_{\Omega} \operatorname{Im} f \mathrm{~d} \mu
$$

(c) Die Menge der $\mu$-integrierbaren $\mathbb{K}$-wertigen Funktionen bezeichnen wir mit $\mathfrak{L}^{1}(\Omega)$ oder $\mathfrak{L}^{1}(\Omega, \mathrm{~d} \mu)$.\\
4.17 Satz. (a) Ist $f: \Omega \rightarrow[0, \infty]$ messbar, so existiert eine Folge $\left(f_{k}\right)_{k}$ in $\mathcal{E}(\Omega)$ mit

$$
0 \leq f_{k}(x) \leq f_{k+1}(x) \leq f(x) \quad \text { für alle } x \in \Omega, k \in \mathbb{N}
$$

und $\lim _{k \rightarrow \infty} f_{k}(x)=f(x)$ für alle $x \in \Omega$.\\
(b) Ist $f \in \mathcal{M}(\Omega, \mathbb{C})$, so existiert eine Folge in $\mathcal{E}(\Omega, \mathbb{C})$ mit $\left|f_{k}\right| \leq|f|$ in $\Omega$ für alle $k$ und $f_{k} \rightarrow f$ für $k \rightarrow \infty$ punktweise in $\Omega$.\\
4.18 Bemerkung. Ist $\Omega=\bigcup_{k \in \mathbb{N}} M_{k}$ mit $\mu$-messbaren Teilmengen $M_{k}$ und $\mu\left(M_{k}\right)<\infty$ für alle $k \in \mathbb{N}$, so können die Folgen $\left(f_{k}\right)_{k}$ in Satz 4.17(a) und (b) derart gewählt werden, dass gilt:

$$
\mu\left(\left\{f_{k} \neq 0\right\}\right)<\infty \quad \text { für alle } k \in \mathbb{N} .
$$

Dies gilt insbesondere im Fall $\Omega=\mathbb{R}^{N}$ und $\mu=\lambda$ (Lebesguemaß).\\
4.19 Satz (von der monotonen Konvergenz). (a) Seien $f_{k} \in \mathcal{M}(\Omega, \mathbb{R}), k \in \mathbb{N}$ nichtnegative Funktionen mit $f_{k} \leq f_{k+1}$ f. ü. auf $\Omega$ für alle $k$, und sei $f \in \mathcal{M}(\Omega, \mathbb{R})$ f.ü. definiert durch $f(x):=\lim _{k \rightarrow \infty} f_{k}(x)$. Dann gilt: $\int_{\Omega} f \mathrm{~d} \mu=\lim _{k \rightarrow \infty} \int_{\Omega} f_{k} \mathrm{~d} \mu$.\\
(b) Seien $f_{k} \in \mathfrak{L}^{1}(\Omega) \mathbb{R}$-wertige Funktionen für $k \in \mathbb{N}$ mit $f_{k} \leq f_{k+1}$ f.ü. auf $\Omega$ für alle $k$, und sei $f \in \mathcal{M}(\Omega, \mathbb{R})$ f.ü. definiert durch $f(x):=\lim _{k \rightarrow \infty} f_{k}(x)$. Dann gilt: Ist die Folge der Integrale $\int_{\Omega} f_{k} \mathrm{~d} \mu$ beschränkt, so ist $f$ integrierbar und $\int_{\Omega} f \mathrm{~d} \mu=$ $\lim _{k \rightarrow \infty} \int_{\Omega} f_{k} \mathrm{~d} \mu$.\\
4.20 Lemma (von Fatou). Sei $\left(f_{k}\right)_{k} \subseteq \mathfrak{L}^{1}(\Omega)$ eine Folge $\mathbb{R}$-wertiger Funktionen.\\
(a) Ist $\liminf _{k \rightarrow \infty} \int_{\Omega} f_{k} \mathrm{~d} \mu<\infty$ und existiert $g \in \mathfrak{L}^{1}(\Omega, \mathbb{R})$ mit $f_{k} \geq g$ f. ü. auf $\Omega$ für alle $k \in \mathbb{N}$, so ist $f:=\liminf _{k \rightarrow \infty} f_{k}$ integrierbar und $\int_{\Omega} f \mathrm{~d} \mu \leq \liminf _{k \rightarrow \infty} \int_{\Omega} f_{k} \mathrm{~d} \mu$.\\
(b) Ist $\lim \sup _{k \rightarrow \infty} \int_{\Omega} f_{k} \mathrm{~d} \mu>-\infty$ und existiert $g \in \mathfrak{L}^{1}(\Omega, \mathbb{R})$ mit $f_{k} \leq g$ f.ü. auf $\Omega$ für alle $k \in \mathbb{N}$, so ist $f:=\limsup _{k \rightarrow \infty} f_{k}$ integrierbar und $\int_{\Omega} f \mathrm{~d} \mu \geq$ $\lim \sup _{k \rightarrow \infty} \int_{\Omega} f_{k} \mathrm{~d} \mu$.\\
4.21 Satz (von der majorisierten Konvergenz (Satz von Lebesgue)). Seien $f_{k} \in \mathfrak{L}^{1}(\Omega)$, $k \in \mathbb{N}$ derart, dass der punktweise Grenzwert $f:=\lim _{k \rightarrow \infty} f_{k}$ punktweise f.ü. auf $\Omega$ existiert. Ferner existiere eine $\mathbb{R}$-wertige Funktion $g \in \mathfrak{L}^{1}(\Omega)$ derart, dass für alle $k \in \mathbb{N}$ die Ungleichung $\left|f_{k}\right| \leq g$ f.ü. auf $\Omega$ gilt. Dann ist $f \in \mathfrak{L}^{1}(\Omega)$, und es gilt $\int_{\Omega} f \mathrm{~d} \mu=\lim _{k \rightarrow \infty} \int_{\Omega} f_{k} \mathrm{~d} \mu$.

Jetzt kommen wir zur Definition der $L^{p}$-Räume.\\
4.22 Definition und Bemerkung. Für eine $\mu$-messbare Funktion $u: \Omega \rightarrow \overline{\mathbb{R}}$ sei

$$
\underset{\Omega}{\operatorname{ess} \sup } u:=\inf \{t \in \mathbb{R} \mid u \leq t \text { f.ü. auf } \Omega\} \quad \text { (wesentliches Supremum von } u \text { ). }
$$

Man sieht leicht, dass stets $u \leq \operatorname{ess} \sup _{\Omega} u$ f.ü. auf $\Omega$ gilt.\\
4.23 Satz und Definition. Sei $1 \leq p \leq \infty$. Die Menge $\mathfrak{L}^{p}(\Omega)$ aller messbaren Funktionen $u: \Omega \rightarrow \mathbb{K}$ mit

$$
\begin{cases}\int_{\Omega}|u|^{p} \mathrm{~d} \mu<\infty, & \text { im Fall } p<\infty \\ \operatorname{ess}_{\sup }|u|<\infty, & \text { im Fall } p=\infty\end{cases}
$$

ist ein $\mathbb{K}$-Vektorraum, und durch

$$
\|u\|_{p}:= \begin{cases}\left(\int_{\Omega}|u|^{p} \mathrm{~d} \mu\right)^{\frac{1}{p}}, & \text { im Fall } p<\infty \\ \operatorname{ess}_{\sup }|u|, & \text { im Fall } p=\infty\end{cases}
$$

wird eine Halbnorm auf $\mathfrak{L}^{p}(\Omega)$ definiert.\\
Beweis. Die Vektorraumeigenschaften sieht man sehr ähnlich wie in Definition und Satz 2.7. Sei nun

$$
A_{p}:= \begin{cases}\left\{\left.u \in \mathfrak{L}^{p}(\Omega)\left|\int_{\Omega}\right| u\right|^{p} \mathrm{~d} \mu \leq 1\right\} & \text { im Fall } p<\infty \\ \left\{u \in \mathfrak{L}^{p}(\Omega)|\underset{\Omega}{\operatorname{ess} \sup }| u \mid \leq 1\right\} & \text { im Fall } p=\infty\end{cases}
$$

Offensichtlich ist $A_{p}$ dann absolut konvex und absorbierend in $\mathfrak{L}^{p}(\Omega)$, und gemäß Satz 2.6 definiert

$$
u \mapsto \inf \left\{\begin{array}{l|l}
t>0 & \frac{u}{t} \in A_{p}
\end{array}\right\}=\|u\|_{p}
$$

eine Halbnorm auf $\mathfrak{L}^{p}(\Omega)$.\\
4.24 Satz (Höldersche Ungleichung). Seien $p, q \in[1, \infty]$ konjugierte Exponenten, d.h. es gelte $\frac{1}{p}+\frac{1}{q}=1$. Seien ferner $u \in \mathfrak{L}^{p}(\Omega)$ und $v \in \mathfrak{L}^{q}(\Omega)$. Dann gilt

$$
\int_{\Omega}|u v| \mathrm{d} \mu \leq\|u\|_{p}\|v\|_{q}
$$

Beweis. Die Ungleichung ist klar im Fall $p \in\{1, \infty\}$. Sei also nun $1<p<\infty$, also auch $1<q<\infty$. Ohne Einschränkung können wir $\|u\|_{p},\|v\|_{q}>0$ annehmen, denn andernfalls ist $u v=0$ f.ü. in $\Omega$ und die Ungleichung ist trivialerweise erfüllt. Mit der Youngschen Ungleichung (vgl. Bemerkung 2.8(b)) folgt dann mit $s=\|u\|_{p}, t=\|v\|_{q}$ :

$$
\begin{aligned}
\int_{\Omega}|u v| \mathrm{d} \mu=s t \int_{\Omega}\left|\frac{u}{s}\right|\left|\frac{v}{t}\right| \mathrm{d} \mu \leq s t \int_{\Omega} & \left(\frac{1}{p}\left|\frac{u}{s}\right|^{p}+\frac{1}{q}\left|\frac{v}{t}\right|^{q}\right) \mathrm{d} \mu \\
& =s t\left(\left.\frac{1}{p}| | \frac{u}{s}\right|_{p} ^{p}+\left.\frac{1}{q}| | \frac{v}{t}\right|_{q} ^{q}\right)=s t\left(\frac{1}{p}+\frac{1}{q}\right)=s t
\end{aligned}
$$

wie behauptet.\\
4.25 Bemerkung (verallgemeinerte Höldersche Ungleichung). Seien $p_{1}, \ldots, p_{n} \in[1, \infty]$ mit $\sum_{i=1}^{n} \frac{1}{p_{i}}=\frac{1}{p}$ für ein $p \in[1, \infty]$. Seien ferner $u_{i} \in \mathfrak{L}^{p_{i}}(\Omega), i=1, \ldots, n$ gegeben. Dann gilt

$$
u:=u_{1} \cdots \cdots u_{n} \in \mathfrak{L}^{p}(\Omega) \quad \text { und } \quad\|u\|_{p} \leq\left\|u_{1}\right\|_{p_{1}} \cdots \cdots\left\|u_{n}\right\|_{p_{n}}
$$

Dies folgt aus Satz 4.24 per Induktion (Übung).\\
4.26 Bemerkung und Beispiel. Sei $1 \leq p<q \leq \infty$.\\
(a) Ist $u \in \mathfrak{L}^{p}(\Omega) \cap \mathfrak{L}^{q}(\Omega)$, so ist auch $u \in \mathfrak{L}^{s}(\Omega)$ für $s \in(p, q)$, wobei gilt:

$$
\|u\|_{s} \leq\|u\|_{p}^{\alpha}\|u\|_{q}^{1-\alpha} \quad \text { mit } \alpha:=\frac{p(q-s)}{s(q-p)}
$$

Merkregel: $\frac{1}{s}=\frac{\alpha}{p}+\frac{1-\alpha}{q}$.\\
(b) Ist $\mu(\Omega)<\infty$, so ist $\mathfrak{L}^{q}(\Omega) \subseteq \mathfrak{L}^{p}(\Omega)$ und

$$
\|u\|_{p} \leq \mu(\Omega)^{\frac{1}{p}-\frac{1}{q}}\|u\|_{q} \quad \text { für alle } u \in \mathfrak{L}^{q}(\Omega) .
$$

(c) Ist $\mu=\lambda$ das Lebesguemaß auf $\mathbb{R}^{N}$, so gilt

$$
\mathfrak{L}^{p}\left(\mathbb{R}^{N}\right) \backslash \mathfrak{L}^{q}\left(\mathbb{R}^{N}\right) \neq \varnothing \neq \mathfrak{L}^{q}\left(\mathbb{R}^{N}\right) \backslash \mathfrak{L}^{p}\left(\mathbb{R}^{N}\right)
$$

Beweis. (a) Nach Voraussetzung ist $|u|^{\alpha} \in \mathfrak{L}^{p / \alpha}(\Omega)$ und $|u|^{1-\alpha} \in \mathfrak{L}^{q /(1-\alpha)}(\Omega)$. Ferner gilt

$$
\frac{1}{p / \alpha}+\frac{1}{q /(1-\alpha)}=\frac{\alpha}{p}+\frac{1-\alpha}{q}=\frac{1}{s}
$$

nach Voraussetzung. Gemäß Bemerkung 4.25 ist also $|u|=|u|^{\alpha}|u|^{1-\alpha} \in \mathfrak{L}^{s}(\Omega)$, also

$$
u \in \mathfrak{L}^{s}(\Omega) \quad \text { mit } \quad\|u\|_{s} \leq\left\||u|^{\alpha}\right\|_{\frac{p}{\alpha}}\left\||u|^{1-\alpha}\right\|_{\frac{q}{1-\alpha}}=\|u\|_{p}^{\alpha}\|u\|_{q}^{1-\alpha}
$$

(b) Sei $u \in \mathfrak{L}^{q}(\Omega)$. Wir schreiben $u=u \cdot 1_{\Omega}$. Mit $\beta:=\left(\frac{1}{p}-\frac{1}{q}\right)^{-1}$ ist dann $1_{\Omega} \in \mathfrak{L}^{\beta}$ mit

$$
\left\|1_{\Omega}\right\|_{\beta}=\left(\int_{\Omega} 1 \mathrm{~d} \mu\right)^{1 / \beta}=(\mu(\Omega))^{\frac{1}{p}-\frac{1}{q}}
$$

Ferner ist $\frac{1}{q}+\frac{1}{\beta}=\frac{1}{p}$, und somit folgt mit Bemerkung 4.25:

$$
u \in \mathfrak{L}^{p}(\Omega) \quad \text { mit } \quad\|u\|_{p} \leq\left\|1_{\Omega}\right\|_{\beta}\|u\|_{q}=(\mu(\Omega))^{\frac{1}{p}-\frac{1}{q}}\|u\|_{q}
$$

(c) Übung.

Ende des Inhalts für 2024-05-14\\
4.27 Definition und Satz ( $L^{p}$-Raum). Sei $Z:=\{u \in \mathcal{M}(\Omega, \mathbb{K}) \mid u=0$ f.ü. auf $\Omega\}$, und sei $1 \leq p \leq \infty$. Dann ist $Z \subseteq \mathfrak{L}^{p}(\Omega)$ ein Unterraum, und der Faktorraum $L^{p}(\Omega)=$ $L^{p}(\Omega, \mathbb{K}):=\mathfrak{L}^{p}(\Omega) / Z$ ist ein normierter Raum mit Norm $\|\cdot\|_{p}$ gegeben durch

$$
\|u+Z\|_{p}:=\|u\|_{p} \quad \text { für } u \in \mathfrak{L}^{p}(\Omega) .
$$

Beweis. Es ist klar, dass $Z$ ein linearer Unterraum von $\mathfrak{L}^{p}$ ist. Die Integrationstheorie\\
liefert für $u \in \mathfrak{L}^{p}(\Omega)$, dass $\|u\|_{p}=0$ genau dann gilt, wenn $u$ in $Z$ liegt. Daraus folgt sofort, dass durch die Definition eine Norm gegeben ist.\\
4.28 Bemerkung (u.a. zur Notation). (a) Man schreibt anstelle von $u+Z$ oder $[u]$ (Äquivalenzklasse von $u$ ) oft nur $u \in L^{p}(\Omega)$. Man betrachtet die Elemente von $L^{p}(\Omega)$ also als Funktionen, die nur bis auf Abänderung auf Nullmengen eindeutig definiert sind.\\
(b) Angenommen, $\Omega$ ist ein metrischer Raum mit der Eigenschaft

$$
\text { Ist } V \subseteq \Omega \text { eine Nullmenge, so ist } \Omega \backslash V \text { dicht in } \Omega
$$

(z.B. $\Omega \subseteq \mathbb{R}^{N}$ offen mit Lebesgue-Maß). Sind dann $u, v \in \mathfrak{L}^{p}(\Omega)$ stetige Funktionen mit $u=v$ f.ü. auf $\Omega$, so gilt $u=v$ auf ganz $\Omega$. Es folgt also, dass dann die Abbildung

$$
\mathfrak{L}^{p}(\Omega) \cap C(\Omega) \rightarrow L^{p}(\Omega), \quad u \mapsto u+Z
$$

injektiv ist.\\
4.29 Satz. Sei $1 \leq p \leq \infty$.\\
(a) $L^{p}(\Omega)$ ist ein Banachraum.\\
(b) Seien $u, u_{k} \in L^{p}(\Omega), k \in \mathbb{N}$ gegeben mit $\left\|u_{k}-u\right\|_{p} \rightarrow 0$ für $k \rightarrow \infty$. Dann existiert eine Teilfolge $\left(u_{k_{j}}\right)_{j}$, welche punktweise f. ü. auf $\Omega$ gegen u konvergiert.

Beweis. (a) Gemäß Satz 2.37 reicht es, zu zeigen, dass jede absolut konvergente Reihe in $L^{p}(\Omega)$ auch konvergiert. Seien also $u_{k} \in L^{p}(\Omega), k \in \mathbb{N}$ gegeben mit

$$
C_{m}:=\sum_{k=m}^{\infty}\left\|u_{k}\right\|_{p} \rightarrow 0 \quad \text { für } m \rightarrow \infty
$$

Wir betrachten im Folgenden die Hilfsfunktionen

$$
h_{m}:=\sum_{k=m}^{\infty}\left|u_{k}\right|: \Omega \rightarrow[0, \infty], \quad m \in \mathbb{N}
$$

Wir zeigen zunächst:

$$
h_{m} \in L^{p}(\Omega) \quad \text { und } \quad\left\|h_{m}\right\|_{p} \leq C_{m} \quad \text { für } m \in \mathbb{N} .
$$

Für festes $m \in \mathbb{N}$ und $j>m$ liefert die Dreiecksungleichung

$$
\left\|\sum_{k=m}^{j}\left|u_{k}\right|\right\|_{p} \leq \sum_{k=m}^{j}\left\|u_{k}\right\|_{p} \leq C_{m}
$$

Im Fall $p<\infty$ folgt dann aus dem Satz von der monotonen Konvergenz

$$
\int_{\Omega}\left|h_{m}\right|^{p} \mathrm{~d} \mu=\lim _{j \rightarrow \infty} \int_{\Omega}\left(\sum_{k=m}^{j}\left|u_{k}\right|\right)^{p} \mathrm{~d} \mu=\lim _{j \rightarrow \infty}\left\|\sum_{k=m}^{j}\left|u_{k}\right|\right\|_{p}^{p} \leq C_{m}^{p}
$$

und somit (4.10). Im Fall $p=\infty$ folgt ebenfalls

$$
h_{m}=\lim _{j \rightarrow \infty} \sum_{k=m}^{j}\left|u_{k}\right| \leq C_{m} \quad \text { f.ü. in } \Omega
$$

und somit wiederum (4.10). Hier haben wir verwendet, dass für alle $k \in \mathbb{N}$ gilt: $\left|u_{k}\right| \leq$ $\left\|u_{k}\right\|_{\infty}$ f.ü. und dass die abzählbare Vereinigung von Nullmengen wieder eine Nullmenge ist. Aus (4.10) folgt insbesondere:

$$
h_{1}=\sum_{k=1}^{\infty}\left|u_{k}\right|<\infty \quad \text { f.ü. auf } \Omega,
$$

denn im Fall von $p<\infty$ gilt $\left|h_{1}\right|^{p} \in L^{1}(\Omega)$. Somit existiert $u(x):=\sum_{k=1}^{\infty} u_{k}(x) \in \mathbb{K}$ für fast alle $x \in \Omega$. Sei $u$ trivial (also durch 0) auf ganz $\Omega$ fortgesetzt; dann ist $u \mu$-messbar und $|u| \leq h_{1}$ auf $\Omega$. Gemäß (4.10) gilt also $u \in L^{p}(\Omega)$ mit $\|u\|_{p} \leq C_{1}$. Sei schließlich $g_{m}:=u-\sum_{k=1}^{m-1} u_{k}$ für $m \geq 2$. Dann gilt $\left|g_{m}\right| \leq h_{m}$ f. ü. in $\Omega$, und wegen (4.10) folgt

$$
\left\|g_{m}\right\|_{p} \leq C_{m} \rightarrow 0 \quad \text { für } m \rightarrow \infty .
$$

Somit konvergiert die Reihe $\sum_{k=1}^{\infty} u_{k}$ bzgl. $\|\cdot\|_{p}$ gegen $u$.\\
(b) Wähle sukzessive $k_{j} \in \mathbb{N}$ mit $k_{j+1}>k_{j}$ und derart, dass für alle $j \in \mathbb{N}$ gilt:

$$
\left\|v_{j}\right\|_{p} \leq 2^{-j} \quad \text { für } v_{j}:=u_{k_{j+1}}-u_{k_{j}} \in L^{p}(\Omega) .
$$

Insbesondere konvergiert die Reihe $\sum_{j \in \mathbb{N}} v_{j}$ absolut in $L^{p}(\Omega)$. Wie im Beweis von (a) folgt daher, dass sie sowohl bzgl. $\|\cdot\|_{p}$ als auch punktweise f.ü. auf $\Omega$ gegen eine Funktion $v \in L^{p}(\Omega)$ konvergiert. Nach Voraussetzung gilt aber auch

$$
\left\|u-u_{k_{1}}-\sum_{j=1}^{m} v_{j}\right\|_{p}=\left\|u-u_{k_{m+1}}\right\|_{p} \rightarrow 0 \quad \text { für } m \rightarrow \infty
$$

und somit folgt $v=u-u_{k_{1}}$. Es gilt also

$$
u_{k_{m}}=\sum_{j=1}^{m-1} v_{j}+u_{k_{1}} \rightarrow v+u_{k_{1}}=u \quad \text { für } m \rightarrow \infty \text { punktweise fast überall auf } \Omega,
$$

d.h. die Teilfolge ( $u_{k_{m}}$ ) hat die gewünschte Eigenschaft.\\
4.30 Beispiel (wandernder Buckel). Der Übergang zu einer Teilfolge in Satz 4.29(b) ist i.A. notwendig. Um dies einzusehen, betrachten wir $1 \leq p<\infty$ und $f_{n}:[0,1] \rightarrow \mathbb{R}$, $n \in \mathbb{N}$ definiert durch

$$
f_{n}(t):= \begin{cases}1 & \text { für } t \in\left[k 2^{-\nu},(k+1) 2^{-\nu}\right] ; \\ 0 & \text { sonst. }\end{cases}
$$

Hier seien $k, \nu \in \mathbb{N}_{0}$ die jeweils eindeutig bestimmten Zahlen mit $n=2^{\nu}+k$ und $k<2^{\nu}$. Dann ist $f_{n} \in L^{p}([0,1])$ für alle $n$ und $\lim _{n \rightarrow \infty}\left\|f_{n}\right\|_{p}=0$; aber in keinem Punkt $x \in[0,1]$ konvergiert die Folge $\left(f_{n}(x)\right)_{n}$ gegen 0 .\\
4.31 Definition und Satz. Auf $L^{2}(\Omega)$ ist ein Skalarprodukt $\langle\cdot, \cdot\rangle$ definiert durch

$$
\langle f, g\rangle=\int_{\Omega} f \bar{g} \mathrm{~d} \mu \quad \text { für } f, g \in L^{2}(\Omega)
$$

(Im Fall $\mathbb{K}=\mathbb{R}$ erübrigt sich hierbei natürlich die komplexe Konjugation). Dabei ist $\|\cdot\|_{2}$ die von $\langle\cdot, \cdot\rangle$ induzierte Norm. Nach Satz 4.29 ist $L^{2}(\Omega)$ also ein Hilbertraum.

Beweis. Die Skalarprodukteigenschaften rechnet man problemlos nach.\\
4.32 Satz. Sei $1 \leq p<\infty$, und sei $u \in L^{p}(\Omega)$. Dann existiert eine Folge $\left(u_{k}\right)_{k} \subseteq \mathcal{E}(\Omega)$ mit $\left|u_{k}\right| \leq|u|$ in $\Omega$ und $\mu\left(\left\{u_{k} \neq 0\right\}\right)<\infty$ für alle $k \in \mathbb{N}$ sowie $\lim _{k \rightarrow \infty}\left\|u_{k}-u\right\|_{p}=0$.

Beweis. Nach Satz 4.17 existieren $u_{k} \in \mathcal{E}(\Omega), k \in \mathbb{N}$, mit $\left|u_{k}\right| \leq|u|$ für alle $k$ und $u_{k} \rightarrow u$ punktweise. Für festes $k$ existiert dabei $c=c(k)>0$ mit $\left|u_{k}\right| \geq c$ auf der Menge $\left\{u_{k} \neq 0\right\}$, und somit folgt

$$
\mu\left(\left\{u_{k} \neq 0\right\}\right) c^{p} \leq \int_{\Omega}\left|u_{k}\right|^{p} \mathrm{~d} \mu \leq \int_{\Omega}|u|^{p} \mathrm{~d} \mu<\infty
$$

Dies zeigt $\mu\left(\left\{u_{k} \neq 0\right\}\right)<\infty$. Sei nun $g_{k}:=\left|u_{k}-u\right|^{p}$ für $k \in \mathbb{N}$. Dann konvergiert die Folge $\left(g_{k}\right)_{k}$ punktweise gegen 0 , wobei

$$
\left|g_{k}\right| \leq\left(\left|u_{k}\right|+|u|\right)^{p} \leq 2^{p}|u|^{p} \quad \text { auf } \Omega \text { für alle } k \in \mathbb{N} \text { gilt. }
$$

Mit dem Satz von Lebesgue (siehe Satz 4.21) folgt daher

$$
\lim _{k \rightarrow \infty} \int_{\Omega} g_{k} \mathrm{~d} \mu=0
$$

und dies liefert $\lim _{k \rightarrow \infty}\left\|u_{k}-u\right\|_{p}=0$.\\
Bisher haben wir noch nicht geklärt, unter welchem Umständen die Räume $L^{p}(\Omega)$ separabel sind. Ist darüber hinaus $\Omega$ ein metrischer Raum, so ist es in Anwendung oft sehr nützlich, zu wissen, inwieweit stetige Funktionen in $L^{p}(\Omega)$ dicht liegen. Diese Fragen wollen wir im Folgenden erörtern. Dazu ist der folgende Satz sehr nützlich; des Weiteren benötigen wir einige Sachverhalte über Borel- und Radonmaße.\\
4.33 Satz (von Egorov (und Severini)). Sei Y ein separabler metrischer Raum, und seien $u_{n}: \Omega \rightarrow Y$ messbare Funktionen für $n \in \mathbb{N}$, welche punktweise f.ü. gegen eine Funktion $u: \Omega \rightarrow Y$ konvergieren. Ist $\mu(\Omega)<\infty$, so existiert zu jedem $\varepsilon>0$ eine messbare Teilmenge $A \subseteq \Omega$ mit $\mu(\Omega \backslash A)<\varepsilon$ und derart, dass $\left(u_{n}\right)_{n}$ auf A gleichmäßig gegen $u$ konvergiert.

Beweis. Sei $\varepsilon>0$, und für $n, k \in \mathbb{N}$ sei

$$
E_{n, k}:=\bigcup_{m=n}^{\infty}\left\{x \in \Omega \left\lvert\, d\left(u_{m}(x), u(x)\right) \geq \frac{1}{k}\right.\right\}
$$

Dann gilt $E_{n+1, k} \subseteq E_{n, k}$ für alle $n, k \in \mathbb{N}$. Da ferner $\left(u_{n}\right)_{n}$ punktweise f.ü. gegen $u$ konvergiert und $\mu\left(E_{1, k}\right) \leq \mu(\Omega)<\infty$ gilt, folgt

$$
\lim _{n \rightarrow \infty} \mu\left(E_{n, k}\right)=\mu\left(\bigcap_{n=1}^{\infty} E_{n, k}\right)=0 \quad \text { für alle } k \in \mathbb{N} \text {. }
$$

Somit existiert zu jedem $k \in \mathbb{N}$ ein $n_{k} \in \mathbb{N}$ mit $\mu\left(E_{n_{k}, k}\right)<\varepsilon 2^{-k}$. Sei nun

$$
A:=\Omega \backslash \bigcup_{k \in \mathbb{N}} E_{n_{k}, k}
$$

Für $x \in A$ und $k \in \mathbb{N}$ gilt dann $x \notin E_{n_{k}, k}$, also

$$
d\left(u_{m}(x), u(x)\right)<\frac{1}{k} \quad \text { für } m \geq n_{k} .
$$

Es folgt, dass die Folge $\left(u_{n}\right)_{n}$ auf $A$ gleichmäßig gegen $u$ konvergiert. Ferner gilt

$$
\mu(\Omega \backslash A) \leq \sum_{k=1}^{\infty} \mu\left(E_{n_{k}, k}\right)<\varepsilon
$$

Bemerkung: Die Separabilität von $Y$ wird im obigen Beweis benötigt, um die Messbarkeit der Mengen $E_{n, k}$ sicherzustellen.

\subsection*{Approximation in $L^{p}$ durch stetige Funktionen}
Wir benötigen zunächst einige ergänzende Begriffsbildungen zur Kompaktheit. Sei im Folgenden ( $\Omega, d$ ) ein metrischer Raum.\\
4.34 Definition. Seien $E$ ein normierter Raum und $U \subseteq E$ eine Teilmenge mit $0 \in U$. Mit $C_{\mathrm{c}}(\Omega, U)$ bezeichnen wir die Menge der stetigen Funktionen $u: \Omega \rightarrow U$ derart, dass der Träger

$$
\operatorname{supp} u:=\overline{\{x \in \Omega \mid u(x) \neq 0\}} \subseteq \Omega
$$

kompakt ist. Im Fall $U=E$ ist diese Menge ein (i.A. nicht abgeschlossener!) Unterraum des normierten Raums $C_{\mathrm{b}}(\Omega, E)$. Kurzschreibweise: $C_{\mathrm{c}}(\Omega)$ anstelle von $C_{\mathrm{c}}(\Omega, \mathbb{K})$.\\
4.35 Definition. $\Omega$ heißt\\
(a) lokal kompakt, wenn jeder Punkt $x \in \Omega$ eine kompakte Umgebung besitzt.\\
(b) $\sigma$-kompakt, wenn kompakte Mengen $K_{j} \subseteq \Omega, j \in \mathbb{N}$ existieren mit $\Omega=\bigcup_{j=1}^{\infty} K_{j}$.\\
4.36 Bemerkungen und Beispiele. (a) Ist $\Omega$ lokal kompakt und $K \subseteq \Omega$ kompakt, so besitzt $K$ eine kompakte Umgebung. Ferner existiert in diesem Fall eine Funktion $u \in C_{\mathrm{c}}(\Omega,[0,1])$ mit $u \equiv 1$ in einer offenen Umgebung von $K$.\\
(b) Ist $\Omega \sigma$-kompakt, so ist $\Omega$ separabel: Sind nämlich $K_{j} \subseteq \Omega, j \in \mathbb{N}$ kompakt mit $\Omega:=\bigcup_{j=1}^{\infty} K_{j}$, so sind die Mengen $K_{j}$ separabel nach Bemerkung 2.42 und somit auch $\Omega$ (als abzählbare Vereinigung separabler Mengen).\\
(c) Ein normierter Raum $E$ ist lokal kompakt genau dann, wenn $E$ endlichdimensional ist (vgl. Satz 2.45).\\
(d) Jeder endlichdimensionale normierte Raum ist $\sigma$-kompakt.\\
(e) Der Raum $\mathcal{F} \subseteq \ell^{\infty}$ der finiten $\mathbb{K}$-wertigen Folgen (versehen mit der Norm $\|\cdot\|_{\infty}$ ) ist $\sigma$-kompakt, denn es ist $\mathcal{F}=\bigcup_{j=1}^{\infty} K_{j}$, wobei $K_{j} \subseteq \mathcal{F}$ die kompakte Teilmenge der Folgen $x=\left(x_{\ell}\right)_{\ell} \in \mathcal{F}$ mit $\|x\|_{\infty} \leq j$ und $x_{\ell}=0$ für $\ell \geq j$ sei. $\mathcal{F}$ ist aber nicht lokal kompakt, da $\operatorname{dim} \mathcal{F}=\infty$ gilt.\\
4.37 Satz. Sei $\Omega$ lokal kompakt und $\sigma$-kompakt. Dann gilt:\\
(a) Es existiert eine Folge kompakter Teilmengen $M_{j}, j \in \mathbb{N}$ von $\Omega$ derart, dass jede kompakte Menge $K \subseteq \Omega$ in einer der Mengen $M_{j}$ enthalten ist.\\
(b) Der Raum ( $C_{\mathrm{c}}(\Omega),\|\cdot\|_{\infty}$ ) ist separabel.

Beweis. (a) Seien $K_{j} \subseteq \Omega, j \in \mathbb{N}$ kompakt mit $\Omega:=\bigcup_{j=1}^{\infty} K_{j}$. Gemäß Bemerkungen und Beispiele 4.36(a) existiert für jedes $j \in \mathbb{N}$ eine kompakte Umgebung $L_{j}$ von $K_{j}$. Setze $M_{j}:=\bigcup_{i=1}^{j} L_{i}$ für $j \in \mathbb{N}$. Dann bilden die Mengen $U_{j}:=\stackrel{\circ}{M}_{j}, j \in \mathbb{N}$ eine offene Überdeckung von $\Omega$ mit $U_{j} \subseteq U_{j+1}$ für $j \in \mathbb{N}$. Für jede kompakte Teilmenge $K \subseteq \Omega$ existiert somit ein $j \in \mathbb{N}$ mit $K \subseteq U_{j} \subseteq M_{j}$, wie behauptet.\\
(b) Sei $M_{j}, j \in \mathbb{N}$ wie in (a), und sei

$$
C_{j}:=\left\{f \in C_{\mathrm{c}}(\Omega) \mid \operatorname{supp} f \subseteq M_{j}\right\} \quad \text { für } j \in \mathbb{N} .
$$

Dabei ist $C_{j}$ also isometrisch isomorph zu einem Unterraum von $C\left(M_{j}\right)$ vermöge der isometrischen Einbettung $C_{j} \rightarrow C\left(M_{j}\right),\left.f \mapsto f\right|_{M_{j}}$. Gemäß Korollar 4.13 ist dabei $C\left(M_{j}\right)$ separabel bzgl. $\|\cdot\|_{\infty}$ und wegen Satz 2.18 auch $C_{j}$. Ferner ist $C_{\mathrm{c}}(\Omega)=\bigcup_{j \in \mathbb{N}} C_{j}$ wegen (a), und somit ist $C_{\mathrm{c}}(\Omega)$ ebenfalls separabel.\\
4.38 Bemerkung. In Satz 4.37 kann die Bedingung der lokalen Kompaktheit nicht weggelassen werden, wie ein mit Hilfe von Bemerkungen und Beispiele 4.36(e) einfach zu konstruierendes Gegenbeispiel zeigt (Übung).\\
4.39 Definition. (a) Für ein Mengenteilsystem $\mathcal{O} \subseteq \mathcal{P}(\Omega)$ definiert man $\langle\mathcal{O}\rangle^{\sigma}$ als den Schnitt aller $\sigma$-Algebren, welche $\mathcal{O}$ enthalten. Man nennt $\langle\mathcal{O}\rangle^{\sigma}$ dann die von $\mathcal{O}$ erzeugte $\sigma$-Algebra.\\
(b) Ist $\mathcal{O} \subseteq \mathcal{P}(\Omega)$ das System aller offenen Teilmengen von $\Omega$, so nennt man $\mathcal{B}(\Omega):=$ $\langle\mathcal{O}\rangle^{\sigma}$ die Borelalgebra von $\Omega$, und die Elemente von $\mathcal{B}(\Omega)$ heißen Borelmengen.\\
(c) Ist $\mu$ ein Maß auf einer $\sigma$-Algebra $\mathcal{M}_{\mu}$ auf $\Omega$ mit $\mathcal{B}(\Omega) \subseteq \mathcal{M}_{\mu}$, so nent man $\mu$ ein Borelma $\beta$ auf $\Omega$. Vorsicht: In Teilen der Literatur wird bei einem Borelmaß $\mu$ noch zusätzlich die Eigenschaft verlangt, dass $\mu$ lokal endlich ist, d.h. das jeder Punkt in $\Omega$ eine offene Umgebung $U$ besitzt mit $\mu(U)<\infty$ (vgl. Definition 4.41(c)).\\
4.40 Beispiele (und Bemerkung zur Notation). (a) Bekanntermaßen ist (nach Konstruktion, s. Vorlesung Integrationstheorie) das Lebesguemaß auf $\mathbb{R}^{N}$ ein Borelmaß. Ferner ist z.B. das Diracmaß $\delta_{x}$ zu $x \in \Omega$, gegeben durch

$$
\delta_{x}=\left\{\begin{array}{ll}
1, & x \in A, \\
0, & x \notin A
\end{array} \quad \text { für } A \subseteq \Omega\right.
$$

ein Borelmaß.\\
(b) Ist im Folgenden von einem Borelmaß $\mu$ auf $\Omega$ die Rede, so betrachten wir $\mu$ stets auf einer gegebenen, zugehörigen $\sigma$-Algebra $\mathcal{M}_{\mu} \supseteq \mathcal{B}(\Omega)$, deren Elemente wir im Folgenden $\mu$-messbar nennen.\\
4.41 Definition. Sei $\mu$ ein Borelmaß auf $\Omega$.\\
(a) $\mu$ heißt von außen regulär, wenn für jede $\mu$-messbare Teilmenge $A \subseteq \Omega$ gilt:

$$
\mu(A)=\inf \{\mu(U) \mid U \text { offen, } A \subseteq U\} .
$$

(b) $\mu$ heißt von innen regulär, wenn für jede $\mu$-messbare Teilmenge $A \subseteq \Omega$ gilt:

$$
\mu(A)=\sup \{\mu(K) \mid K \text { kompakt, } K \subseteq A\} .
$$

(c) $\mu$ heißt lokal endlich, wenn jeder Punkt $x \in \Omega$ eine offene Umgebung $U$ besitzt mit $\mu(U)<\infty$.\\
(d) $\mu$ heißt Radon-Maß auf $\Omega$, wenn $\mu$ von innen regulär und lokal endlich ist.

Bemerkung: Ist $\Omega$ lokal kompakt, so ist (c) offensichtlich äquivalent zur Eigenschaft

$$
\mu(K)<\infty \quad \text { für jede kompakte Teilmenge } K \subseteq \Omega \text {. }
$$

4.42 Satz. Sei $\Omega$-kompakt, und sei $\mu$ ein von außen reguläres und lokal endliches Borelmaß auf $\Omega$. Dann gilt:\\
(a) $\mu$ ist auch von innen regulär und somit ein Radon-Maß.\\
(b) Für jede $\mu$-messbare Teilmenge $A \subseteq \Omega$ und jedes $\varepsilon>0$ existieren eine offene Menge $U \subseteq \Omega$ und eine abgeschlossene Menge $V \subseteq \Omega$ mit $V \subseteq A \subseteq U$ und $\mu(U \backslash V)<\varepsilon$.

Beweis. Sei $\Omega=\bigcup_{j \in \mathbb{N}} K_{j}$ mit kompakten Teilmengen $K_{j} \subseteq \Omega$. Ohne Einschränkung können wir dabei $K_{j} \subseteq K_{j+1}$ für $j \in \mathbb{N}$ annehmen. Sei $A \subseteq \Omega \mu$-messbar. Wir zeigen zunächst (b) für $A$.

\begin{enumerate}
  \item Behauptung: Zu jedem $\varepsilon>0$ existiert $U \subseteq \Omega$ offen mit $A \subseteq U$ und $\mu(U \backslash A)<\frac{\varepsilon}{2}$. Zum Beweis sei $A_{j}:=A \cap K_{j}$ für $j \in \mathbb{N}$. Dann ist $A_{j} \subseteq K_{j}$ und somit $\mu\left(A_{j}\right) \leq \mu\left(K_{j}\right)<\infty$ gemäß (4.11). Da $\mu$ von außen regulär ist, existieren nun offene Mengen $U_{j} \subseteq \Omega$ mit $A_{j} \subseteq U_{j}$ und $\mu\left(U_{j}\right)<\mu\left(A_{j}\right)+\varepsilon 2^{-j-1}$, somit also
\end{enumerate}

$$
\mu\left(U_{j} \backslash A_{j}\right)=\mu\left(U_{j}\right)-\mu\left(A_{j}\right)<\varepsilon 2^{-j-1} \quad \text { für } j \in \mathbb{N} \text {. }
$$

Mit $U:=\bigcup_{j \in \mathbb{N}} U_{j}$ folgt $U \backslash A \subseteq \bigcup_{j \in \mathbb{N}}\left(U_{j} \backslash A_{j}\right)$ und

$$
\mu(U \backslash A) \leq \sum_{j \in \mathbb{N}} \mu\left(U_{j} \backslash A_{j}\right)<\frac{\varepsilon}{2}
$$

wie behauptet. Anwendung der 1. Behauptung auf die $\mu$-messbare Menge $\Omega \backslash A$ liefert nun auch $W \subseteq \Omega$ offen mit $\Omega \backslash A \subseteq W$ und $\mu(W \backslash(\Omega \backslash A))<\frac{\varepsilon}{2}$. Sei $V:=\Omega \backslash W$. Dann ist $V$ abgeschlossen in $\Omega$ mit $V \subseteq A$; ferner gilt $W \backslash(\Omega \backslash A)=W \cap A=A \backslash V$ und daher

$$
\mu(U \backslash V)=\mu(U \backslash A)+\mu(A \backslash V)<\frac{\varepsilon}{2}+\frac{\varepsilon}{2}=\varepsilon .
$$

Somit ist (b) gezeigt.\\
Nun zu (a). Wir müssen zeigen:

$$
\mu(A) \leq \sup \{\mu(K) \mid K \text { kompakt, } K \subseteq A\}
$$

die umgekehrte Ungleichung ist offensichtlich. Für die oben definierten Mengen $A_{j}:=$ $A \cap K_{j}$ gilt $A_{j} \subseteq A_{j+1}$ für $j \in \mathbb{N}$ und $A=\bigcup_{j=1}^{\infty} A_{j}$, somit also

$$
\mu(A)=\lim _{j \rightarrow \infty} \mu\left(A_{j}\right)=\sup _{j \in \mathbb{N}} \mu\left(A_{j}\right) .
$$

Sei nun $\varepsilon>0$. Anwendung von (b) liefert abgeschlossene Mengen $V_{j} \subseteq \Omega, j \in \mathbb{N}$ mit $V_{j} \subseteq A_{j}$ und $\mu\left(A_{j} \backslash V_{j}\right)<\varepsilon$, d.h $\mu\left(A_{j}\right) \leq \mu\left(V_{j}\right)+\varepsilon$. Ferner ist $V_{j}$ kompakt, da $V_{j} \subseteq K_{j}$\\
gilt. Es folgt

$$
\mu(A) \leq \sup _{j \in \mathbb{N}} \mu\left(V_{j}\right)+\varepsilon \leq \sup \{\mu(K) \mid K \text { kompakt, } K \subseteq A\}+\varepsilon
$$

$\mathrm{Da} \varepsilon>0$ beliebig gewählt war, folgt die Behauptung.\\
4.43 Beispiel. Bekannterweise ist das Lebesguemaß $\lambda$ auf $\mathbb{R}^{N}$ ein von außen reguläres und lokal endliches Borelmaß. Somit ist $\lambda$ ein Radonmaß.\\
4.44 Hilfssatz. Sei $\Omega$ lokal kompakt und $\mu$ ein lokal endliches Borelmaß auf $\Omega$. Dann existiert zu jeder kompakten Teilmenge $K \subseteq \Omega$ und zu jedem $\varepsilon>0$ eine kompakte Umgebung $L$ von $K$ mit $\mu(L \backslash K)<\varepsilon$.

Beweis. Übung. Man zeige zunächst, dass $B_{\frac{1}{n}}(K)$ für ausreichend großes $n$ kompakt ist.\\
4.45 Satz (von Lusin). Sei $\mu$ ein Radon-Maß auf $\Omega$. Sei ferner $u: \Omega \rightarrow \mathbb{K} \mu$-messbar. Dann gilt:\\
(a) Ist $A \subseteq \Omega \mu$-messbar mit $\mu(A)<\infty$, so existiert zu jedem $\varepsilon>0$ eine kompakte Teilmenge $K \subseteq A$ mit $\mu(A \backslash K)<\varepsilon$ und derart, dass $\left.u\right|_{K}$ stetig ist.\\
(b) Ist $\Omega$ lokal kompakt und $\mu(\{x \in \Omega \mid u(x) \neq 0\})<\infty$, so existiert zu jedem $\varepsilon>0$ eine Funktion $u_{\varepsilon} \in C_{\mathrm{c}}(\Omega)$ mit

$$
\mu\left(\left\{x \in \Omega \mid u(x) \neq u_{\varepsilon}(x)\right\}\right)<\varepsilon \quad \text { und } \quad \sup _{x \in \Omega}\left|u_{\varepsilon}(x)\right| \leq \sup _{x \in \Omega}|u(x)| .
$$

Beweis. (a) Ohne Einschränkung sei $u \equiv 0$ in $\Omega \backslash A$ (sonst betrachte $u 1_{A}$ anstelle von $u$ ). 1. Spezialfall: $u$ ist eine Elementarfunktion, d.h. $u=\sum_{k=1}^{n} a_{k} 1_{B_{k}}$ mit $\mu$-messbaren disjunkten Teilmengen $B_{k} \subseteq A$ und $a_{k} \in \mathbb{K} \backslash\{0\}, k=1, \ldots, n$. Sei ferner $B_{0}:=$ $\{x \in A \mid u(x)=0\}$. Wegen $\mu\left(B_{k}\right) \leq \mu(A)<\infty$ und der inneren Regularität von $\mu$ existieren dann kompakte Teilmengen $C_{k} \subseteq B_{k}$ mit $\mu\left(B_{k} \backslash C_{k}\right)<\frac{\varepsilon}{n+1}$ für $k=0, \ldots, n$. Dabei gilt $\min _{k \neq j} \operatorname{dist}\left(C_{k}, C_{j}\right)>0$, da es sich um disjunkte kompakte Mengen handelt. Sei nun $K:=\bigcup_{k=0}^{n} C_{k} \subseteq A$. Dann ist $K$ kompakt. Ferner ist $u$ lokal konstant und somit stetig auf $K$. Schließlich ist $A \backslash K$ die disjunkte Vereinigung der Mengen $B_{k} \backslash C_{k}$ und somit

$$
\mu(A \backslash K)=\sum_{k=0}^{n} \mu\left(B_{k} \backslash C_{k}\right)<\varepsilon
$$

\begin{enumerate}
  \setcounter{enumi}{1}
  \item Allgemeiner Fall: Gemäß Satz 4.17 existiert eine Folge von Elementarfunktionen $u_{n}$ mit $\left|u_{n}\right| \leq|u|$ für alle $n$ und $u_{n} \rightarrow u$ punktweise in $\Omega$ für $n \rightarrow \infty$. Nach dem Spezialfall existieren ferner kompakte Teilmengen $K_{n} \subseteq A$ mit $\mu\left(A \backslash K_{n}\right)<\varepsilon 2^{-n-1}$ und derart, dass $\left.u_{n}\right|_{K_{n}}$ stetig ist. Ferner existiert gemäß Satz 4.33 eine $\mu$-messbare Teilmenge $L \subseteq A$ mit\\
$\mu(A \backslash L)<\frac{\varepsilon}{4}$ und $u_{n} \rightarrow u$ gleichmäßig auf $L$ für $n \rightarrow \infty$. Schließlich existiert aufgrund der inneren Regularität von $\mu$ eine kompakte Teilmenge $K_{0} \subseteq L$ mit $\mu\left(L \backslash K_{0}\right)<\frac{\varepsilon}{4}$, also
\end{enumerate}

$$
\mu\left(A \backslash K_{0}\right) \leq \mu(A \backslash L)+\mu\left(L \backslash K_{0}\right)<\frac{\varepsilon}{2}
$$

Sei nun $K:=\bigcap_{n=0}^{\infty} K_{n}$. Dann ist $K \subseteq A$ kompakt und

$$
\mu(A \backslash K) \leq \sum_{n=0}^{\infty} \mu\left(A \backslash K_{n}\right)<\varepsilon \sum_{n=0}^{\infty} 2^{-n-1}=\varepsilon
$$

In $K$ sind ferner alle Funktionen $u_{n}, n \in \mathbb{N}$ stetig; ferner gilt $u_{n} \rightarrow u$ gleichmäßig auf $K$. Es folgt, dass $\left.u\right|_{K}$ stetig ist.\\
(b) Anwendung von (a) auf $A:=\{x \in \Omega \mid u(x) \neq 0\}$ liefert eine kompakte Teilmenge $K \subseteq A$ mit $\mu(A \backslash K)<\frac{\varepsilon}{2}$ und derart, dass $\left.u\right|_{K}$ stetig ist. Der Fortsetzungssatz von Tietze (aus der mengentheoretischen Topologie) besagt, dass sich $\left.u\right|_{K}$ zu einer stetigen Funktion $\tilde{u}: \Omega \rightarrow \mathbb{K}$ fortsetzen lässt mit $\|\tilde{u}\|_{\infty} \leq \sup _{x \in K}|u(x)|$. Gemäß Hilfssatz 4.44 existiert ferner eine kompakte Umgebung $L$ von $K$ mit $\mu(L \backslash K)<\frac{\varepsilon}{2}$. Sei nun

$$
\xi \in C_{\mathrm{c}}(\Omega,[0,1]) \quad \text { definiert durch } \quad \xi(x)=\frac{\operatorname{dist}(x, \Omega \backslash L)}{\operatorname{dist}(x, K)+\operatorname{dist}(x, \Omega \backslash L)},
$$

und sei $u_{\varepsilon}:=\xi \tilde{u} \in C_{\mathrm{c}}(\Omega)$. Dann ist $\xi \equiv 1$ auf $K$ und $\xi \equiv 0$ auf $\Omega \backslash L$. Nach Konstruktion ist also

$$
\left\{x \in \Omega \mid u(x) \neq u_{\varepsilon}(x)\right\} \subseteq(A \cup L) \backslash K=(A \backslash K) \cup(L \backslash K)
$$

und somit

$$
\mu\left(\left\{x \in \Omega \mid u(x) \neq u_{\varepsilon}(x)\right\}\right) \leq \frac{\varepsilon}{2}+\frac{\varepsilon}{2}=\varepsilon .
$$

Ferner ist

$$
\sup _{x \in \Omega}\left|u_{\varepsilon}(x)\right| \leq\|\tilde{u}\|_{\infty} \leq \sup _{x \in K}|u(x)| \leq \sup _{x \in \Omega}|u(x)|,
$$

wie gewünscht.\\
4.46 Beispiel. Sei $\Omega=[0,1] \subseteq \mathbb{R}$, versehen mit der Betragsmetrik. Das Lebesguemaß $\lambda$ auf $\Omega$ erfüllt dann die Voraussetzungen des Satzes von Lusin. Wir betrachten die Dirichletfunktion

$$
u:[0,1] \rightarrow \mathbb{R}, \quad u(x)= \begin{cases}1, & x \in \mathbb{Q} \\ 0, & x \notin \mathbb{Q}\end{cases}
$$

geben uns $\varepsilon>0$ vor und konstruieren eine kompakte Menge $K \subseteq[0,1]$ mit $\lambda([0,1] \backslash$ $K)<\varepsilon$ so, dass $\left.u\right|_{K}$ stetig ist: Sei $\left\{q_{n} \mid n \in \mathbb{N}\right\}$ eine Abzählung von $\mathbb{Q} \cap[0,1]$ und\\
$U_{n}:=U_{\varepsilon 2^{-n-2}}\left(q_{n}\right)$ für $n \in \mathbb{N}$. Sei ferner

$$
K:=[0,1] \backslash \bigcup_{n \in \mathbb{N}} U_{n}
$$

Dann ist $K$ kompakt, $u \equiv 0$ auf $K$ (und somit dort stetig) sowie

$$
\lambda([0,1] \backslash K) \leq \sum_{n \in \mathbb{N}} \lambda\left(U_{n}\right)=\varepsilon \sum_{n \in \mathbb{N}} 2^{-n-1}=\frac{\varepsilon}{2}<\varepsilon .
$$

4.47 Satz. Seien $\Omega$ ein lokal kompakter metrischer Raum und $\mu$ ein Radon-Maß auf $\Omega$.

Sei ferner $1 \leq p<\infty$. Dann gilt:\\
(a) $C_{\mathrm{c}}(\Omega) \subseteq L^{p}(\Omega)$ ist dicht.\\
(b) Ist $\Omega$ ferner $\sigma$-kompakt, so ist $L^{p}(\Omega)$ separabel.

Beweis. (a) Sei $u \in L^{p}(\Omega)$ und $\varepsilon>0$ gegeben. Nach Satz 4.32 existiert $v \in \mathcal{E}(\Omega)$ mit $|v| \leq|u|$ in $\Omega, \mu(\{v \neq 0\})<\infty$ und $\|u-v\|_{p}<\frac{\varepsilon}{2}$. Ohne Einschränkung sei dabei $v \neq 0$, d.h. $\|v\|_{\infty}>0$. Nach dem Satz von Lusin (Satz 4.45(b)) existiert nun $w \in C_{\mathrm{c}}(\Omega)$ mit $\|w\|_{\infty} \leq\|v\|_{\infty}$ und $\mu(\{w \neq v\}) \leq\left(\frac{\varepsilon}{4\|v\|_{\infty}}\right)^{p}$, also

$$
\|w-v\|_{p}^{p}=\int_{\Omega}|w-v|^{p} \mathrm{~d} \mu \leq \mu(\{w \neq v\}) 2^{p}\|v\|_{\infty}^{p} \leq\left(\frac{\varepsilon}{2}\right)^{p}
$$

Es folgt $\|v-w\|_{p} \leq \frac{\varepsilon}{2}$ und somit $\|u-w\|_{p} \leq\|u-v\|_{p}+\|v-w\|_{p}<\varepsilon$.\\
(b) Gemäß Satz 4.37(a) existieren kompakte Mengen $M_{j} \subseteq \Omega, j \in \mathbb{N}$ mit $\Omega=\bigcup_{j=1}^{\infty} M_{j}$ und derart, dass jede kompakte Menge $K \subseteq \Omega$ in einer der Mengen $M_{j}$ enthalten ist. Wie im Beweis von Satz 4.37(b) setzen wir

$$
C_{j}:=\left\{f \in C_{\mathrm{c}}(\Omega) \mid \operatorname{supp} f \subseteq M_{j}\right\} \subseteq C\left(M_{j}\right) \quad \text { für } j \in \mathbb{N} .
$$

Gemäß Korollar 4.13 ist $C\left(M_{j}\right)$ und somit auch $C_{j}$ separabel bzgl. $\|\cdot\|_{\infty}$. Daher existieren abzählbare und dichte Teilmengen $A_{j} \subseteq C_{j}$ für $j \in \mathbb{N}$. Wir zeigen: Die abzählbare Menge $A:=\bigcup_{j \in \mathbb{N}} A_{j}$ liegt dicht in $L^{p}(\Omega)$. Sei dazu $u \in L^{p}(\Omega)$ und $\varepsilon>0$ beliebig. Gemäß (a) existiert $v \in C_{\mathrm{c}}(\Omega)$ mit $\|u-v\|_{p}<\frac{\varepsilon}{2}$; dabei ist $v \in C_{j}$ für ein $j \in \mathbb{N}$. Somit existiert $w \in A_{j}$ mit $\mu\left(M_{j}\right)^{\frac{1}{p}}\|v-w\|_{\infty} \leq \frac{\varepsilon}{2}$. Dann ist

$$
\|w-v\|_{p}^{p}=\int_{\Omega}|w-v|^{p} \mathrm{~d} \mu \leq \mu\left(M_{j}\right)\|v-w\|_{\infty}^{p} \leq\left(\frac{\varepsilon}{2}\right)^{p}
$$

und somit $\|w-v\|_{p} \leq \frac{\varepsilon}{2}$. Es folgt $\|u-w\|_{p} \leq\|u-v\|_{p}+\|v-w\|_{p}<\varepsilon$. Also ist $A$ dicht in $L^{p}(\Omega)$, wie behauptet.\\
4.48 Bemerkung. Ist $\Omega \subseteq \mathbb{R}^{N}$ Lebesgue-messbar, so ist $L^{p}(\Omega)$ nach Satz 2.18 separabel, da man $L^{p}(\Omega)$ als eine Teilmenge des gemäß Satz 4.47 separablen Raums $L^{p}\left(\mathbb{R}^{N}\right)$ auffassen kann (vermöge trivialer Fortsetzung).\\
4.49 Satz. Ist $\Omega \subseteq \mathbb{R}^{N}$ Lebesgue-messbar mit $\lambda(\Omega)>0$, so ist $L^{\infty}(\Omega)$ (definiert bzgl. des Lebesgue-Maßes) nicht separabel.

Beweisskizze. Für $r \geq 0$ sei $\Omega_{r}:=B_{r}(0) \cap \Omega$. Mit dem Satz von Lebesgue (siehe Satz 4.21) sieht man, dass die monoton wachsende Funktion

$$
g:[0, \infty) \rightarrow[0, \infty), \quad g(r)=\lambda\left(\Omega_{r}\right)
$$

stetig ist. Ferner ist $g(0)=0$ und $\lim _{r \rightarrow \infty} g(r)=\lambda(\Omega)$ nach dem Satz von der monotonen Konvergenz (siehe Satz 4.19). Für $s \in(0, \lambda(\Omega))$ existiert nach dem Zwischenwertsatz also ein $r_{s} \in[0, \infty)$ mit $g\left(r_{s}\right)=s$. Betrachte nun die Menge

$$
M:=\left\{1_{\Omega_{r_{s}}} \mid 0<s<\lambda(\Omega)\right\} \subseteq L^{\infty}(\Omega) .
$$

Für $0<s<t<\lambda(\Omega)$ ist dann $r_{s}<r_{t}$, also $\Omega_{r_{s}} \subseteq \Omega_{r_{t}}$ sowie $\lambda\left(\Omega_{r_{t}} \backslash \Omega_{r_{s}}\right)=t-s>0$; somit also $\left\|1_{\Omega_{r_{s}}}-1_{\Omega_{r_{t}}}\right\|_{\infty}=1$. Daher besteht $U_{1 / 3}(M)$ aus überabzählbar vielen paarweise disjunkten offenen Teilmengen von $L^{\infty}(\Omega)$. Es folgt, dass $L^{\infty}(\Omega)$ nicht separabel ist.\\
4.50 Definition und Satz. Sei $\Omega \subseteq \mathbb{R}^{N}$ offen und

$$
C_{\mathrm{c}}^{\infty}(\Omega):=\left\{f \in C^{\infty}(\Omega, \mathbb{K}) \mid \operatorname{supp} f \subseteq \Omega \text { kompakt }\right\}
$$

Dann liegt $C_{\mathrm{c}}^{\infty}(\Omega)$ dicht in $L^{p}(\Omega)$ für $1 \leq p<\infty$, aber (offensichtlich) nicht dicht in $L^{\infty}(\Omega)$. Hierbei sei das Lebesguemaß auf $\Omega$ zugrunde gelegt.

Beweisskizze. Sei $u \in L^{p}(\Omega), \varepsilon>0$. Nach Satz 4.47 existiert $v \in C_{\mathrm{c}}(\Omega)$ mit $\|u-v\|_{p}<\frac{\varepsilon}{2}$. Wir setzen $v$ trivial auf $\mathbb{R}^{N}$ fort. Wir betrachten ferner den Glättungskern

$$
\rho \in C_{\mathrm{c}}^{\infty}\left(\mathbb{R}^{N}\right), \quad \rho(x):= \begin{cases}c \exp \left(\frac{1}{|x|^{2}-1}\right), & |x|<1 \\ 0, & |x| \geq 1\end{cases}
$$

wobei $c>0$ so gewählt sei, dass $\int_{\mathbb{R}^{N}} \rho=1$ gilt. Für $h>0$ sei ferner

$$
\rho_{h} \in C_{\mathrm{c}}^{\infty}\left(\mathbb{R}^{N}\right), \quad \rho_{h}(x)=\frac{1}{h^{N}} \rho\left(\frac{x}{h}\right)
$$

Dann ist supp $\rho_{h}=B_{h}(0)$ und es gilt ebenfalls $\int_{\mathbb{R}^{N}} \rho_{h}=1$ für $h>0$. Nun definieren wir

$$
v_{h}: \mathbb{R}^{N} \rightarrow \mathbb{R}, \quad v_{h}(x):=\int_{\mathbb{R}^{N}} \rho_{h}(x-y) v(y) \mathrm{d} y
$$

für $0<h<\operatorname{dist}(\operatorname{supp} v, \partial \Omega)$. Dann ist $v_{h} \in C_{\mathrm{c}}^{\infty}\left(\mathbb{R}^{N}\right)$ mit $\operatorname{supp} v_{h} \subseteq B_{h}(\operatorname{supp} v) \subseteq \Omega$. Ferner gilt mit

$$
\varepsilon_{h}:=\sup \{|v(x)-v(y)||x, y \in \Omega,|x-y| \leq h\}
$$

die Abschätzung

$$
\left|v_{h}(x)-v(x)\right| \leq \int_{\mathbb{R}^{N}} \rho_{h}(x-y)|v(y)-v(x)| \mathrm{d} y \leq \varepsilon_{h} \int_{\mathbb{R}^{N}} \rho_{h}(x-y) \mathrm{d} y=\varepsilon_{h} \quad \text { für } x \in \Omega
$$

Da $v$ gleichmäßig stetig ist, gilt $\lim _{h \rightarrow 0} \varepsilon_{h}=0$. Es folgt $v_{h} \rightarrow v$ bzgl. $\|\cdot\|_{\infty}$ in $\Omega$ und somit auch $\left\|v_{h}-v\right\|_{L^{p}(\Omega)} \rightarrow 0$ für $h \rightarrow 0$, da $\operatorname{supp} v_{h}$ für $h \leq 1$ in der kompakten Menge $B_{1}(\operatorname{supp} v)$ enthalten ist (vgl. Bemerkung und Beispiel 4.26(b)). Die Einschränkung $w_{h}:=\left.v_{h}\right|_{\Omega} \in C_{\mathrm{c}}^{\infty}(\Omega)$ erfüllt also $\left\|w_{h}-u\right\|_{p} \leq\left\|w_{h}-v\right\|_{p}+\|v-u\|_{p}<\frac{\varepsilon}{2}+\frac{\varepsilon}{2}=\varepsilon$ für $0<h<\operatorname{dist}(\operatorname{supp} v, \partial \Omega)$ genügend klein gewählt.

\section*{Hilberträume}
\subsection*{Orthogonalität}
Stets sei $H$ ein Hilbertraum über $\mathbb{K}$ mit Skalarprodukt $\langle\cdot, \cdot\rangle$ und induzierter Norm $\|\cdot\|$.\\
5.1 Satz und Definition. Sei $C \subseteq H$ konvex, abgeschlossen und nichtleer. Dann existiert zu jedem $x \in H$ genau ein $y \in C$ mit $\|x-y\|=\operatorname{dist}(x, C)$. Wir setzen $y=P(x)$ und nennen $P: H \rightarrow H$ die (abstandsminimierende) Projektion auf $C$.

Beweis. Zur Existenz: Sei $x \in H$ und $\alpha:=\operatorname{dist}(x, C)$. Sei ferner $\left(y_{n}\right)_{n} \subseteq C$ eine Folge mit

$$
\left\|x-y_{n}\right\| \rightarrow \alpha \quad \text { für } n \rightarrow \infty .
$$

Für $n, m \in \mathbb{N}$ gilt $\frac{y_{n}+y_{m}}{2} \in C$ aufgrund der Konvexität von $C$, also

$$
\begin{aligned}
4 \alpha^{2} \leq 4\left\|x-\frac{y_{n}+y_{m}}{2}\right\|^{2}=\left\|x-y_{m}+x-y_{n}\right\|^{2} & \\
& =2\left(\left\|x-y_{m}\right\|^{2}+\left\|x-y_{n}\right\|^{2}\right)-\|\underbrace{x-y_{m}-\left(x-y_{n}\right)}_{y_{n}-y_{m}}\|^{2}
\end{aligned}
$$

und somit wegen (5.1)

$$
\sup _{m \geq n}\left\|y_{m}-y_{n}\right\|^{2} \leq 2 \sup _{m \geq n}\left(\left\|x-y_{m}\right\|^{2}+\left\|x-y_{n}\right\|^{2}\right)-4 \alpha^{2} \rightarrow 0 \quad \text { für } n \rightarrow \infty
$$

Es folgt, dass $\left(y_{n}\right)_{n} \subseteq C$ eine Cauchyfolge ist. Ferner ist $C$ als abgeschlossene Teilmenge des Hilbertraums $H$ nach Satz 2.32 vollständig, also existiert $y:=\lim _{n} y_{n} \in C$. Ferner gilt $\|x-y\|=\lim _{n \rightarrow \infty}\left\|x-y_{n}\right\|=\alpha$.

Zur Eindeutigkeit: Sei $y^{\prime} \in C$ mit $\left\|x-y^{\prime}\right\|=\alpha=\|x-y\|$. Dann erfüllt die Folge ( $y, y^{\prime}, y, y^{\prime}, \ldots$ ) Bedingung (5.1) und ist daher nach obigem Argument eine Cauchyfolge. Dies liefert $y=y^{\prime}$.\\
5.2 Satz. Sind $C \subseteq H$ und $P: H \rightarrow H$ wie in Satz und Definition 5.1 und $Q:=$ id $-P: H \rightarrow H$, so gilt:\\
(a) $\left.P\right|_{C}=\mathrm{id}_{C}$.\\
(b) Für alle $x \in H, y \in C$ ist $\operatorname{Re}\langle Q(x), P(x)-y\rangle \geq 0$.\\
(c) $\left\|P\left(x_{1}\right)-P\left(x_{2}\right)\right\| \leq\left\|x_{1}-x_{2}\right\|$ für alle $x_{1}, x_{2} \in H$, insbesondere ist $P$ Lipschitzstetig.

Beweis. (a) ist klar nach Definition.\\
(b) Für $\varepsilon \in(0,1)$ ist $z_{\varepsilon}:=P(x)+\varepsilon(y-P(x)) \in C$ (da $C$ konvex), also

$$
\begin{aligned}
0 \leq\left\|x-z_{\varepsilon}\right\|^{2}-\operatorname{dist}(x, C)^{2}=\| Q(x) & +\varepsilon(P(x)-y)\left\|^{2}-\right\| Q(x) \|^{2} \\
& =2 \varepsilon \operatorname{Re}\langle Q(x), P(x)-y\rangle+\varepsilon^{2}\|P(x)-y\|^{2}=: t(\varepsilon)
\end{aligned}
$$

und somit

$$
0 \leq \lim _{\varepsilon \rightarrow 0^{+}} \frac{t(\varepsilon)}{2 \varepsilon}=\operatorname{Re}\langle Q(x), P(x)-y\rangle
$$

(c) Es ist wegen (b)

$$
\begin{aligned}
& \left\|x_{1}-x_{2}\right\|^{2} \\
& \quad=\left\|Q\left(x_{1}\right)-Q\left(x_{2}\right)+P\left(x_{1}\right)-P\left(x_{2}\right)\right\|^{2} \\
& \quad=\left\|Q\left(x_{1}\right)-Q\left(x_{2}\right)\right\|^{2}+\left\|P\left(x_{1}\right)-P\left(x_{2}\right)\right\|^{2}+2 \operatorname{Re}\left\langle Q\left(x_{1}\right)-Q\left(x_{2}\right), P\left(x_{1}\right)-P\left(x_{2}\right)\right\rangle \\
& \quad \geq\left\|P\left(x_{1}\right)-P\left(x_{2}\right)\right\|^{2}+2 \operatorname{Re}\left\langle Q\left(x_{1}\right), P\left(x_{1}\right)-P\left(x_{2}\right)\right\rangle+2 \operatorname{Re}\left\langle Q\left(x_{2}\right), P\left(x_{2}\right)-P\left(x_{1}\right)\right\rangle \\
& \quad \geq\left\|P\left(x_{1}\right)-P\left(x_{2}\right)\right\|^{2}
\end{aligned}
$$

5.3 Bemerkung. Satz und Definition 5.1 und Satz 5.2 gelten auch, wenn $(H,\langle\cdot, \cdot\rangle)$ ein Prähilbertraum und $C \subseteq H$ nichtleer, konvex und vollständig ist, insbesondere also dann, wenn die abgeschlossene, nichtleere und konvexe Teilmenge $C$ in einem endlichdimensionalen Teilraum des Prähilbertraums $H$ enthalten ist.\\
5.4 Beispiel. Sei $\Omega \subseteq \mathbb{R}^{N}$ Lebesgue-messbar, $H=L^{2}(\Omega, \mathbb{R})$ sowie $C:=$ $\left\{u \in L^{2}(\Omega) \mid u \geq 0\right.$ f.ü. auf $\left.\Omega\right\}$. Man sieht leicht, dass $C$ eine nichtleere, abgeschlossene und konvexe Teilmenge von $H$ ist. Wir zeigen nun, dass die Projektion $P$ von $H$ auf $C$ gegeben ist durch $P(u)=u^{+}=\max \{0, u\}$ für $u \in H$. Tatsächlich ist $u^{+} \in C$ für alle $u \in H$, und für $v \in C$ gilt

$$
\begin{aligned}
\|v-u\|_{2}^{2}=\int_{\Omega}(v-u)^{2} \geq \int_{\{u<0\}}(v-u)^{2} & \\
& \geq \int_{\{u<0\}} u^{2}=\int_{\Omega}\left(u^{+}-u\right)^{2}=\left\|u^{+}-u\right\|_{2}^{2}
\end{aligned}
$$

Also ist $u^{+}$das zu $u$ abstandsminimierende Element in $C$, d.h. $P(u)=u^{+}$gemäß Satz und Definition 5.1.

Ende des Inhalts für 2024-05-24\\
5.5 Definition und Satz. Für $M \subseteq H$ sei $M^{\perp}:=\{x \in H \mid\langle x, y\rangle=0$ für alle $y \in M\}$. Dann gilt:\\
(a) $M^{\perp}$ ist ein abgeschlossener Unterraum von $H$\\
(b) $M^{\perp}=(\overline{\operatorname{span} M})^{\perp}$

Beweis. (a) folgt direkt aus der Linearität und Stetigkeit des Skalarprodukts in der ersten Komponente.

Zu (b): Offensichtlich gilt für Teilmengen $A \subseteq B \subseteq H$ die Inklusionsumkehrung $B^{\perp} \subseteq A^{\perp}$. Wegen $M \subseteq \overline{\operatorname{span} M}$ folgt also $(\overline{\operatorname{span} M})^{\perp} \subseteq M^{\perp}$. Sei umgekehrt $x \in M^{\perp}$, und sei zunächst $y=\alpha_{1} y_{1}+\cdots+\alpha_{n} y_{n} \in \operatorname{span} M$ mit $y_{i} \in M, \alpha_{i} \in K$ gegeben. Dann ist

$$
\langle x, y\rangle=\sum_{i=1}^{n} \bar{\alpha}_{i}\left\langle x, y_{i}\right\rangle=0
$$

Ist schließlich $y \in \overline{\operatorname{span} M}$ und $\left(y^{n}\right)_{n}$ eine Folge in span $M$ mit $\lim _{n \rightarrow \infty} y^{n}=y$, so folgt $\langle x, y\rangle=\lim _{n \rightarrow \infty}\left\langle x, y^{n}\right\rangle=0$. Somit ist $x \in(\overline{\operatorname{span} M})^{\perp}$, und insgesamt folgt $M^{\perp}=$ $(\overline{\operatorname{span} M})^{\perp}$.\\
5.6 Satz und Definition. Sind $V, W \subseteq H$ Unterräume mit $H=V+W$ und $\langle v, w\rangle=0$ für alle $v \in V, w \in W$, so gilt:\\
(a) $\|v+w\|^{2}=\|v\|^{2}+\|w\|^{2}$ für alle $v \in V, w \in W$ ("Satz des Pythagoras").\\
(b) $W=V^{\perp}$ und $V=W^{\perp}$. Insbesondere sind $V, W$ abgeschlossen in $H$ nach Definition und Satz 5.5 (a).\\
(c) $H=V \oplus_{\text {top }} W$.

Wir schreiben dann $H=V \oplus_{\perp} W$, nennen diese Zerlegung orthogonal und die Räume $V$ und $W$ orthogonale Komplementärräume.

Beweis. (a) Für alle $v \in V, w \in W$ ist

$$
\|v+w\|^{2}=\|v\|^{2}+\|w\|^{2}+2 \operatorname{Re} \underbrace{\langle v, w\rangle}_{=0}=\|v\|^{2}+\|w\|^{2} .
$$

(b) Nach Voraussetzung ist $W \subseteq V^{\perp}$ und $V \subseteq W^{\perp}$. Sei nun $w \in V^{\perp}$. Da $H=V+W$ ist, existieren $v \in V$ und $\tilde{w} \in W$ mit $w=v+\tilde{w}$, also

$$
\|v\|^{2}=\langle v, v\rangle=\langle\underbrace{w-\tilde{w}}_{\in V^{\perp}}, v\rangle=0
$$

Es folgt $w=\tilde{w} \in W$. Dies zeigt $W=V^{\perp}$. Genauso folgt $V=W^{\perp}$.\\
(c) Es gilt $V \cap W=\{0\}$, da $\|x\|^{2}=\langle x, x\rangle=0$ für alle $x \in V \cap W$. Also gilt: $H=V \oplus W$. Da $V, W$ gemäß (b) in $H$ abgeschlossen sind, folgt $H=V \oplus_{\text {top }} W$ nach Satz 3.24(b).\\
5.7 Satz und Definition. Sei $V$ ein abgeschlossener Unterraum von H. Dann gilt $H=V \oplus_{\perp} V^{\perp}$, und die (abstandsminimierende) Projektion $P$ auf $V$ aus Satz und Definition 5.1 erfüllt:\\
(a) $P \in \mathcal{L}(H), P^{2}=P$ und Bild $P=V$, Kern $P=V^{\perp}$\\
(b) $\|P\|=1$ falls $V \neq\{0\}$

Man nennt $P$ die orthogonale Projektion auf $V$.\\
Beweis. Wir zeigen zunächst: Für $x \in H$ ist

$$
Q(x):=x-P(x) \in V^{\perp}
$$

Nach Satz 5.2(b) gilt für alle $y \in V, \alpha \in \mathbb{K}$ :

$$
0 \leq \operatorname{Re}\langle Q(x), P(x)-(\underbrace{-\bar{\alpha} y+P(x)}_{\in V})\rangle=\operatorname{Re}\langle Q(x), \bar{\alpha} y\rangle=\operatorname{Re}(\alpha\langle Q(x), y\rangle) .
$$

Durch Wahl von $\alpha=1$ und $\alpha=-1$ folgt $\operatorname{Re}\langle Q(x), y\rangle=0$ für alle $y \in V$. Ist $\mathbb{K}=\mathbb{C}$, so folgt durch Wahl von $\alpha= \pm i$ auch

$$
0 \leq \operatorname{Re}( \pm \mathrm{i}\langle Q(x), y\rangle)=\mp \operatorname{Im}\langle Q(x), y\rangle, \quad \text { also } \quad \operatorname{Im}\langle Q(x), y\rangle=0 \quad \text { für alle } y \in V .
$$

Insgesamt folgt $\langle Q(x), y\rangle=0$ für alle $y \in V$, und somit ist $Q(x) \in V^{\perp}$, wie behauptet. Aus (5.2) folgt nun $x=P(x)+Q(x) \in V+V^{\perp}$ für alle $x \in H$, also $H=V+V^{\perp}$. Nach Satz und Definition 5.6 ist also $H=V \oplus_{\perp} V^{\perp}$. Sei nun $\widetilde{P} \in \mathcal{L}(H)$ die Projektion auf $V$ längs $V^{\perp}$ im Sinne von Definition 3.22. Dann gilt für alle $x \in H$ :

$$
\begin{aligned}
& \operatorname{dist}(x, V)^{2}=\|x-P(x)\|^{2}=\|\underbrace{x-\widetilde{P}(x)}_{\in V^{\perp}}+\underbrace{\widetilde{P}(x)-P(x)}_{\in V}\|^{2} \\
& \quad \text { Satz und Definition 5.6(a) }\|x-\widetilde{P}(x)\|^{2}+\|\widetilde{P}(x)-P(x)\|^{2} \geq\|x-\widetilde{P}(x)\|^{2} \geq \operatorname{dist}(x, V)^{2} .
\end{aligned}
$$

Dies erzwingt $\|\widetilde{P}(x)-P(x)\|^{2}=0$, also $\widetilde{P}(x)=P(x)$ für alle $x \in H$. Insbesondere ist $P \in \mathcal{L}(H), P^{2}=P$ sowie Bild $P=V$ und $\operatorname{Kern} P=V^{\perp}$. Weiterhin ist $\|P x\| \leq\|x\|$ für $x \in H$ nach Satz 5.2(c) und $\|P x\|=\|x\|$ für $x \in V$. Dies liefert $\|P\|=1$, falls $V \neq\{0\}$.\\
5.8 Korollar. Für $M \subseteq H$ gelten:\\
(a) $M^{\perp \perp}=\overline{\operatorname{span} M}$.\\
(b) $M$ total in $H \Longleftrightarrow M^{\perp}=\{0\}$.

Beweis. (a) Setze $V=\overline{\operatorname{span} M}$. Dann ist $H=V \oplus_{\perp} V^{\perp}$ nach Satz und Definition 5.7, also $V=V^{\perp \perp}=M^{\perp \perp}$ wegen Definition und Satz 5.5(b).\\
(b) Ist $M$ total, so ist $M^{\perp}=\overline{\operatorname{span} M^{\perp}}=H^{\perp}=\{0\}$ wegen Definition und Satz 5.5(b). Ist umgekehrt $M^{\perp}=\{0\}$, so ist $H=\{0\}^{\perp}=M^{\perp \perp} \stackrel{(a)}{=} \stackrel{\text { span } M}{\text { span }}$. Also ist $M$ total.\\
5.9 Definition und Bemerkung. Für $x \in H$ sei

$$
\varphi_{x} \in H^{\prime}=\mathcal{L}(H, \mathbb{K}) \quad \text { definiert durch } \quad \varphi_{x}(y):=\langle y, x\rangle \quad \text { für } y \in H .
$$

Die Linearität von $\varphi_{x}$ ist offensichtlich. Da ferner

$$
\left|\varphi_{x}(y)\right|=|\langle y, x\rangle| \leq\|y\|\|x\| \quad \text { für alle } y \in H \quad \text { und } \quad\left|\varphi_{x}(x)\right|=\|x\|^{2},
$$

ist $\varphi_{x}$ in der Tat auch stetig mit

$$
\left\|\varphi_{x}\right\|=\|x\| .
$$

Es folgt, dass die Abbildung $J: H \rightarrow H^{\prime}, x \mapsto J(x):=\varphi_{x}$ eine Isometrie ist und damit insbesondere injektiv. Man beachte ferner: $J$ ist antilinear, d.h. es gilt

$$
J\left(x_{1}+x_{2}\right)=J\left(x_{1}\right)+J\left(x_{2}\right) \quad \text { und } \quad J(\lambda x)=\bar{\lambda} J(x) \quad \text { für alle } x, x_{1}, x_{2} \in H, \lambda \in \mathbb{K} .
$$

5.10 Satz (Rieszscher Darstellungssatz). (a) Für alle $\varphi \in H^{\prime}$ existiert $x \in H$ mit $\varphi=\varphi_{x}$. Die Abbildung $J: H \rightarrow H^{\prime}$ aus Definition und Bemerkung 5.9 ist also bijektiv.\\
(b) Die Operatornorm auf $H^{\prime}$ wird induziert von einem Skalarprodukt $\langle\cdot, \cdot\rangle_{*}$, gegeben durch

$$
\langle\varphi, \psi\rangle_{*}:=\left\langle J^{-1}(\psi), J^{-1}(\varphi)\right\rangle \quad \text { für alle } \varphi, \psi \in H^{\prime} .
$$

Insbesondere ist ( $H^{\prime},\langle\cdot, \cdot\rangle_{*}$ ) ein Hilbertraum.\\
Beweis. (a) Sei $\varphi \in H^{\prime} \backslash\{0\}$. Dann ist $V:=\operatorname{Kern} \varphi$ ein abgeschlossener Unterraum von $H$. Nach Satz und Definition 5.7 ist also $H=V \oplus_{\perp} V^{\perp}$, und $\left.\varphi\right|_{V^{\perp}}: V^{\perp} \rightarrow \mathbb{K}$ ist linear, injektiv, hat $\operatorname{Bild}\left(\left.\varphi\right|_{V^{\perp}}\right)=\operatorname{Bild} \varphi \neq\{0\}$. Daher gilt $\operatorname{dim} \operatorname{Bild}\left(\left.\varphi\right|_{V^{\perp}}\right)=1,\left.\varphi\right|_{V^{\perp}}$ ist bijektiv und daher $\operatorname{dim} V^{\perp}=1$. Es existiert also $x_{0} \in V^{\perp} \backslash\{0\}$ mit $V^{\perp}=\operatorname{span}\left\{x_{0}\right\}$. Setze $x:=\frac{\overline{\varphi\left(x_{0}\right)}}{\left\|x_{0}\right\|^{2}} x_{0}$. Dann gilt für alle $y=v+\lambda x \in H$ mit $v \in V$ :

$$
\begin{aligned}
\varphi_{x}(y)=\langle y, x\rangle=\langle v, x\rangle+\lambda\|x\|^{2}=\lambda\|x\|^{2} & \\
& \text { und } \quad \varphi(y)=\lambda \varphi(x)=\lambda \frac{\left|\varphi\left(x_{0}\right)\right|^{2}}{\left\|x_{0}\right\|^{2}}=\lambda\|x\|^{2}
\end{aligned}
$$

Es folgt $\varphi=\varphi_{x}$.\\
(b) Die Skalarprodukt-Eigenschaften von $\langle\cdot, \cdot\rangle_{*}$ folgen direkt aus denen von $\langle\cdot, \cdot\rangle$. Ferner gilt für $\varphi \in H^{\prime}$ mit $x=J^{-1}(\varphi)$ :

$$
\langle\varphi, \varphi\rangle_{*} \stackrel{\text { Def. }}{=}\langle x, x\rangle=\|x\|^{2} \stackrel{(5.3)}{=}\|\varphi\|^{2}
$$

Somit induziert $\langle\cdot, \cdot\rangle_{*}$ also die Operatornorm auf $H^{\prime}$.

\subsection*{Anwendung auf ein Neumann-Randwertproblem}
Wir betrachten stets $\mathbb{K}=\mathbb{R}$ und $\Omega=(-1,1) \subseteq \mathbb{R}$. Gegeben seien Funktionen $a \in$ $L^{\infty}(\Omega), f \in L^{2}(\Omega)$, wobei $\operatorname{ess}_{\inf } a>0$ sei. Wir suchen eine Lösung des NeumannRandwertproblems

$$
\left\{\begin{aligned}
-u^{\prime \prime}+a(x) u & =f(x), \quad x \in \Omega, \\
u^{\prime}(1)=u^{\prime}(-1) & =0 .
\end{aligned}\right.
$$

Dazu benötigen wir einige Vorbereitungen.\\
5.11 Definition. (a) Sei $u \in L^{2}(\Omega)$. Wir nennen $v \in L^{2}(\Omega)$ eine schwache Ableitung von $u$, wenn gilt:

$$
\int_{\Omega} v w=-\int_{\Omega} u w^{\prime} \quad \text { für alle } w \in C_{\mathrm{c}}^{1}(\Omega) .
$$

Dadurch ist $v$ eindeutig bestimmt, denn gilt (5.4) auch für $\tilde{v} \in L^{2}(\Omega)$, so ist

$$
\langle v-\tilde{v}, w\rangle_{L^{2}(\Omega)}=\int_{\Omega}(v-\tilde{v}) w=-\int_{\Omega}(u-u) w^{\prime}=0 \quad \text { für alle } w \in C_{\mathrm{c}}^{1}(\Omega) .
$$

Da $C_{\mathrm{c}}^{1}(\Omega)$ in $L^{2}(\Omega)$ dicht liegt, folgt $v-\tilde{v} \in C_{\mathrm{c}}^{1}(\Omega)^{\perp}=\{0\}$. Wir schreiben $u^{\prime}$ anstelle von $v$.\\
(b) Wir setzen

$$
H^{1}(\Omega):=\left\{u \in L^{2}(\Omega) \mid u \text { besitzt eine schwache Ableitung } u^{\prime} \in L^{2}(\Omega)\right\}
$$

und versehen $H^{1}(\Omega)$ mit dem Skalarprodukt

$$
\left\langle u_{1}, u_{2}\right\rangle_{H^{1}}:=\left\langle u_{1}, u_{2}\right\rangle_{L^{2}(\Omega)}+\left\langle u_{1}^{\prime}, u_{2}^{\prime}\right\rangle_{L^{2}(\Omega)} .
$$

Die induzierte Norm sei mit $\|\cdot\|_{H^{1}}$ bezeichnet. $H^{1}$ gehört zu der Klasse der sogenannten Sobolevräume.\\
5.12 Bemerkung und Beispiel. (a) Ist $u \in C^{1}(\Omega) \cap L^{2}(\Omega)$ mit klassischer Ableitung $u^{\prime} \in L^{2}(\Omega)$, so ist $u^{\prime}$ auch die schwache Ableitung von $u$. Dies folgt aus der gewöhnlichen Regel zur partiellen Integration. Dies gilt insbesondere dann, wenn $u \in C^{1}(\bar{\Omega})$ ist.\\
(b) Sei $u \in L^{2}(\Omega)$ definiert durch $u(x)=|x|^{\alpha}$ für ein $\alpha>\frac{1}{2}$. Dann ist $u \in H^{1}(\Omega)$ mit schwacher Ableitung

$$
v=u^{\prime} \in L^{2}(\Omega), \quad v(x)=\alpha|x|^{\alpha-2} x \quad \text { für } x \neq 0 .
$$

In der Tat ist $v \in L^{2}(\Omega)$, denn wegen $2(\alpha-1)>-1$ ist

$$
\int_{\Omega}|v|^{2}=\alpha^{2} \int_{\Omega}|x|^{2(\alpha-1)} \mathrm{d} x<\infty
$$

Ferner ist $v(x)=u^{\prime}(x)$ im klassischen Sinne für $x \neq 0$, und es folgt für $w \in C_{\mathrm{c}}^{1}(\Omega)$ mit dem Satz von Lebesgue (Satz 4.21) wegen $v w, u w^{\prime} \in L^{1}(\Omega)$ :

$$
\begin{aligned}
\int_{\Omega} v w & =\lim _{\varepsilon \rightarrow 0^{+}}\left(\int_{-1}^{-\varepsilon} u^{\prime} w+\int_{\varepsilon}^{1} u^{\prime} w\right) \\
& =\lim _{\varepsilon \rightarrow 0^{+}}\left((u w)(-\varepsilon)-\int_{-1}^{-\varepsilon} u w^{\prime}-(u w)(\varepsilon)-\int_{\varepsilon}^{1} u w^{\prime}\right) \\
& =-\lim _{\varepsilon \rightarrow 0^{+}}\left(\int_{-1}^{-\varepsilon} u w^{\prime}+\int_{\varepsilon}^{1} u w^{\prime}\right)=-\int_{\Omega} u w^{\prime}
\end{aligned}
$$

Dies zeigt $u^{\prime}=v$ im schwachen Sinne. Wir werden weiter unten in Satz 5.14 sehen, dass die oben definierte Funktion im Fall $\alpha<\frac{1}{2}$ nicht mehr in $H^{1}(\Omega)$ liegt.

Ende des Inhalts für 2024-05-28\\
5.13 Satz. $H^{1}(\Omega)$ ist ein Hilbertraum.

Beweis. Die Skalarprodukteigenschaften sind klar. Sei $\left(u_{k}\right)_{k}$ eine Cauchyfolge in $H^{1}(\Omega)$. Nach Definition des Skalarprodukts ist $\left(u_{k}\right)_{k}$ dann eine Cauchyfolge in $L^{2}(\Omega)$ und $\left(u_{k}^{\prime}\right)_{k}$ ist eine Cauchyfolge in $L^{2}(\Omega)$. Somit gilt

$$
u_{k} \rightarrow u \in L^{2}(\Omega) \quad \text { und } \quad u_{k}^{\prime} \rightarrow v \in L^{2}(\Omega) \quad \text { für } k \rightarrow \infty .
$$

Ferner gilt für alle $w \in C_{\mathrm{c}}^{1}(\Omega)$ :

$$
\int_{\Omega} v w=\lim _{k \rightarrow \infty} \int_{\Omega} u_{k}^{\prime} w=-\lim _{k \rightarrow \infty} \int_{\Omega} u_{k} w^{\prime}=-\int_{\Omega} u w^{\prime} .
$$

Somit ist also $u \in H^{1}(\Omega)$ mit $u^{\prime}=v$, und

$$
\left\|u_{k}-u\right\|_{H^{1}}^{2}=\left\|u_{k}-u\right\|_{2}^{2}+\left\|u_{k}^{\prime}-v\right\|_{2}^{2} \rightarrow 0 \quad \text { für } n \rightarrow \infty .
$$

Dies zeigt die Vollständigkeit von $H^{1}(\Omega)$.\\
5.14 Satz. Sei $u \in H^{1}(\Omega)$. Dann gilt:\\
(a) Ist $u^{\prime} \equiv 0$, so besitzt $u$ einen konstanten Repräsentanten.\\
(b) Die Funktion u lässt sich (nach Wahl eines geeigneten Repräsentanten) zu einer stetigen Funktion auf $\bar{\Omega}$ fortsetzen, welche

$$
u(y)-u(x)=\int_{x}^{y} u^{\prime} \quad \text { für alle } x, y \in \bar{\Omega}
$$

erfüllt.\\
(c) Es gilt $|u(x)-u(y)| \leq\left\|u^{\prime}\right\|_{2}|x-y|^{\frac{1}{2}}$ für $x, y \in \bar{\Omega}$ und $\|u\|_{\infty} \leq \sqrt{2}\|u\|_{H^{1}}$.

Beweis. Zu (a): Sei $\psi \in C_{\mathrm{c}}(\Omega)$ mit $\int_{\Omega} \psi=1$ fest gewählt. Sei ferner $w \in C_{\mathrm{c}}(\Omega)$ beliebig, und sei

$$
f_{w}:=w-\left(\int_{\Omega} w\right) \psi \in C_{\mathrm{c}}(\Omega)
$$

Diese Funktion besitzt eine Stammfunktion gegeben durch

$$
F_{w} \in C_{\mathrm{c}}^{1}(\Omega), \quad F_{w}(x)=\int_{-1}^{x} f_{w}
$$

und somit gilt (nach Voraussetzung und Definition der schwachen Ableitung)

$$
\begin{aligned}
0=-\int_{\Omega} u^{\prime} F_{w}=\int_{\Omega} u f_{w}=\int_{\Omega} u w-\left(\int_{\Omega} w\right) & \left(\int_{\Omega} u \psi\right) \\
& =\int_{\Omega}\left(u-c_{u}\right) w \quad \text { mit } c_{u}:=\int_{\Omega} u \psi
\end{aligned}
$$

Da $w \in C_{\mathrm{c}}(\Omega)$ beliebig gewählt war, folgt $u-c_{u} \in\left(C_{\mathrm{c}}(\Omega)\right)^{\perp}$ bzgl. des $L^{2}$-Skalarprodukts, also $u-c_{u}=0$ in $L^{2}(\Omega)$ aufgrund der Dichtheit von $C_{\mathrm{c}}(\Omega)$ in $L^{2}(\Omega)$. Somit gilt $u \equiv c_{u}$ f.ü. in $\Omega$.\\
(b): Sei $u \in H^{1}(\Omega)$, und sei $v: \Omega \rightarrow \mathbb{R}$ definiert durch $v(x)=\int_{0}^{x} u^{\prime}$. Für $x, y \in \Omega$ gilt dann

$$
|v(x)-v(y)|=\left|\int_{y}^{x} u^{\prime}\right| \leq|x-y|^{\frac{1}{2}}\left(\int_{\min \{x, y\}}^{\max \{x, y\}}\left(u^{\prime}\right)^{2}\right)^{\frac{1}{2}} \leq|x-y|^{\frac{1}{2}}\left\|u^{\prime}\right\|_{2}
$$

Insbesondere ist $v$ gleichmäßig stetig und lässt sich somit zu einer stetigen Funktion auf $\bar{\Omega}$ fortsetzen, welche dann (5.6) sogar für $x, y \in \bar{\Omega}$ erfüllt. Insbesondere ist $v \in L^{2}(\Omega)$. Wir zeigen nun, dass $v$ die schwache Ableitung $u^{\prime}$ besitzt. Sei dazu $\varphi \in C_{\mathrm{c}}^{1}(\Omega)$. Dann gilt mit dem Satz von Fubini

$$
\begin{aligned}
\int_{\Omega} \varphi^{\prime} v & =\int_{-1}^{1} \varphi^{\prime}(x) \int_{0}^{x} u^{\prime}(t) \mathrm{d} t \mathrm{~d} x \\
& =-\int_{-1}^{0} \varphi^{\prime}(x) \int_{x}^{0} u^{\prime}(t) \mathrm{d} t \mathrm{~d} x+\int_{0}^{1} \varphi^{\prime}(x) \int_{0}^{x} u^{\prime}(t) \mathrm{d} t \mathrm{~d} x \\
& =-\int_{-1}^{0} u^{\prime}(t) \int_{-1}^{t} \varphi^{\prime}(x) \mathrm{d} x \mathrm{~d} t+\int_{0}^{1} u^{\prime}(t) \int_{t}^{1} \varphi^{\prime}(x) \mathrm{d} x \mathrm{~d} t \\
& =-\int_{-1}^{0} u^{\prime} \varphi-\int_{0}^{1} u^{\prime} \varphi=-\int_{-1}^{1} u^{\prime} \varphi
\end{aligned}
$$

Dies zeigt $v^{\prime}=u^{\prime}$ im schwachen Sinne. Somit ist die Funktion $u-v$ gemäß (a) nach

Übergang zu einem geeigneten Repräsentanten konstant, und dies liefert die Behauptung; insbesondere ist

$$
u(y)-u(x)=v(y)-v(x)=\int_{x}^{y} u^{\prime} \quad \text { für alle } x, y \in \bar{\Omega} .
$$

Zu (c): Sei $u$ der stetige Repräsentant nach (b), also $u \in C(\bar{\Omega})$. Die erste Ungleichung folgt direkt aus (5.6). Ist zudem $y \in \bar{\Omega}$ mit $|u(y)|=\min _{\bar{\Omega}}|u|$ gewählt, so folgt mit (b):

$$
|u(x)| \leq|u(x)-u(y)|+|u(y)| \leq|x-y|^{\frac{1}{2}}\left\|u^{\prime}\right\|_{2}+|u(y)| \leq \sqrt{2}\left\|u^{\prime}\right\|_{2}+\min _{\bar{\Omega}}|u|
$$

wobei

$$
\left(\min _{\bar{\Omega}}|u|\right)^{2} \leq \frac{1}{2} \int_{\Omega}|u|^{2}=\frac{\|u\|_{2}^{2}}{2}
$$

gilt. Es folgt $\|u\|_{\infty} \leq \sqrt{2}\left\|u^{\prime}\right\|_{2}+\frac{1}{\sqrt{2}}\|u\|_{2} \leq \sqrt{2}\|u\|_{H^{1}}$, wie behauptet.\\
5.15 Bemerkung. (a) Aus Satz 5.14(c) folgt, dass $H^{1}(\Omega)$ stetig in $C(\bar{\Omega})$ eingebettet ist. Wegen der Faktorisierung

$$
H^{1}(\Omega) \hookrightarrow C^{\frac{1}{2}}(\bar{\Omega}) \hookrightarrow C(\bar{\Omega})
$$

mit stetigen Einbettungen, wobei hier $C^{\frac{1}{2}}(\bar{\Omega})$ den Raum der $\frac{1}{2}$-Hölder-stetigen Funktionen auf $\bar{\Omega}$ bezeichne, gilt sogar, dass jede beschränkte Folge in $H^{1}(\Omega)$ eine in $C(\bar{\Omega})$ konvergente Teilfolge besitzt (Übung).\\
(b) Besitzt $u \in H^{1}(\Omega)$ eine schwache Ableitung $u^{\prime} \in C(\bar{\Omega})$, so ist $u \in C^{1}(\bar{\Omega})$ (nach Wahl eines geeigneten Repräsentanten). Dies folgt unmittelbar aus (5.5) und dem klassischen Hauptsatz der Integral- und Differentialrechnung.\\
5.16 Definition und Satz. Seien Funktionen $a \in L^{\infty}(\Omega), f \in L^{2}(\Omega)$ gegeben.\\
(a) Wir nennen $u \in H^{1}(\Omega)$ eine schwache Lösung des Randwertproblems (NP), wenn

$$
\langle u, w\rangle_{*}=\int_{\Omega} f w \quad \text { für alle } w \in H^{1}(\Omega)
$$

gilt, wobei

$$
\langle u, w\rangle_{*}:=\int_{\Omega}\left(u^{\prime} w^{\prime}+a u w\right)
$$

sei.\\
(b) Ist $a_{0}:=\operatorname{essinf}_{\Omega} a>0$, so ist durch

$$
(u, w) \mapsto\langle u, w\rangle_{*}
$$

ein Skalarprodukt definiert mit der Eigenschaft, dass die induzierte Norm $\|\cdot\|_{*}$ äquivalent zur Norm $\|\cdot\|_{H^{1}}$ ist. Insbesondere ist ( $H^{1}(\Omega),\langle\cdot, \cdot\rangle_{*}$ ) ebenfalls ein Hilbertraum.

Beweis von (b). Offensichtlich ist $\langle\cdot, \cdot\rangle_{*}$ eine symmetrische Bilinearform auf $H^{1}(\Omega)$, und für $u \in H^{1}(\Omega)$ gilt

$$
a_{0}\|u\|_{2}^{2} \leq \int_{\Omega} a u^{2} \leq\|a\|_{\infty}\|u\|_{2}^{2}
$$

also

$$
\min \left\{a_{0}, 1\right\}\|u\|_{H^{1}}^{2} \leq\langle u, u\rangle_{*} \leq \max \left\{\|a\|_{\infty}, 1\right\}\|u\|_{H^{1}}^{2}
$$

Somit ist $\langle\cdot, \cdot\rangle_{*}$ ein Skalarprodukt derart, dass die induzierte Norm $\|\cdot\|_{*}$ äquivalent zur Norm $\|\cdot\|_{H^{1}}$ ist.\\
5.17 Satz. Sind $a \in L^{\infty}(\Omega), f \in L^{2}(\Omega)$ mit $a_{0}:=\operatorname{essinf}_{\Omega} a>0$ gegeben, so besitzt das Randwertproblem (NP) genau eine schwache Lösung.

Beweis. Sei $H:=H^{1}(\Omega)$, versehen mit dem Skalarprodukt $\langle\cdot, \cdot\rangle_{*}$, und sei $\varphi \in H^{\prime}$ definiert durch $\varphi(w):=\int_{\Omega} f w$. Die Stetigkeit von $\varphi$ folgt aus der Abschätzung

$$
|\varphi(w)|=\left|\langle f, w\rangle_{L^{2}(\Omega)}\right| \leq\|f\|_{2}\|w\|_{2} \leq\|f\|_{2}\left(\frac{1}{a_{0}} \int_{\Omega} a w^{2}\right)^{\frac{1}{2}} \leq \frac{1}{\sqrt{a_{0}}}\|f\|_{2}\|w\|_{*} .
$$

Nach dem Rieszschen Darstellungssatz (Satz 5.10) existiert nun genau ein $u \in H^{1}(\Omega)$ mit

$$
\langle u, w\rangle_{*}=\varphi(w)=\int_{\Omega} f w \quad \text { für alle } w \in H^{1}(\Omega)
$$

Also ist $u \in H^{1}(\Omega)$ die gesuchte eindeutige schwache Lösung von (NP).\\
5.18 Bemerkung. Seien Funktionen $a \in L^{\infty}(\Omega), f \in L^{2}(\Omega)$ gegeben.\\
(a) Ist $u$ eine schwache Lösung von (NP), so folgt

$$
\int_{\Omega} u^{\prime} w^{\prime}=-\int_{\Omega}(a u-f) w \quad \text { für alle } w \in C^{1}(\bar{\Omega}) \subseteq H^{1}(\Omega)
$$

also insbesondere für $w \in C_{\mathrm{c}}^{1}(\Omega)$. Somit ist

$$
u^{\prime} \in H^{1}(\Omega) \quad \text { mit schwacher Ableitung } \quad u^{\prime \prime}=a u-f .
$$

Insbesondere folgt mit Bemerkung 5.15(a) (nach Wahl eines geeigneten Repräsentanten) $u^{\prime} \in C(\bar{\Omega})$, und mit Bemerkung 5.15(b) folgt $u \in C^{1}(\bar{\Omega})$.\\
(b) Sind sogar $a, f \in C(\bar{\Omega})$, so folgt $u^{\prime \prime}=a u-f \in C(\bar{\Omega})$ für die schwache Ableitung von $u^{\prime}$, d.h. $u^{\prime} \in C^{1}(\bar{\Omega})$ gemäß Bemerkung 5.15(b). In diesem Fall ist also $u \in C^{2}(\bar{\Omega})$,\\
und die Differentialgleichung $u^{\prime \prime}=a u-f$ ist im klassischen Sinne in $\Omega$ erfüllt. Aus (5.7) folgt dann durch partielle Integration für alle $w \in C^{1}(\bar{\Omega})$ die Gleichung

$$
u^{\prime}(1) w(1)-u^{\prime}(-1) w(-1)=\int_{\Omega}\left(u^{\prime} w\right)^{\prime}=\int_{\Omega}\left(u^{\prime} w^{\prime}+u^{\prime \prime} w\right) \stackrel{(5.7)}{=} 0,
$$

und dies liefert die Neumann-Randbedingungen

$$
u^{\prime}(1)=u^{\prime}(-1)=0 .
$$

(c) Die Bedingungen (NR) sind auch ohne die Zusatzvoraussetzung $a, f \in C(\bar{\Omega})$ erfüllt. Tatsächlich gilt (Übung): Ist $\tilde{u} \in L^{2}(\Omega)$ und existiert $C>0$ mit

$$
\left|\int_{\Omega} \tilde{u} w^{\prime}\right| \leq C\|w\|_{2} \quad \text { für alle } w \in C^{1}(\bar{\Omega}),
$$

so ist $\tilde{u} \in H^{1}(\Omega)$, und für den Repräsentanten $\tilde{u} \in C(\bar{\Omega})$ gilt $\tilde{u}(1)=\tilde{u}(-1)=0$.

\subsection*{Orthonormalsysteme und abstrakte Fourierreihen}
Nun sei wieder ( $H,\langle\cdot, \cdot\rangle$ ) ein beliebiger Hilbertraum.\\
5.19 Definition. (a) Eine Menge $A \subseteq H$ heißt Orthonormalsystem (kurz: ONS) falls gilt:

$$
\langle x, y\rangle=\left\{\begin{array}{ll}
0, & x \neq y, \\
1, & x=y,
\end{array} \quad \text { für alle } x, y \in A .\right.
$$

(b) Ein totales Orthonormalsystem $A$ heißt Orthonormalbasis (kurz: ONB) oder Hilbertbasis. Man beachte, dass dies nicht bedeutet, dass $A$ eine Basis im Sinne der Linearen Algebra ist.\\
5.20 Satz. Sei $A=\left\{e_{1}, e_{2}, \ldots\right\}$ ein abzählbar unendliches ONS in H. Dann gelten:\\
(a) Sind $c_{n} \in \mathbb{K}, n \in \mathbb{N}$ derart, dass $x=\sum_{n=1}^{\infty} c_{n} e_{n} \in H$ existiert, so gilt $c_{n}=\left\langle x, e_{n}\right\rangle$ für alle $n \in \mathbb{N}$.\\
(b) Ist $\left(c_{n}\right)_{n} \in \ell^{2}$ gegeben, so konvergiert $\sum_{n=1}^{\infty} c_{n} e_{n}$ in $H$.

Beweis. (a) Aus der Stetigkeit und Linearität des Skalarprodukts im ersten Argument folgt

$$
\left\langle x, e_{n}\right\rangle=\left\langle\lim _{k \rightarrow \infty} \sum_{j=1}^{k} c_{j} e_{j}, e_{n}\right\rangle=\lim _{k \rightarrow \infty} \sum_{j=1}^{k} c_{j} \underbrace{\left\langle e_{j}, e_{n}\right\rangle}_{=\delta_{j n}}=c_{n} .
$$

(b) Sei $s_{n}:=\sum_{k=1}^{n} c_{k} e_{k}$. Dann gilt

$$
\left\|s_{m}-s_{n}\right\|^{2}=\left\|\sum_{k=n+1}^{m} c_{k} e_{k}\right\|^{2}=\sum_{i, j=n+1}^{m} c_{i} \overline{c_{j}}\langle\underbrace{e_{i}, e_{j}}_{\delta_{i j}}\rangle=\sum_{k=n+1}^{m}\left|c_{k}\right|^{2}
$$

für alle $n, m \in \mathbb{N}, m>n$. Da ferner $\left(c_{k}\right)_{k} \in \ell^{2}$ ist, folgt

$$
\sup _{m \geq n}\left\|s_{m}-s_{n}\right\|^{2} \leq \sum_{k=n+1}^{\infty}\left|c_{k}\right|^{2} \rightarrow 0 \quad \text { für } n \rightarrow \infty
$$

Also ist $\left(s_{n}\right)_{n}$ eine Cauchyfolge in $H$ und die Behauptung folgt.\\
5.21 Satz. Sei $A=\left\{e_{1}, \ldots, e_{n}\right\}$ ein endliches $O N S, V=\operatorname{span} A$ und $P \in \mathcal{L}(H)$ die orthogonale Projektion auf V. Dann gelten für alle $x, y \in H$ :\\
(a) $P x=\sum_{k=1}^{n}\left\langle x, e_{k}\right\rangle e_{k}$.\\
(b) $\langle P x, P y\rangle=\sum_{k=1}^{n}\left\langle x, e_{k}\right\rangle \overline{\left\langle y, e_{k}\right\rangle}$.

Beweis. (a) Sei $z:=\sum_{k=1}^{n}\left\langle x, e_{k}\right\rangle e_{k}$. Dann ist

$$
\left\langle x-z, e_{j}\right\rangle=\left\langle x, e_{j}\right\rangle-\sum_{k=1}^{n}\left\langle x, e_{k}\right\rangle\left\langle e_{k}, e_{j}\right\rangle=\left\langle x, e_{j}\right\rangle-\left\langle x, e_{j}\right\rangle=0
$$

für alle $j \in\{1, \ldots, n\}$. Also gilt $x-z \in A^{\perp}=V^{\perp}=\operatorname{Kern} P$. Es folgt

$$
P x=P z+P(x-z)=P z=z
$$

da $z$ in $V$ liegt.\\
(b) Es ist

$$
\begin{aligned}
\langle P x, P y\rangle & =\left\langle\sum_{k=1}^{n}\left\langle x, e_{k}\right\rangle e_{k}, \sum_{j=1}^{n}\left\langle y, e_{j}\right\rangle e_{j}\right\rangle \\
& =\sum_{k, j=1}^{n}\left\langle x, e_{k}\right\rangle \overline{\left\langle y, e_{j}\right\rangle}\left\langle e_{k}, e_{j}\right\rangle=\sum_{k=1}^{n}\left\langle x, e_{k}\right\rangle \overline{\left\langle y, e_{k}\right\rangle}
\end{aligned}
$$

5.22 Definition. Für ein fest gewähltes, abzählbares ONS $A=\left\{e_{1}, e_{2}, \ldots\right\}$ in $H$ und $x \in H$ nennen wir\\
(a) $\hat{x}(n):=\left\langle x, e_{n}\right\rangle$ den $n$-ten Fourierkoeffizienten von $x$ (bzgl. A),\\
(b) $\sum_{n} \hat{x}(n) e_{n}$ die (abstrakte) Fourierreihe von $x$ (bzgl. A).\\
5.23 Satz. Sei $A=\left\{e_{1}, e_{2}, \ldots\right\}$ ein abzählbar unendliches ONS in $H$. Sei ferner $V:=$ $\overline{\operatorname{span} A}$ und $P \in \mathcal{L}(H)$ die orthogonale Projektion auf $V$. Dann gelten für alle $x, y \in H$ :\\
(a) $P x=\sum_{n=1}^{\infty} \hat{x}(n) e_{n}$,\\
(b) $\langle P x, P y\rangle=\sum_{n=1}^{\infty} \hat{x}(n) \overline{\hat{y}(n)}$,\\
(c) $\sum_{n=1}^{\infty}|\hat{x}(n)|^{2} \leq\|x\|^{2} \quad$ (Besselsche Ungleichung).

Beweis. Für $n \in \mathbb{N}$ sei $P_{n} \in \mathcal{L}(H)$ die orthogonale Projektion auf $\operatorname{span}\left\{e_{1}, \ldots, e_{n}\right\}$. Dann gilt

$$
\sum_{k=1}^{n}|\hat{x}(k)|^{2} \stackrel{\text { Satz }}{=} \stackrel{5.21(\mathrm{~b})}{=}\left\|P_{n} x\right\|^{2} \stackrel{\text { Satz und Definition 5.7(b) }}{\leq}\|x\|^{2} \quad \text { für alle } n \in \mathbb{N} \text {. }
$$

Es folgt (c) und $(\hat{x}(n))_{n} \in \ell^{2}$. Gemäß Satz 5.20(b) existiert somit $z:=\sum_{n=1}^{\infty} \hat{x}(n) e_{n} \in H$ (d.h. die Reihe konvergiert). Aus der Stetigkeit des Skalarprodukts folgt

$$
\left\langle x-z, e_{j}\right\rangle=\hat{x}(j)-\sum_{n=1}^{\infty} \hat{x}(n)\left\langle e_{n}, e_{j}\right\rangle=\hat{x}(j)-\hat{x}(j)=0 \quad \text { für alle } j \in \mathbb{N},
$$

also $x-z \in A^{\perp}=V^{\perp}$ und somit $P x=P z+P(x-z)=P z=z$. Dies liefert (a).\\
Zu (b):

$$
\langle P x, P y\rangle=\left\langle\sum_{n=1}^{\infty} \hat{x}(n) e_{n}, \sum_{k=1}^{\infty} \hat{y}(k) e_{k}\right\rangle=\sum_{n, k=1}^{\infty} \hat{x}(n) \overline{\hat{y}(k)}\left\langle e_{n}, e_{k}\right\rangle=\sum_{n=1}^{\infty} \hat{x}(n) \overline{\hat{y}(n)}
$$

Hier haben wir wieder die Stetigkeit des Skalarprodukts verwendet.\\
5.24 Korollar. Sei $A=\left\{e_{1}, e_{2}, \ldots\right\}$ ein abzählbar unendliches ONS in H. Dann sind äquivalent:\\
(i) A ist eine ONB.\\
(ii) Für alle $x \in H$ gilt $x=\sum_{k=1}^{\infty} \hat{x}(k) e_{k}$.\\
(iii) Für alle $x \in H$ gilt $\|x\|^{2}=\sum_{k=1}^{\infty}|\hat{x}(k)|^{2}$.\\
(iv) Für alle $x, y \in H$ gilt $\langle x, y\rangle=\sum_{k=1}^{\infty} \hat{x}(k) \overline{\hat{y}(k)}$ (Parsevalsche Gleichung).

Beweis. Sei $V:=\overline{\operatorname{span} A}$ und $P \in \mathcal{L}(H)$ die orthogonale Projektion auf $V$. Dann gilt:

$$
A \text { total } \Longleftrightarrow V=H \quad \Longleftrightarrow \quad P=I .
$$

Somit erhalten wir die Implikationen:

$$
(\mathrm{i}) \stackrel{\text { Satz } 5.23}{\Longrightarrow}(\mathrm{ii}) \stackrel{\text { Stet. d. SKP }}{\Longrightarrow}(\mathrm{iv}) \stackrel{\text { trivial }}{\Longrightarrow}(\mathrm{iii}) \text {. }
$$

Gilt (iii), so folgt für alle $x \in H$ :

$$
\|P x\|^{2}+\|x-P x\|^{2} \stackrel{\text { Satz und Definition 5.6(a) }}{=}\|x\|^{2}=\sum_{k=1}^{\infty}|\hat{x}(k)|^{2} \stackrel{\text { Satz }}{=} \stackrel{5.23(\mathrm{~b})}{=}\|P x\|^{2}
$$

also $x=P x \in V$. Also ist $H=V$ und somit $A$ total, d.h. (i) gilt.\\
5.25 Bemerkung. Ist $A$ ein abzählbar unendliches ONS und existiert eine totale Teilmenge $M \subseteq H$ mit

$$
\|x\|^{2}=\sum_{k=1}^{\infty}|\hat{x}(k)|^{2} \quad \text { für alle } x \in M,
$$

so ist $A$ bereits eine ONB.\\
Beweis. Sei wieder $V:=\overline{\operatorname{span} A}$. Wie im Beweisschritt (iii) auf (i) in Korollar 5.24 erhält man mit (5.8) die Implikation: $x \in M \Longrightarrow x \in V$. Somit gilt $M \subseteq V$ und damit $H=\overline{\operatorname{span} M} \subseteq V$. Es folgt, dass $A$ eine ONB ist.\\
5.26 Satz. Äquivalent sind:\\
(i) $H$ ist unendlich-dimensional und separabel.\\
(ii) $H$ besitzt eine abzählbar unendliche ONB.\\
(iii) $H$ ist isometrisch isomorph $z u \ell^{2}$.

Man beachte: Die Aussagen gelten insbesondere im Fall $H=L^{2}(\Omega)$, wobei $\Omega \subseteq \mathbb{R}^{N}$ eine Lebesgue-messbare Menge mit $\lambda(\Omega)>0$ sei (vgl. Bemerkung 4.48).

Beweis. Aus (ii) folgt (iii): Wir betrachten die lineare Abbildung $\mathcal{F}: H \rightarrow \ell^{2}, \mathcal{F} x:=$ $(\hat{x}(k))_{k}$. Aus Korollar 5.24(iii) folgt

$$
\|\mathcal{F} x\|_{2}=\|x\|_{H} \quad \text { für alle } x \in H .
$$

Die Surjektivität von $\mathcal{F}$ folgt aus Satz 5.20(b). Somit ist $\mathcal{F}$ ein isometrischer Isomorphismus.\\
Aus (iii) folgt (i), da $\ell^{2}$ separabel ist. Noch zu zeigen: (i) $\Longrightarrow$ (ii). Sei dazu $M=$ $\left\{a_{1}, a_{2}, \ldots\right\} \subseteq H$ eine abzählbar unendliche totale Menge. Wir erzeugen hieraus eine ONB mittels des Gram-Schmidtschen Orthogonalisierungsverfahrens: Ohne Einschränkung gelte $a_{1} \neq 0$ und $a_{k+1} \notin W_{k}:=\operatorname{span}\left\{a_{1}, \ldots, a_{k}\right\}$ für $k \in \mathbb{N}$. Sei $P_{k} \in \mathcal{L}(H)$ die orthogonale Projektion auf $W_{k}$ für $k \in \mathbb{N}$. Wir setzen $e_{1}:=\frac{a_{1}}{\left\|a_{1}\right\|}$ und

$$
e_{k}:=\frac{a_{k}-P_{k-1} a_{k}}{\left\|a_{k}-P_{k-1} a_{k}\right\|} \in W_{k} \cap W_{k-1}^{\perp} \quad \text { für } k \geq 2 .
$$

Nach Konstruktion ist $A:=\left\{e_{1}, e_{2}, \ldots\right\}$ ein ONS in $H$. Per Induktion sieht man ferner: $W_{k}=\operatorname{span}\left\{e_{1}, \ldots, e_{k}\right\}$ für alle $k \in \mathbb{N}$. Es folgt

$$
\operatorname{span} M=\bigcup_{k \in \mathbb{N}} W_{k}=\operatorname{span} A
$$

also ist $A$ total und somit eine ONB. Es folgt (ii).

\subsection*{Klassische Fourierreihen}
Im Folgenden betrachten wir den Hilbertraum $H:=L^{2}[0,2 \pi]=L^{2}([0,2 \pi], \mathbb{C})$ (bezüglich des Lebesguemaßes) mit Skalarprodukt $\langle\cdot, \cdot\rangle$ definiert durch

$$
\langle f, g\rangle=\int_{0}^{2 \pi} f \bar{g} \quad \text { für } f, g \in H
$$

5.27 Definition und Satz. Für $k \in \mathbb{Z}$ sei $e_{k} \in H$ definiert durch $e_{k}(t):=\frac{1}{\sqrt{2 \pi}} \mathrm{e}^{\mathrm{i} k t}$. Dann ist $A:=\left\{e_{k} \mid k \in \mathbb{Z}\right\}$ eine ONB von $H$. Die Fourierkoeffizienten von $f \in H$ bzgl. $A$ sind dann gegeben durch

$$
\hat{f}(k):=\frac{1}{\sqrt{2 \pi}} \int_{0}^{2 \pi} f(t) \mathrm{e}^{-\mathrm{i} k t} \mathrm{~d} t, \quad \text { für } k \in \mathbb{Z}
$$

Beweis. Es ist

$$
\left\langle e_{k}, e_{k}\right\rangle=\frac{1}{2 \pi} \int_{0}^{2 \pi} \underbrace{e^{\mathrm{i} k t} e^{-\mathrm{i} k t}}_{=1} \mathrm{~d} t=1 \quad \text { für alle } k \in \mathbb{Z}
$$

und

$$
\left\langle e_{k}, e_{j}\right\rangle=\frac{1}{2 \pi} \int_{0}^{2 \pi} \mathrm{e}^{\mathrm{i}(k-j) t} \mathrm{~d} t=\left.\frac{1}{2 \pi} \frac{1}{\mathrm{i}(k-j)} \mathrm{e}^{\mathrm{i}(k-j) t}\right|_{0} ^{2 \pi}=0 \quad \text { für } k \neq j .
$$

Es folgt, dass $A$ ein ONS in $H=L^{2}([0,2 \pi])$ ist. Ferner ist $A$ total in $H$, da die Menge der trigonometrischen Polynome dicht in $H$ liegt. Die letzte Aussage ist eine leichte Folgerung aus Definition und Korollar 4.12 und Satz 4.47 (Übung).\\
5.28 Bemerkung. Aus Definition und Satz 5.27 folgt, dass für jedes $f \in H$ die Fourierreihe $\sum_{k \in \mathbb{Z}} \hat{f}(k) e_{k}$ bzgl. $\|\cdot\|_{2}$ gegen $f$ konvergiert. Weiterhin gelten:\\
(a) Die Fourierreihe konvergiert punktweise f.ü. gegen $f$ (Satz von Carleson, 1966).\\
(b) Ist $(\hat{f}(k))_{k} \in \ell^{1}(\mathbb{Z})$, so konvergiert die Fourierreihe in $\left(C_{2 \pi}(\mathbb{R}, \mathbb{C}),\|\cdot\|_{\infty}\right)$, also gleichmäßig. Es folgt dann, dass $f$ (nach Wahl eines geeigneten Repräsentanten) zu einem Element in $C_{2 \pi}(\mathbb{R}, \mathbb{C})$ fortgesetzt werden kann.\\
(c) Ist $f \in C_{2 \pi}(\mathbb{R}, \mathbb{C})$, so muss die Fourierreihe i.A. trotzdem nicht punktweise gegen $f$ konvergieren.\\
Beweis. (b) Weil die Folge $(\hat{f}(k))_{k}$ in $\ell^{1}(\mathbb{Z})$ liegt, gilt für alle $t \in[0,2 \pi]$ :

$$
\sum_{k \in \mathbb{Z}}\left|\hat{f}(k) e_{k}(t)\right|=\frac{1}{\sqrt{2 \pi}} \sum_{k \in \mathbb{Z}}|\hat{f}(k)|<\infty
$$

d.h. die Funktionenreihe $\tilde{f}:=\sum_{k} \hat{f}(k) e_{k}$ konvergiert nach dem Majorantenkriterium absolut und gleichmäßig. Zudem liegen alle Funktionen $e_{k}$ in $C_{2 \pi}(\mathbb{R}, \mathbb{C})$; dies gilt also auch für $\tilde{f}$. Wegen (a) ist dann $\tilde{f}$ ein Repräsentant von $f$ in $C_{2 \pi}(\mathbb{R}, \mathbb{C})$.

Skizze des Beweises von (c): Sei $t_{0} \in[0,2 \pi]$, und sei im Folgenden $C_{2 \pi}:=C_{2 \pi}(\mathbb{R}, \mathbb{C})$. Wir zeigen, dass ein $f \in C_{2 \pi}$ existiert derart, dass die Fourierreihe von $f$ in $t_{0}$ nicht gegen $f$ konvergiert. Wir verwenden den Satz von der gleichmäßigen Beschränktheit. Dazu definieren wir die Operatoren

$$
S_{n} \in \mathcal{L}\left(C_{2 \pi}\right), \quad S_{n} f:=\sum_{k=-n}^{n} \hat{f}(k) e_{k}
$$

mit $e_{k}$ wie in Definition und Satz 5.27. Dann gilt für alle $t \in[0,2 \pi]$ :

$$
\left(S_{n} f\right)(t)=\frac{1}{2 \pi} \int_{0}^{2 \pi} f(s)\left(\sum_{k=-n}^{n} \mathrm{e}^{\mathrm{i} k(t-s)}\right) \mathrm{d} s=\frac{1}{2 \pi} \int_{0}^{2 \pi} d_{n}(t-s) f(s) \mathrm{d} s
$$

mit dem $n$-ten Dirichletkern

$$
\tau \mapsto d_{n}(\tau):=\sum_{k=-n}^{n} \mathrm{e}^{\mathrm{i} k \tau}= \begin{cases}\frac{\sin \left(\left(n+\frac{1}{2}\right) \tau\right)}{\sin \frac{\tau}{2}}, & \tau \in \mathbb{R} \backslash\{2 k \pi \mid k \in \mathbb{Z}\} \\ 2 n+1, & \tau \in\{2 k \pi \mid k \in \mathbb{Z}\}\end{cases}
$$

Die letztere Darstellung des Kerns kann man dabei z.B. mit der geometrischen Summenformel herleiten (Übung). Wir betrachten nun die stetigen Linearformen

$$
T_{n} \in C_{2 \pi}^{\prime}, \quad T_{n}(f):=\left(S_{n} f\right)\left(t_{0}\right)=\frac{1}{2 \pi} \int_{0}^{2 \pi} d_{n}\left(t_{0}-s\right) f(s) \mathrm{d} s
$$

Ähnlich wie in Beispiel 2.29(c) sieht man, dass für die zugehörige Operatornorm dann

$$
\left\|T_{n}\right\|=\frac{1}{2 \pi} \int_{0}^{2 \pi}\left|d_{n}\right|
$$

gilt (Übung). Allerdings gilt wegen (5.9) auch

$$
\int_{0}^{2 \pi}\left|d_{n}\right| \rightarrow \infty \quad \text { für } n \rightarrow \infty
$$

(Übung). Somit ist die Folge der Linearformen $T_{n}, n \in \mathbb{N}$, nicht beschränkt. Nach Hauptsatz 3.8 ist sie damit auch nicht punktweise beschränkt, d.h. es existiert ein $f \in C_{2 \pi}$ derart, dass die im Punkt $t_{0}$ ausgewerteten Partialsummen

$$
T_{n} f=\left(S_{n} f\right)\left(t_{0}\right), \quad n \in \mathbb{N},
$$

der Fourierreihe von $f$ eine unbeschränkte Folge in $\mathbb{C}$ bilden und somit nicht konvergieren.

\subsection*{Der Satz von Lax-Milgram mit Anwendungen}
Stets sei $(H,\langle\cdot, \cdot\rangle)$ ein Hilbertraum über $\mathbb{K}$.\\
5.29 Definition. Eine Abbildung $a: H \times H \rightarrow \mathbb{K}$ heißt sesquilinear, falls gilt:\\
(i) $a(\cdot, y): H \rightarrow \mathbb{K}$ ist linear für alle $y \in H$.\\
(ii) $a(x, \cdot): H \rightarrow \mathbb{K}$ ist antilinear für alle $x \in H$, d.h. $a(x, \lambda y+z)=\bar{\lambda} a(x, y)+a(x, z)$ für alle $x, y, z \in H, \lambda \in \mathbb{K}$.

Wir setzen $\operatorname{Sql}(H):=\{a: H \times H \rightarrow \mathbb{K} \mid a$ sesquilinear $\}$.\\
5.30 Satz und Definition. Zu $T \in \mathcal{L}(H)$ sei $a_{T} \in \operatorname{Sql}(H)$ definiert durch $a_{T}(x, y):=$ $\langle x, T y\rangle$ für alle $x, y \in H$. Dann gelten:\\
(a) $\left|a_{T}(x, y)\right| \leq\|T\|\|x\|\|y\|$ für alle $x, y \in H$.\\
(b) Existiert $c_{1}>0$ mit $\left|a_{T}(x, x)\right| \geq c_{1}\|x\|^{2}$ für alle $x \in H$, so ist $T$ ein topologischer Isomorphismus.

Beweis. (a) $\left|a_{T}(x, y)\right|=|\langle x, T y\rangle| \leq\|x\|\|T y\| \leq\|x\|\|T\|\|y\|$ für alle $x, y \in H$.\\
(b) Für alle $x \in H$ ist $c_{1}\|x\|^{2} \leq|\langle x, T x\rangle| \leq\|T x\|\|x\|$, also

$$
\|x\| \leq \frac{1}{c_{1}}\|T x\|
$$

Somit ist $T$ injektiv. Sei $V:=$ Bild $T$. Wir behaupten, dass $V$ abgeschlossen ist. Um das zu zeigen, seien $\left(x_{n}\right)_{n} \subseteq H$ und $y \in H$ mit $T x_{n} \rightarrow y$ für $n \rightarrow \infty$. Dann gilt

$$
\sup _{m \geq n}\left\|x_{m}-x_{n}\right\| \stackrel{(5.10)}{\leq} \frac{1}{c_{1}} \sup _{m \geq n}\left\|T x_{m}-T x_{n}\right\| \rightarrow 0 \quad \text { für } n \rightarrow \infty .
$$

Also ist $\left(x_{n}\right)_{n}$ eine Cauchyfolge, und somit existiert $x:=\lim _{n \rightarrow \infty} x_{n} \in H$. Da $T$ stetig ist, gilt $y=T x \in V$. Es folgt die Behauptung. Für $y \in V^{\perp}$ gilt nun $c_{1}\|y\|^{2} \leq|\langle y, T y\rangle|=0$, also $V^{\perp}=\{0\}$. Es folgt $V=\bar{V}=V^{\perp \perp}=\{0\}^{\perp}=H$, d.h. $T$ ist surjektiv. Schließlich gilt wegen (5.10) $\left\|T^{-1} y\right\| \leq \frac{1}{c_{1}}\left\|T\left(T^{-1} y\right)\right\|=\frac{1}{c_{1}}\|y\|$ für alle $y \in H \Longrightarrow$, d.h. $T^{-1} \in \mathcal{L}(H)$ und $\left\|T^{-1}\right\| \leq \frac{1}{c_{1}}$.\\
5.31 Satz (von Lax-Milgram). Seien $a \in \operatorname{Sql}(H)$ und $c_{0}>0$ mit

$$
|a(x, y)| \leq c_{0}\|x\|\|y\| \quad \text { für alle } x, y \in H
$$

gegeben. Dann gelten:\\
(a) Es existiert genau ein $T \in \mathcal{L}(H)$ mit $a=a_{T}$, d.h. mit $a(x, y)=\langle x, T y\rangle$ für alle $x, y \in H$.\\
(b) Existiert ferner eine Konstante $c_{1}>0$ mit

$$
|a(x, x)| \geq c_{1}\|x\|^{2} \quad \text { für alle } x \in H,
$$

so ist $T$ ein topologischer Isomorphismus.\\
(c) Sind (5.11) und (5.12) erfüllt, so existiert zu jedem $f \in H$ genau ein $y \in H$ mit

$$
a(x, y)=\langle x, f\rangle \quad \text { für alle } x \in H
$$

(d) Sind (5.11) und (5.12) erfüllt, so existiert zu jedem $\varphi \in H^{\prime}$ genau ein $y \in H$ mit

$$
a(x, y)=\varphi(x) \quad \text { für alle } x \in H .
$$

Beweis. (a) Sei $J: H \rightarrow H^{\prime}$ die antilineare isometrische Bijektion aus dem Rieszschen Darstellungssatz (Satz 5.10). Für $y \in H$ sei ferner $\psi_{y}: H \rightarrow \mathbb{K}$ definiert durch $\psi_{y}(x):=$ $a(x, y)$. Dann ist $\psi_{y}$ linear und stetig mit $\left\|\psi_{y}\right\| \leq c_{0}\|y\|$, denn wegen (5.11) gilt $\left|\psi_{y}(x)\right| \leq$ $c_{0}\|x\|\|y\|$ für alle $x \in H$. Gilt $a=a_{T}$ für ein $T \in \mathcal{L}(H)$, so folgt

$$
\langle x, T y\rangle=a_{T}(x, y)=\psi_{y}(x) \quad \text { für alle } y, x \in H,
$$

also notwendigerweise

$$
T y=J^{-1} \psi_{y} \quad \text { für alle } y \in H .
$$

Insbesondere ist $T$ eindeutig bestimmt. Ist $T$ durch (5.14) definiert, so gilt für $y_{1}, y_{2} \in H$, $\lambda \in \mathbb{K}$ :

$$
\begin{aligned}
T\left(\lambda y_{1}+y_{2}\right)=J^{-1} \psi_{\lambda y_{1}+y_{2}} \stackrel{a \in \operatorname{Sql}(H)}{=} J^{-1}\left(\bar{\lambda} \psi_{y_{1}}+\psi_{y_{2}}\right) & \\
& =\lambda J^{-1} \psi_{y_{1}}+J^{-1} \psi_{y_{2}}=\lambda T y_{1}+T y_{2}
\end{aligned}
$$

Also ist $T$ ist linear. Ferner gilt

$$
\|T y\|=\left\|J^{-1} \psi_{y}\right\| \stackrel{\text { Satz } 5.10}{=}\left\|\psi_{y}\right\| \leq c_{0}\|y\| \quad \text { für alle } y \in H,
$$

also $T \in \mathcal{L}(H)$.\\
(b) Dies folgt direkt aus (a) und Satz und Definition 5.30(b).\\
(c) Ist $T \in \mathcal{L}(H)$ gemäß (a) gewählt, so ist (5.13) äquivalent zu

$$
\langle x, T y\rangle=\langle x, f\rangle \quad \text { für alle } x \in H
$$

wobei $T$ gemäß (b) ein topologischer Isomorphismus ist. Also existiert genau ein $y \in H$, welches (5.15) erfüllt, und dieses ist gegeben durch $y:=T^{-1} f$.\\
(d) folgt aus (c) und dem Darstellungssatz von Riesz, Satz 5.10.\\
5.32 Bemerkung. Der Vorteil der Aussage Satz 5.31(d) gegenüber dem Rieszschen Darstellungssatz (wo die betrachtete Sesquilinearform a das Skalarprodukt ist) liegt darin, dass hier keine Symmetrie von $a$ verlangt wird.\\
5.33 Beispiel. Seien $\mathbb{K}=\mathbb{R}$ und $\Omega=(-1,1) \subseteq \mathbb{R}$. Gegeben seien Funktionen $h, b \in$ $L^{\infty}(\Omega), f \in L^{2}(\Omega)$, wobei $\operatorname{ess} \inf _{\Omega} h>0$ und

$$
\frac{\|b\|_{\infty}}{2}<\min \{1, \underset{\Omega}{\operatorname{ess} \inf } h\}
$$

sei. Wir suchen eine Lösung des Neumann-Randwertproblems

$$
\left\{\begin{aligned}
-u^{\prime \prime}+b(x) u^{\prime}+h(x) u & =f(x), \quad x \in \Omega \\
u^{\prime}(1)=u^{\prime}(-1) & =0
\end{aligned}\right.
$$

Ähnlich wie in Definition und Satz 5.16 nennen wir $u \in H^{1}(\Omega)$ eine schwache Lösung von (5.17), wenn gilt:

$$
a(w, u):=\int_{\Omega}\left(u^{\prime} w^{\prime}+b u^{\prime} w+h u w\right)=\int_{\Omega} f w \quad \text { für alle } w \in H^{1}(\Omega)
$$

Dabei ist $a: H^{1}(\Omega) \times H^{1}(\Omega) \rightarrow \mathbb{R}$ eine Bilinearform (wegen $\mathbb{K}=\mathbb{R}$ also auch eine Sesquilinearform) mit

$$
\begin{aligned}
|a(u, w)| \leq\left\|u^{\prime}\right\|_{2}\left\|w^{\prime}\right\|_{2}+\|b\|_{\infty}\left\|u^{\prime}\right\|_{2}\|w\|_{2}+\|h\|_{\infty}\|u\|_{2}\|w\|_{2} & \\
\leq & \left(1+\|b\|_{\infty}+\|h\|_{\infty}\right)\|u\|_{H^{1}}\|w\|_{H^{1}}
\end{aligned}
$$

für alle $u, w \in H^{1}(\Omega)$. Umgekehrt haben wir

$$
\begin{aligned}
|a(u, u)| & \geq\left\|u^{\prime}\right\|_{2}^{2}-\|b\|_{\infty}\left\|u^{\prime}\right\|_{2}\|u\|_{2}+(\underset{\Omega}{\operatorname{ess} \inf } h)\|u\|_{2}^{2} \\
& \geq\left(1-\frac{\|b\|_{\infty}}{2}\right)\left\|u^{\prime}\right\|_{2}^{2}+\left(\underset{\Omega}{\operatorname{ess} \inf } h-\frac{\|b\|_{\infty}}{2}\right)\|u\|_{2}^{2} \geq c_{1}\|u\|_{H^{1}}^{2}
\end{aligned}
$$

für $u \in H^{1}(\Omega)$ mit $c_{1}:=\min \left\{1, \operatorname{essinf}_{\Omega} h\right\}-\frac{\|b\|_{\infty}}{2}>0$ wegen (5.16). Somit erfüllt die Bilinearform $a$ auf $H^{1}(\Omega)$ die Voraussetzungen des Satzes von Lax-Milgram. Ähnlich wie im Beweis von Satz 5.17 betrachten wir nun die Linearform

$$
\varphi \in H^{1}(\Omega)^{\prime}, \quad \varphi(w):=\int_{\Omega} f w
$$

Nach Satz 5.31(d) existiert genau ein $u \in H^{1}(\Omega)$ mit $a(w, u)=\varphi(w)$ für alle $w \in H^{1}(\Omega)$, d.h. genau eine schwache Lösung $u \in H^{1}(\Omega)$ von (5.17). Für diese schwache Lösung kann man analoge Regularitätsaussagen wie in Bemerkung 5.18 erhalten. Man beachte: Die hier definierte Bilinearform $a$ ist nicht symmetrisch; man kann die (eindeutige) Existenz schwacher Lösungen daher nicht allein aus dem Rieszschen Darstellungssatz folgern.\\
5.34 Bemerkung. Selbst das einfachere Problem

$$
\left\{\begin{aligned}
-u^{\prime \prime}+h(x) u & =f(x), \quad x \in \Omega \\
u^{\prime}(1)=u^{\prime}(-1) & =0
\end{aligned}\right.
$$

(vgl. Abschnitt 5.2) ohne den Term erster Ordnung kann auf eine nicht symmetrische Bilinearform führen, wenn man eine komplexwertige Funktion $x \mapsto h(x)$ betrachtet. Analog wie in Satz 5.17 kann man dann die Existenz genau einer schwachen Lösung aus dem Satz von Lax-Milgram folgern, wenn $\operatorname{ess}_{\inf } \operatorname{Re} h>0$ gilt. Hierzu betrachtet man $L^{2}(\Omega)=L^{2}(\Omega, \mathbb{C})$ und einen analog definierten komplexwertigen Sobolevraum $H^{1}(\Omega)=H^{1}(\Omega, \mathbb{C})$.\\
5.35 Beispiel. Seien $h, f \in C([0,1], \mathbb{C})$ Funktionen mit $\operatorname{Re} h>0$ auf $[0,1]$. Wir zeigen, dass das Dirichletproblem

$$
-u^{\prime \prime}+h u=f \quad \text { in }(0,1), \quad u(0)=u(1)=0
$$

eine Lösung $u \in C^{2}([0,1], \mathbb{C})$ besitzt. Man kann diese Gleichung ähnlich wie das Neumannsche Randwertproblem lösen, muss dann aber im Abschluss des Raumes $C_{\mathrm{c}}^{\infty}(\Omega)$ in $H^{1}(\Omega)$ statt in $H^{1}(\Omega)$ (wie in Abschnitt 5.2) arbeiten, um die Randbedingungen zu erfüllen.

Wir beschreiten hier einen anderen Weg. Dazu setzen wir $H:=L^{2}([0,1], \mathbb{C})$ und schreiben das Problem in der Form

$$
u=K(f-h u) \quad \Longleftrightarrow \quad K \tilde{u}+\frac{\tilde{u}}{h}=K f, \tilde{u}=h u
$$

mit dem Integraloperator

$$
K: H \rightarrow H, \quad(K v)(t):=\int_{0}^{1} G(t, s) v(s) \mathrm{d} s \quad \text { für } t \in[0,1]
$$

zur Greenfunktion

$$
G \in C\left([0,1]^{2}\right), \quad G(t, s)= \begin{cases}(1-s) t, & t \leq s \\ (1-t) s, & t>s\end{cases}
$$

aus Kapitel 1. Da $G$ beschränkt ist, liegt $K$ in $\mathcal{L}(H)$. Ferner definieren wir $a \in \operatorname{Sql}(H)$ durch

$$
a(w, v):=\int_{0}^{1} w \overline{\left(K v+\frac{v}{h}\right)} \quad \text { für } v, w \in H
$$

Aufgrund der Beschränktheit der Funktionen $G$ und $\frac{1}{h}$ sieht man leicht, dass dann die Voraussetzung (5.11) aus Satz 5.31 erfüllt ist. Betrachtet man ferner $v \in C_{\mathrm{c}}(0,1)$ und setzt $w:=K v$, so ist $w \in C^{2}[0,1]$ eine Lösung von

$$
-w^{\prime \prime}=v \quad \text { in }(0,1), \quad w(0)=w(1)=0
$$

(Übung) und somit folgt mit partieller Integration

$$
\int_{0}^{1} v \overline{(K v)}=-\int_{0}^{1} w^{\prime \prime} \bar{w}=\int_{0}^{1} w^{\prime} \bar{w}^{\prime}=\int_{0}^{1}\left|w^{\prime}\right|^{2} \geq 0
$$

also

$$
\operatorname{Re} a(v, v) \geq \operatorname{Re} \int_{0}^{1} v \overline{\left(\frac{v}{h}\right)} \geq c_{1} \int_{0}^{1}|v|^{2}=c_{1}\|v\|_{2}^{2}
$$

mit

$$
c_{1}:=\min _{t \in[0,1]} \operatorname{Re} \frac{1}{h(t)}>0
$$

nach Voraussetzung. Da $C_{\mathrm{c}}(0,1)$ gemäß Satz 4.47 in $H$ dicht liegt, erhält man durch Approximation auch

$$
|a(v, v)| \geq \operatorname{Re} a(v, v) \geq c_{1}\|v\|_{2}^{2} \quad \text { für alle } v \in H .
$$

Also erfüllt $a$ die Voraussetzungen (5.11) und (5.12) von Satz 5.31, und somit existiert gemäß Satz 5.31(c) zu $K f \in H$ genau ein $\tilde{u} \in H$ mit

$$
a(w, \tilde{u})=\langle w, K f\rangle \quad \text { für alle } w \in H,
$$

d.h.

$$
\int_{0}^{1} w \overline{\left(K \tilde{u}+\frac{\tilde{u}}{h}-K f\right)}=0 \quad \text { für } w \in H .
$$

Dies liefert $K \tilde{u}+\frac{\tilde{u}}{h}-K f \in H^{\perp}=\{0\}$, d.h. $u:=\frac{\tilde{u}}{h}$ erfüllt (5.20). Aus Eigenschaften der Greenfunktion $G$ (Übung) erhält man dann, dass $u \in C^{2}([0,1])$ eine Lösung von (5.19) ist.

Ende des Inhalts für 2024-06-07

\subsection*{Der Adjungierte Operator}
5.36 Satz und Definition. Zu jedem $T \in \mathcal{L}(H)$ existiert genau ein $T^{*} \in \mathcal{L}(H)$ mit $\langle T x, y\rangle=\left\langle x, T^{*} y\right\rangle$ für alle $x, y \in H$. Man nennt $T^{*}$ den zu $T$ adjungierten Operator.

Beweis. Definiere $a \in \operatorname{Sql}(H)$ durch $a(x, y)=\langle T x, y\rangle$. Dann ist

$$
|a(x, y)| \leq\|T x\|\|y\| \leq\|T\|\|x\|\|y\| \quad \text { für } x, y \in H,
$$

also gilt (5.11) von Satz 5.31 mit $c_{0}:=\|T\|$. Somit existiert genau ein $T^{*} \in \mathcal{L}(H)$ mit

$$
\langle T x, y\rangle=a(x, y)=\left\langle x, T^{*} y\right\rangle \quad \text { für alle } x, y \in H .
$$

5.37 Satz. Für alle $S, T \in \mathcal{L}(H)$ gilt:\\
(a) $T^{* *}=T$\\
(b) $(\lambda S+\mu T)^{*}=\bar{\lambda} S^{*}+\bar{\mu} T^{*}$ für $\lambda, \mu \in \mathbb{K}$.\\
(c) $(S T)^{*}=T^{*} S^{*}$\\
(d) Ist $T$ ein topologischer Isomorphismus, so auch $T^{*}$; und dann gilt $\left(T^{*}\right)^{-1}=\left(T^{-1}\right)^{*}$.\\
(e) $\left\|T^{*}\right\|=\|T\|$\\
(f) $\left\|T^{*} T\right\|=\|T\|^{2}$

Beweis. (a)-(c): Leichte Übung, z.B. (c): Für alle $x, y \in H$ ist $\langle S T x, y\rangle=\left\langle T x, S^{*} y\right\rangle=$ $\left\langle x, T^{*} S^{*} y\right\rangle$, und mit Satz und Definition 5.36 (Eindeutigkeit) folgt $(S T)^{*}=T^{*} S^{*}$.

Zu (d): Gemäß (c) ist $\left(T^{-1}\right)^{*} T^{*}=\left(T T^{-1}\right)^{*}=I^{*}=I$ und analog $T^{*}\left(T^{-1}\right)^{*}=I$. Somit folgt die Behauptung.\\
$\mathrm{Zu}(\mathrm{e})$ und (f): Für alle $x \in H$ ist

$$
\|T x\|^{2}=\left\langle x, T^{*} T x\right\rangle \leq\left\|T^{*} T x\right\|\|x\| \leq\left\|T^{*} T\right\|\|x\|^{2}
$$

also

$$
\|T\|^{2}=\left(\sup _{\|x\| \leq 1}\|T x\|\right)^{2}=\sup _{\|x\| \leq 1}\|T x\|^{2} \leq\left\|T^{*} T\right\| \leq\left\|T^{*}\right\|\|T\|
$$

Es folgt $\|T\| \leq\left\|T^{*}\right\|$, und dann auch $\left\|T^{*}\right\| \leq\left\|T^{* *}\right\|=\|T\|$. Somit ist $\|T\|=\left\|T^{*}\right\|$, und aus (5.21) folgt nun $\left\|T^{*} T\right\|=\|T\|^{2}$.\\
5.38 Definition. $T \in \mathcal{L}(H)$ heißt

\begin{itemize}
  \item normal, falls $T T^{*}=T^{*} T$ gilt.
  \item hermitesch oder selbstadjungiert, falls $T=T^{*}$ gilt.
  \item unitär, falls $T T^{*}=I=T^{*} T$, also $T^{*}=T^{-1}$ gilt.\\
5.39 Beispiel (Multiplikationsoperatoren). Seien $\Omega \subseteq \mathbb{R}^{N}$ (Lebesgue-)messbar, $H:=$ $L^{2}(\Omega), q \in L^{\infty}(\Omega)$ und $T \in \mathcal{L}(H)$ definiert durch $(T f)(x):=q(x) f(x)$ für $f \in H$ und $x \in \Omega$. Wegen
\end{itemize}

$$
\|T f\|_{2}^{2}=\int_{\Omega}|q f|^{2} \leq\|q\|_{\infty}^{2}\|f\|_{2}^{2} \quad \text { für alle } f \in H
$$

ist $T$ tatsächlich stetig mit $\|T\| \leq\|q\|_{\infty}$. Dabei gilt:

$$
\left\langle f, T^{*} g\right\rangle=\langle T f, g\rangle=\int_{\Omega} q f \bar{g}=\int_{\Omega} f \overline{\bar{q}} \bar{g}=\langle f, \bar{q} g\rangle
$$

Also ist $T^{*} \in \mathcal{L}(H)$ gegeben durch $T^{*} g=\bar{q} g$. Beachte: $T$ ist normal, denn $T T^{*} g=$ $T^{*} T g=|q|^{2} g$ für alle $g \in H$. Ferner gilt:

\begin{itemize}
  \item $T$ unitär $\Longleftrightarrow|q|=1$ f.ü. auf $\Omega$;
  \item $T$ selbstadjungiert $\Longleftrightarrow q$ f.ü. auf $\Omega$ reellwertig.
\end{itemize}

Randbemerkung: $T$ definiert ebenfalls einen stetigen Operator $L^{p}(\Omega) \rightarrow L^{p}(\Omega)$ für $1 \leq p \leq \infty$, wobei auch dann $\|T\| \leq\|q\|_{\infty}$ gilt.\\
5.40 Beispiel (Integraloperatoren). Sei $\Omega \subseteq \mathbb{R}^{N}$ (Lebesgue-)messbar und $H:=L^{2}(\Omega)$. Sei ferner $k \in L^{2}(\Omega \times \Omega)$. Der Integraloperator $K \in \mathcal{L}(H)$ zur Kernfunktion $k$ ist definiert durch

$$
(K f)(x):=\int_{\Omega} k(x, y) f(y) \mathrm{d} y \quad \text { für } f \in H \text { und } x \in \Omega \text { (f.ü.). }
$$

$K$ ist wohldefiniert, denn: Da $\int_{\Omega \times \Omega}|k(x, y)|^{2} d(x, y)<\infty$ ist, gilt nach dem Satz von Fubini:

\begin{itemize}
  \item $\int_{\Omega}|k(x, y)|^{2} \mathrm{~d} y<\infty$ für fast alle $x \in \Omega$;
  \item $x \mapsto \int_{\Omega}|k(x, y)|^{2} \mathrm{~d} y$ ist integrierbar.
\end{itemize}

Aufgrund der Cauchy-Schwarz-Ungleichung ist somit ( $K f$ )( $x$ ) für $f \in H$ und fast alle $x \in \Omega$ wohldefiniert, und es gilt

$$
|(K f)(x)|^{2}=\left|\int_{\Omega} k(x, y) f(y) \mathrm{d} y\right|^{2} \leq\|f\|_{2}^{2} \int_{\Omega}|k(x, y)|^{2} \mathrm{~d} y
$$

Es folgt

$$
\|K f\|_{2}^{2}=\int_{\Omega}|(K f)(x)|^{2} \mathrm{~d} x \leq\|f\|_{2}^{2} \int_{\Omega} \int_{\Omega}|k(x, y)|^{2} \mathrm{~d} y \mathrm{~d} x=\|f\|_{2}^{2}\|k\|_{L^{2}(\Omega \times \Omega)}^{2}
$$

also $K \in \mathcal{L}(H)$ mit $\|K\| \leq\|k\|_{L^{2}(\Omega \times \Omega)}$.\\
Es gilt dann: $K^{*} \in \mathcal{L}(H)$ ist der Integraloperator zur Kernfunktion

$$
k^{*} \in L^{2}(\Omega \times \Omega), \quad k^{*}(x, y):=\overline{k(y, x)},
$$

denn für $f, g \in H$ gilt nach Tonelli und Fubini:

$$
\begin{aligned}
\langle K f, g\rangle & =\int_{\Omega}\left(\int_{\Omega} k(x, y) f(y) \mathrm{d} y\right) \overline{g(x)} \mathrm{d} x=\int_{\Omega} f(y) \overline{\left(\int_{\Omega} \overline{k(x, y)} g(x) \mathrm{d} x\right)} \mathrm{d} y \\
& =\int_{\Omega} f(y) \overline{\left(\int_{\Omega} k^{*}(y, x) g(x) \mathrm{d} x\right)} \mathrm{d} y=\left\langle f, K^{*} g\right\rangle
\end{aligned}
$$

5.41 Satz. Für alle $T \in \mathcal{L}(H)$ gilt $\operatorname{Kern} T^{*}=(\operatorname{Bild} T)^{\perp}$ und $\overline{\operatorname{Bild} T^{*}}=(\operatorname{Kern} T)^{\perp}$.

Beweis. Für $y \in H$ gilt

$$
y \in \operatorname{Kern} T^{*} \quad \Longleftrightarrow \quad\left\langle x, T^{*} y\right\rangle=0 \quad \text { für alle } x \in H
$$

$$
\Longleftrightarrow \quad\langle T x, y\rangle=0 \quad \text { für alle } x \in H \quad \Longleftrightarrow \quad y \in(\operatorname{Bild} T)^{\perp} .
$$

Damit folgt $\operatorname{Kern} T^{*}=(\operatorname{Bild} T)^{\perp}$. Weiterhin gilt

$$
\operatorname{Kern} T=\operatorname{Kern} T^{* *}=\left(\operatorname{Bild} T^{*}\right)^{\perp}, \text { also }(\operatorname{Kern} T)^{\perp}=\left(\operatorname{Bild} T^{*}\right)^{\perp \perp}=\overline{\operatorname{Bild} T^{*}}
$$

5.42 Satz. $P \in \mathcal{L}(H)$ ist eine orthogonale Projektion genau dann, wenn $P^{2}=P=P^{*}$ gilt.

Beweis. „": Da $P^{2}=P$ gilt, ist $P$ eine Projektion. Ferner ist Bild $P=\operatorname{Kern}(I-P)$ abgeschlossen, da $P$ stetig ist. Da schließlich $P^{*}=P$ ist, folgt aus Satz 5.41 Kern $P=$ $\operatorname{Kern} P^{*}=(\operatorname{Bild} P)^{\perp}$. Somit ist $P$ eine orthogonale Projektion.\\
„"": Nach Satz und Definition 5.7 gilt $P^{2}=P$, und für $Q=I-P$ gilt Bild $Q=(\operatorname{Bild} P)^{\perp}$. Also gilt für alle $x, y \in H$ :

$$
\langle P x, y\rangle=\langle P x, P y+Q y\rangle=\langle P x, P y\rangle=\langle P x+Q x, P y\rangle=\langle x, P y\rangle
$$

Mit Satz und Definition 5.36 folgt $P=P^{*}$.\\
5.43 Bemerkung. Für $T \in \mathcal{L}(H)$ ist $T$ genau dann unitär, wenn $T$ eine surjektive Isometrie ist.

Beweis. Es gilt

$$
\begin{aligned}
& T \text { surjektive Isometrie } \Longleftrightarrow \\
& T \text { surjektiv und }\|T x\|=\|x\| \text { für alle } x \in H \quad \Longleftrightarrow \\
& T \text { bijektiv und }\langle T x, T y\rangle=\langle x, y\rangle \text { für alle } x, y \in H \quad \Longleftrightarrow \\
& T \text { bijektiv und } T^{*} T=I \text {, also } T^{-1}=T^{*} \Longleftrightarrow \\
& T \text { unitär. }
\end{aligned}
$$

Bei der zweiten Äquivalenz haben wir die Polarisationsgleichungen

$$
\begin{aligned}
& \langle x, y\rangle=\frac{1}{2}\left(\|x+y\|^{2}-\|x\|^{2}-\|y\|^{2}\right) \quad \text { falls } \mathbb{K}=\mathbb{R} \\
& \langle x, y\rangle=\frac{1}{2}\left(\|x+y\|^{2}-\|x\|^{2}-\|y\|^{2}\right)+\frac{i}{2}\left(\|x+i y\|^{2}-\|x\|^{2}-\|y\|^{2}\right) \quad \text { falls } \mathbb{K}=\mathbb{C}
\end{aligned}
$$

(Übung) verwendet.\\
5.44 Bemerkung. Ist $\mathbb{K}=\mathbb{C}$ und $T \in \mathcal{L}(H)$, so ist $T^{*}$ bereits durch

$$
\langle T x, x\rangle=\left\langle x, T^{*} x\right\rangle \quad \text { für alle } x \in H
$$

eindeutig festgelegt.\\
Beweis. Seien $x, y \in H$. Anwendung von (5.22) auf $x+\mathrm{i} y$ und $x+y$ liefert die Gleichung $\langle T x, y\rangle=\left\langle x, T^{*} y\right\rangle$ (Übung!).


\end{document}